\input{preamble}


% OK, start here.
%
\begin{document}

\title{Catégories}


\maketitle

\phantomsection
\label{section-phantom}

\tableofcontents

\section{Introduction}
\label{section-introduction}

\noindent
Les catégories ont été introduites dans \cite{GenEqui}.
La catégorie des catégories (qui est une classe propre)
est une $2$-catégorie. De même, la catégorie des champs
forme une $2$-catégorie. Si vous connaissez déjà
les catégories, mais pas les $2$-catégories, vous devriez lire
la section \ref{section-formal-cat-cat}
comme introduction aux définitions formelles données plus loin.

\section{Définitions}
\label{section-definition-categories}

\noindent
Nous rappelons les définitions, en partie pour fixer les notations.

\begin{definition}
\label{definition-category}
Une {\it catégorie} $\mathcal{C}$ est constituée des données suivantes :
\begin{enumerate}
\item Un ensemble d'objets $\Ob(\mathcal{C})$.
\item Pour chaque paire $x, y \in \Ob(\mathcal{C})$, un ensemble de morphismes
$\Mor_\mathcal{C}(x, y)$.
\item Pour chaque triplet $x, y, z\in \Ob(\mathcal{C})$, une application de
composition $ \Mor_\mathcal{C}(y, z) \times \Mor_\mathcal{C}(x, y)
\to \Mor_\mathcal{C}(x, z) $, notée $(\phi, \psi) \mapsto
\phi \circ \psi$.
\end{enumerate}
Ces données doivent satisfaire aux règles suivantes :
\begin{enumerate}
\item Pour tout élément $x\in \Ob(\mathcal{C})$, il existe un
morphisme $\text{id}_x\in \Mor_\mathcal{C}(x, x)$ tel que
$\text{id}_x \circ \phi = \phi$ et $\psi \circ \text{id}_x = \psi $ dès que
ces compositions ont un sens.
\item La composition est associative, c'est-à-dire
$(\phi \circ \psi) \circ \chi =
\phi \circ ( \psi \circ \chi)$ dès que ces compositions ont un sens.
\end{enumerate}
\end{definition}

\noindent
Il est d'usage d'exiger que tous les ensembles de morphismes
$\Mor_\mathcal{C}(x, y)$ soient disjoints.
Ainsi, un morphisme $\phi : x \to y$ possède une {\it source} $x$
unique et un {\it but} $y$ unique. Cela n'est pas strictement nécessaire,
bien qu'il faille alors prendre garde à la formulation de la condition (2)
ci-dessus. Cette convention est commode et nous la supposerons souvent
satisfaite. Dans ce cas, nous disons que $\phi$ et $\psi$ sont
{\it composables} si la source de $\phi$ est égale au
but de $\psi$ ; alors $\phi \circ \psi$ est défini.
Une définition équivalente consisterait à définir une catégorie
comme un quintuplet $(\text{Ob}, \text{Flèches}, s, t, \circ)$
formé d'un ensemble d'objets, d'un ensemble de morphismes (flèches),
d'une source, d'un but et d'une composition, soumis à une longue liste
d'axiomes. Nous adopterons occasionnellement ce point de vue.

\begin{remark}
\label{remark-big-categories}
Grandes catégories. Dans certains textes, une catégorie peut avoir une classe
propre d'objets. Nous l'autoriserons également dans ces notes, mais seulement
dans les cas de la liste suivante (qui sera complétée au fur et à mesure).
En particulier, lorsque nous disons : « Soit $\mathcal{C}$ une catégorie »,
il est entendu que $\Ob(\mathcal{C})$ est un ensemble.
\begin{enumerate}
\item La catégorie $\textit{Ensembles}$ des ensembles.
\item La catégorie $\textit{Ab}$ des groupes abéliens.
\item La catégorie $\textit{Groups}$ des groupes.
\item Étant donné un groupe $G$, la catégorie $G\textit{-Ensembles}$ des
ensembles munis d'une action à gauche de $G$.
\item Étant donné un anneau $R$, la catégorie $\text{Mod}_R$ des $R$-modules.
\item Étant donné un corps $k$, la catégorie des espaces vectoriels sur $k$.
\item La catégorie des anneaux.
\item La catégorie des anneaux à puissances divisées, voir
Algèbre à puissances divisées, section \ref{dpa-section-divided-power-rings}.
\item La catégorie des schémas.
\item La catégorie $\textit{Top}$ des espaces topologiques.
\item Étant donné un espace topologique $X$, la catégorie
$\textit{PSh}(X)$ des préfaisceaux d'ensembles sur $X$.
\item Étant donné un espace topologique $X$, la catégorie
$\Sh(X)$ des faisceaux d'ensembles sur $X$.
\item Étant donné un espace topologique $X$, la catégorie
$\textit{PAb}(X)$ des préfaisceaux de groupes abéliens sur $X$.
\item Étant donné un espace topologique $X$, la catégorie
$\textit{Ab}(X)$ des faisceaux de groupes abéliens sur $X$.
\item Étant donné une petite catégorie $\mathcal{C}$, la catégorie des foncteurs
de $\mathcal{C}$ dans $\textit{Ensembles}$.
\item Étant donné une catégorie $\mathcal{C}$, la catégorie des préfaisceaux
d'ensembles sur $\mathcal{C}$.
\item Étant donné un site $\mathcal{C}$, la catégorie des faisceaux
d'ensembles sur $\mathcal{C}$.
\end{enumerate}
L'une des raisons de cette énumération est d'éviter de travailler
avec quelque chose comme la « collection » des « grandes » catégories,
ce qui reviendrait à travailler avec la collection de toutes les classes
et constituerait, à notre avis, un objet définitivement métamathématique.
\end{remark}

\begin{remark}
\label{remark-unique-identity}
Il découle directement de la définition que deux morphismes identités
d'un objet $x$ de $\mathcal{A}$ sont égaux. Nous pouvons donc parler,
et parlerons, {\it du} morphisme identité $\text{id}_x$ de $x$.
\end{remark}

\begin{definition}
\label{definition-isomorphism}
Un morphisme $\phi : x \to y$ est un {\it isomorphisme} de la catégorie
$\mathcal{C}$ s'il existe un morphisme $\psi : y \to x$
tel que $\phi \circ \psi = \text{id}_y$ et
$\psi \circ \phi = \text{id}_x$.
\end{definition}

\noindent
Un isomorphisme $\phi$ est aussi parfois appelé morphisme
{\it inversible}, et le morphisme $\psi$ de la définition est appelé
l'{\it inverse} et noté $\phi^{-1}$. Il est unique s'il existe. Remarquons que,
pour un objet $x$ d'une catégorie $\mathcal{A}$, l'ensemble des éléments
inversibles $\text{Aut}_\mathcal{A}(x)$
de $\Mor_\mathcal{A}(x, x)$ forme un groupe pour la composition.
Ce groupe est appelé le groupe des {\it automorphismes} de $x$ dans
$\mathcal{A}$.

\begin{definition}
\label{definition-groupoid}
Un {\it groupoïde} est une catégorie dans laquelle tout morphisme est
un isomorphisme.
\end{definition}

\begin{example}
\label{example-group-groupoid}
Un groupe $G$ définit un groupoïde ayant un unique objet $x$
et pour morphismes $\Mor(x, x) = G$, la loi de composition
étant donnée par la loi de groupe de $G$. Tout groupoïde ayant un unique
objet est de cette forme.
\end{example}

\begin{example}
\label{example-set-groupoid}
Un ensemble $C$ définit un groupoïde $\mathcal{C}$ de la manière suivante :
on pose $\Ob(\mathcal{C}) := C$ et, pour les morphismes,
on prend $\Mor(x, y)$ vide si $x\neq y$ et égal à
$\{\text{id}_x\}$ si $x = y$.
\end{example}

\begin{definition}
\label{definition-functor}
Un {\it foncteur} $F : \mathcal{A} \to \mathcal{B}$
entre deux catégories $\mathcal{A}, \mathcal{B}$ est défini par les
données suivantes :
\begin{enumerate}
\item Une application $F : \Ob(\mathcal{A}) \to \Ob(\mathcal{B})$.
\item Pour tous $x, y \in \Ob(\mathcal{A})$, une application
$F : \Mor_\mathcal{A}(x, y) \to \Mor_\mathcal{B}(F(x), F(y))$,
notée $\phi \mapsto F(\phi)$.
\end{enumerate}
Ces données doivent être compatibles avec la composition et les morphismes
identités de la manière suivante : $F(\phi \circ \psi) =
F(\phi) \circ F(\psi)$ pour toute paire composable $(\phi, \psi)$ de
morphismes de $\mathcal{A}$, et $F(\text{id}_x) = \text{id}_{F(x)}$.
\end{definition}

\noindent
Toute catégorie $\mathcal{A}$ possède un foncteur
{\it identité} $\text{id}_\mathcal{A}$.
De plus, étant donnés un foncteur $G : \mathcal{B} \to \mathcal{C}$
et un foncteur $F : \mathcal{A} \to \mathcal{B}$, il existe
un foncteur {\it composé} $G \circ F : \mathcal{A} \to \mathcal{C}$
défini de manière évidente.

\begin{definition}
\label{definition-faithful}
Soit $F : \mathcal{A} \to \mathcal{B}$ un foncteur.
\begin{enumerate}
\item Nous disons que $F$ est {\it fidèle} si,
pour tous objets $x, y \in \Ob(\mathcal{A})$, l'application
$$
F : \Mor_\mathcal{A}(x, y) \to \Mor_\mathcal{B}(F(x), F(y))
$$
est injective.
\item Si toutes ces applications sont bijectives, $F$ est dit
{\it pleinement fidèle}.
\item
Le foncteur $F$ est dit {\it essentiellement surjectif} si, pour tout
objet $y \in \Ob(\mathcal{B})$, il existe un objet
$x \in \Ob(\mathcal{A})$ tel que $F(x)$ soit isomorphe à $y$ dans
$\mathcal{B}$.
\end{enumerate}
\end{definition}

\begin{definition}
\label{definition-subcategory}
Une {\it sous-catégorie} d'une catégorie $\mathcal{B}$ est une catégorie
$\mathcal{A}$ dont les objets et les flèches forment des sous-ensembles des
objets et des flèches de $\mathcal{B}$, telle que la source, le but et la
composition dans $\mathcal{A}$ coïncident avec ceux de $\mathcal{B}$ et que
le morphisme identité d'un objet de $\mathcal{A}$ coïncide avec celui de
$\mathcal{B}$. Nous disons que $\mathcal{A}$ est une
{\it sous-catégorie pleine} de $\mathcal{B}$ si
$\Mor_\mathcal{A}(x, y)
= \Mor_\mathcal{B}(x, y)$ pour tous
$x, y \in \Ob(\mathcal{A})$.
Nous disons que $\mathcal{A}$ est une sous-catégorie
{\it strictement pleine} de $\mathcal{B}$ si elle est pleine et si,
pour tout $x \in \Ob(\mathcal{A})$, tout objet de $\mathcal{B}$
isomorphe à $x$ appartient également à $\mathcal{A}$.
\end{definition}

\noindent
Si $\mathcal{A} \subset \mathcal{B}$ est une sous-catégorie, l'application
identité est un foncteur de $\mathcal{A}$ dans $\mathcal{B}$.
De plus, une sous-catégorie $\mathcal{A} \subset \mathcal{B}$
est pleine si et seulement si le foncteur d'inclusion est pleinement fidèle.
Remarquons que, pour une catégorie $\mathcal{B}$, l'ensemble des
sous-catégories pleines de $\mathcal{B}$ s'identifie à l'ensemble des sous-ensembles de
$\Ob(\mathcal{B})$.

\begin{remark}
\label{remark-functor-into-sets}
Supposons que $\mathcal{A}$ soit une catégorie.
Un foncteur $F$ de $\mathcal{A}$ dans $\textit{Ensembles}$
est un objet mathématique (c'est-à-dire un ensemble, et non une classe
ni une formule de la théorie des ensembles ; voir
Ensembles, section \ref{sets-section-sets-everything}),
bien que la catégorie des ensembles soit « grande ».
En effet, l'image de $F$ sur les objets est
un ensemble $F(\Ob(\mathcal{A}))$ et nous pouvons alors
considérer $F$ comme un foncteur entre
$\mathcal{A}$ et la sous-catégorie pleine
de la catégorie des ensembles dont les objets
sont les éléments de $F(\Ob(\mathcal{A}))$.
\end{remark}

\begin{example}
\label{example-group-homomorphism-functor}
Un homomorphisme de groupes $p : G\to H$ définit un foncteur entre
les groupoïdes associés de l'exemple \ref{example-group-groupoid}. Il est
fidèle (resp.\ pleinement fidèle) si et seulement si $p$ est injectif
(resp.\ est un isomorphisme).
\end{example}

\begin{example}
\label{example-category-over-X}
Étant donnés une catégorie $\mathcal{C}$ et un objet
$X\in \Ob(\mathcal{C})$, nous définissons la
{\it catégorie des objets au-dessus de $X$},
notée $\mathcal{C}/X$, comme suit.
Les objets de $\mathcal{C}/X$ sont les morphismes $Y\to X$,
où $Y\in \Ob(\mathcal{C})$. Les morphismes entre les objets
$Y\to X$ et $Y'\to X$ sont les morphismes $Y\to Y'$ de $\mathcal{C}$
qui rendent commutatif le diagramme évident. Il existe un foncteur
$p_X : \mathcal{C}/X\to \mathcal{C}$ qui oublie simplement le morphisme.
De plus, tout morphisme $f : X'\to X$ de $\mathcal{C}$ induit un foncteur
$F : \mathcal{C}/X' \to \mathcal{C}/X$ obtenu par composition avec $f$,
et $p_X\circ F = p_{X'}$.
\end{example}

\begin{example}
\label{example-category-under-X}
Étant donnés une catégorie $\mathcal{C}$ et un objet
$X\in \Ob(\mathcal{C})$, nous définissons la
{\it catégorie des objets au-dessous de $X$},
notée $X/\mathcal{C}$, comme suit.
Les objets de $X/\mathcal{C}$ sont les morphismes $X\to Y$,
où $Y\in \Ob(\mathcal{C})$. Les morphismes entre les objets
$X\to Y$ et $X\to Y'$ sont les morphismes $Y\to Y'$ de $\mathcal{C}$
qui rendent commutatif le diagramme évident. Il existe un foncteur
$p_X : X/\mathcal{C}\to \mathcal{C}$ qui oublie simplement le morphisme.
De plus, tout morphisme $f : X'\to X$ de $\mathcal{C}$ induit un foncteur
$F : X/\mathcal{C} \to X'/\mathcal{C}$
obtenu par composition avec $f$,
et $p_{X'}\circ F = p_X$.
\end{example}




\begin{definition}
\label{definition-transformation-functors}
Soient $F, G : \mathcal{A} \to \mathcal{B}$ des foncteurs.
Une {\it transformation naturelle}, ou un {\it morphisme de foncteurs},
$t : F \to G$, est une famille $\{t_x\}_{x\in \Ob(\mathcal{A})}$
telle que
\begin{enumerate}
\item $t_x : F(x) \to G(x)$ soit un morphisme de la catégorie $\mathcal{B}$, et
\item pour tout morphisme $\phi : x \to y$ de $\mathcal{A}$, le diagramme
suivant soit commutatif :
$$
\xymatrix{
F(x) \ar[r]^{t_x} \ar[d]_{F(\phi)} & G(x) \ar[d]^{G(\phi)} \\
F(y) \ar[r]^{t_y} & G(y) }
$$
\end{enumerate}
\end{definition}

\noindent
Nous utilisons parfois le diagramme
$$
\xymatrix{
\mathcal{A}
\rtwocell^F_G{t}
&
\mathcal{B}
}
$$
pour indiquer que $t$ est un morphisme de $F$ vers $G$.

\medskip\noindent
Tout foncteur $F$ est muni de la transformation
{\it identité} $\text{id}_F : F \to F$. De plus, étant donnés un morphisme de
foncteurs $t : F \to G$ et un morphisme de foncteurs $s : E \to F$,
leur {\it composé} $t \circ s$ est défini par la règle
$$
(t \circ s)_x = t_x \circ s_x : E(x) \to G(x)
$$
pour $x \in \Ob(\mathcal{A})$.
On vérifie facilement qu'il s'agit bien d'un morphisme de foncteurs
de $E$ vers $G$.
Ainsi, à partir de deux catégories
$\mathcal{A}$ et $\mathcal{B}$, nous obtenons une nouvelle catégorie,
à savoir la catégorie des foncteurs de $\mathcal{A}$ dans $\mathcal{B}$.

\begin{remark}
\label{remark-functors-sets-sets}
C'est l'un des cas où le même énoncé ne vaut pas si
$\mathcal{A}$ est une « grande » catégorie. Considérons par exemple les
foncteurs $\textit{Ensembles} \to \textit{Ensembles}$. Avec notre définition actuelle,
un tel foncteur est une classe et non un ensemble. Autrement dit, il est
donné par une formule de la théorie des ensembles (dont certaines variables
sont égales à des ensembles prescrits) ! Il n'est pas judicieux d'essayer
d'intégrer toutes les formules possibles de la théorie des ensembles à la
définition d'un objet mathématique. Le même problème se présente, par exemple,
lorsqu'on considère les faisceaux sur la catégorie des schémas.
Nous reviendrons plus tard sur ce point.
\end{remark}

\begin{definition}
\label{definition-equivalence-categories}
Une {\it équivalence de catégories}
$F : \mathcal{A} \to \mathcal{B}$ est un foncteur tel qu'il existe
un foncteur $G : \mathcal{B} \to \mathcal{A}$ pour lequel
les composés $F \circ G$ et $G \circ F$ sont respectivement isomorphes aux
foncteurs identité $\text{id}_\mathcal{B}$ et
$\text{id}_\mathcal{A}$.
Dans ce cas, nous disons que $G$ est un {\it quasi-inverse} de $F$.
\end{definition}

\begin{lemma}
\label{lemma-construct-quasi-inverse}
Soit $F : \mathcal{A} \to \mathcal{B}$ un foncteur pleinement fidèle.
Supposons donnés, pour tout $X \in \Ob(\mathcal{B})$, un
objet $j(X)$ de $\mathcal{A}$ et un isomorphisme $i_X : X \to F(j(X))$.
Il existe alors un unique foncteur $j : \mathcal{B} \to \mathcal{A}$
tel que l'action de $j$ sur les objets soit celle qui vient d'être prescrite et
pour lequel les isomorphismes
$i_X$ définissent un isomorphisme de foncteurs
$\text{id}_\mathcal{B} \to F \circ j$. De plus, $j$ et $F$
sont des équivalences de catégories quasi-inverses.
\end{lemma}

\begin{proof}
La construction de $j : \mathcal{B} \to \mathcal{A}$ se fait en deux étapes.
Premièrement, nous définissons l'application
$j : \Ob(\mathcal{B}) \to \Ob(\mathcal{A})$ qui associe
$j(X)$ à $X \in \mathcal{B}$. Deuxièmement, si
$X,Y \in \Ob(\mathcal{B})$ et $\phi : X \to Y$, nous considérons
$\phi' := i_Y \circ\phi\circ i_X^{-1}$. Il existe un unique
$\varphi$ tel que $F(\varphi) = \phi'$, puisque $F$ est pleinement fidèle.
Nous posons $j(\phi) = \varphi$. Nous omettons la vérification que $j$
est un foncteur. Par construction, le diagramme
$$
\xymatrix{
X \ar[r]_-{i_X} \ar[d]_{\phi}
&
F(j(X)) \ar[d]^{F(j(\phi))}
\\
Y \ar[r]^-{i_Y}
&
F(j(Y))
}
$$
est commutatif. Puisque chaque $i_X$ est un isomorphisme, la famille
$\{i_X\}_X$ est donc un isomorphisme de foncteurs
$\text{id}_\mathcal{B} \to F\circ j$. Pour conclure, il reste à montrer que
$j \circ F$ est isomorphe à $\text{id}_\mathcal{A}$.
Cependant, puisque $F$ est pleinement fidèle, il suffit pour cela de le
démontrer après composition à gauche par $F$, c'est-à-dire qu'il suffit de
montrer que $F \circ j \circ F$ est isomorphe à
$F \circ \text{id}_\mathcal{A}$
(nous omettons ce petit détail). Cela est clair puisque
$F \circ j \cong \text{id}_\mathcal{B}$.
\end{proof}

\begin{lemma}
\label{lemma-equivalence-categories}
Un foncteur est une équivalence de catégories si et seulement s'il est
pleinement fidèle et essentiellement surjectif.
\end{lemma}

\begin{proof}
Soit $F : \mathcal{A} \to \mathcal{B}$ essentiellement surjectif et pleinement
fidèle. Puisque, par convention, toutes les catégories sont petites et que
$F$ est essentiellement surjectif, l'axiome du choix permet de choisir, pour
tout $X \in \Ob(\mathcal{B})$, un objet $j(X)$ de $\mathcal{A}$ et un
isomorphisme $i_X : X \to F(j(X))$. Nous appliquons alors le
lemme \ref{lemma-construct-quasi-inverse},
en utilisant le fait que $F$ est pleinement fidèle.
\end{proof}

\begin{definition}
\label{definition-product-category}
Soient $\mathcal{A}$ et $\mathcal{B}$ des catégories.
Nous définissons la {\it catégorie produit}
$\mathcal{A} \times \mathcal{B}$ comme la catégorie dont les objets sont
$\Ob(\mathcal{A} \times \mathcal{B}) =
\Ob(\mathcal{A}) \times \Ob(\mathcal{B})$
et telle que
$$
\Mor_{\mathcal{A} \times \mathcal{B}}((x, y), (x', y'))
:=
\Mor_\mathcal{A}(x, x')\times
\Mor_\mathcal{B}(y, y').
$$
La composition est définie composante par composante.
\end{definition}


\section{Catégories opposées et lemme de Yoneda}
\label{section-opposite}

\begin{definition}
\label{definition-opposite}
Étant donnée une catégorie $\mathcal{C}$, la {\it catégorie opposée}
$\mathcal{C}^{opp}$ est la catégorie qui a les mêmes objets
que $\mathcal{C}$, mais dont tous les morphismes sont renversés.
\end{definition}

\noindent
Autrement dit,
$\Mor_{\mathcal{C}^{opp}}(x, y) = \Mor_\mathcal{C}(y, x)$.
La composition dans $\mathcal{C}^{opp}$ est celle de $\mathcal{C}$,
mais prise à l'envers : si $\phi : y \to z$ et $\psi : x \to y$
sont des morphismes de $\mathcal{C}^{opp}$, c'est-à-dire des flèches
$z \to y$ et $y \to x$ dans $\mathcal{C}$,
alors $\phi \circ^{opp} \psi$ est le morphisme $x \to z$
de $\mathcal{C}^{opp}$ qui correspond au composé
$z \to y \to x$ dans $\mathcal{C}$.

\begin{definition}
\label{definition-contravariant}
Soient $\mathcal{C}$ et $\mathcal{S}$ des catégories.
Un foncteur {\it contravariant} $F$
de $\mathcal{C}$ dans $\mathcal{S}$
est un foncteur $\mathcal{C}^{opp}\to \mathcal{S}$.
\end{definition}

\noindent
Concrètement, un foncteur contravariant $F$ est donné
par une application $F : \Ob(\mathcal{C}) \to
\Ob(\mathcal{S})$ et, pour tout morphisme
$\psi : x \to y$ dans $\mathcal{C}$, par un morphisme
$F(\psi) : F(y) \to F(x)$. Ces données doivent vérifier la propriété
suivante : étant donné un autre morphisme
$\phi : y \to z$, on a $F(\phi \circ \psi)
= F(\psi) \circ F(\phi)$ comme morphismes $F(z) \to F(x)$.
(Remarquer l'inversion de l'ordre.)

\begin{definition}
\label{definition-presheaf}
Soit $\mathcal{C}$ une catégorie.
\begin{enumerate}
\item Un {\it préfaisceau d'ensembles sur $\mathcal{C}$},
ou simplement un {\it préfaisceau}, est un foncteur contravariant
$F$ de $\mathcal{C}$ dans $\textit{Ensembles}$.
\item La catégorie des préfaisceaux est notée $\textit{PSh}(\mathcal{C})$.
\end{enumerate}
\end{definition}

\noindent
Bien entendu, la catégorie des préfaisceaux est une classe propre.

\begin{example}
\label{example-hom-functor}
Foncteur des points.
Pour tout $U\in \Ob(\mathcal{C})$, il existe un foncteur contravariant
$$
\begin{matrix}
h_U & : & \mathcal{C}
&
\longrightarrow
&
\textit{Ensembles} \\
& &
X
&
\longmapsto
&
\Mor_\mathcal{C}(X, U)
\end{matrix}
$$
qui associe à un objet $X$ l'ensemble
$\Mor_\mathcal{C}(X, U)$. Autrement dit, $h_U$ est un préfaisceau.
À un morphisme $f : X\to Y$, l'application correspondante
$h_U(f) : \Mor_\mathcal{C}(Y, U)\to \Mor_\mathcal{C}(X, U)$
envoie $\phi$ sur $\phi\circ f$. Nous noterons toujours ce préfaisceau
$h_U : \mathcal{C}^{opp} \to \textit{Ensembles}$.
On l'appelle le {\it préfaisceau représentable} associé à $U$.
Si $\mathcal{C}$ est la catégorie des schémas, ce foncteur est
parfois appelé le \emph{foncteur des points} de $U$.
\end{example}

\noindent
Remarquons qu'à un morphisme $\phi : U \to V$ de $\mathcal{C}$ correspond
une transformation naturelle de foncteurs $h(\phi) : h_U \to h_V$,
définie par composition avec le morphisme $U \to V$. Ainsi,
la composition des morphismes dans $\mathcal{C}$ devient la composition
des transformations de foncteurs. Autrement dit, nous obtenons un foncteur
$$
h :
\mathcal{C}
\longrightarrow
\textit{PSh}(\mathcal{C})
$$
Remarquons que le but est une « grande » catégorie ; voir la
remarque \ref{remark-big-categories}. En revanche,
$h$ est bien un objet mathématique (c'est-à-dire\ un ensemble) ; comparer
la remarque \ref{remark-functor-into-sets}.

\begin{lemma}[Lemme de Yoneda]
\label{lemma-yoneda}
\begin{reference}
Est apparu sous une certaine forme dans \cite{Yoneda-homology}. Utilisé par
Grothendieck sous une forme généralisée dans \cite{Gr-II}.
\end{reference}
Soient $U, V \in \Ob(\mathcal{C})$.
Pour tout morphisme de foncteurs $s : h_U \to h_V$,
il existe un unique morphisme $\phi : U \to V$
tel que $h(\phi) = s$. Autrement dit, le
foncteur $h$ est pleinement fidèle. Plus généralement,
pour tout foncteur contravariant $F$ et tout objet
$U$ de $\mathcal{C}$, on a une bijection naturelle
$$
\Mor_{\textit{PSh}(\mathcal{C})}(h_U, F) \longrightarrow F(U),
\quad
s \longmapsto s_U(\text{id}_U).
$$
\end{lemma}

\begin{proof}
Pour le premier énoncé, il suffit de prendre
$\phi = s_U(\text{id}_U) \in \Mor_\mathcal{C}(U, V)$.
Pour le second, étant donné $\xi \in F(U)$, on définit
$s$ par $s_V : h_U(V) \to F(V)$ en envoyant l'élément $f : V \to U$
de $h_U(V) = \Mor_\mathcal{C}(V, U)$ sur $F(f)(\xi)$.
\end{proof}

\begin{definition}
\label{definition-representable-functor}
Un foncteur contravariant $F : \mathcal{C}\to \textit{Ensembles}$ est dit
{\it représentable} s'il est isomorphe au foncteur des points
$h_U$ pour un certain objet $U$ de $\mathcal{C}$.
\end{definition}

\noindent
Soit $\mathcal{C}$ une catégorie et soit
$F : \mathcal{C}^{opp} \to \textit{Ensembles}$ un foncteur représentable.
Choisissons un objet $U$ de $\mathcal{C}$ et un isomorphisme $s : h_U \to F$.
Le lemme de Yoneda garantit que le couple $(U, s)$
est unique à isomorphisme unique près. L'objet
$U$ est appelé un objet {\it représentant} $F$.
Par le lemme de Yoneda, la transformation $s$ correspond à un unique
élément $\xi \in F(U)$. Cet élément est appelé l'{\it objet universel}.
Il possède la propriété que, pour $V \in \Ob(\mathcal{C})$, l'application
$$
\Mor_\mathcal{C}(V, U) \longrightarrow F(V),\quad
(f : V \to U) \longmapsto F(f)(\xi)
$$
est une bijection. Ainsi, $\xi$ est universel en ce sens que tout élément
de $F(V)$ est l'image de $\xi$ par $F(f)$ pour un unique morphisme
$f : V \to U$ dans $\mathcal{C}$.






\section{Produits de deux objets}
\label{section-products-pairs}

\begin{definition}
\label{definition-products}
Soient $x, y\in \Ob(\mathcal{C})$.
Un {\it produit} de $x$ et $y$ est
un objet $x \times y \in \Ob(\mathcal{C})$
muni de morphismes
$p\in \Mor_{\mathcal C}(x \times y, x)$ et
$q\in\Mor_{\mathcal C}(x \times y, y)$
satisfaisant la propriété universelle suivante : pour tout
$w\in \Ob(\mathcal{C})$ et tous morphismes
$\alpha \in \Mor_{\mathcal C}(w, x)$ et
$\beta \in \Mor_\mathcal{C}(w, y)$,
il existe un unique
$\gamma\in \Mor_{\mathcal C}(w, x \times y)$ qui rend le diagramme
$$
\xymatrix{
w \ar[rrrd]^\beta \ar@{-->}[rrd]_\gamma \ar[rrdd]_\alpha & & \\
& & x \times y \ar[d]_p \ar[r]_q & y \\
& & x &
}
$$
commutatif.
\end{definition}

\noindent
S'il existe, un produit est unique à isomorphisme unique près.
Cela découle du lemme de Yoneda, car la définition exige que
$x \times y$ soit un objet de $\mathcal{C}$ tel que
$$
h_{x \times y}(w) = h_x(w) \times h_y(w)
$$
fonctoriellement en $w$. Autrement dit, le produit $x \times y$
est un objet représentant le foncteur
$w \mapsto h_x(w) \times h_y(w)$.

\begin{definition}
\label{definition-has-products-of-pairs}
Nous disons que la catégorie $\mathcal{C}$ {\it admet les produits de deux
objets} si un produit $x \times y$ existe pour tous
$x, y \in \Ob(\mathcal{C})$.
\end{definition}

\noindent
Nous employons cette terminologie pour distinguer cette notion de celles
d'« admettre les produits » ou d'« admettre les produits finis », qui ont
habituellement un autre sens (elles impliquent notamment toujours l'existence
d'un objet final).





\section{Coproduits de deux objets}
\label{section-coproducts-pairs}

\begin{definition}
\label{definition-coproducts}
Soient $x, y \in \Ob(\mathcal{C})$.
Un {\it coproduit}, ou une {\it somme}, de $x$ et $y$ est
un objet $x \amalg y \in \Ob(\mathcal{C})$
muni de morphismes
$i \in \Mor_{\mathcal C}(x, x \amalg y)$ et
$j \in \Mor_{\mathcal C}(y, x \amalg y)$
satisfaisant la propriété universelle suivante : pour tout
$w \in \Ob(\mathcal{C})$ et tous morphismes
$\alpha \in \Mor_{\mathcal C}(x, w)$ et
$\beta \in \Mor_\mathcal{C}(y, w)$,
il existe un unique
$\gamma \in \Mor_{\mathcal C}(x \amalg y, w)$ qui rend le diagramme
$$
\xymatrix{
& y \ar[d]^j \ar[rrdd]^\beta \\
x \ar[r]^i \ar[rrrd]_\alpha & x \amalg y \ar@{-->}[rrd]^\gamma \\
& & & w
}
$$
commutatif.
\end{definition}

\noindent
S'il existe, un coproduit est unique à isomorphisme unique près.
Cela découle du lemme de Yoneda (appliqué à la catégorie opposée),
car la définition exige que $x \amalg y$ soit un objet
de $\mathcal{C}$ tel que
$$
\Mor_\mathcal{C}(x \amalg y, w) =
\Mor_\mathcal{C}(x, w) \times \Mor_\mathcal{C}(y, w)
$$
fonctoriellement en $w$.

\begin{definition}
\label{definition-has-coproducts-of-pairs}
Nous disons que la catégorie $\mathcal{C}$ {\it admet les coproduits de deux
objets} si un coproduit $x \amalg y$ existe pour tous
$x, y \in \Ob(\mathcal{C})$.
\end{definition}

\noindent
Nous employons cette terminologie pour distinguer cette notion de celles
d'« admettre les coproduits » ou d'« admettre les coproduits finis », qui ont
habituellement un autre sens (elles impliquent notamment toujours l'existence
d'un objet initial dans $\mathcal{C}$).





\section{Produits fibrés}
\label{section-fibre-products}

\begin{definition}
\label{definition-fibre-products}
Soient $x, y, z\in \Ob(\mathcal{C})$,
$f\in \Mor_\mathcal{C}(x, y)$
et $g\in \Mor_{\mathcal C}(z, y)$.
Un {\it produit fibré} de $f$ et $g$ est
un objet $x \times_y z\in \Ob(\mathcal{C})$
muni de morphismes
$p \in \Mor_{\mathcal C}(x \times_y z, x)$ et
$q \in \Mor_{\mathcal C}(x \times_y z, z)$ qui rendent le diagramme
$$
\xymatrix{
x \times_y z \ar[r]_q \ar[d]_p & z \ar[d]^g \\
x \ar[r]^f & y
}
$$
commutatif et qui satisfont la propriété universelle suivante : pour tout
$w\in \Ob(\mathcal{C})$ et tous morphismes
$\alpha \in \Mor_{\mathcal C}(w, x)$ et
$\beta \in \Mor_\mathcal{C}(w, z)$ tels que
$f \circ \alpha = g \circ \beta$,
il existe un unique
$\gamma \in \Mor_{\mathcal C}(w, x \times_y z)$ qui rend le diagramme
$$
\xymatrix{
w \ar[rrrd]^\beta \ar@{-->}[rrd]_\gamma \ar[rrdd]_\alpha & & \\
& & x \times_y z \ar[d]^p \ar[r]_q & z \ar[d]^g \\
& & x \ar[r]^f & y
}
$$
commutatif.
\end{definition}

\noindent
S'il existe, un produit fibré est unique à isomorphisme unique près.
Cela découle du lemme de Yoneda, car la définition exige que
$x \times_y z$ soit un objet de $\mathcal{C}$ tel que
$$
h_{x \times_y z}(w) = h_x(w) \times_{h_y(w)} h_z(w)
$$
fonctoriellement en $w$. Autrement dit, le produit fibré $x \times_y z$
est un objet représentant le foncteur
$w \mapsto h_x(w) \times_{h_y(w)} h_z(w)$.

\begin{definition}
\label{definition-cartesian}
Nous disons qu'un diagramme commutatif
$$
\xymatrix{
w \ar[r] \ar[d] &
z \ar[d] \\
x \ar[r] &
y
}
$$
dans une catégorie est {\it cartésien} si $w$ et les morphismes
$w \to x$ et $w \to z$ forment un produit fibré des morphismes
$x \to y$ et $z \to y$.
\end{definition}

\begin{definition}
\label{definition-has-fibre-products}
Nous disons que la catégorie $\mathcal{C}$ {\it admet les produits fibrés}
si le produit fibré existe pour tous
$f\in \Mor_{\mathcal C}(x, y)$ et
$g\in \Mor_{\mathcal C}(z, y)$.
\end{definition}

\begin{definition}
\label{definition-representable-morphism}
Un morphisme $f : x \to y$ d'une catégorie $\mathcal{C}$ est dit
{\it représentable} si, pour tout morphisme $z \to y$
dans $\mathcal{C}$, le produit fibré $x \times_y z$ existe.
\end{definition}

\begin{lemma}
\label{lemma-composition-representable}
Soit $\mathcal{C}$ une catégorie.
Soient $f : x \to y$ et $g : y \to z$ des morphismes représentables.
Alors $g \circ f : x \to z$ est représentable.
\end{lemma}

\begin{proof}
Soient $t \in \Ob(\mathcal C)$ et
$ \varphi \in \Mor_{\mathcal C}(t,z) $.
Puisque $g$ et $f$ sont représentables, nous obtenons des diagrammes
commutatifs
$$
\xymatrix{
y \times_z t \ar[r]_q \ar[d]_p &
t \ar[d]^{\varphi} \\
y \ar[r]^{g} & z
}
\quad\quad
\xymatrix{
x \times_y (y\! \times_z\! t) \ar[r]_{q'} \ar[d]_{p'} &
y \times_z t \ar[d]^p \\
x \ar[r]^f & y
}
$$
possédant la propriété universelle de la
définition \ref{definition-fibre-products}.
Nous affirmons que $x \times_z t = x \times_y (y \times_z t)$,
muni des morphismes $q \circ q' : x \times_z t \to t$ et
$p' : x \times_z t \to x$, est un produit fibré.
La commutativité des diagrammes ci-dessus donne d'abord
$\varphi \circ q \circ q' = g \circ f \circ p'$.
Pour vérifier la propriété universelle, soient
$w \in \Ob(\mathcal C)$ et des morphismes $\alpha : w \to x$ et
$\beta : w \to t$ tels que
$\varphi \circ \beta = g \circ f \circ \alpha$.
Par définition du produit fibré, il existe des morphismes uniques
$\delta$ et $\gamma$ tels que
$$
\xymatrix{
w \ar[rrrd]^\beta \ar@{-->}[rrd]_\delta \ar[rrdd]_{f\circ\alpha} & & \\
& & y \times_z t \ar[d]_p \ar[r]_q & t \ar[d]^{\varphi} \\
& & y \ar[r]^{g} & z
}
$$
et
$$
\xymatrix{
w \ar[rrrd]^\delta \ar@{-->}[rrd]_\gamma \ar[rrdd]_{\alpha} & & \\
& & x \times_y (y\!\times_z\! t) \ar[d]_{p'} \ar[r]_{q'} &
y \times_z t \ar[d]^{p} \\
& & x \ar[r]^{f} & y
}
$$
soient commutatifs. Alors $\gamma$ rend commutatif le diagramme
$$
\xymatrix{
w \ar[rrrd]^\beta \ar@{-->}[rrd]_\gamma \ar[rrdd]_{\alpha} & & \\
& & x \times_z t \ar[d]_{p'} \ar[r]_{q\circ q'} & t \ar[d]^{\varphi} \\
& & x \ar[r]^{g\circ f} & z
}
$$
Pour montrer son unicité, soit $\gamma'$ tel que
$ q\circ q'\circ\gamma' =
\beta $ et
$ p'\circ \gamma' = \alpha $. Par l'unicité de $\gamma$, il suffit de
prouver que $ q'\circ\gamma' = \delta $ et
$ p'\circ\gamma' = \alpha $. La seconde égalité est supposée.
Pour la première, nous utilisons aussi l'unicité de $\delta$.
En effet, $\delta$ est l'unique morphisme vérifiant
$ q\circ\delta = \beta $ et
$ p\circ\delta = f\circ\alpha $. Nous avons déjà supposé
$ q\circ (q'\circ\gamma') = \beta $. De plus, par définition du produit
fibré, $ f\circ p' = p\circ q' $. Par conséquent :
$$
p\circ (q'\circ\gamma') =
(p\circ q')\circ\gamma' =
(f\circ p')\circ\gamma' =
f\circ (p'\circ\gamma') = f\circ\alpha.
$$
Ainsi $ q'\circ\gamma' = \delta $, ce qui achève la preuve.
\end{proof}

\begin{lemma}
\label{lemma-base-change-representable}
Soit $\mathcal{C}$ une catégorie.
Soit $f : x \to y$ un morphisme représentable.
Soit $y' \to y$ un morphisme de $\mathcal{C}$.
Alors le morphisme $x' := x \times_y y' \to y'$ est lui aussi représentable.
\end{lemma}

\begin{proof}
Soit $z \to y'$ un morphisme. Le produit fibré
$x' \times_{y'} z$ doit représenter le foncteur
\begin{eqnarray*}
w & \mapsto & h_{x'}(w)\times_{h_{y'}(w)} h_z(w) \\
& = & (h_x(w) \times_{h_y(w)} h_{y'}(w)) \times_{h_{y'}(w)} h_z(w) \\
& = & h_x(w) \times_{h_y(w)} h_z(w)
\end{eqnarray*}
qui est représentable par hypothèse.
\end{proof}

\section{Exemples de produits fibrés}
\label{section-example-fibre-products}

\noindent
Dans cette section, nous énumérons des exemples de produits fibrés
et les décrivons.

\medskip\noindent
Comme premier exemple tout à fait trivial, observons
que la catégorie des ensembles admet les produits fibrés et que tout
morphisme y est donc représentable. En effet, si $f : X \to Y$
et $g : Z \to Y$ sont des applications d'ensembles, nous définissons
$X \times_Y Z$ comme le sous-ensemble de $X \times Z$ formé
des couples $(x, z)$ tels que $f(x) = g(z)$. Les morphismes
$p : X \times_Y Z \to X$ et $q : X \times_Y Z \to Z$ sont
les projections $(x, z) \mapsto x$ et $(x, z) \mapsto z$.
Enfin, si $\alpha : W \to X$ et $\beta : W \to Z$
sont des morphismes tels que $f \circ \alpha = g \circ \beta$,
alors l'application $W \to X \times Z$, $w\mapsto (\alpha(w), \beta(w))$,
est manifestement à valeurs dans $X \times_Y Z$, comme voulu.

\medskip\noindent
Dans de nombreuses catégories dont les objets sont des ensembles munis
de certains types de structures algébriques, le produit fibré des ensembles
sous-jacents fournit aussi le produit fibré dans la catégorie. Supposons,
par exemple, que $X$, $Y$ et $Z$ ci-dessus soient des groupes et que
$f$, $g$ soient des homomorphismes de groupes. Le produit fibré ensembliste
$X \times_Y Z$ hérite alors d'une structure de groupe, simplement en
définissant le produit de deux couples par la formule
$(x, z) \cdot (x', z') = (xx', zz')$. Voici les catégories
pour lesquelles un raisonnement analogue s'applique.
\begin{enumerate}
\item La catégorie $\textit{Groups}$ des groupes.
\item La catégorie $G\textit{-Ensembles}$ des ensembles
munis d'une action à gauche de $G$, pour un groupe fixé $G$.
\item La catégorie des anneaux.
\item La catégorie des $R$-modules, pour un anneau $R$ donné.
\end{enumerate}




\section{Produits fibrés et représentabilité}
\label{section-representable-map-presheaves}

\noindent
Dans cette section, nous déterminons les produits fibrés dans la
catégorie des foncteurs contravariants d'une catégorie
dans la catégorie des ensembles. Cette discussion sera remplacée plus loin
par celle des sites, des préfaisceaux et des faisceaux. La notion de
« morphisme représentable » entre de tels foncteurs présente un certain
intérêt.

\begin{lemma}
\label{lemma-fibre-product-presheaves}
Soit $\mathcal{C}$ une catégorie.
Soient $F, G, H : \mathcal{C}^{opp} \to \textit{Ensembles}$
des foncteurs. Soient $a : F \to G$ et $b : H \to G$ des
transformations de foncteurs. Alors le produit fibré
$F \times_{a, G, b} H$ dans la catégorie
$\textit{PSh}(\mathcal{C})$
existe et est donné par la formule
$$
(F \times_{a, G, b} H)(X) =
F(X) \times_{a_X, G(X), b_X} H(X)
$$
pour tout objet $X$ de $\mathcal{C}$.
\end{lemma}

\begin{proof}
Omis.
\end{proof}

\noindent
Comme cas particulier, supposons donnés un morphisme
$a : F \to G$, un objet $U \in \Ob(\mathcal{C})$
et un élément $\xi \in G(U)$. D'après le
lemme de Yoneda \ref{lemma-yoneda}, cela définit une transformation
$\xi : h_U \to G$. Dans ce cas, le produit fibré est décrit par la règle
$$
(h_U \times_{\xi, G, a} F)(X) =
\{ (f, \xi') \mid f : X \to U, \ \xi' \in F(X), \ G(f)(\xi) = a_X(\xi')\}
$$
Si $F$ et $G$ sont eux aussi représentables, ce foncteur représente le
produit fibré, lorsque celui-ci existe ; voir la
section \ref{section-fibre-products}.
L'analogie avec la définition \ref{definition-representable-morphism}
nous conduit à définir une notion de transformations représentables.

\begin{definition}
\label{definition-representable-map-presheaves}
Soit $\mathcal{C}$ une catégorie.
Soient $F, G : \mathcal{C}^{opp} \to \textit{Ensembles}$
des foncteurs. Nous disons qu'un morphisme $a : F \to G$ est
{\it représentable}, ou que {\it $F$ est relativement représentable
au-dessus de $G$}, si, pour tout $U \in \Ob(\mathcal{C})$
et tout $\xi \in G(U)$, le foncteur
$h_U \times_{\xi, G, a} F$ est représentable.
\end{definition}

\begin{lemma}
\label{lemma-representable-over-representable}
Soit $\mathcal{C}$ une catégorie.
Soit $a : F \to G$ un morphisme de foncteurs contravariants
de $\mathcal{C}$ dans $\textit{Ensembles}$. Si $a$ est représentable
et si $G$ est un foncteur représentable, alors $F$ est représentable.
\end{lemma}

\begin{proof}
Omis.
\end{proof}

\begin{lemma}
\label{lemma-representable-diagonal}
Soit $\mathcal{C}$ une catégorie.
Soit $F : \mathcal{C}^{opp} \to \textit{Ensembles}$ un foncteur.
Supposons que $\mathcal{C}$ admette les produits de deux objets et les
produits fibrés. Les assertions suivantes sont équivalentes :
\begin{enumerate}
\item la diagonale $\Delta : F \to F \times F$ est représentable ;
\item pour tout $U$ dans $\mathcal{C}$
et tout $\xi \in F(U)$, le morphisme $\xi : h_U \to F$ est représentable ;
\item pour toute paire $U, V$ dans $\mathcal{C}$
et tous $\xi \in F(U)$, $\xi' \in F(V)$, le produit fibré
$h_U \times_{\xi, F, \xi'} h_V$ est représentable.
\end{enumerate}
\end{lemma}

\begin{proof}
Nous continuerons d'utiliser le lemme de Yoneda pour identifier $F(U)$
aux transformations de foncteurs $h_U \to F$.

\medskip\noindent
Équivalence de (2) et (3). Soient $U, \xi, V, \xi'$ comme dans (3).
Les assertions (2) et (3) disent toutes deux exactement que
$h_U \times_{\xi, F, \xi'} h_V$ est représentable ; la seule différence
est que l'énoncé (3) est symétrique en $U$ et $V$, tandis que (2) ne l'est pas.

\medskip\noindent
Supposons la condition (1). Soient $U, \xi, V, \xi'$
comme dans (3). Remarquons que
$h_U \times h_V = h_{U \times V}$ est représentable.
Notons $\eta : h_{U \times V} \to F \times F$ le morphisme
correspondant au produit $\xi \times \xi' : h_U \times h_V \to F \times F$.
Le produit fibré $F \times_{\Delta, F \times F, \eta} h_{U \times V}$
est alors représentable par hypothèse. Cela signifie qu'il existe
$W \in \Ob(\mathcal{C})$, des morphismes
$W \to U$, $W \to V$ et $h_W \to F$ tels que
$$
\xymatrix{
h_W \ar[d] \ar[r] & h_U \times h_V \ar[d]^{\xi \times \xi'} \\
F \ar[r] & F \times F
}
$$
soit cartésien. La description explicite des produits fibrés donnée dans le
lemme \ref{lemma-fibre-product-presheaves} montre alors que
$h_W = h_U \times_{\xi, F, \xi'} h_V$, comme voulu.

\medskip\noindent
Supposons les conditions équivalentes (2) et (3). Soit $U$ un objet
de $\mathcal{C}$ et soit $(\xi, \xi') \in (F \times F)(U)$.
Par (3), le produit fibré $h_U \times_{\xi, F, \xi'} h_U$ est
représentable. Choisissons un objet $W$ et un isomorphisme
$h_W \to h_U \times_{\xi, F, \xi'} h_U$. Les deux projections
$\text{pr}_i : h_U \times_{\xi, F, \xi'} h_U \to h_U$
correspondent, par Yoneda, à des morphismes $p_i : W \to U$.
Considérons
$W' = W \times_{(p_1, p_2), U \times U} U$. Un argument formel montre
que $W'$ représente $F \times_{\Delta, F \times F} h_U$, car
$$
h_{W'} = h_W \times_{h_U \times h_U} h_U
= (h_U \times_{\xi, F, \xi'} h_U) \times_{h_U \times h_U} h_U
= F \times_{F \times F} h_U.
$$
Ainsi, $\Delta$ est représentable, ce qui achève la preuve.
\end{proof}








\section{Sommes amalgamées}
\label{section-pushouts}

\noindent
La notion duale de celle de produit fibré est celle de somme amalgamée.

\begin{definition}
\label{definition-pushouts}
Soient $x, y, z\in \Ob(\mathcal{C})$,
$f\in \Mor_\mathcal{C}(y, x)$
et $g\in \Mor_{\mathcal C}(y, z)$.
Une {\it somme amalgamée} de $f$ et $g$ est
un objet $x\amalg_y z\in \Ob(\mathcal{C})$
muni de morphismes
$p\in \Mor_{\mathcal C}(x, x\amalg_y z)$ et
$q\in\Mor_{\mathcal C}(z, x\amalg_y z)$ qui rendent le diagramme
$$
\xymatrix{
y \ar[r]_g \ar[d]_f & z \ar[d]^q \\
x \ar[r]^p & x\amalg_y z
}
$$
commutatif et qui satisfont la propriété universelle suivante :
pour tout $w\in \Ob(\mathcal{C})$ et tous morphismes
$\alpha \in \Mor_{\mathcal C}(x, w)$ et
$\beta \in \Mor_\mathcal{C}(z, w)$ tels que
$\alpha \circ f = \beta \circ g$, il existe un unique
$\gamma\in \Mor_{\mathcal C}(x\amalg_y z, w)$ qui rend le diagramme
$$
\xymatrix{
y \ar[r]_g \ar[d]_f & z \ar[d]^q \ar[rrdd]^\beta & & \\
x \ar[r]^p \ar[rrrd]^\alpha & x \amalg_y z \ar@{-->}[rrd]^\gamma & & \\
& & & w
}
$$
commutatif.
\end{definition}

\noindent
On peut démontrer directement, et sans difficulté, l'unicité du triplet
$(x\amalg_y z, p, q)$ à isomorphisme unique près (lorsqu'il existe).
On peut aussi considérer la somme amalgamée comme le
produit fibré dans la catégorie opposée ; son unicité découle alors
immédiatement de la discussion de la section \ref{section-fibre-products}.

\begin{definition}
\label{definition-cocartesian}
Nous disons qu'un diagramme commutatif
$$
\xymatrix{
y \ar[r] \ar[d] & z \ar[d] \\
x \ar[r] & w
}
$$
dans une catégorie est {\it cocartésien} si $w$ et les morphismes
$x \to w$ et $z \to w$ forment une somme amalgamée des morphismes
$y \to x$ et $y \to z$.
\end{definition}


\section{Égalisateurs}
\label{section-equalizers}

\begin{definition}
\label{definition-equalizers}
Supposons que $X$ et $Y$ soient des objets d'une catégorie $\mathcal{C}$
et que $a, b : X \to Y$ soient des morphismes. Un morphisme
$e : Z \to X$ est appelé un {\it égalisateur} de la paire $(a, b)$ si
$a \circ e = b \circ e$ et si $(Z, e)$ satisfait la propriété universelle
suivante : pour tout morphisme $t : W \to X$
dans $\mathcal{C}$ tel que $a \circ t = b \circ t$, il existe
un unique morphisme $s : W \to Z$ tel que $t = e \circ s$.
\end{definition}

\noindent
Comme pour les produits fibrés ci-dessus, lorsqu'ils existent, les
égalisateurs sont uniques à isomorphisme unique près. Cette définition
se généralise immédiatement au cas où l'on a plus de $2$ morphismes.

\section{Coégalisateurs}
\label{section-coequalizers}

\begin{definition}
\label{definition-coequalizers}
Supposons que $X$ et $Y$ soient des objets d'une catégorie $\mathcal{C}$
et que $a, b : X \to Y$ soient des morphismes. Un morphisme
$c : Y \to Z$ est appelé un {\it coégalisateur} de la paire $(a, b)$ si
$c \circ a = c \circ b$ et si $(Z, c)$ satisfait la propriété universelle
suivante : pour tout morphisme $t : Y \to W$
dans $\mathcal{C}$ tel que $t \circ a = t \circ b$, il existe
un unique morphisme $s : Z \to W$ tel que $t = s \circ c$.
\end{definition}

\noindent
Comme pour les sommes amalgamées ci-dessus, lorsqu'ils existent, les
coégalisateurs sont uniques à isomorphisme unique près ; cela résulte
de l'unicité des égalisateurs en considérant la catégorie opposée.
Cette définition se généralise immédiatement au cas où l'on a plus de
$2$ morphismes.

\section{Objets initiaux et finaux}
\label{section-initial-final}

\begin{definition}
\label{definition-initial-final}
Soit $\mathcal{C}$ une catégorie.
\begin{enumerate}
\item Un objet $x$ de la catégorie $\mathcal{C}$ est dit
{\it initial} si, pour tout objet $y$ de $\mathcal{C}$,
il existe exactement un morphisme $x \to y$.
\item Un objet $x$ de la catégorie $\mathcal{C}$ est dit
{\it final} si, pour tout objet $y$ de $\mathcal{C}$,
il existe exactement un morphisme $y \to x$.
\end{enumerate}
\end{definition}

\noindent
Dans la catégorie des ensembles, l'ensemble vide $\emptyset$ est un
objet initial, et c'est même le seul objet initial.
Tout {\it singleton}, c'est-à-dire tout ensemble à un élément,
est par ailleurs un objet final (il n'est donc pas unique).




\section{Monomorphismes et épimorphismes}
\label{section-mono-epi}

\begin{definition}
\label{definition-mono-epi}
Soit $\mathcal{C}$ une catégorie et soit $f : X \to Y$
un morphisme de $\mathcal{C}$.
\begin{enumerate}
\item Nous disons que $f$ est un {\it monomorphisme} si, pour tout objet
$W$ et toute paire de morphismes $a, b : W \to X$ tels que
$f \circ a = f \circ b$, on a $a = b$.
\item Nous disons que $f$ est un {\it épimorphisme} si, pour tout objet
$W$ et toute paire de morphismes $a, b : Y \to W$ tels que
$a \circ f = b \circ f$, on a $a = b$.
\end{enumerate}
\end{definition}

\begin{example}
\label{example-mono-epi-sets}
Dans la catégorie des ensembles, les monomorphismes sont les applications
injectives et les épimorphismes sont les applications surjectives.
\end{example}

\begin{lemma}
\label{lemma-characterize-mono-epi}
Soit $\mathcal{C}$ une catégorie et soit $f : X \to Y$
un morphisme de $\mathcal{C}$. Alors
\begin{enumerate}
\item $f$ est un monomorphisme si et seulement si la diagonale
$X \to X \times_Y X$ est un isomorphisme, et
\item $f$ est un épimorphisme si et seulement si la codiagonale
$Y \amalg_X Y \to Y$ est un isomorphisme.
\end{enumerate}
\end{lemma}

\begin{proof}
Supposons que $f$ soit un monomorphisme. Soit $W$ un objet de $\mathcal C$
et soient $\alpha, \beta \in \Mor_\mathcal C(W,X)$ tels que
$f\circ \alpha = f\circ
\beta$. On a donc $\alpha = \beta$,
puisque $f$ est un monomorphisme. De plus, on a le diagramme commutatif
$$
\xymatrix{
X \ar[r]^{\text{id}_X} \ar[d]_{\text{id}_X} & X \ar[d]^{f} \\
X \ar[r]^{f} & Y
}
$$
qui vérifie la propriété universelle avec $\gamma := \alpha = \beta$.
Ainsi, $X$ est bien le produit fibré $X\times_Y X$.

\medskip\noindent
Supposons que la diagonale $X \to X \times_Y X$ soit un isomorphisme.
Le diagramme
$$
\xymatrix{
X \ar[r]^{\text{id}_X} \ar[d]_{\text{id}_X} & X \ar[d]^{f} \\
X \ar[r]^{f} & Y }
$$
est commutatif et, si $W \in \Ob(\mathcal C)$ et
$\alpha, \beta : W \to X$ sont tels que
$f \circ \alpha = f \circ \beta$, il existe un unique
$\gamma$ vérifiant
$$
\gamma = \text{id}_X\circ\gamma = \alpha = \beta
$$
ce qui montre que $\alpha = \beta$.

\medskip\noindent
La preuve du second point est exactement la même, en utilisant la codiagonale
$Y\amalg_X Y \to Y$.
\end{proof}





\section{Limites et colimites}
\label{section-limits}

\noindent
Soit $\mathcal{C}$ une catégorie. Un {\it diagramme} dans $\mathcal{C}$ est
simplement un foncteur $M : \mathcal{I} \to \mathcal{C}$. Nous disons que
$\mathcal{I}$ est la {\it catégorie d'indices}, ou que $M$ est un
diagramme indexé par $\mathcal{I}$. Nous noterons $M_i$ l'image de l'objet
$i$ de $\mathcal{I}$. Ainsi, pour un morphisme $\phi : i \to i'$
dans $\mathcal{I}$, on a $M(\phi) : M_i \to M_{i'}$.

\begin{definition}
\label{definition-limit}
Une {\it limite} du diagramme indexé par $\mathcal{I}$ que définit $M$ dans la
catégorie
$\mathcal{C}$ est la donnée d'un objet $\lim_\mathcal{I} M$ de $\mathcal{C}$
et de morphismes $p_i : \lim_\mathcal{I} M \to M_i$ tels que
\begin{enumerate}
\item pour tout morphisme $\phi : i \to i'$
dans $\mathcal{I}$, on ait $p_{i'} = M(\phi) \circ p_i$, et
\item pour tout objet $W$ de $\mathcal{C}$ et toute famille de
morphismes $q_i : W \to M_i$ (indexée par $i \in \Ob(\mathcal{I})$)
telle que, pour tout $\phi : i \to i'$
dans $\mathcal{I}$, on ait $q_{i'} = M(\phi) \circ q_i$, il existe
un unique morphisme $q : W \to \lim_\mathcal{I} M$ tel que
$q_i = p_i \circ q$ pour tout objet $i$ de $\mathcal{I}$.
\end{enumerate}
\end{definition}

\noindent
Lorsqu'elles existent, les limites
$(\lim_\mathcal{I} M, (p_i)_{i\in \Ob(\mathcal{I})})$
sont uniques à isomorphisme unique près, par la condition d'unicité
de la définition. Les produits de deux objets, les produits fibrés et les
égalisateurs sont des exemples de limites. La limite du diagramme vide
est un objet final de $\mathcal{C}$.
Dans la catégorie des ensembles, toutes les limites existent.
La notion duale est celle de colimite.

\begin{definition}
\label{definition-colimit}
Une {\it colimite} du diagramme indexé par $\mathcal{I}$ que définit $M$ dans la
catégorie
$\mathcal{C}$ est la donnée d'un objet $\colim_\mathcal{I} M$ de $\mathcal{C}$
et de morphismes $s_i : M_i \to \colim_\mathcal{I} M$ tels que
\begin{enumerate}
\item pour tout morphisme $\phi : i \to i'$
dans $\mathcal{I}$, on ait $s_i = s_{i'} \circ M(\phi)$, et
\item pour tout objet $W$ de $\mathcal{C}$ et toute famille de
morphismes $t_i : M_i \to W$ (indexée par $i \in \Ob(\mathcal{I})$)
telle que, pour tout $\phi : i \to i'$
dans $\mathcal{I}$, on ait $t_i = t_{i'} \circ M(\phi)$, il existe
un unique morphisme $t : \colim_\mathcal{I} M \to W$ tel que
$t_i = t \circ s_i$ pour tout objet $i$ de $\mathcal{I}$.
\end{enumerate}
\end{definition}

\noindent
Lorsqu'elles existent, les colimites
$(\colim_\mathcal{I} M, (s_i)_{i\in \Ob(\mathcal{I})})$
sont uniques à isomorphisme unique près, par la condition d'unicité
de la définition. Les coproduits de deux objets, les sommes amalgamées et
les coégalisateurs sont des exemples de colimites. La colimite du diagramme
vide est un objet initial de $\mathcal{C}$. Dans la catégorie des ensembles,
toutes les colimites existent.

\begin{remark}
\label{remark-diagram-small}
La catégorie d'indices d'une (co)limite ne pourra jamais avoir une classe
propre d'objets. Dans ce projet, cela signifie qu'elle ne peut être l'une
des catégories énumérées dans la remarque \ref{remark-big-categories}.
\end{remark}

\begin{remark}
\label{remark-limit-colim}
Nous écrivons souvent $\lim_i M_i$, $\colim_i M_i$,
$\lim_{i\in \mathcal{I}} M_i$ ou $\colim_{i\in \mathcal{I}} M_i$
à la place des versions indexées par $\mathcal{I}$.
Avec cette notation et la description des limites et colimites d'ensembles
donnée dans la section \ref{section-limit-sets} ci-dessous,
nous pouvons énoncer ce qui suit.
Soit $M : \mathcal{I} \to \mathcal{C}$ un diagramme.
\begin{enumerate}
\item S'il existe, l'objet $\lim_i M_i$ vérifie la propriété suivante :
$$
\Mor_\mathcal{C}(W, \lim_i M_i)
=
\lim_i \Mor_\mathcal{C}(W, M_i)
$$
où la limite du membre de droite est prise dans la catégorie des ensembles.
\item S'il existe, l'objet $\colim_i M_i$ vérifie la propriété suivante :
$$
\Mor_\mathcal{C}(\colim_i M_i, W)
=
\lim_{i \in \mathcal{I}^{opp}} \Mor_\mathcal{C}(M_i, W)
$$
où la limite du membre de droite, prise sur la catégorie opposée, est calculée
dans la catégorie des ensembles.
\end{enumerate}
Par le lemme de Yoneda (et son dual), cette formule détermine complètement
la limite, respectivement la colimite.
\end{remark}

\begin{remark}
\label{remark-cones-and-cocones}
Soit $M : \mathcal{I} \to \mathcal{C}$ un diagramme. Dans ce contexte,
un {\it cône} sur $M$ est la donnée d'un objet $W$ et d'une famille de
morphismes $q_i : W \to M_i$, $i \in \Ob(\mathcal{I})$, tels que, pour tout
morphisme $\phi : i \to i'$ de $\mathcal{I}$, le diagramme
$$
\xymatrix{
& W \ar[dl]_{q_i} \ar[dr]^{q_{i'}} \\
M_i \ar[rr]^{M(\phi)} & & M_{i'}
}
$$
soit commutatif. Les cônes forment une catégorie dont la notion de morphisme
est évidente. Manifestement, la limite de $M$, si elle existe, est un objet
final de la catégorie des cônes. Dualement, un {\it cocône} sur $M$ est la
donnée d'un objet $W$ et d'une famille de morphismes $t_i : M_i \to W$ tels
que, pour tout morphisme $\phi : i \to i'$ dans $\mathcal{I}$, le diagramme
$$
\xymatrix{
M_i \ar[rr]^{M(\phi)} \ar[dr]_{t_i} & & M_{i'} \ar[dl]^{t_{i'}} \\
& W
}
$$
soit commutatif. Les cocônes forment une catégorie dont la notion de
morphisme est évidente. Comme ci-dessus, la colimite de $M$ existe si et
seulement si la catégorie des cocônes possède un objet initial.
\end{remark}

\noindent
Comme application des notions de limite et de colimite,
nous définissons les produits et les coproduits.

\begin{definition}
\label{definition-product}
Supposons que $I$ soit un ensemble et que, pour tout $i \in I$, soit donné
un objet $M_i$ de la catégorie $\mathcal{C}$. Un {\it produit}
$\prod_{i\in I} M_i$ est, par définition, $\lim_\mathcal{I} M$
(s'il existe), où $\mathcal{I}$ est la catégorie dont les objets
sont les éléments de $I$ et dont les seuls morphismes sont les identités.
\end{definition}

\noindent
Un cas particulier important est $I = \emptyset$ ; le
produit est alors un objet final de la catégorie.
Les morphismes $p_i : \prod M_i \to M_i$ sont appelés les
{\it morphismes de projection}.

\begin{definition}
\label{definition-coproduct}
Supposons que $I$ soit un ensemble et que, pour tout $i \in I$, soit donné
un objet $M_i$ de la catégorie $\mathcal{C}$. Un {\it coproduit}
$\coprod_{i\in I} M_i$ est, par définition, $\colim_\mathcal{I} M$
(s'il existe), où $\mathcal{I}$ est la catégorie dont les objets
sont les éléments de $I$ et dont les seuls morphismes sont les identités.
\end{definition}

\noindent
Un cas particulier important est $I = \emptyset$ ; le
coproduit est alors un objet initial de la catégorie.
Remarquons que le coproduit est muni de morphismes
$M_i \to \coprod M_i$. Ceux-ci sont parfois appelés les
{\it coprojections}.

\begin{lemma}
\label{lemma-functorial-colimit}
Supposons que $M : \mathcal{I} \to \mathcal{C}$
et $N : \mathcal{J} \to \mathcal{C}$ soient des diagrammes
dont les colimites existent. Supposons que
$H : \mathcal{I} \to \mathcal{J}$ soit
un foncteur et que $t : M \to N \circ H$
soit une transformation de foncteurs.
Il existe alors un unique morphisme
$$
\theta :
\colim_\mathcal{I} M
\longrightarrow
\colim_\mathcal{J} N
$$
tel que tous les diagrammes
$$
\xymatrix{
M_i \ar[d]_{t_i} \ar[r]
&
\colim_\mathcal{I} M \ar[d]^{\theta}
\\
N_{H(i)} \ar[r]
&
\colim_\mathcal{J} N
}
$$
soient commutatifs.
\end{lemma}

\begin{proof}
Omis.
\end{proof}

\begin{lemma}
\label{lemma-functorial-limit}
Supposons que $M : \mathcal{I} \to \mathcal{C}$
et $N : \mathcal{J} \to \mathcal{C}$ soient des diagrammes
dont les limites existent. Supposons que
$H : \mathcal{I} \to \mathcal{J}$ soit un foncteur et que
$t : N \circ H \to M$ soit une transformation de foncteurs.
Il existe alors un unique morphisme
$$
\theta :
\lim_\mathcal{J} N
\longrightarrow
\lim_\mathcal{I} M
$$
tel que tous les diagrammes
$$
\xymatrix{
\lim_\mathcal{J} N \ar[d]^{\theta} \ar[r]
&
N_{H(i)} \ar[d]_{t_i}
\\
\lim_\mathcal{I} M \ar[r]
&
M_i
}
$$
soient commutatifs.
\end{lemma}

\begin{proof}
Omis.
\end{proof}


\begin{lemma}
\label{lemma-colimits-commute}
Soient $\mathcal{I}$ et $\mathcal{J}$ des catégories d'indices. Soit
$M : \mathcal{I} \times \mathcal{J} \to \mathcal{C}$ un foncteur.
Supposons que $M_{i, \infty} = \colim_j M_{i,j}$
existe pour tout $i$. Alors le foncteur ainsi obtenu
$M_{-, \infty} : \mathcal{I} \to \mathcal{C}$
admet une colimite si et seulement si $M$ en admet une, et les deux
colimites coïncident. En particulier,
$$
\colim_i \colim_j M_{i, j}
=
\colim_{i, j} M_{i, j}
=
\colim_j \colim_i M_{i, j}
$$
pourvu que toutes les colimites indiquées existent. Il en va de même
pour les limites.
\end{lemma}

\begin{proof}
Omis.
\end{proof}

\begin{lemma}
\label{lemma-limits-products-equalizers}
Soit $M : \mathcal{I} \to \mathcal{C}$ un diagramme.
Écrivons $I = \Ob(\mathcal{I})$ et $A = \text{Flèches}(\mathcal{I})$.
Notons $s, t : A \to I$ les applications source et but.
Supposons que $\prod_{i \in I} M_i$ et $\prod_{a \in A} M_{t(a)}$
existent. Supposons que l'égalisateur de
$$
\xymatrix{
\prod_{i \in I} M_i
\ar@<1ex>[r]^\phi \ar@<-1ex>[r]_\psi
&
\prod_{a \in A} M_{t(a)}
}
$$
existe, les morphismes étant déterminés par leurs composantes
comme suit : $p_a \circ \psi = M(a) \circ p_{s(a)}$
et $p_a \circ \phi = p_{t(a)}$. Cet égalisateur est alors la
limite du diagramme.
\end{lemma}

\begin{proof}
Omis.
\end{proof}


\begin{lemma}
\label{lemma-colimits-coproducts-coequalizers}
\begin{slogan}
Si tous les coproduits et tous les coégalisateurs existent, toutes les
colimites existent.
\end{slogan}
Soit $M : \mathcal{I} \to \mathcal{C}$ un diagramme.
Écrivons $I = \Ob(\mathcal{I})$ et $A = \text{Flèches}(\mathcal{I})$.
Notons $s, t : A \to I$ les applications source et but.
Supposons que $\coprod_{i \in I} M_i$ et
$\coprod_{a \in A} M_{s(a)}$ existent. Supposons que le coégalisateur de
$$
\xymatrix{
\coprod_{a \in A} M_{s(a)}
\ar@<1ex>[r]^\phi \ar@<-1ex>[r]_\psi
&
\coprod_{i \in I} M_i
}
$$
existe, les morphismes étant déterminés par leurs composantes
comme suit : la composante $M_{s(a)}$ s'envoie par $\psi$
dans la composante $M_{t(a)}$ au moyen du morphisme $M(a)$.
La composante $M_{s(a)}$ s'envoie par $\phi$ dans la composante
$M_{s(a)}$ au moyen du morphisme identité. Ce coégalisateur est alors la
colimite du diagramme.
\end{lemma}

\begin{proof}
Omis.
\end{proof}








\section{Limites et colimites dans la catégorie des ensembles}
\label{section-limit-sets}

\noindent
Non seulement les limites et les colimites existent dans $\textit{Ensembles}$,
mais elles sont également faciles à décrire. Soit en effet
$M : \mathcal{I}
\to \textit{Ensembles}$, $i \mapsto M_i$, un diagramme
d'ensembles. Notons $I = \Ob(\mathcal{I})$.
La limite s'écrit
$$
\lim_\mathcal{I} M
=
\{
(m_i)_{i\in I} \in \prod\nolimits_{i\in I} M_i
\mid
\forall \phi : i \to i' \text{ dans }\mathcal{I},
M(\phi)(m_i) = m_{i'}
\}.
$$
Nous considérons donc un élément de la limite comme un système compatible
d'éléments de tous les ensembles $M_i$.

\medskip\noindent
D'autre part, la colimite est
$$
\colim_\mathcal{I} M
=
(\coprod\nolimits_{i\in I} M_i)/\sim
$$
où la relation d'équivalence $\sim$ est engendrée en imposant
$m_i \sim m_{i'}$ si $m_i \in M_i$,
$m_{i'} \in M_{i'}$ et $M(\phi)(m_i) = m_{i'}$ pour un certain
$\phi : i \to i'$. Autrement dit, $m_i \in M_i$
et $m_{i'} \in M_{i'}$ sont équivalents s'il existe une
chaîne de morphismes dans $\mathcal{I}$
$$
\xymatrix{
&
i_1 \ar[ld] \ar[rd] & &
i_3 \ar[ld] & &
i_{2n-1} \ar[rd] & \\
i = i_0 & &
i_2 & &
\ldots & &
i_{2n} = i'
}
$$
et des éléments $m_{i_j} \in M_{i_j}$ dont les images coïncident par
les applications $M_{i_{2k-1}} \to M_{i_{2k-2}}$ et
$M_{i_{2k-1}}
\to M_{i_{2k}}$ induites par les flèches de
$\mathcal{I}$ ci-dessus.

\medskip\noindent
Ce type d'objet n'est pas très commode à manipuler.
Mais si le diagramme est filtrant, sa description est beaucoup plus simple.
Nous l'expliquerons dans la section \ref{section-directed-colimits}.



\section{Limites connexes}
\label{section-connected-limits}

\noindent
Une (co)limite est dite connexe si sa catégorie d'indices est connexe.

\begin{definition}
\label{definition-category-connected}
Nous disons qu'une catégorie $\mathcal{I}$ est {\it connexe}
si la relation d'équivalence engendrée par
$x \sim y \Leftrightarrow \Mor_\mathcal{I}(x, y) \not = \emptyset$
possède exactement une classe d'équivalence.
\end{definition}

\noindent
Nous suivons ici la convention de
Topologie, définition \ref{topology-definition-connected-components},
selon laquelle les espaces connexes sont non vides.
Le résultat suivant caractérise, en un sens quelque peu vague,
les limites connexes.

\begin{lemma}
\label{lemma-connected-limit-over-X}
Soit $\mathcal{C}$ une catégorie.
Soit $X$ un objet de $\mathcal{C}$.
Soit $M : \mathcal{I} \to \mathcal{C}/X$ un diagramme
dans la catégorie des objets au-dessus de $X$.
Si la catégorie d'indices $\mathcal{I}$ est connexe
et si la limite de $M$ existe dans $\mathcal{C}/X$,
alors la limite du composé
$\mathcal{I} \to \mathcal{C}/X \to \mathcal{C}$
existe et lui est égale.
\end{lemma}

\begin{proof}
Soit $L \to X$ un objet représentant la limite dans $\mathcal{C}/X$.
Considérons le foncteur
$$
W \longmapsto \lim_i \Mor_\mathcal{C}(W, M_i).
$$
Soit $(\varphi_i)$ un élément de l'ensemble de droite.
Comme chaque $M_i$ est muni d'un morphisme $s_i : M_i \to X$, nous
obtenons des morphismes $f_i = s_i \circ \varphi_i : W \to X$.
La connexité de $\mathcal{I}$ montre que tous les $f_i$ sont égaux.
Puisque $\mathcal{I}$ est non vide, il existe au moins un $f_i$.
Leur valeur commune $W \to X$ confère donc à $W$ une structure
d'objet de $\mathcal{C}/X$, et $(\varphi_i)$ définit un
élément de $\lim_i \Mor_{\mathcal{C}/X}(W, M_i)$.
Nous obtenons ainsi un unique morphisme $\phi : W \to L$ tel que
$\varphi_i$ soit le composé de $\phi$ et de $L \to M_i$, comme voulu.
\end{proof}

\begin{lemma}
\label{lemma-connected-colimit-under-X}
Soit $\mathcal{C}$ une catégorie.
Soit $X$ un objet de $\mathcal{C}$.
Soit $M : \mathcal{I} \to X/\mathcal{C}$ un diagramme
dans la catégorie des objets au-dessous de $X$.
Si la catégorie d'indices $\mathcal{I}$ est connexe
et si la colimite de $M$ existe dans $X/\mathcal{C}$,
alors la colimite du composé
$\mathcal{I} \to X/\mathcal{C} \to \mathcal{C}$
existe et lui est égale.
\end{lemma}

\begin{proof}
Omis. Indication : ce lemme est dual du
lemme \ref{lemma-connected-limit-over-X}.
\end{proof}




\section{Catégories cofinales et initiales}
\label{section-cofinal}

\noindent
Dans la littérature, le mot « final » est parfois employé à la place de
« cofinal » dans la définition suivante.

\begin{definition}
\label{definition-cofinal}
Soit $H : \mathcal{I} \to \mathcal{J}$ un foncteur entre catégories.
Nous disons que {\it $\mathcal{I}$ est cofinale dans $\mathcal{J}$},
ou que $H$ est {\it cofinal}, si
\begin{enumerate}
\item pour tout $y \in \Ob(\mathcal{J})$, il existe
$x \in \Ob(\mathcal{I})$ et un morphisme $y \to H(x)$, et
\item étant donnés $y \in \Ob(\mathcal{J})$,
$x, x' \in \Ob(\mathcal{I})$ et des morphismes
$y \to H(x)$ et $y \to H(x')$, il existe une suite de morphismes
$$
x = x_0 \leftarrow x_1 \rightarrow x_2 \leftarrow x_3 \rightarrow \ldots
\rightarrow x_{2n} = x'
$$
dans $\mathcal{I}$ et des morphismes $y \to H(x_i)$ dans $\mathcal{J}$,
les morphismes $y \to H(x_0)$ et $y \to H(x_{2n})$ étant ceux donnés,
tels que les diagrammes
$$
\xymatrix{
& y \ar[ld] \ar[d] \ar[rd] \\
H(x_{2k}) & H(x_{2k + 1}) \ar[l] \ar[r] & H(x_{2k + 2})
}
$$
soient commutatifs pour $k = 0, \ldots, n - 1$.
\end{enumerate}
Autrement dit, fixons un objet $y$ de $\mathcal{J}$ et considérons
l'ensemble $S$ des couples $(x, b)$, où $x$ est un objet de $\mathcal{I}$
et $b : y \to H(x)$ un morphisme. Considérons sur $S$ la relation
d'équivalence engendrée par $(x, b) \sim (x', b')$ s'il existe un
morphisme $a : x \to x'$ tel que $b' = H(a) \circ b$. Alors $S$ doit
être formé d'une seule classe d'équivalence.
\end{definition}

\begin{lemma}
\label{lemma-cofinal}
Soit $H : \mathcal{I} \to \mathcal{J}$ un foncteur entre catégories.
Supposons $\mathcal{I}$ cofinale dans $\mathcal{J}$. Alors, pour tout
diagramme $M : \mathcal{J} \to \mathcal{C}$, on a un isomorphisme canonique
$$
\colim_\mathcal{I} M \circ H
=
\colim_\mathcal{J} M
$$
dès que l'un des deux membres existe.
\end{lemma}

\begin{proof}
Omis.
\end{proof}

\begin{definition}
\label{definition-initial}
Soit $H : \mathcal{I} \to \mathcal{J}$ un foncteur entre catégories.
Nous disons que {\it $\mathcal{I}$ est initiale dans $\mathcal{J}$},
ou que $H$ est {\it initial}, si
\begin{enumerate}
\item pour tout $y \in \Ob(\mathcal{J})$, il existe
$x \in \Ob(\mathcal{I})$ et un morphisme $H(x) \to y$ ;
\item pour tout $y \in \Ob(\mathcal{J})$, tous
$x , x' \in \Ob(\mathcal{I})$ et tous morphismes
$H(x) \to y$, $H(x') \to y$ dans $\mathcal{J}$,
il existe une suite de morphismes
$$
x = x_0 \leftarrow x_1 \rightarrow x_2 \leftarrow x_3 \rightarrow \ldots
\rightarrow x_{2n} = x'
$$
dans $\mathcal{I}$ et des morphismes $H(x_i) \to y$ dans $\mathcal{J}$
tels que les diagrammes
$$
\xymatrix{
H(x_{2k}) \ar[rd] &
H(x_{2k + 1}) \ar[l] \ar[r] \ar[d] &
H(x_{2k + 2}) \ar[ld] \\
& y
}
$$
soient commutatifs pour $k = 0, \ldots, n - 1$.
\end{enumerate}
\end{definition}

\noindent
C'est simplement la notion duale de celle de foncteur « cofinal ».

\begin{lemma}
\label{lemma-initial}
Soit $H : \mathcal{I} \to \mathcal{J}$ un foncteur de catégories.
Supposons $\mathcal{I}$ initiale dans $\mathcal{J}$.
Alors, pour tout diagramme $M : \mathcal{J} \to \mathcal{C}$,
on a un isomorphisme canonique
$$
\lim_\mathcal{I} M \circ H = \lim_\mathcal{J} M
$$
dès que l'un des deux membres existe.
\end{lemma}

\begin{proof}
Omis.
\end{proof}

\begin{lemma}
\label{lemma-colimit-constant-connected-fibers}
Soit $F : \mathcal{I} \to \mathcal{I}'$ un foncteur.
Supposons que
\begin{enumerate}
\item les catégories fibres (voir la
définition \ref{definition-fibre-category})
de $\mathcal{I}$ au-dessus de $\mathcal{I}'$ soient toutes connexes, et que
\item pour tout morphisme $\alpha' : x' \to y'$ dans $\mathcal{I}'$, il
existe un morphisme $\alpha : x \to y$ dans $\mathcal{I}$ tel que
$F(\alpha) = \alpha'$.
\end{enumerate}
Alors, pour tout diagramme $M : \mathcal{I}' \to \mathcal{C}$,
la colimite $\colim_\mathcal{I} M \circ F$ existe si et seulement si
$\colim_{\mathcal{I}'} M$ existe ; lorsqu'elles existent, ces colimites
coïncident.
\end{lemma}

\begin{proof}
On peut le démontrer en établissant que $\mathcal{I}$ est cofinale dans
$\mathcal{I}'$, puis en appliquant le lemme \ref{lemma-cofinal}.
Mais on peut aussi donner la preuve directe suivante.
Il suffit de montrer que, pour tout objet $T$ de $\mathcal{C}$, on a
$$
\lim_{\mathcal{I}^{opp}} \Mor_\mathcal{C}(M_{F(i)}, T)
=
\lim_{(\mathcal{I}')^{opp}} \Mor_\mathcal{C}(M_{i'}, T)
$$
Si $(g_{i'})_{i' \in \Ob(\mathcal{I}')}$ est un élément du
membre de droite, alors, en posant $f_i = g_{F(i)}$, on obtient un
élément $(f_i)_{i \in \Ob(\mathcal{I})}$ du membre de gauche.
Réciproquement, soit $(f_i)_{i \in \Ob(\mathcal{I})}$ un élément du
membre de gauche. Sur chaque catégorie fibre (connexe)
$\mathcal{I}_{i'}$, le foncteur $M \circ F$
est constant de valeur $M_{i'}$. Les morphismes
$f_i$, pour $i \in \Ob(\mathcal{I})$ tel que $F(i) = i'$,
sont donc tous égaux et déterminent un morphisme bien défini
$g_{i'} : M_{i'} \to T$. Par l'hypothèse (2), la famille
$(g_{i'})_{i' \in \Ob(\mathcal{I}')}$ définit un élément
du membre de droite.
\end{proof}

\begin{lemma}
\label{lemma-product-with-connected}
Soient $\mathcal{I}$ et $\mathcal{J}$ des catégories et notons
$p : \mathcal{I} \times \mathcal{J} \to \mathcal{J}$ la projection.
Si $\mathcal{I}$ est connexe, alors, pour un diagramme
$M : \mathcal{J} \to \mathcal{C}$, la colimite
$\colim_\mathcal{J} M$ existe si et seulement si
$\colim_{\mathcal{I} \times \mathcal{J}} M \circ p$ existe et,
dans ce cas, ces colimites sont égales.
\end{lemma}

\begin{proof}
C'est un cas particulier du
lemme \ref{lemma-colimit-constant-connected-fibers}.
\end{proof}





\section{Limites et colimites finies}
\label{section-finite-limits}

\noindent
Une (co)limite est dite {\it finie} si sa catégorie d'indices est finie,
c'est-à-dire si celle-ci possède un nombre fini d'objets et un nombre fini
de morphismes. Une (co)limite est dite {\it non vide} si sa catégorie
d'indices est non vide. Une (co)limite est dite {\it connexe} si sa catégorie
d'indices est connexe ; voir la
définition \ref{definition-category-connected}.
Il existe en fait « suffisamment » de catégories d'indices finies.

\begin{lemma}
\label{lemma-finite-diagram-category}
Soit $\mathcal{I}$ une catégorie telle que
\begin{enumerate}
\item $\Ob(\mathcal{I})$ soit fini, et
\item il existe un nombre fini de morphismes
$f_1, \ldots, f_m \in \text{Flèches}(\mathcal{I})$ tels
que tout morphisme de $\mathcal{I}$ soit un composé
$f_{j_1} \circ f_{j_2} \circ \ldots \circ f_{j_k}$.
\end{enumerate}
Il existe alors un foncteur $F : \mathcal{J} \to \mathcal{I}$ tel que
\begin{enumerate}
\item[(a)] $\mathcal{J}$ soit une catégorie finie, et
\item[(b)] pour tout diagramme $M : \mathcal{I} \to \mathcal{C}$, la
(co)limite de $M$ sur $\mathcal{I}$ existe si et seulement si
la (co)limite de $M \circ F$ sur $\mathcal{J}$ existe ; dans ce cas,
les (co)limites sont canoniquement isomorphes.
\end{enumerate}
De plus, $\mathcal{J}$ est connexe (resp.\ non vide) si et seulement si
$\mathcal{I}$ l'est.
\end{lemma}

\begin{proof}
Écrivons $\Ob(\mathcal{I}) = \{x_1, \ldots, x_n\}$.
Notons $s, t : \{1, \ldots, m\} \to \{1, \ldots, n\}$ les applications
telles que $f_j : x_{s(j)} \to x_{t(j)}$.
Posons $\Ob(\mathcal{J}) = \{y_1, \ldots, y_n, z_1, \ldots, z_n\}$.
Outre les identités, introduisons des morphismes
$g_j : y_{s(j)} \to z_{t(j)}$, $j = 1, \ldots, m$, et des morphismes
$h_i : y_i \to z_i$, $i = 1, \ldots, n$. Comme tous les morphismes de
$\mathcal{J}$ autres que les identités vont d'un $y$ vers un $z$, il n'y a
aucune composition à définir ni aucune associativité à vérifier.
Posons $F(y_i) = F(z_i) = x_i$. Posons
$F(g_j) = f_j$ et $F(h_i) = \text{id}_{x_i}$.
Il est clair que $F$ est un foncteur.
Il est clair que $\mathcal{J}$ est finie.
Il est clair que $\mathcal{J}$ est connexe, resp.\ non vide,
si et seulement si $\mathcal{I}$ l'est.

\medskip\noindent
Soit $M : \mathcal{I} \to \mathcal{C}$ un diagramme.
Considérons un objet $W$ de $\mathcal{C}$ et des morphismes
$q_i : W \to M(x_i)$ comme dans la
définition \ref{definition-limit}.
En prenant $q_i : W \to M(F(y_i)) = M(F(z_i)) = M(x_i)$, nous obtenons
une famille de morphismes comme dans la
définition \ref{definition-limit}
pour le diagramme $M \circ F$.
Réciproquement, supposons donnés des morphismes
$q_{y_i} : W \to M(F(y_i))$ et $q_{z_i} : W \to M(F(z_i))$
comme dans la définition \ref{definition-limit}
pour le diagramme $M \circ F$. Puisque
$$
M(F(h_i)) = \text{id} : M(F(y_i)) = M(x_i) \longrightarrow M(x_i) = M(F(z_i))
$$
on en déduit que $q_{y_i} = q_{z_i}$ pour tout $i$. Posons $q_i$ égal à
cette valeur commune. La compatibilité de
$q_{s(j)} = q_{y_{s(j)}}$ et de $q_{t(j)} = q_{z_{t(j)}}$ avec le morphisme
$M(f_j)$ garantit que la famille $q_i$ est compatible avec tous les
morphismes de $\mathcal{I}$, puisque, par hypothèse, chacun d'eux est un
composé de morphismes $f_j$. Nous avons ainsi obtenu une bijection canonique
$$
\lim_{B \in \Ob(\mathcal{J})} \Mor_\mathcal{C}(W, M(F(B)))
=
\lim_{A \in \Ob(\mathcal{I})} \Mor_\mathcal{C}(W, M(A))
$$
qui entraîne l'assertion du lemme concernant les limites. L'assertion
concernant les colimites se démontre de la même manière (preuve omise).
\end{proof}

\begin{lemma}
\label{lemma-fibre-products-equalizers-exist}
Soit $\mathcal{C}$ une catégorie.
Les assertions suivantes sont équivalentes :
\begin{enumerate}
\item Les limites finies connexes existent dans $\mathcal{C}$.
\item Les égalisateurs et les produits fibrés existent dans $\mathcal{C}$.
\end{enumerate}
\end{lemma}

\begin{proof}
Puisque les égalisateurs et les produits fibrés sont des limites finies
connexes, (1) implique (2). Réciproquement, soit $\mathcal{I}$
une catégorie d'indices finie et connexe. Soit
$F : \mathcal{J} \to \mathcal{I}$
le foncteur de catégories d'indices construit dans la preuve du
lemme \ref{lemma-finite-diagram-category}.
Nous pouvons alors remplacer $\mathcal{I}$ par $\mathcal{J}$.
Nous pouvons donc supposer que
$\Ob(\mathcal{I}) = \{x_1, \ldots, x_n\} \amalg \{y_1, \ldots, y_m\}$,
avec $n, m \geq 1$, et que tous les morphismes de $\mathcal{I}$ autres que les
identités sont des morphismes $f : x_i \to y_j$ pour certains
$i$ et $j$.

\medskip\noindent
Supposons $n > 1$. Puisque $\mathcal{I}$ est connexe, il existe
des indices $i_1, i_2$ et $j_0$, ainsi que des morphismes
$a : x_{i_1} \to y_{j_0}$ et $b : x_{i_2} \to y_{j_0}$.
Considérons la catégorie
$$
\mathcal{I}' =
\{x\} \amalg \{x_1, \ldots, \hat x_{i_1}, \ldots, \hat x_{i_2}, \ldots, x_n\}
\amalg \{y_1, \ldots, y_m\}
$$
avec
$$
\Mor_{\mathcal{I}'}(x, y_j) = \Mor_\mathcal{I}(x_{i_1}, y_j)
\amalg \Mor_\mathcal{I}(x_{i_2}, y_j)
$$
et tous les autres ensembles de morphismes égaux à ceux de
$\mathcal{I}$. Pour tout foncteur
$M : \mathcal{I} \to \mathcal{C}$, nous pouvons construire un foncteur
$M' : \mathcal{I}' \to \mathcal{C}$ en posant
$$
M'(x) = M(x_{i_1}) \times_{M(a), M(y_{j_0}), M(b)} M(x_{i_2})
$$
et, pour un morphisme $f' : x \to y_j$ correspondant, disons, à
$f : x_{i_1} \to y_j$, en posant $M'(f') = M(f) \circ \text{pr}_1$.
Le foncteur $M$ admet alors une limite si et seulement si le foncteur
$M'$ en admet une (preuve omise). Par récurrence, on se ramène donc
au cas $n = 1$.

\medskip\noindent
Si $n = 1$, la limite de tout $M : \mathcal{I} \to \mathcal{C}$ est
l'égalisateur successif des paires de morphismes $x_1 \to y_j$ ;
elle existe donc par hypothèse.
\end{proof}

\begin{lemma}
\label{lemma-almost-finite-limits-exist}
Soit $\mathcal{C}$ une catégorie.
Les assertions suivantes sont équivalentes :
\begin{enumerate}
\item Les limites finies non vides existent dans $\mathcal{C}$.
\item Les produits de deux objets et les égalisateurs existent dans
$\mathcal{C}$.
\item Les produits de deux objets et les produits fibrés existent dans
$\mathcal{C}$.
\end{enumerate}
\end{lemma}

\begin{proof}
Puisque les produits de deux objets, les produits fibrés et les égalisateurs
sont des limites dont les catégories d'indices sont non vides, (1) implique
(2) et (3). Supposons (2). Les produits finis non vides et les égalisateurs
existent alors. Le lemme \ref{lemma-limits-products-equalizers}
montre donc que les limites finies non vides existent, c'est-à-dire que
(1) est vraie. Supposons (3).
Si $a, b : A \to B$ sont des morphismes de $\mathcal{C}$, l'égalisateur
de $a, b$ est
$$
(A \times_{a, B, b} A)\times_{(\text{pr}_1, \text{pr}_2), A \times A, \Delta} A.
$$
Ainsi, (3) implique (2), ce qui démontre le lemme.
\end{proof}

\begin{lemma}
\label{lemma-finite-limits-exist}
Soit $\mathcal{C}$ une catégorie.
Les assertions suivantes sont équivalentes :
\begin{enumerate}
\item Les limites finies existent dans $\mathcal{C}$.
\item Les produits finis et les égalisateurs existent.
\item La catégorie possède un objet final et les produits fibrés existent.
\end{enumerate}
\end{lemma}

\begin{proof}
Puisque les produits finis, les produits fibrés, les égalisateurs et les
objets finaux sont des limites sur des catégories d'indices finies, (1)
implique (2) et (3). Le lemme \ref{lemma-limits-products-equalizers}
ci-dessus montre que (2) implique (1). Supposons (3).
Le produit $A \times B$ est le produit fibré au-dessus de l'objet final.
Si $a, b : A \to B$ sont des morphismes de $\mathcal{C}$, l'égalisateur
de $a, b$ est
$$
(A \times_{a, B, b} A)\times_{(\text{pr}_1, \text{pr}_2), A \times A, \Delta} A.
$$
Ainsi, (3) implique (2), ce qui démontre le lemme.
\end{proof}

\begin{lemma}
\label{lemma-push-outs-coequalizers-exist}
Soit $\mathcal{C}$ une catégorie.
Les assertions suivantes sont équivalentes :
\begin{enumerate}
\item Les colimites finies connexes existent dans $\mathcal{C}$.
\item Les coégalisateurs et les sommes amalgamées existent dans $\mathcal{C}$.
\end{enumerate}
\end{lemma}

\begin{proof}
Omis. Indication : c'est l'énoncé dual du
lemme \ref{lemma-fibre-products-equalizers-exist}.
\end{proof}

\begin{lemma}
\label{lemma-almost-finite-colimits-exist}
Soit $\mathcal{C}$ une catégorie.
Les assertions suivantes sont équivalentes :
\begin{enumerate}
\item Les colimites finies non vides existent dans $\mathcal{C}$.
\item Les coproduits de deux objets et les coégalisateurs existent dans
$\mathcal{C}$.
\item Les coproduits de deux objets et les sommes amalgamées existent dans
$\mathcal{C}$.
\end{enumerate}
\end{lemma}

\begin{proof}
Omis. Indication : c'est l'énoncé dual du
lemme \ref{lemma-almost-finite-limits-exist}.
\end{proof}

\begin{lemma}
\label{lemma-colimits-exist}
Soit $\mathcal{C}$ une catégorie.
Les assertions suivantes sont équivalentes :
\begin{enumerate}
\item Les colimites finies existent dans $\mathcal{C}$.
\item Les coproduits finis et les coégalisateurs existent dans $\mathcal{C}$.
\item La catégorie possède un objet initial et les sommes amalgamées existent.
\end{enumerate}
\end{lemma}

\begin{proof}
Omis. Indication : c'est l'énoncé dual du
lemme \ref{lemma-finite-limits-exist}.
\end{proof}





\section{Limites inductives filtrantes}
\label{section-directed-colimits}

\noindent
Les colimites sont plus faciles à calculer ou à décrire lorsqu'elles
portent sur un diagramme filtrant. En voici la définition.

\begin{definition}
\label{definition-directed}
Nous disons qu'un diagramme $M : \mathcal{I} \to \mathcal{C}$ est
{\it dirigé}, ou {\it filtrant}, si les conditions suivantes sont satisfaites :
\begin{enumerate}
\item la catégorie $\mathcal{I}$ possède au moins un objet ;
\item pour toute paire d'objets $x, y$ de $\mathcal{I}$,
il existe un objet $z$ et des morphismes $x \to z$,
$y \to z$ ; et
\item pour toute paire d'objets $x, y$ de $\mathcal{I}$
et toute paire de morphismes $a, b : x \to y$ de $\mathcal{I}$,
il existe un morphisme $c : y \to z$ de $\mathcal{I}$
tel que $M(c \circ a) = M(c \circ b)$ comme morphismes dans $\mathcal{C}$.
\end{enumerate}
Nous disons qu'une catégorie d'indices $\mathcal{I}$ est
{\it dirigée}, ou {\it filtrante}, si
$\text{id} : \mathcal{I} \to \mathcal{I}$ est filtrant
(autrement dit, on supprime le $M$ dans la partie (3) ci-dessus).
\end{definition}

\noindent
Tout diagramme dont la catégorie d'indices est filtrante est filtrant ;
c'est ainsi que les limites inductives filtrantes se présentent habituellement.
En fait, si $M : \mathcal{I} \to \mathcal{C}$ est un diagramme filtrant,
on peut factoriser $M$ sous la forme
$\mathcal{I} \to \mathcal{I}' \to \mathcal{C}$,
où $\mathcal{I}'$ est une catégorie d'indices
filtrante\footnote{On prend pour $\mathcal{I}'$ les mêmes objets que pour
$\mathcal{I}$, mais $\Mor_{\mathcal{I}'}(x, y)$ est le quotient de
$\Mor_\mathcal{I}(x, y)$ par la relation d'équivalence qui identifie
$a, b : x \to y$ si $M(a) = M(b)$.},
de telle sorte que $\colim_\mathcal{I} M$ existe si et seulement si
$\colim_{\mathcal{I}'} M'$ existe ; dans ce cas, les colimites sont
canoniquement isomorphes.

\medskip\noindent
Supposons que $M : \mathcal{I} \to \textit{Ensembles}$ soit un diagramme filtrant.
Dans ce cas, la relation d'équivalence de la formule
$$
\colim_\mathcal{I} M
=
(\coprod\nolimits_{i\in I} M_i)/\sim
$$
se décrit simplement comme suit :
$$
m_i \sim m_{i'}
\Leftrightarrow
\exists i'', \phi : i \to i'', \phi': i' \to i'',
M(\phi)(m_i) = M(\phi')(m_{i'}).
$$
Autrement dit, deux éléments sont égaux dans la colimite si et seulement
s'ils « finissent par devenir égaux ».

\begin{lemma}
\label{lemma-directed-commutes}
Soient $\mathcal{I}$ et $\mathcal{J}$ des catégories d'indices.
Supposons que $\mathcal{I}$ soit filtrante et que $\mathcal{J}$ soit finie.
Soit $M : \mathcal{I} \times \mathcal{J} \to \textit{Ensembles}$,
$(i, j) \mapsto M_{i, j}$, un diagramme de diagrammes d'ensembles.
Alors
$$
\colim_i \lim_j M_{i, j}
=
\lim_j \colim_i M_{i, j}.
$$
En particulier, les colimites sur $\mathcal{I}$ commutent aux produits finis,
aux produits fibrés et aux égalisateurs d'ensembles.
\end{lemma}

\begin{proof}
Omis. C'est en fait un exercice agréable de montrer qu'une catégorie est
filtrante si et seulement si les colimites sur cette catégorie commutent
aux limites finies (à valeurs dans la catégorie des ensembles).
\end{proof}

\noindent
Donnons un contre-exemple au lemme lorsque $\mathcal{J}$ est infinie.
Prenons pour $\mathcal{I}$ la catégorie dont les objets sont
$\mathbf{N} = \{1, 2, 3, \ldots\}$ et dans laquelle il existe un unique
morphisme $i \to i'$ dès que $i \leq i'$.
Prenons pour $\mathcal{J}$ la catégorie discrète
$\mathbf{N} = \{1, 2, 3, \ldots\}$ (ses seuls morphismes sont les identités).
Posons $M_{i, j} = \{1, 2, \ldots, i\}$, avec les applications d'inclusion
évidentes $M_{i, j} \to M_{i', j}$ lorsque $i \leq i'$. Dans ce cas,
$\colim_i M_{i, j} = \mathbf{N}$, et donc
$$
\lim_j \colim_i M_{i, j}
=
\prod\nolimits_j \mathbf{N}
=
\mathbf{N}^\mathbf{N}
$$
D'autre part, $\lim_j M_{i, j} = \prod\nolimits_j M_{i, j}$, et donc
$$
\colim_i \lim_j M_{i, j}
=
\bigcup\nolimits_i \{1, 2, \ldots, i\}^{\mathbf{N}}
$$
qui est plus petit que l'autre limite.

\begin{lemma}
\label{lemma-cofinal-in-filtered}
\begin{reference}
\cite[Proposition 3.2.4]{KS}
\end{reference}
Soit $\mathcal{I}$ une catégorie. Soit $\mathcal{J}$ une sous-catégorie
pleine. Supposons que $\mathcal{I}$ soit filtrante. Supposons en outre que,
pour tout objet $i$ de $\mathcal{I}$, il existe un morphisme $i \to j$
vers un objet $j$ de $\mathcal{J}$. Alors $\mathcal{J}$
est filtrante et cofinale dans $\mathcal{I}$.
\end{lemma}

\begin{proof}
Omis. Exercice agréable sur les notions en jeu.
\end{proof}

\noindent
Nous avons parfois besoin d'un contrôle plus fin des
conditions que l'on peut imposer aux catégories d'indices. Nous ajoutons
donc quelques lemmes sur les diverses conditions possibles.

\begin{lemma}
\label{lemma-preserve-products}
Soit $\mathcal{I}$ une catégorie d'indices, c'est-à-dire une catégorie.
Supposons que, pour toute paire d'objets $x, y$ de $\mathcal{I}$,
il existe un objet $z$ et des morphismes $x \to z$ et $y \to z$.
Alors
\begin{enumerate}
\item si $M$ et $N$ sont des diagrammes d'ensembles sur $\mathcal{I}$,
alors $\colim (M_i \times N_i) \to \colim M_i \times \colim N_i$
est surjective ;
\item en général, les colimites de diagrammes d'ensembles sur $\mathcal{I}$
ne commutent pas aux produits finis non vides.
\end{enumerate}
\end{lemma}

\begin{proof}
Preuve de (1). Soit $(\overline{m}, \overline{n})$
un élément de $\colim M_i \times \colim N_i$.
Il existe alors $m \in M_x$ et $n \in N_y$, pour certains
$x, y \in \Ob(\mathcal{I})$, tels que $m$ s'envoie sur
$\overline{m}$ et $n$ sur $\overline{n}$. Voir la
section \ref{section-limit-sets}.
Choisissons $a : x \to z$ et $b : y \to z$
dans $\mathcal{I}$. Alors $(M(a)(m), N(b)(n))$ est un élément de
$(M \times N)_z$ dont l'image dans $\colim (M_i \times N_i)$
s'envoie sur $(\overline{m}, \overline{n})$, comme voulu.

\medskip\noindent
Preuve de (2). Soit $G$ un groupe non trivial et
soit $\mathcal{I}$ la catégorie à un objet dont le monoïde des
endomorphismes est $G$. Alors $\mathcal{I}$ satisfait trivialement la
condition du lemme. Faisons agir $G$ sur lui-même par translations et
considérons le $G$-ensemble $G$ comme un diagramme indexé par $\mathcal{I}$
à valeurs dans les ensembles. Alors
$$
\colim_\mathcal{I} G \times \colim_\mathcal{I} G \cong G/G \times G/G
$$
n'est pas isomorphe à
$$
\colim_\mathcal{I} (G \times G) \cong (G \times G)/G
$$
Cet exemple montre qu'on ne peut pas simplement supprimer la condition
supplémentaire du lemme \ref{lemma-directed-commutes},
même si l'on ne s'intéresse qu'aux produits finis.
\end{proof}

\begin{lemma}
\label{lemma-colimits-abelian-as-sets}
Soit $\mathcal{I}$ une catégorie d'indices, c'est-à-dire une catégorie.
Supposons que, pour toute paire d'objets $x, y$ de $\mathcal{I}$,
il existe un objet $z$ et des morphismes $x \to z$ et $y \to z$.
Soit $M : \mathcal{I} \to \textit{Ab}$ un diagramme de groupes abéliens
sur $\mathcal{I}$. Alors l'application naturelle de la colimite de $M$
dans la catégorie des ensembles vers la colimite de $M$ dans la catégorie
des groupes abéliens est surjective.
\end{lemma}

\begin{proof}
Rappelons que la colimite dans la catégorie des ensembles est le quotient
de la réunion disjointe $\coprod M_i$ par une relation d'équivalence ; voir la
section \ref{section-limit-sets}.
De même, la colimite dans la catégorie des groupes abéliens est un quotient
de la somme directe $\bigoplus M_i$.
L'hypothèse du lemme signifie que, étant donnés
$i, j \in \Ob(\mathcal{I})$, $m \in M_i$ et $n \in M_j$, on peut trouver
un objet $k$ et des morphismes $a : i \to k$ et $b : j \to k$.
Ainsi, $m + n$ est représenté dans la colimite par l'élément
$M(a)(m) + M(b)(n)$ de $M_k$. L'application naturelle de la réunion
disjointe $\coprod M_i$ vers la colimite est donc surjective.
\end{proof}

\begin{lemma}
\label{lemma-split-into-connected}
Soit $\mathcal{I}$ une catégorie d'indices, c'est-à-dire une catégorie.
Supposons que, pour tout diagramme en traits pleins
$$
\xymatrix{
x \ar[d] \ar[r] & y \ar@{..>}[d] \\
z \ar@{..>}[r] & w
}
$$
dans $\mathcal{I}$, il existe un objet $w$ et des flèches pointillées
qui rendent le diagramme commutatif. Alors $\mathcal{I}$ est soit vide,
soit une réunion disjointe non vide de catégories connexes possédant
la même propriété.
\end{lemma}

\begin{proof}
Si $\mathcal{I}$ est la catégorie vide, le lemme est vrai.
Sinon, définissons une relation sur les objets de $\mathcal{I}$ en posant
$x \sim y$ s'il existe $z$ et des morphismes
$x \to z$ et $y \to z$. C'est une relation d'équivalence
par l'hypothèse du lemme. Ainsi, $\Ob(\mathcal{I})$
est une réunion disjointe de classes d'équivalence. Soient
$\mathcal{I}_j$ les sous-catégories pleines correspondant à ces classes.
Alors $\mathcal{I} = \coprod \mathcal{I}_j$, où les $\mathcal{I}_j$
sont non vides, comme voulu.
\end{proof}

\begin{lemma}
\label{lemma-preserve-injective-maps}
Soit $\mathcal{I}$ une catégorie d'indices, c'est-à-dire une catégorie.
Supposons que, pour tout diagramme en traits pleins
$$
\xymatrix{
x \ar[d] \ar[r] & y \ar@{..>}[d] \\
z \ar@{..>}[r] & w
}
$$
dans $\mathcal{I}$, il existe un objet $w$ et des flèches pointillées
qui rendent le diagramme commutatif. Alors
\begin{enumerate}
\item tout morphisme injectif $M \to N$ de diagrammes d'ensembles sur
$\mathcal{I}$ induit une application injective
$\colim M_i \to \colim N_i$ d'ensembles ;
\item en général, il n'en va pas de même pour les diagrammes de groupes
abéliens et leurs colimites.
\end{enumerate}
\end{lemma}

\begin{proof}
Si $\mathcal{I}$ est la catégorie vide, le lemme est vrai.
Nous pouvons donc supposer $\mathcal{I}$ non vide. Dans ce cas,
on peut écrire $\mathcal{I} = \coprod \mathcal{I}_j$, où chaque
$\mathcal{I}_j$ est non vide et satisfait la même propriété ; voir le
lemme \ref{lemma-split-into-connected}. Puisque
$\colim_\mathcal{I} M = \coprod_j \colim_{\mathcal{I}_j} M|_{\mathcal{I}_j}$,
la preuve de (1) se ramène au cas connexe.

\medskip\noindent
Supposons $\mathcal{I}$ connexe et $M \to N$ injectif, c'est-à-dire
que toutes les applications $M_i \to N_i$ sont injectives.
Identifions $M_i$ à l'image de $M_i \to N_i$ ; autrement dit,
considérons $M_i$ comme un sous-ensemble de $N_i$.
Nous utiliserons sans plus de mention la description des colimites donnée
dans la section \ref{section-limit-sets}.
Soient $s, s' \in \colim M_i$ ayant la même image dans $\colim N_i$.
Disons que $s$ provient d'un élément $m$ de $M_i$ et $s'$ d'un
élément $m'$ de $M_{i'}$. Il existe alors une suite
$i = i_0, i_1, \ldots, i_{2n} = i'$ d'objets de $\mathcal{I}$
et des morphismes
$$
\xymatrix{
&
i_1 \ar[ld] \ar[rd] & &
i_3 \ar[ld] & &
i_{2n-1} \ar[rd] & \\
i = i_0 & &
i_2 & &
\ldots & &
i_{2n} = i'
}
$$
ainsi que des éléments $n_{i_j} \in N_{i_j}$ dont les images coïncident
par les applications $N_{i_{2k-1}} \to N_{i_{2k-2}}$ et
$N_{i_{2k-1}}
\to N_{i_{2k}}$ induites par les flèches de
$\mathcal{I}$ ci-dessus, avec $n_{i_0} = m$ et $n_{i_{2n}} = m'$.
Nous allons montrer par récurrence sur $n$ que cela entraîne $s = s'$.
Le cas initial $n = 0$ est trivial. Supposons $n \geq 1$.
L'hypothèse sur $\mathcal{I}$ fournit un diagramme commutatif
$$
\xymatrix{
& i_1 \ar[ld] \ar[rd] \\
i_0 \ar[rd] & & i_2 \ar[ld] \\
& w
}
$$
Nous en déduisons que $m$ et $n_{i_2}$ ont la même image dans $N_w$,
puisqu'ils sont tous deux l'image de l'élément $n_{i_1}$.
En particulier, cet élément est un élément $m'' \in M_w$ qui
définit le même élément que $s$ dans $\colim M_i$.
Nous obtenons alors la chaîne
$$
\xymatrix{
&
i_3 \ar[ld] \ar[rd] & &
i_5 \ar[ld] & &
i_{2n-1} \ar[rd] & \\
w & &
i_4 & &
\ldots & &
i_{2n} = i'
}
$$
et les éléments $n_{i_j}$ pour $j \geq 3$ ; cette chaîne est plus courte
que la chaîne initiale. Cela démontre l'étape de récurrence et achève
la preuve de (1).

\medskip\noindent
Soit $G$ un groupe et soit $\mathcal{I}$ la catégorie à un objet dont
le monoïde des endomorphismes est $G$. Alors $\mathcal{I}$ satisfait
la condition du lemme, car, étant donnés $g_1, g_2 \in G$, on peut trouver
$h_1, h_2 \in G$ tels que $h_1 g_1 = h_2 g_2$.
Un diagramme $M$ sur $\mathcal{I}$ à valeurs dans $\textit{Ab}$ est la même
chose qu'un groupe abélien $M$ muni d'une action de $G$, et
$\colim_\mathcal{I} M$ est le groupe des coinvariants $M_G$ de $M$.
Prenons pour $G$ le groupe d'ordre $2$ agissant trivialement sur
$M = \mathbf{Z}/2\mathbf{Z}$, plongé dans le premier facteur de
$N = \mathbf{Z}/2\mathbf{Z} \times \mathbf{Z}/2\mathbf{Z}$,
où l'élément non trivial de $G$ agit par
$(x, y) \mapsto (x + y, y)$. Alors $M_G \to N_G$ est nul.
\end{proof}

\begin{lemma}
\label{lemma-split-into-directed}
Soit $\mathcal{I}$ une catégorie d'indices, c'est-à-dire une catégorie.
Supposons que
\begin{enumerate}
\item pour toute paire de morphismes $a : w \to x$ et $b : w \to y$
dans $\mathcal{I}$, il existe un objet $z$ et des morphismes
$c : x \to z$ et $d : y \to z$ tels que $c \circ a = d \circ b$, et
\item pour toute paire de morphismes $a, b : x \to y$, il existe
un morphisme $c : y \to z$ tel que $c \circ a = c \circ b$.
\end{enumerate}
Alors $\mathcal{I}$ est une réunion (éventuellement vide)
de catégories d'indices filtrantes disjointes $\mathcal{I}_j$.
\end{lemma}

\begin{proof}
Si $\mathcal{I}$ est la catégorie vide, le lemme est vrai.
Sinon, définissons une relation sur les objets de $\mathcal{I}$ en posant
$x \sim y$ s'il existe $z$ et des morphismes
$x \to z$ et $y \to z$. C'est une relation d'équivalence
par la première hypothèse du lemme. Ainsi, $\Ob(\mathcal{I})$
est une réunion disjointe de classes d'équivalence. Soient
$\mathcal{I}_j$ les sous-catégories pleines correspondant à ces classes.
Le reste découle immédiatement des définitions.
\end{proof}

\begin{lemma}
\label{lemma-almost-directed-commutes-equalizers}
Soit $\mathcal{I}$ une catégorie d'indices satisfaisant les hypothèses du
lemme \ref{lemma-split-into-directed} ci-dessus. Alors les colimites sur
$\mathcal{I}$ commutent aux produits fibrés et aux égalisateurs d'ensembles
(et, plus généralement, aux limites finies connexes).
\end{lemma}

\begin{proof}
D'après le lemme \ref{lemma-split-into-directed},
nous pouvons écrire $\mathcal{I} = \coprod \mathcal{I}_j$, chaque
$\mathcal{I}_j$ étant filtrante. D'après le
lemme \ref{lemma-directed-commutes}, les colimites sur $\mathcal{I}_j$
commutent aux égalisateurs et aux produits fibrés. Il suffit donc de montrer
que les égalisateurs et les produits fibrés commutent aux coproduits dans la
catégorie des ensembles (y compris aux coproduits vides).
Autrement dit, étant donnés un ensemble $J$, des ensembles
$A_j, B_j, C_j$ et des applications
$A_j \to B_j$, $C_j \to B_j$ pour $j \in J$, il faut montrer que
$$
(\coprod\nolimits_{j \in J} A_j)
\times_{(\coprod\nolimits_{j \in J} B_j)}
(\coprod\nolimits_{j \in J} C_j)
=
\coprod\nolimits_{j \in J} A_j \times_{B_j} C_j
$$
et, étant données $a_j, a'_j : A_j \to B_j$, que
$$
\text{Égalisateur}(
\coprod\nolimits_{j \in J} a_j,
\coprod\nolimits_{j \in J} a'_j)
=
\coprod\nolimits_{j \in J}
\text{Égalisateur}(a_j, a'_j)
$$
Cela reste vrai si $J = \emptyset$. Les détails sont omis.
\end{proof}



\section{Limites projectives cofiltrantes}
\label{section-codirected-limits}

\noindent
Les limites projectives sont plus faciles à calculer ou à décrire lorsqu'elles
portent sur un diagramme cofiltrant. En voici la définition.

\begin{definition}
\label{definition-codirected}
On dit qu'un diagramme $M : \mathcal{I} \to \mathcal{C}$ est
{\it codirigé} ou {\it cofiltrant} si les conditions suivantes sont satisfaites :
\begin{enumerate}
\item la catégorie $\mathcal{I}$ possède au moins un objet ;
\item pour toute paire d'objets $x, y$ de $\mathcal{I}$, il existe un objet
$z$ et des morphismes $z \to x$, $z \to y$ ;
\item pour toute paire d'objets $x, y$ de $\mathcal{I}$ et toute paire de
morphismes $a, b : x \to y$ de $\mathcal{I}$, il existe un morphisme
$c : w \to x$ de $\mathcal{I}$ tel que
$M(a \circ c) = M(b \circ c)$ comme morphismes dans $\mathcal{C}$.
\end{enumerate}
On dit qu'une catégorie d'indices $\mathcal{I}$ est {\it codirigée}, ou
{\it cofiltrante}, si $\text{id} : \mathcal{I} \to \mathcal{I}$ est
cofiltrant (autrement dit, on efface le $M$ dans la partie (3) ci-dessus).
\end{definition}

\noindent
On observe que tout diagramme dont la catégorie d'indices est cofiltrante
est cofiltrant ; c'est sous cette forme que la situation se présente
habituellement.

\medskip\noindent
Pour illustrer en quoi les limites projectives cofiltrantes d'ensembles sont
« plus faciles » que les limites projectives générales, signalons qu'un
diagramme cofiltrant d'ensembles finis non vides a une limite non vide
(lemme \ref{lemma-nonempty-limit}). Ce résultat n'est pas vrai pour une
limite générale d'ensembles finis non vides.











\section{Limites projectives et inductives sur des ensembles préordonnés}
\sectionmark{Limites sur les ensembles préordonnés}
\label{section-posets-limits}

\noindent
Les systèmes sur des ensembles préordonnés constituent un cas particulier
de diagrammes.

\begin{definition}
\label{definition-directed-set}
Soit $I$ un ensemble et soit $\leq$ une relation binaire sur $I$.
\begin{enumerate}
\item On dit que $\leq$ est un {\it préordre} si cette relation est
transitive (si $i \leq j$ et $j \leq k$, alors $i \leq k$) et réflexive
($i \leq i$ pour tout $i \in I$).
\item Un {\it ensemble préordonné} est un ensemble muni d'un préordre.
\item Un {\it ensemble ordonné filtrant} est un ensemble préordonné
$(I, \leq)$ tel que $I$ soit non vide et tel que $\forall i, j \in I$,
il existe $k \in I$ vérifiant $i \leq k, j \leq k$.
\item On dit que $\leq$ est un {\it ordre partiel} si c'est un préordre
antisymétrique (si $i \leq j$ et $j \leq i$, alors $i = j$).
\item Un {\it ensemble partiellement ordonné} est un ensemble muni d'un
ordre partiel.
\item Un {\it ensemble partiellement ordonné filtrant} est un ensemble
ordonné filtrant dont l'ordre est partiel.
\end{enumerate}
\end{definition}

\noindent
Il est d'usage d'omettre $\leq$ de la notation lorsqu'on parle d'ensembles
préordonnés, c'est-à-dire de parler de l'ensemble préordonné $I$ plutôt que
de l'ensemble préordonné $(I, \leq)$. Étant donné un ensemble préordonné
$I$, le symbole $\geq$ est défini par la règle
$i \geq j \Leftrightarrow j \leq i$ pour tous $i, j \in I$.
L'abréviation anglaise « poset » est parfois employée pour désigner un
ensemble partiellement ordonné.

\medskip\noindent
À tout ensemble préordonné $I$, on peut associer une catégorie : ses objets
sont les éléments de $I$, il existe exactement un morphisme $i \to i'$ si
$i \leq i'$, et il n'en existe aucun sinon. Réciproquement, étant donnée une
catégorie $\mathcal{C}$ ayant au plus une flèche entre deux objets quelconques,
l'ensemble $\Ob(\mathcal{C})$ est muni du préordre défini par la règle
$x \leq y \Leftrightarrow \Mor_\mathcal{C}(x, y) \not = \emptyset$.

\begin{definition}
\label{definition-system-over-poset}
Soit $(I, \leq)$ un ensemble préordonné. Soit $\mathcal{C}$ une catégorie.
\begin{enumerate}
\item Un {\it système sur $I$ à valeurs dans $\mathcal{C}$}, parfois appelé
{\it système inductif sur $I$ à valeurs dans $\mathcal{C}$}, est la donnée d'objets
$M_i$ de $\mathcal{C}$ et, pour chaque $i \leq i'$, d'un morphisme
$f_{ii'} : M_i \to M_{i'}$ tel que $f_{ii} = \text{id}$ et
$f_{ii''} = f_{i'i''} \circ f_{i i'}$ dès que $i \leq i' \leq i''$.
\item Un {\it système inverse sur $I$ à valeurs dans $\mathcal{C}$}, parfois appelé
{\it système projectif sur $I$ à valeurs dans $\mathcal{C}$}, est la donnée d'objets
$M_i$ de $\mathcal{C}$ et, pour chaque $i' \leq i$, d'un morphisme
$f_{ii'} : M_i \to M_{i'}$ tel que $f_{ii} = \text{id}$ et
$f_{ii''} = f_{i'i''} \circ f_{i i'}$ dès que $i'' \leq i' \leq i$.
(Remarquer le renversement des inégalités.)
\end{enumerate}
On écrira $(M_i, f_{ii'})$ pour désigner un système de l'un ou l'autre type sur $I$.
Les morphismes $f_{ii'}$ sont parfois appelés les
{\it morphismes de transition}.
\end{definition}

\noindent
Autrement dit, un système sur $I$ n'est autre qu'un diagramme
$M : \mathcal{I} \to \mathcal{C}$, où $\mathcal{I}$ est la catégorie
associée ci-dessus à $I$ : ses objets sont les éléments de $I$ et il existe
une unique flèche $i \to i'$ dans $\mathcal{I}$ si et seulement si
$i \leq i'$. Un système inverse est un diagramme
$M : \mathcal{I}^{opp} \to \mathcal{C}$. De ce point de vue, on pourrait
prendre les limites et colimites de tout système, inverse ou non, sur $I$. Toutefois, il est
d'usage de ne prendre {\it que les limites inductives des systèmes sur $I$}
et {\it que les limites projectives des systèmes inverses sur $I$}.
Plus précisément, étant donné un système $(M_i, f_{ii'})$ sur $I$, la limite
inductive du système $(M_i, f_{ii'})$ est définie par
$$
\colim_{i \in I} M_i = \colim_\mathcal{I} M,
$$
c'est-à-dire comme la limite inductive du diagramme correspondant.
Étant donné un système inverse $(M_i, f_{ii'})$ sur $I$, la limite projective
du système inverse $(M_i, f_{ii'})$ est définie par
$$
\lim_{i \in I} M_i = \lim_{\mathcal{I}^{opp}} M,
$$
c'est-à-dire comme la limite projective du diagramme correspondant.

\begin{remark}
\label{remark-preorder-versus-partial-order}
Soit $I$ un ensemble préordonné. À partir de $I$, on peut construire un
ensemble partiellement ordonné canonique $\overline{I}$ et une application
croissante $\pi : I \to \overline{I}$. En effet, on peut définir une relation
d'équivalence $\sim$ sur $I$ par la règle
$$
i \sim j \Leftrightarrow (i \leq j\text{ et }j \leq i).
$$
On pose $\overline{I} = I/\sim$ et l'on note
$\pi : I \to \overline{I}$ l'application quotient. Enfin, $\overline{I}$
est muni d'un unique ordre partiel tel que
$\pi(i) \leq \pi(j) \Leftrightarrow i \leq j$.
Remarquons que, si $I$ est un ensemble ordonné filtrant, alors
$\overline{I}$ est un ensemble partiellement ordonné filtrant.
Étant donné un système, inverse ou non, $N$ sur $\overline{I}$, on obtient
un système du même type $M$ sur $I$ en posant $M_i = N_{\pi(i)}$.
Cette construction définit un foncteur entre les catégories de systèmes,
inverses ou non, sur $I$ et sur $\overline{I}$. C'est en fait une équivalence.
La raison en est que, si $i \sim j$, alors, pour tout système $M$ sur $I$,
les morphismes $M_i \to M_j$ et $M_j \to M_i$ sont des isomorphismes
réciproques. Plus précisément, si l'on choisit une section
$s : \overline{I} \to I$ de $\pi$, un quasi-inverse du foncteur précédent
envoie $M$ sur $N$ défini par
$N_{\overline{i}} = M_{s(\overline{i})}$.
Enfin, cette correspondance est compatible avec les limites inductives des
systèmes : si $M$ et $N$ sont reliés comme ci-dessus et si l'une des limites
$\colim_{\overline{I}} N$ ou $\colim_I M$ existe, alors l'autre existe aussi et
$\colim_{\overline{I}} N = \colim_I M$.
Des résultats analogues valent pour les systèmes inverses et leurs limites
projectives.
\end{remark}

\noindent
Il résulte de la remarque \ref{remark-preorder-versus-partial-order} que,
lorsqu'on calcule la limite inductive d'un système ou la limite projective
d'un système inverse, on peut toujours supposer que le préordre est un ordre
partiel.

\begin{definition}
\label{definition-directed-system}
Soit $I$ un ensemble préordonné. On dit qu'un système (resp.\ un système
inverse) $(M_i, f_{ii'})$ est un {\it système filtrant} (resp.\ un
{\it système inverse filtrant}) si $I$ est un ensemble ordonné filtrant
(définition \ref{definition-directed-set}) : $I$ est non vide et, pour tous
$i_1, i_2 \in I$, il existe $i\in I$ tel que
$i_1 \leq i$ et $i_2 \leq i$.
\end{definition}

\noindent
Dans ce cas, la limite inductive est parfois appelée (malheureusement)
« limite directe ». Nous n'emploierons pas cette dernière terminologie.
En fait, les diagrammes sur une catégorie filtrante ne sont pas plus
généraux que les systèmes filtrants, au sens suivant.

\begin{lemma}
\label{lemma-directed-category-system}
Soit $\mathcal{I}$ une catégorie d'indices filtrante. Il existe un ensemble
ordonné filtrant $I$ et un système $(x_i, \varphi_{ii'})$ sur $I$ à valeurs dans
$\mathcal{I}$ possédant les propriétés suivantes :
\begin{enumerate}
\item Pour toute catégorie $\mathcal{C}$ et tout diagramme
$M : \mathcal{I} \to \mathcal{C}$ à valeurs dans $\mathcal{C}$, notons
$(M(x_i), M(\varphi_{ii'}))$ le système correspondant sur $I$. Si
$\colim_{i \in I} M(x_i)$ existe, alors $\colim_\mathcal{I} M$ existe aussi
et la transformation
$$
\theta :
\colim_{i \in I} M(x_i)
\longrightarrow
\colim_\mathcal{I} M
$$
du lemme \ref{lemma-functorial-colimit} est un isomorphisme.
\item Pour toute catégorie $\mathcal{C}$ et tout diagramme
$M : \mathcal{I}^{opp} \to \mathcal{C}$ à valeurs dans $\mathcal{C}$, notons
$(M(x_i), M(\varphi_{ii'}))$ le système inverse correspondant sur $I$.
Si $\lim_{i \in I} M(x_i)$ existe, alors $\lim_{\mathcal{I}^{opp}} M$
existe aussi et la transformation
$$
\theta :
\lim_{\mathcal{I}^{opp}} M
\longrightarrow
\lim_{i \in I} M(x_i)
$$
du lemme \ref{lemma-functorial-limit} est un isomorphisme.
\end{enumerate}
\end{lemma}

\begin{proof}
Comme expliqué dans le texte qui suit la
définition \ref{definition-system-over-poset}, on peut considérer les
ensembles préordonnés comme des catégories et les systèmes comme des
foncteurs. Tout au long de la preuve, nous passerons librement d'un point de
vue à l'autre. Nous démontrons le premier énoncé en construisant une catégorie
$\mathcal{I}_0$ correspondant à un ensemble ordonné
filtrant\footnote{En fait, notre construction produira un ensemble
partiellement ordonné filtrant.}, et un foncteur cofinal
$M_0 : \mathcal{I}_0 \to \mathcal{I}$. Alors, d'après le
lemme \ref{lemma-cofinal}, la limite inductive d'un diagramme
$M : \mathcal{I} \to \mathcal{C}$ coïncide avec celle du diagramme
$M \circ M_0 : \mathcal{I}_0 \to \mathcal{C}$, d'où l'énoncé. Le second
énoncé est dual du premier et peut se démontrer en interprétant une limite
projective dans $\mathcal{C}$ comme une limite inductive dans
$\mathcal{C}^{opp}$. Nous omettons les détails.

\medskip\noindent
Une catégorie $\mathcal{F}$ est dite {\em finiment engendrée} s'il existe un
ensemble fini $F$ de flèches de $\mathcal{F}$ tel que toute flèche de
$\mathcal{F}$ s'obtienne en composant des flèches de $F$. Cela implique en
particulier que $\mathcal{F}$ possède un nombre fini d'objets. Commençons par
nous ramener au cas où $\mathcal{I}$ possède la propriété suivante : toute
sous-catégorie finiment engendrée de $\mathcal{I}$ est contenue dans une
sous-catégorie finiment engendrée ayant un unique objet final.

\medskip\noindent
Notons $\omega$ l'ensemble ordonné filtrant des ordinaux finis, que nous
considérons comme une catégorie filtrante. On vérifie aisément que la
catégorie produit $\mathcal{I}\times \omega$ est elle aussi filtrante et que
la projection
$\Pi : \mathcal{I} \times \omega \to \mathcal{I}$
est cofinale.

\medskip\noindent
Soit maintenant $\mathcal{F}$ une sous-catégorie finiment engendrée
quelconque de $\mathcal{I}\times \omega$. À l'aide des axiomes d'une catégorie
filtrante et d'une récurrence simple sur un ensemble fini de générateurs de
$\mathcal{F}$, on peut construire un cocône $(\{f_i\}, i_\infty)$ dans
$\mathcal{I} \times \omega$ pour le diagramme
$\mathcal{F} \to \mathcal{I} \times \omega$. Autrement dit, on dispose d'un
morphisme $f_i : i \to i_\infty$ pour chaque objet $i$ de $\mathcal{F}$,
tel que, pour toute flèche $f : i \to i'$ de $\mathcal{F}$, on ait
$f_i = f_{i'} \circ f$. On peut en outre choisir $i_\infty$ de sorte qu'il
n'existe aucune flèche de $i_\infty$ vers un objet de $\mathcal{F}$. Cela est
possible puisqu'on peut toujours composer à gauche les flèches $f_i$ avec
une flèche qui est l'identité sur la composante $\mathcal{I}$ et strictement
croissante sur la composante $\omega$.
Notons maintenant $\mathcal{F}^+$ la catégorie formée de tous les objets et
de toutes les flèches de $\mathcal{F}$, auxquels on adjoint l'objet
$i_\infty$, la flèche identité $\text{id}_{i_\infty}$ et les flèches $f_i$.
Comme aucune flèche de $i_\infty$ dans $\mathcal{F}^+$ n'aboutit à un objet de
$\mathcal{F}$, l'ensemble des flèches de $\mathcal{F}^+$ est stable par
composition ; $\mathcal{F}^+$ est donc bien une catégorie. Par construction,
c'est une sous-catégorie finiment engendrée de $\mathcal{I}\times\omega$ dont
$i_\infty$ est l'unique objet final. Puisque, d'après le
lemme \ref{lemma-cofinal}, la limite inductive de tout diagramme
$M : \mathcal{I} \to \mathcal{C}$ coïncide avec celle de $M\circ\Pi$, on
obtient ainsi la réduction voulue.

\medskip\noindent
L'ensemble de toutes les sous-catégories finiment engendrées de
$\mathcal{I}$ ayant un unique objet final est naturellement ordonné par
inclusion. On prend pour $\mathcal{I}_0$ la catégorie correspondant à cet
ensemble. On dispose aussi d'un foncteur
$M_0 : \mathcal{I}_0 \to \mathcal{I}$ qui envoie une flèche
$\mathcal{F} \subset \mathcal{F'}$ de $\mathcal{I}_0$ sur l'unique morphisme
de l'objet final de $\mathcal{F}$ vers l'objet final de $\mathcal{F}'$.
Étant données deux sous-catégories finiment engendrées quelconques de
$\mathcal{I}$, la catégorie qu'elles engendrent est elle aussi finiment
engendrée. D'après l'hypothèse faite sur $\mathcal{I}$, elle est donc contenue
dans une sous-catégorie finiment engendrée de $\mathcal{I}$ ayant un unique
objet final. Cela montre que $\mathcal{I}_0$ est filtrante.

\medskip\noindent
Vérifions enfin que $M_0$ est cofinal. Tout objet de $\mathcal{I}$ étant
l'objet final de la sous-catégorie formée de ce seul objet et de sa flèche
identité, le foncteur $M_0$ est surjectif sur les objets. La condition (1) de
la définition \ref{definition-cofinal} est donc satisfaite. Étant donnés un
objet $i$ de $\mathcal{I}$, des objets $\mathcal{F}_1, \mathcal{F}_2$ de
$\mathcal{I}_0$ et des morphismes
$\varphi_1 : i \to M_0(\mathcal{F}_1)$ et
$\varphi_2 : i \to M_0(\mathcal{F}_2)$ dans $\mathcal{I}$, on peut prendre
pour $\mathcal{F}_{12}$ une catégorie finiment engendrée ayant un unique
objet final et contenant $\mathcal{F}_1$, $\mathcal{F}_2$ ainsi que les
morphismes $\varphi_1, \varphi_2$. Le diagramme obtenu est commutatif :
$$
\xymatrix{
& M_0(\mathcal{F}_{12}) & \\
M_0(\mathcal{F}_{1}) \ar[ru] & & M_0(\mathcal{F}_{2}) \ar[lu] \\
& i \ar[lu] \ar[ru]
}
$$
car il se trouve dans la catégorie $\mathcal{F}_{12}$ et
$M_0(\mathcal{F}_{12})$ est final dans cette catégorie. La condition (2) est
donc elle aussi satisfaite, ce qui achève la preuve.
\end{proof}

\begin{remark}
\label{remark-trick-needed}
Remarquons qu'un ensemble ordonné filtrant fini $(I, \leq)$ possède toujours
un plus grand élément $i_\infty$. Toute limite inductive d'un système
$(M_i, f_{ii'})$ sur un tel ensemble est donc triviale, au sens où elle est
égale à $M_{i_\infty}$. En revanche, une limite inductive indexée par une
catégorie filtrante finie n'est pas nécessairement triviale. Par exemple,
soit $\mathcal{I}$ la catégorie ayant un unique objet $i$ et un unique
morphisme non trivial $e$ satisfaisant $e = e \circ e$. La limite inductive
d'un diagramme $M : \mathcal{I} \to Ensembles$ est l'image de l'idempotent $M(e)$.
Cela montre qu'un artifice tel que le passage à
$\mathcal{I}\times \omega$ dans la preuve du
lemme \ref{lemma-directed-category-system} est essentiel.
\end{remark}

\begin{lemma}
\label{lemma-nonempty-limit}
Si $S : \mathcal{I} \to \textit{Ensembles}$ est un diagramme cofiltrant
d'ensembles et si tous les $S_i$ sont finis et non vides, alors
$\lim_i S_i$ est non vide. Autrement dit, la limite projective d'un système
inverse filtrant d'ensembles finis non vides est non vide.
\end{lemma}

\begin{proof}
Les deux énoncés sont équivalents d'après le
lemme \ref{lemma-directed-category-system}. Soit $I$ un ensemble ordonné
filtrant et soit $(S_i)_{i \in I}$ un système inverse d'ensembles finis non
vides sur $I$. Appelons {\it sous-système} $T$ une famille
$T = (T_i)_{i \in I}$ de parties non vides $T_i \subset S_i$ telle que
$T_{i'}$ soit envoyé dans $T_i$ par le morphisme de transition
$S_{i'} \to S_i$ pour tout $i' \geq i$. Notons $\mathcal{T}$ l'ensemble des
sous-systèmes et ordonnons $\mathcal{T}$ par inclusion. Supposons que les
$T_\alpha$, $\alpha \in A$, forment une famille totalement ordonnée d'éléments
de $\mathcal{T}$. Écrivons
$T_\alpha = (T_{\alpha, i})_{i \in I}$. On peut alors trouver un minorant
$T = (T_i)_{i \in I}$ en posant
$T_i = \bigcap_{\alpha \in A} T_{\alpha, i}$ ; c'est manifestement une partie
finie non vide de $S_i$, puisque tous les $T_{\alpha, i}$ sont non vides et
que les $T_\alpha$ forment une famille totalement ordonnée. Le lemme de Zorn
montre donc que $\mathcal{T}$ possède des éléments minimaux.

\medskip\noindent
Analysons la forme d'un élément minimal $T \in \mathcal{T}$. Remarquons
d'abord que toutes les applications $T_{i'} \to T_i$ sont surjectives.
En effet, puisque $I$ est un ensemble ordonné filtrant et que $T_i$ est fini,
l'intersection $T'_i = \bigcap_{i' \geq i} \Im(T_{i'} \to T_i)$ est non vide.
Ainsi, $T' = (T'_i)$ est un sous-système contenu dans $T$ et, par minimalité,
$T' = T$. Enfin, affirmons que $T_i$ est un singleton pour tout $i$. En effet,
si $x \in T_i$, on peut définir
$T'_{i'} = (T_{i'} \to T_i)^{-1}(\{x\})$ pour $i' \geq i$ et
$T'_j = T_j$ si $j \not \geq i$. Il s'agit encore d'un sous-système, car on
a vu ci-dessus que les morphismes de transition du sous-système $T$ sont
surjectifs. Par minimalité, $T = T'$, ce qui implique bien que $T_i$ est un
singleton. Cela vaut pour chaque $i \in I$ ; on a donc
$T_i = \{x_i\}$ pour un certain $x_i \in S_i$, avec $x_{i'} \mapsto x_i$
par l'application $S_{i'} \to S_i$ pour tout $i' \geq i$. Autrement dit,
$(x_i) \in \lim S_i$, ce qui démontre le lemme.
\end{proof}






\section{Systèmes essentiellement constants}
\label{section-essentially-constant}

\noindent
Soit $M : \mathcal{I} \to \mathcal{C}$ un diagramme à valeurs dans une catégorie
$\mathcal{C}$. Supposons la catégorie d'indices $\mathcal{I}$ filtrante.
Il existe alors trois notions, de plus en plus fortes, qui distinguent un
objet $X$ de $\mathcal{C}$. La première est simplement
$$
X = \colim_{i \in \mathcal{I}} M_i.
$$
L'objet $X$ est alors muni des coprojections $M_i \to X$. Une condition plus
forte consisterait à exiger que $X$ soit la limite inductive et qu'il existe
$i \in \mathcal{I}$ et un morphisme $X \to M_i$ tels que la composée
$X \to M_i \to X$ soit $\text{id}_X$. Voici une condition encore plus forte.

\begin{definition}
\label{definition-essentially-constant-diagram}
Soit $M : \mathcal{I} \to \mathcal{C}$ un diagramme à valeurs dans une catégorie
$\mathcal{C}$.
\begin{enumerate}
\item Supposons la catégorie d'indices $\mathcal{I}$ filtrante et soit
$(X, \{M_i \to X\}_i)$ un cocône de $M$, voir la
remarque \ref{remark-cones-and-cocones}. On dit que $M$ est
{\it essentiellement constant} de {\it valeur} $X$ s'il existe un
$i \in \mathcal{I}$ et un morphisme $X \to M_i$ tels que
\begin{enumerate}
\item $X \to M_i \to X$ soit $\text{id}_X$ ;
\item pour tout $j$, il existe $k$ et des morphismes $i \to k$ et
$j \to k$ tels que le morphisme $M_j \to M_k$ soit égal à la composée
$M_j \to X \to M_i \to M_k$.
\end{enumerate}
\item Supposons la catégorie d'indices $\mathcal{I}$ cofiltrante et soit
$(X, \{X \to M_i\}_i)$ un cône de $M$, voir la
remarque \ref{remark-cones-and-cocones}. On dit que $M$ est
{\it essentiellement constant} de {\it valeur} $X$ s'il existe un
$i \in \mathcal{I}$ et un morphisme $M_i \to X$ tels que
\begin{enumerate}
\item $X \to M_i \to X$ soit $\text{id}_X$ ;
\item pour tout $j$, il existe $k$ et des morphismes $k \to i$ et
$k \to j$ tels que le morphisme $M_k \to M_j$ soit égal à la composée
$M_k \to M_i \to X \to M_j$.
\end{enumerate}
\end{enumerate}
Lorsqu'on utilise cette définition, il faut garder à l'esprit le
lemme \ref{lemma-essentially-constant-is-limit-colimit}.
\end{definition}

\noindent
Le contexte indiquera laquelle des deux versions est considérée. En cas
d'ambiguïté, on les distinguera en disant, dans le premier cas, que $M$ est
essentiellement constant en tant qu'{\it ind-objet} et, dans le second, qu'il
est essentiellement constant en tant que {\it pro-objet}. Cette terminologie
est expliquée plus avant dans les remarques \ref{remark-ind-category} et
\ref{remark-pro-category}. En fait, nous emploierons souvent l'expression
« système essentiellement constant », laquelle, à strictement parler, n'est
définie que pour les systèmes sur des ensembles ordonnés filtrants.

\begin{definition}
\label{definition-essentially-constant-system}
Soit $\mathcal{C}$ une catégorie. Un système filtrant
$(M_i, f_{ii'})$ est un {\it système essentiellement constant} si $M$,
considéré comme un foncteur $I \to \mathcal{C}$, définit un diagramme
essentiellement constant. Un système inverse filtrant $(M_i, f_{ii'})$ est
un {\it système inverse essentiellement constant} si $M$, considéré comme un
foncteur $I^{opp} \to \mathcal{C}$, définit un diagramme inverse
essentiellement constant.
\end{definition}

\noindent
Si $(M_i, f_{ii'})$ est un système essentiellement constant et si les
morphismes $f_{ii'}$ sont des monomorphismes, alors, pour tous
$i \leq i'$ suffisamment grands, les morphismes $f_{ii'}$ sont des
isomorphismes. Considérons en
revanche le système
$$
\mathbf{Z}^2 \to \mathbf{Z}^2 \to \mathbf{Z}^2 \to \ldots
$$
dont les applications sont données par $(a, b) \mapsto (a + b, 0)$. Ce
système est essentiellement constant de valeur $\mathbf{Z}$, mais tout
morphisme de transition possède un noyau.

\medskip\noindent
Voici un exemple de système qui n'est pas essentiellement constant. Posons
$M = \bigoplus_{n \geq 0} \mathbf{Z}$ et soit $S : M \to M$ l'opérateur de
décalage $(a_0, a_1, \ldots) \mapsto (a_1, a_2, \ldots)$. Alors le système
$M \to M \to M \to \ldots$ dont les morphismes de transition sont $S$ a
pour limite inductive $0$, et la composée $0 \to M \to 0$ est l'identité,
mais ce système n'est pas essentiellement constant.

\medskip\noindent
Le lemme suivant est un contrôle de cohérence élémentaire.

\begin{lemma}
\label{lemma-essentially-constant-is-limit-colimit}
Soit $M : \mathcal{I} \to \mathcal{C}$ un diagramme.
Si $\mathcal{I}$ est filtrante et si $M$ est essentiellement constant en tant
qu'ind-objet, alors $X = \colim M_i$ existe et $M$ est essentiellement
constant de valeur $X$.
Si $\mathcal{I}$ est cofiltrante et si $M$ est essentiellement constant en
tant que pro-objet, alors $X = \lim M_i$ existe et $M$ est essentiellement
constant de valeur $X$.
\end{lemma}

\begin{proof}
Omis. C'est un bon exercice sur les définitions.
\end{proof}

\begin{remark}
\label{remark-ind-category}
Soit $\mathcal{C}$ une catégorie. Il existe une grande catégorie
$\text{Ind-}\mathcal{C}$ des {\it ind-objets} de $\mathcal{C}$. En effet, si
$F : \mathcal{I} \to \mathcal{C}$ et $G : \mathcal{J} \to \mathcal{C}$
sont des diagrammes filtrants à valeurs dans $\mathcal{C}$, on peut définir
$$
\Mor_{\text{Ind-}\mathcal{C}}(F, G) =
\lim_i \colim_j \Mor_\mathcal{C}(F(i), G(j)).
$$
Il existe un foncteur canonique $\mathcal{C} \to \text{Ind-}\mathcal{C}$ qui
envoie $X$ sur le {\it système constant} de valeur $X$. C'est un plongement
pleinement fidèle. Dans ce langage, on voit qu'un diagramme $F$ est
essentiellement constant si et seulement si $F$ est isomorphe à un système
constant. Si nous avons un jour besoin de ce résultat, nous l'énoncerons sous
forme de lemme et le démontrerons ici.
\end{remark}

\begin{remark}
\label{remark-pro-category}
Soit $\mathcal{C}$ une catégorie. Il existe une grande catégorie
$\text{Pro-}\mathcal{C}$ des {\it pro-objets} de $\mathcal{C}$. En effet, si
$F : \mathcal{I} \to \mathcal{C}$ et $G : \mathcal{J} \to \mathcal{C}$
sont des diagrammes cofiltrants à valeurs dans $\mathcal{C}$, on peut définir
$$
\Mor_{\text{Pro-}\mathcal{C}}(F, G) =
\lim_j \colim_i \Mor_\mathcal{C}(F(i), G(j)).
$$
Il existe un foncteur canonique $\mathcal{C} \to \text{Pro-}\mathcal{C}$ qui
envoie $X$ sur le {\it système constant} de valeur $X$. C'est un plongement
pleinement fidèle. Dans ce langage, on voit qu'un diagramme $F$ est
essentiellement constant si et seulement si $F$ est isomorphe à un système
constant. Si nous avons un jour besoin de ce résultat, nous l'énoncerons sous
forme de lemme et le démontrerons ici.
\end{remark}

\begin{example}
\label{example-pro-morphism-inverse-systems}
Soit $\mathcal{C}$ une catégorie. Soient $(X_n)$ et $(Y_n)$ des systèmes
inverses dans $\mathcal{C}$ sur $\mathbf{N}$ muni de l'ordre usuel.
Schématiquement :
$$
\ldots \to X_3 \to X_2 \to X_1
\quad\text{et}\quad
\ldots \to Y_3 \to Y_2 \to Y_1
$$
Soit $a : (X_n) \to (Y_n)$ un morphisme de pro-objets de $\mathcal{C}$.
Quelles données représente $a$ ? Pour chaque $n \in \mathbf{N}$, il devrait
exister un $m(n)$ et un morphisme $a_n : X_{m(n)} \to Y_n$. Ces morphismes
devraient être compatibles au sens suivant : pour tout $n' \geq n$, il existe
un $m(n', n) \geq m(n'), m(n)$ tel que le diagramme
$$
\xymatrix{
X_{m(n', n)} \ar[rr] \ar[d] & & X_{m(n)} \ar[d]^{a_n} \\
X_{m(n')} \ar[r]^{a_{n'}} & Y_{n'} \ar[r] & Y_n
}
$$
soit commutatif. Après avoir remplacé $m(n)$ par
$\max_{l \leq k \leq n}\{m(n, k), m(k, l)\}$, on voit que l'on obtient
$\ldots \geq m(3) \geq m(2) \geq m(1)$ et un diagramme commutatif
$$
\xymatrix{
\ldots \ar[r] &
X_{m(3)} \ar[d]^{a_3} \ar[r] &
X_{m(2)} \ar[d]^{a_2} \ar[r] &
X_{m(1)} \ar[d]^{a_1} \\
\ldots \ar[r] &
Y_3 \ar[r] &
Y_2 \ar[r] &
Y_1
}
$$
Étant donnée une application croissante
$m' : \mathbf{N} \to \mathbf{N}$ telle que $m' \geq m$, et en posant
$a'_i : X_{m'(i)} \to X_{m(i)} \to Y_i$, la paire $(m', a')$ définit le même
morphisme de pro-systèmes. Réciproquement, deux paires $(m_1, a_1)$ et
$(m_2, a_2)$ comme ci-dessus définissent le même morphisme de pro-objets si et
seulement s'il existe $m' \geq m_1, m_2$ tel que $a'_1 = a'_2$.
\end{example}

\begin{remark}
\label{remark-pro-category-copresheaves}
Soit $\mathcal{C}$ une catégorie. Soient
$F : \mathcal{I} \to \mathcal{C}$ et $G : \mathcal{J} \to \mathcal{C}$
des diagrammes cofiltrants à valeurs dans $\mathcal{C}$. Considérons les foncteurs
$A, B : \mathcal{C} \to \textit{Ensembles}$ définis par
$$
A(X) = \colim_i \Mor_\mathcal{C}(F(i), X)
\quad\text{et}\quad
B(X) = \colim_j \Mor_\mathcal{C}(G(j), X)
$$
Nous affirmons qu'un morphisme de pro-systèmes de $F$ vers $G$ équivaut à une
transformation de foncteurs $t : B \to A$. En effet, étant donné $t$, on peut
appliquer $t$ à la classe de $\text{id}_{G(j)}$ dans $B(G(j))$ et obtenir un
système compatible d'éléments
$\xi_j \in A(G(j)) = \colim_i \Mor_\mathcal{C}(F(i), G(j))$ ; c'est exactement
notre définition d'un morphisme dans $\text{Pro-}\mathcal{C}$ donnée dans la
remarque \ref{remark-pro-category}. Nous omettons la construction d'une
transformation $B \to A$ à partir d'un morphisme de pro-objets de $F$ vers $G$.
\end{remark}

\begin{lemma}
\label{lemma-image-essentially-constant}
Soit $\mathcal{C}$ une catégorie. Soit $M : \mathcal{I} \to \mathcal{C}$
un diagramme dont la catégorie d'indices $\mathcal{I}$ est filtrante
(resp.\ cofiltrante). Soit $F : \mathcal{C} \to \mathcal{D}$ un foncteur.
Si $M$ est essentiellement constant en tant qu'ind-objet (resp.\ pro-objet),
alors il en va de même de $F \circ M : \mathcal{I} \to \mathcal{D}$.
\end{lemma}

\begin{proof}
Si $X$ est une valeur de $M$, il découle immédiatement de la définition que
$F(X)$ est une valeur de $F \circ M$.
\end{proof}

\begin{lemma}
\label{lemma-characterize-essentially-constant-ind}
Soit $\mathcal{C}$ une catégorie. Soit $M : \mathcal{I} \to \mathcal{C}$
un diagramme dont la catégorie d'indices $\mathcal{I}$ est filtrante.
Les assertions suivantes sont équivalentes :
\begin{enumerate}
\item $M$ est un ind-objet essentiellement constant ;
\item il existe un cocône $(X, \{M_i \to X\}_i)$ tel que, pour tout $W$ dans
$\mathcal{C}$, l'application
$\colim_i \Mor_\mathcal{C}(W, M_i) \to \Mor_\mathcal{C}(W, X)$
soit bijective ;
\item $X = \colim_i M_i$ existe et, pour tout $W$ dans $\mathcal{C}$,
l'application
$\colim_i \Mor_\mathcal{C}(W, M_i) \to \Mor_\mathcal{C}(W, X)$
est bijective ;
\item il existe un objet $X$ de $\mathcal{C}$, un $i$ dans $\mathcal{I}$ et un
morphisme $X \to M_i$ tels que, pour tout $W$ dans $\mathcal{C}$, l'application
$\Mor_\mathcal{C}(W, X) \to
\colim_{j \in \mathcal{I}} \Mor_\mathcal{C}(W, M_j)$
soit bijective.
\end{enumerate}
Dans les cas (2), (3) et (4), la valeur du système essentiellement constant
est $X$.
\end{lemma}

\begin{proof}
Il est clair que (3) implique (2). Supposons (2). Alors
$\text{id}_X \in \Mor_\mathcal{C}(X, X)$ provient d'un morphisme
$X \to M_i$ pour un certain $i$, c'est-à-dire que $X \to M_i \to X$ est
l'identité. Les deux applications
$$
\Mor_\mathcal{C}(W, X)
\longrightarrow
\colim_i \Mor_\mathcal{C}(W, M_i)
\longrightarrow
\Mor_\mathcal{C}(W, X)
$$
sont alors bijectives pour tout $W$ ; la première est induite par le morphisme
$X \to M_i$ trouvé ci-dessus et leur composée est l'identité. Cela signifie
que la composée
$$
\colim_i \Mor_\mathcal{C}(W, M_i)
\longrightarrow
\Mor_\mathcal{C}(W, X)
\longrightarrow
\colim_i \Mor_\mathcal{C}(W, M_i)
$$
est elle aussi l'identité. En prenant $W = M_j$ et en partant de
$\text{id}_{M_j}$ dans la limite inductive, on voit que
$M_j \to X \to M_i \to M_k$ est égal à $M_j \to M_k$ pour un certain $k$
suffisamment loin dans le système. Cela démontre (1).

\medskip\noindent
Supposons (4). Soit $k$ un objet de $\mathcal{I}$. En prenant $W = M_k$,
on déduit qu'il existe un unique morphisme $M_k \to X$ tel qu'il existe un
$j$ et des morphismes $k \to j$ et $i \to j$ dans $\mathcal{I}$ pour lesquels
$M_k \to X \to M_i \to M_j$ est égal à $M_k \to M_j$. L'unicité garantit
que l'on obtient ainsi un cocône $(X, \{M_k \to X\})$. On voit de cette
manière que (4) implique (2) ; certains détails sont omis.

\medskip\noindent
Nous omettons la preuve que (1) implique les autres assertions.
\end{proof}

\begin{lemma}
\label{lemma-characterize-essentially-constant-pro}
Soit $\mathcal{C}$ une catégorie. Soit $M : \mathcal{I} \to \mathcal{C}$
un diagramme dont la catégorie d'indices $\mathcal{I}$ est cofiltrante.
Les assertions suivantes sont équivalentes :
\begin{enumerate}
\item $M$ est un pro-objet essentiellement constant ;
\item il existe un cône $(X, \{X \to M_i\})$ tel que, pour tout $W$ dans
$\mathcal{C}$, l'application
$\colim_{i \in \mathcal{I}^{opp}} \Mor_\mathcal{C}(M_i, W) \to
\Mor_\mathcal{C}(X, W)$ soit bijective ;
\item $X = \lim_i M_i$ existe et, pour tout $W$ dans $\mathcal{C}$,
l'application
$\colim_{i \in \mathcal{I}^{opp}} \Mor_\mathcal{C}(M_i, W) \to
\Mor_\mathcal{C}(X, W)$ est bijective ;
\item il existe un objet $X$ de $\mathcal{C}$, un $i$ dans $\mathcal{I}$ et un
morphisme $M_i \to X$ tels que, pour tout $W$ dans $\mathcal{C}$, l'application
$\Mor_\mathcal{C}(X, W) \to
\colim_{j \in \mathcal{I}^{opp}} \Mor_\mathcal{C}(M_j, W)$ soit bijective.
\end{enumerate}
Dans les cas (2), (3) et (4), la valeur du système essentiellement constant
est $X$.
\end{lemma}

\begin{proof}
Ce lemme est dual du lemme \ref{lemma-characterize-essentially-constant-ind}.
\end{proof}

\begin{lemma}
\label{lemma-cofinal-essentially-constant}
Soit $\mathcal{C}$ une catégorie. Soit $H : \mathcal{I} \to \mathcal{J}$ un
foncteur entre catégories d'indices filtrantes. Si $H$ est cofinal, alors tout
diagramme $M : \mathcal{J} \to \mathcal{C}$ est essentiellement constant si
et seulement si $M \circ H$ est essentiellement constant.
\end{lemma}

\begin{proof}
Cela résulte formellement des
lemmes \ref{lemma-characterize-essentially-constant-ind} et
\ref{lemma-cofinal}.
\end{proof}

\begin{lemma}
\label{lemma-essentially-constant-over-product}
Soient $\mathcal{I}$ et $\mathcal{J}$ des catégories filtrantes et notons
$p : \mathcal{I} \times \mathcal{J} \to \mathcal{J}$ la projection. Alors
$\mathcal{I} \times \mathcal{J}$ est filtrante, et un diagramme
$M : \mathcal{J} \to \mathcal{C}$ est essentiellement constant si et
seulement si $M \circ p : \mathcal{I} \times \mathcal{J} \to \mathcal{C}$
est essentiellement constant.
\end{lemma}

\begin{proof}
Nous omettons la vérification que $\mathcal{I} \times \mathcal{J}$ est
filtrante. L'équivalence résulte du
lemme \ref{lemma-cofinal-essentially-constant}, car $p$ est cofinal
(vérification omise).
\end{proof}

\begin{lemma}
\label{lemma-initial-essentially-constant}
Soit $\mathcal{C}$ une catégorie. Soit $H : \mathcal{I} \to \mathcal{J}$ un
foncteur entre catégories d'indices cofiltrantes. Si $H$ est initial, alors
tout diagramme $M : \mathcal{J} \to \mathcal{C}$ est essentiellement constant
si et seulement si $M \circ H$ est essentiellement constant.
\end{lemma}

\begin{proof}
Cela résulte formellement des
lemmes \ref{lemma-characterize-essentially-constant-pro},
\ref{lemma-initial}, \ref{lemma-cofinal}, et du fait que, si $\mathcal{I}$ est
initiale dans $\mathcal{J}$, alors $\mathcal{I}^{opp}$ est cofinale dans
$\mathcal{J}^{opp}$.
\end{proof}
\section{Foncteurs exacts}
\label{section-exact-functor}

\noindent
Dans cette section, nous définissons les foncteurs exacts.

\begin{definition}
\label{definition-exact}
Soit $F : \mathcal{A} \to \mathcal{B}$ un foncteur.
\begin{enumerate}
\item Supposons que toutes les limites finies existent dans $\mathcal{A}$.
Nous disons que $F$ est {\it exact à gauche} s'il commute avec toutes
les limites finies.
\item Supposons que toutes les colimites finies existent dans $\mathcal{A}$.
Nous disons que $F$ est {\it exact à droite} s'il commute avec toutes les
colimites finies.
\item Nous disons que $F$ est {\it exact} s'il est exact à la fois à gauche et
à droite.
\end{enumerate}
\end{definition}

\begin{lemma}
\label{lemma-characterize-left-exact}
Soit $F : \mathcal{A} \to \mathcal{B}$ un foncteur. Supposons que toutes les
limites finies existent dans $\mathcal{A}$, voir le lemme
\ref{lemma-finite-limits-exist}. Les assertions suivantes sont équivalentes :
\begin{enumerate}
\item $F$ est exact à gauche,
\item $F$ commute avec des produits finis et des égalisateurs, et
\item $F$ transforme un objet final de $\mathcal{A}$ en un objet final de
$\mathcal{B}$ et commute avec des produits fibrés.
\end{enumerate}
\end{lemma}

\begin{proof}
Le lemme \ref{lemma-limits-products-equalizers} montre que (2) implique (1).
Supposons que (3) soit vrai. Le produit fibré sur l'objet final est le
produit. Si $a, b : A \to B$ sont des morphismes de $\mathcal{A}$, alors
l'égalisateur de $a, b$ est
$$
(A \times_{a, B, b} A)\times_{(\text{pr}_1, \text{pr}_2), A \times A, \Delta} A.
$$
Ainsi (3) implique (2). Enfin (1) implique (3) car la limite vide est un objet
final et les produits fibrés sont des limites.
\end{proof}

\begin{lemma}
\label{lemma-characterize-right-exact}
Soit $F : \mathcal{A} \to \mathcal{B}$ un foncteur. Supposons que toutes les
colimites finies existent dans $\mathcal{A}$, voir le lemme
\ref{lemma-colimits-exist}. Les assertions suivantes sont équivalentes :
\begin{enumerate}
\item $F$ est exact à droite,
\item $F$ commute avec des coproduits finis et des coégalisateurs, et
\item $F$ transforme un objet initial de $\mathcal{A}$ en un objet initial de
$\mathcal{B}$ et commute avec les sommes amalgamées.
\end{enumerate}
\end{lemma}

\begin{proof}
Dual du lemme \ref{lemma-characterize-left-exact}.
\end{proof}




\section{Foncteurs adjoints}
\label{section-adjoint}

\begin{definition}
\label{definition-adjoint}
Soient $\mathcal{C}$, $\mathcal{D}$ des catégories. Soient $u : \mathcal{C} \to \mathcal{D}$ et $v : \mathcal{D} \to \mathcal{C}$ des foncteurs. On dit que
$u$ est un {\it adjoint à gauche} de $v$, ou que $v$ est un {\it adjoint à droite}
de $u$ s'il existe des bijections
$$
\Mor_\mathcal{D}(u(X), Y)
\longrightarrow
\Mor_\mathcal{C}(X, v(Y))
$$
fonctorielles en $X \in \Ob(\mathcal{C})$ et $Y \in \Ob(\mathcal{D})$.
\end{definition}

\noindent
En d'autres termes, cela signifie que l'on s'est {\it donné} un isomorphisme entre
les foncteurs $\mathcal{C}^{opp} \times \mathcal{D} \to \textit{Ensembles}$, allant de
$\Mor_\mathcal{D}(u(-), -)$ à $\Mor_\mathcal{C}(-, v(-))$. Pour tout objet $X$
de $\mathcal{C}$ on obtient un morphisme $X \to v(u(X))$ correspondant à
$\text{id}_{u(X)}$. De même, pour tout objet $Y$ de $\mathcal{D}$ on obtient
un morphisme $u(v(Y)) \to Y$ correspondant à $\text{id}_{v(Y)}$. Ces
morphismes sont appelés {\it morphismes d'adjonction}. Les morphismes
d'adjonction sont fonctoriels en $X$ et $Y$ ; on obtient donc des morphismes
de foncteurs
$$
\eta : \text{id}_\mathcal{C} \to v \circ u\quad (\text{unité})
\quad\text{et}\quad
\epsilon : u \circ v \to \text{id}_\mathcal{D}\quad (\text{counité}).
$$
De plus, si $\alpha : u(X) \to Y$ et $\beta : X \to v(Y)$ sont des morphismes,
alors les conditions suivantes sont équivalentes
\begin{enumerate}
\item $\alpha$ et $\beta$ se correspondent via la bijection de la définition,
\item $\beta$ est la composition $X \to v(u(X)) \xrightarrow{v(\alpha)} v(Y)$,
et
\item $\alpha$ est la composition $u(X) \xrightarrow{u(\beta)} u(v(Y)) \to Y$.
\end{enumerate}
On peut ainsi reformuler la notion de foncteurs adjoints en termes de
morphismes d'adjonction.

\begin{lemma}
\label{lemma-adjoint-exists}
Soit $u : \mathcal{C} \to \mathcal{D}$ un foncteur entre catégories. Si pour
chaque $y \in \Ob(\mathcal{D})$ le foncteur $x \mapsto \Mor_\mathcal{D}(u(x), y)$ est représentable, alors $u$ a un adjoint à droite.
\end{lemma}

\begin{proof}
Pour chaque $y$, choisissons un objet $v(y)$ et un isomorphisme
$\Mor_\mathcal{C}(-, v(y)) \to \Mor_\mathcal{D}(u(-), y)$ de foncteurs. Par le
lemme de Yoneda (lemme \ref{lemma-yoneda}) pour tout morphisme $g : y \to y'$
la transformation de foncteurs
$$
\Mor_\mathcal{C}(-, v(y)) \to \Mor_\mathcal{D}(u(-), y) \to
\Mor_\mathcal{D}(u(-), y') \to \Mor_\mathcal{C}(-, v(y'))
$$
correspond à un morphisme unique $v(g) : v(y) \to v(y')$. Nous omettons de
vérifier que $v$ est un foncteur et qu'il est adjoint à droite de $u$.
\end{proof}

\begin{lemma}
\label{lemma-left-adjoint-composed-fully-faithful}
\begin{reference}
Bhargav Bhatt, communication privée.
\end{reference}
Soit $u$ un adjoint à gauche de $v$ comme dans la définition
\ref{definition-adjoint}.
\begin{enumerate}
\item Si $v \circ u$ est pleinement fidèle, alors $u$ est pleinement fidèle.
\item Si $u \circ v$ est pleinement fidèle, alors $v$ est pleinement fidèle.
\end{enumerate}
\end{lemma}

\begin{proof}
Preuve de (2). Supposons que $u \circ v$ soit pleinement fidèle. Soient
$X$ et $Y$ des objets de $\mathcal{D}$. Alors l'application composite
naturelle
$$
\Mor(X,Y) \to \Mor(v(X),v(Y)) \to \Mor(u(v(X)), u(v(Y)))
$$
est une bijection, donc $v$ est au moins fidèle. Pour montrer qu'il est
pleinement fidèle, il faut montrer que la deuxième application ci-dessus est injective.
Mais l'adjonction entre $u$ et $v$ dit que
$$
\Mor(v(X), v(Y)) \to \Mor(u(v(X)), u(v(Y))) \to \Mor(u(v(X)), Y)
$$
est une bijection, où la première application est naturelle et la deuxième
est induite par la counité $u(v(Y)) \to Y$ de l'adjonction. Cela signifie
donc que $\Mor(v(X), v(Y)) \to \Mor(u(v(X)), u(v(Y)))$ est également injectif,
comme souhaité. La preuve de (1) est duale.
\end{proof}

\begin{lemma}
\label{lemma-adjoint-fully-faithful}
Soit $u$ un adjoint à gauche de $v$ comme dans la définition
\ref{definition-adjoint}. Alors
\begin{enumerate}
\item $u$ est pleinement fidèle $\Leftrightarrow$ $\text{id} \cong v \circ u$
$\Leftrightarrow$ $\eta : \text{id} \to v \circ u$ est un isomorphisme,
\item $v$ est pleinement fidèle $\Leftrightarrow$ $u \circ v \cong \text{id}$
$\Leftrightarrow$ $\epsilon : u \circ v \to \text{id}$ est un isomorphisme.
\end{enumerate}
\end{lemma}

\begin{proof}
Preuve de (1). Supposons que $u$ soit pleinement fidèle. Nous allons montrer
que $\eta_X : X \to v(u(X))$ est un isomorphisme. Considérons un morphisme
quelconque $X' \to v(u(X))$. Par adjonction cela correspond à un morphisme $u(X') \to u(X)$. Par la pleine fidélité de $u$, cela correspond à un morphisme unique $X' \to X$. On voit ainsi que la post-composition par $\eta_X$ définit une
bijection $\Mor(X', X) \to \Mor(X', v(u(X)))$. Donc $\eta_X$ est un
isomorphisme. S'il existe un isomorphisme $\text{id} \cong v \circ u$ de
foncteurs, alors $v \circ u$ est pleinement fidèle. Par le lemme
\ref{lemma-left-adjoint-composed-fully-faithful} nous voyons que $u$ est
pleinement fidèle. D'après ce qui précède, cela implique que $\eta$ est un
isomorphisme. Ainsi, les $3$ conditions sont équivalentes (et ces
conditions sont également équivalentes au fait que $v \circ u$ soit
pleinement fidèle).

\medskip\noindent
La partie (2) est duale de la partie (1).
\end{proof}

\begin{lemma}
\label{lemma-adjoint-exact}
Soit $u$ un adjoint à gauche de $v$ comme dans la définition
\ref{definition-adjoint}.
\begin{enumerate}
\item Supposons que $M : \mathcal{I} \to \mathcal{C}$ soit un diagramme et
supposons que $\colim_\mathcal{I} M$ existe dans $\mathcal{C}$. Alors
$u(\colim_\mathcal{I} M) =
\colim_\mathcal{I} u \circ M$. En d'autres termes,
$u$ commute avec les colimites existantes.
\item Supposons que $M : \mathcal{I} \to \mathcal{D}$ soit un diagramme et
supposons que $\lim_\mathcal{I} M$ existe dans $\mathcal{D}$. Alors
$v(\lim_\mathcal{I} M) =
\lim_\mathcal{I} v \circ M$. En d'autres termes, $v$
commute avec les limites existantes.
\end{enumerate}
\end{lemma}

\begin{proof}
La donnée d'un morphisme d'une colimite vers un objet équivaut à celle d'un
système compatible de morphismes des constituants de la colimite vers l'objet, voir
remarque \ref{remark-limit-colim}. Alors
$$
\begin{matrix}
\Mor_\mathcal{D}(u(\colim_{i \in \mathcal{I}} M_i), Y) &
= & \Mor_\mathcal{C}(\colim_{i \in \mathcal{I}} M_i, v(Y)) \\
& = &
\lim_{i \in \mathcal{I}^{opp}}
\Mor_\mathcal{C}(M_i, v(Y)) \\
& = &
\lim_{i \in \mathcal{I}^{opp}}
\Mor_\mathcal{D}(u(M_i), Y)
\end{matrix}
$$
prouve que $u(\colim_{i \in \mathcal{I}} M_i)$ est la colimite que nous
recherchons. Un argument analogue s'applique au second énoncé.
\end{proof}

\begin{lemma}
\label{lemma-exact-adjoint}
Soit $u$ un adjoint à gauche de $v$ comme dans la définition
\ref{definition-adjoint}.
\begin{enumerate}
\item Si $\mathcal{C}$ a des colimites finies, alors $u$ est exact à droite.
\item Si $\mathcal{D}$ a des limites finies, alors $v$ est exact à gauche.
\end{enumerate}
\end{lemma}

\begin{proof}
Évident d’après les définitions et le lemme \ref{lemma-adjoint-exact}.
\end{proof}

\begin{lemma}
\label{lemma-unit-counit-relations}
Soit $u : \mathcal{C} \to \mathcal{D}$ un adjoint à gauche du foncteur $v : \mathcal{D} \to \mathcal{C}$. Soit $\eta_X : X \to v(u(X))$ l'unité et
$\epsilon_Y : u(v(Y)) \to Y$ la counité. Alors
$$
u(X) \xrightarrow{u(\eta_X)} u(v(u(X)))
\xrightarrow{\epsilon_{u(X)}} u(X)
\quad\text{et}\quad
v(Y) \xrightarrow{\eta_{v(Y)}} v(u(v(Y))) \xrightarrow{v(\epsilon_Y)}
v(Y)
$$
sont les identités.
\end{lemma}

\begin{proof}
Omis.
\end{proof}



\begin{lemma}
\label{lemma-transformation-between-functors-and-adjoints}
Soient $u_1, u_2 : \mathcal{C} \to \mathcal{D}$ des foncteurs avec adjoints
droits $v_1, v_2 : \mathcal{D} \to \mathcal{C}$. Soit $\beta : u_2 \to u_1$
une transformation de foncteurs. Soit $\beta^\vee : v_1 \to v_2$ la
transformation correspondante des foncteurs adjoints. Alors
$$
\xymatrix{
u_2 \circ v_1 \ar[r]_\beta \ar[d]_{\beta^\vee} &
u_1 \circ v_1 \ar[d] \\
u_2 \circ v_2 \ar[r] & \text{id}
}
$$
est commutatif, où les flèches sans étiquette sont les transformations de
counité.
\end{lemma}

\begin{proof}
Cela est vrai car $\beta^\vee_D : v_1D \to v_2D$ est le morphisme unique tel
que l'application induite $\Mor(C, v_1D) \to \Mor(C, v_2D)$ soit
l'application $\Mor(u_1C, D) \to \Mor(u_2C, D)$ induite par $\beta_C : u_2C \to u_1C$. Autrement dit, cela signifie que l'application
$$
\Mor(u_1 v_1 D, D') \to \Mor(u_2 v_1 D, D')
$$
induite par $\beta_{v_1 D}$ est la même que l'application
$$
\Mor(v_1 D, v_1 D') \to \Mor(v_1 D, v_2 D')
$$
induite par $\beta^\vee_{D'}$. En prenant $D' = D$, nous constatons que la counité
$u_1 v_1 D \to D$ précomposée par $\beta_{v_1D}$ correspond à $\beta^\vee_D$
sous adjonction. Cela signifie exactement que le diagramme est commutatif
lorsqu'il est évalué sur $D$.
\end{proof}

\begin{lemma}
\label{lemma-compose-counits}
Soient $\mathcal{A}$, $\mathcal{B}$ et $\mathcal{C}$ des catégories. Soient $v : \mathcal{A} \to \mathcal{B}$ et $v' : \mathcal{B} \to \mathcal{C}$ des
foncteurs ayant pour adjoints à gauche $u$ et $u'$ respectivement. Alors
\begin{enumerate}
\item Le foncteur $v'' = v' \circ v$ a un adjoint à gauche égal à $u'' = u \circ u'$.
\item Étant donné $X$ dans $\mathcal{A}$, nous avons
\begin{equation}
\label{equation-compose-counits}
\epsilon_X^v \circ u(\epsilon^{v'}_{v(X)}) = \epsilon^{v''}_X :
u''(v''(X)) \to X
\end{equation}
où $\epsilon$ est la counité des adjonctions.
\end{enumerate}
\end{lemma}

\begin{proof}
Déroulons la formule de (2) car cela prouvera aussi immédiatement (1).
Premièrement, les counités des adjonctions correspondant aux paires $(u, v)$ et $(u', v')$
sont les morphismes $\epsilon_X^v : u(v(X)) \to X$ et $\epsilon_Y^{v'} : u'(v'(Y)) \to Y$, voir la discussion qui suit la définition
\ref{definition-adjoint}. Avec $u''$ et $v''$ comme en (1) on déroule tout
$$
u''(v''(X)) = u(u'(v'(v(X)))) \xrightarrow{u(\epsilon_{v(X)}^{v'})}
u(v(X)) \xrightarrow{\epsilon_X^v} X
$$
pour obtenir le morphisme du membre de gauche de
(\ref{equation-compose-counits}). Notons-le $\epsilon_X^{v''}$ pour
l'instant. Pour voir qu'il s'agit de la counité d'une paire de foncteurs adjoints $(u'', v'')$
il faut montrer qu'étant donné $Z$ dans $\mathcal{C}$ la règle qui envoie un
morphisme $\beta : Z \to v''(X)$ à $\alpha = \epsilon_X^{v''} \circ u''(\beta) : u''(Z) \to X$ est une bijection entre ensembles de morphismes. C'est vrai
car il s'agit de la composition de la règle envoyant $\beta$ à
$\epsilon_{v(X)}^{v'} \circ u'(\beta)$ qui est une bijection par hypothèse sur
$(u', v')$ puis envoyant celle-ci à $\epsilon_X^v \circ u(\epsilon_{v(X)}^{v'} \circ u'(\beta))$ qui est une bijection par hypothèse sur $(u, v)$.
\end{proof}






\section{Un critère de représentabilité}
\label{section-representable}

\noindent
Le lemme suivant est souvent utile pour prouver l'existence d'objets
universels dans de grandes catégories ; voir la discussion dans
la remarque \ref{remark-how-to-use-it}.

\begin{lemma}
\label{lemma-a-version-of-brown}
Soit $\mathcal{C}$ une grande\footnote{Voir remarque
\ref{remark-big-categories}.} catégorie qui a des limites. Soit $F : \mathcal{C} \to \textit{Ensembles}$ un foncteur. Supposons que
\begin{enumerate}
\item $F$ commute avec les limites,
\item il existe une famille $\{x_i\}_{i \in I}$ d'objets de $\mathcal{C}$ et
pour chaque $i \in I$ un élément $f_i \in F(x_i)$ tel que pour $y \in \Ob(\mathcal{C})$ et $g \in F(y)$ il existe un $i$ et un morphisme $\varphi : x_i \to y$ avec $F(\varphi)(f_i) = g$.
\end{enumerate}
Alors $F$ est représentable, c'est-à-dire qu'il existe un objet $x$ de
$\mathcal{C}$ tel que
$$
F(y) = \Mor_\mathcal{C}(x, y)
$$
fonctoriellement en $y$.
\end{lemma}

\begin{proof}
Soit $\mathcal{I}$ la catégorie dont les objets sont les couples $(x_i, f_i)$
et dont les morphismes $(x_i, f_i) \to (x_{i'}, f_{i'})$ sont des morphismes
$\varphi : x_i \to x_{i'}$ dans $\mathcal{C}$ tels que $F(\varphi)(f_i) = f_{i'}$. Posons
$$
x = \lim_{(x_i, f_i) \in \mathcal{I}} x_i
$$
(ce ne sera pas le $x$ que nous recherchons, voir ci-dessous). La limite
existe par hypothèse. Comme $F$ commute avec les limites, nous avons
$$
F(x) = \lim_{(x_i, f_i) \in \mathcal{I}} F(x_i).
$$
Il existe donc un élément universel $f \in F(x)$ dont l'image est
$f_i \in F(x_i)$ lorsque l'on applique $F$ au morphisme de projection
$x \to x_i$. En
utilisant $f$ on obtient une transformation de foncteurs
$$
\xi : \Mor_\mathcal{C}(x, - ) \longrightarrow F(-)
$$
voir la section \ref{section-opposite}. Soit $y$ un objet quelconque de
$\mathcal{C}$ et soit $g \in F(y)$. Choisissons $x_i \to y$ de telle sorte que
$f_i$ corresponde à $g$, ce qui est possible par hypothèse. Alors $F$ appliqué
aux morphismes
$$
x \longrightarrow x_i \longrightarrow y
$$
(le premier étant le morphisme de projection de la limite définissant $x$)
envoie $f$ sur $g$. La transformation $\xi$ est donc surjective.

\medskip\noindent
Afin de trouver l'objet représentant $F$, prenons pour $e : x' \to x$
l'égalisateur de tous les endomorphismes $\varphi : x \to x$ avec
$F(\varphi)(f) = f$. Puisque $F$ commute avec les limites, il commute avec les
égalisateurs, et nous voyons qu'il existe un élément $f' \in F(x')$ dont l'image
est $f$ dans $F(x)$. Puisque $\xi$ est surjective et que $f'$
correspond à $f$, nous voyons que $\xi' : \Mor_\mathcal{C}(x', -) \to F(-)$
est également surjective. Enfin, supposons que $a, b : x' \to y$ soient deux
morphismes tels que $F(a)(f') = F(b)(f')$. Nous devons montrer $a = b$.
Considérons l'égalisateur $e' : x'' \to x'$. De nouveau, nous trouvons un
$f'' \in F(x'')$ dont l'image est $f'$. Choisissons un morphisme $\psi : x \to x''$ tel que $F(\psi)(f) = f''$. On voit alors que $e \circ e' \circ \psi : x \to x$ est un morphisme avec $F(e \circ e' \circ \psi)(f) = f$. D'où $e \circ e' \circ \psi \circ e = e$. Puisque $e$ est un monomorphisme, cela
implique que $e'$ est un épimorphisme, donc $a = b$ comme souhaité.
\end{proof}

\begin{remark}
\label{remark-how-to-use-it}
Le lemme ci-dessus est souvent utilisé pour construire quelque chose de libre
sur quelque chose. Par exemple le groupe abélien libre sur un ensemble, le
groupe libre sur un ensemble, etc. L'idée, disons dans le cas du groupe libre
sur un ensemble $E$, est de considérer le foncteur
$$
F : \textit{Groups} \to \textit{Ensembles},\quad
G \longmapsto \text{Map}(E, G)
$$
Ce foncteur commute avec les limites. Comme famille d'objets, nous pouvons
prendre une famille de structures $E \to G_i$ constituées de groupes $G_i$
de cardinalité au plus $\max(\aleph_0, |E|)$ et d'applications ensemblistes
$E \to G_i$, de telle sorte que chaque classe d'isomorphisme d'une telle
structure soit représentée
au moins une fois. En effet, si $E \to G$ est une application de $E$ vers un
groupe $G$, alors le sous-groupe $G'$ généré par l'image a une cardinalité au
plus $\max(\aleph_0, |E|)$. Le lemme nous dit que le foncteur est
représentable, il existe donc un groupe $F_E$ tel que
$\Mor_{\textit{Groups}}(F_E, G) = \text{Map}(E, G)$. En particulier, le
morphisme identité de $F_E$ correspond à une application $E \to F_E$ et on
peut montrer que $F_E$ est généré par l'image sans imposer de relations.

\medskip\noindent
Une autre application typique est que nous pouvons utiliser le lemme pour
construire des colimites dès que l'on sait que des limites existent. Nous
l'illustrons à l'aide de la catégorie des espaces topologiques, qui admet des
limites d'après Topologie, Lemme \ref{topology-lemma-limits}. Supposons
que $\mathcal{I} \to \textit{Top}$, $i \mapsto X_i$ soit un foncteur. On peut
alors considérer
$$
F : \textit{Top} \longrightarrow \textit{Ensembles},\quad
Y \longmapsto \lim_\mathcal{I} \Mor_{\textit{Top}}(X_i, Y)
$$
Ce foncteur commute avec les limites. De plus, étant donné tout espace
topologique $Y$ et un élément $(\varphi_i : X_i \to Y)$ de $F(Y)$, il existe
un sous-espace $Y' \subset Y$ de cardinalité au plus $|\coprod X_i|$ tel que
les morphismes $\varphi_i$ se factorisent par $Y'$. En effet, nous pouvons prendre
la topologie induite sur l'union des images des $\varphi_i$. Ainsi il est clair
que les hypothèses du lemme sont satisfaites et on retrouve un espace
topologique $X$ représentant le foncteur $F$, ce qui signifie précisément que
$X$ est la colimite du diagramme $i \mapsto X_i$.
\end{remark}

\begin{theorem}[Théorème du foncteur adjoint]
\label{theorem-adjoint-functor}
Soit $G : \mathcal{C} \to \mathcal{D}$ un foncteur entre grandes catégories.
Supposons que $\mathcal{C}$ admette des limites, que $G$ commute avec elles et que,
pour chaque objet $y$ de $\mathcal{D}$, il existe un ensemble de paires $(x_i, f_i)_{i \in I}$ avec $x_i \in \Ob(\mathcal{C})$, $f_i \in \Mor_\mathcal{D}(y, G(x_i))$ tel que pour toute paire $(x, f)$ avec $x \in \Ob(\mathcal{C})$, $f \in \Mor_\mathcal{D}(y, G(x))$ il existe un $i$ et un morphisme $h : x_i \to x$ tel que $f = G(h) \circ f_i$. Alors $G$ a un adjoint à gauche $F$.
\end{theorem}

\begin{proof}
Les hypothèses impliquent que pour chaque objet $y$ de $\mathcal{D}$ le
foncteur $x \mapsto \Mor_\mathcal{D}(y, G(x))$ satisfait les hypothèses du
lemme \ref{lemma-a-version-of-brown}. Il est donc représentable par un objet,
appelons-le $F(y)$. Une application du lemme de Yoneda (Lemme
\ref{lemma-yoneda}) permet de faire de la règle $y \mapsto F(y)$ un foncteur qui
par construction est un adjoint à gauche de $G$. Nous omettons les détails.
\end{proof}





\section{Objets catégoriquement compacts}
\label{section-compact}

\noindent
Un peu sur les « petits » objets d'une catégorie.

\begin{definition}
\label{definition-compact-object}
Soit $\mathcal{C}$ une grande catégorie \footnote{Voir la remarque
\ref{remark-big-categories}.}. Un objet $X$ de $\mathcal{C}$ est dit
{\it catégoriquement compact} si on a
$$
\Mor_\mathcal{C}(X, \colim_i M_i) =
\colim_i \Mor_\mathcal{C}(X, M_i)
$$
pour chaque diagramme filtrant $M : \mathcal{I} \to \mathcal{C}$ tel que
$\colim_i M_i$ existe.
\end{definition}

\noindent
Souvent, cette définition n'est donnée que sous l'hypothèse que $\mathcal{C}$
admet toutes les limites inductives filtrantes.

\begin{lemma}
\label{lemma-extend-functor-by-colim}
Soient $\mathcal{C}$ et $\mathcal{D}$ deux grandes catégories admettant des limites
inductives filtrantes. Soit $\mathcal{C}' \subset \mathcal{C}$ une petite sous-catégorie
pleine composée d'objets catégoriquement compacts de $\mathcal{C}$, telle
que chaque objet de $\mathcal{C}$ soit une limite inductive filtrante d'objets
de $\mathcal{C}'$. Alors chaque foncteur $F' : \mathcal{C}' \to \mathcal{D}$ a une extension unique $F : \mathcal{C} \to \mathcal{D}$ qui commute
avec les limites inductives filtrantes.
\end{lemma}

\begin{proof}
Pour chaque objet $X$ de $\mathcal{C}$, nous pouvons écrire $X$ comme limite
inductive filtrante $X = \colim X_i$ avec $X_i \in \Ob(\mathcal{C}')$. Posons
$$
F(X) = \colim F'(X_i)
$$
dans $\mathcal{D}$. Nous montrerons ci-dessous que cette construction ne
dépend pas du choix d'une telle présentation de $X$.

\medskip\noindent
Soit $\alpha : X \to Y$ un morphisme de $\mathcal{C}$, et supposons que $X = \colim_{i \in I} X_i$ et $Y = \colim_{j \in J} Y_j$ soient écrits comme
limites inductives filtrantes d'objets de $\mathcal{C}'$. Pour chaque $i \in I$, puisque
$X_i$ est un objet catégoriquement compact de $\mathcal{C}$, nous pouvons
trouver un $j \in J$ et un diagramme commutatif
$$
\xymatrix{
X_i \ar[r] \ar[d] & X \ar[d]^\alpha \\
Y_j \ar[r] & Y
}
$$
On obtient alors un morphisme $F'(X_i) \to F'(Y_j) \to F(Y)$ où le deuxième
morphisme est la coprojection dans $F(Y) = \colim F'(Y_j)$. La flèche $\beta_i : F'(X_i) \to F(Y)$ ne dépend pas du choix de $j$. Pour $i \leq i'$, la
composition
$$
F'(X_i) \to F'(X_{i'}) \xrightarrow{\beta_{i'}} F(Y)
$$
est égale à $\beta_i$. On obtient ainsi une flèche bien définie
$$
F(\alpha) : F(X) = \colim F(X_i) \to F(Y)
$$
par la propriété universelle de la colimite. Si $\alpha' : Y \to Z$ est un
deuxième morphisme de $\mathcal{C}$ et que $Z = \colim Z_k$ s'écrit également
comme limite inductive filtrante d'objets de $\mathcal{C}'$, alors c'est un exercice
agréable de montrer que les morphismes induits $F(\alpha) : F(X) \to F(Y)$ et
$F(\alpha') : F(Y) \to F(Z)$ ont pour composé le morphisme $F(\alpha' \circ \alpha)$.
Détails omis.

\medskip\noindent
En particulier, si on nous donne deux présentations $X = \colim X_i$ et $X = \colim X'_{i'}$ comme limites inductives filtrantes de systèmes dans $\mathcal{C}'$,
alors nous obtenons des flèches mutuellement inverses $\colim F'(X_i) \to \colim F'(X'_{i'})$ et $\colim F'(X'_{i'}) \to \colim F'(X_i)$. Autrement dit,
la valeur $F(X)$ est bien définie indépendamment du choix de la présentation
de $X$ comme limite inductive filtrante d'objets de $\mathcal{C}'$. Avec la
fonctorialité de $F$ discutée dans le paragraphe précédent, nous constatons
que $F$ est un foncteur. En outre, il est clair que $F(X) = F'(X)$ si $X \in \Ob(\mathcal{C}')$.

\medskip\noindent
L'énoncé d'unicité du lemme est clair, à condition de montrer que
$F$ commute avec les limites inductives filtrantes (car cet énoncé n'a
aucun sens autrement). Pour cela, supposons que $X = \colim_{\lambda \in \Lambda} X_\lambda$ soit une limite inductive filtrante dans $\mathcal{C}$. Puisque
$F$ est un foncteur, nous obtenons un morphisme naturel
$$
\colim_\lambda F(X_\lambda) \longrightarrow F(X)
$$
D'autre part, écrivons $X = \colim X_i$ comme limite inductive filtrante d'objets de
$\mathcal{C}'$. Comme ci-dessus, pour chaque $i \in I$ on peut choisir un
$\lambda \in \Lambda$ et un diagramme commutatif
$$
\xymatrix{
X_i \ar[rr] \ar[rd] & & X_\lambda \ar[ld] \\
& X
}
$$
Comme ci-dessus cela détermine un morphisme bien défini $F'(X_i) \to \colim_\lambda F(X_\lambda)$ compatible avec les morphismes de transition et
donc un morphisme
$$
F(X) = \colim_i F'(X_i) \longrightarrow \colim_\lambda F(X_\lambda)
$$
Ce morphisme est inverse du morphisme ci-dessus (détails omis) et prouve que
$F(X) = \colim_\lambda F(X_\lambda)$ comme souhaité.
\end{proof}







\section{Localisation dans les catégories}
\label{section-localization}

\noindent
L'idée de base de cette section est la suivante : étant donnés une catégorie $\mathcal{C}$ et
un ensemble de flèches $S$, construire un foncteur $F : \mathcal{C} \to S^{-1}\mathcal{C}$ tel que tous les éléments de $S$ deviennent inversibles
dans $S^{-1}\mathcal{C}$ et tel que $F$ soit universel parmi tous les foncteurs
possédant cette propriété. Les références pour cette section sont
\cite[chapitre I, section 2]{GZ} et \cite[chapitre II, section 2]{Verdier}.

\begin{definition}
\label{definition-multiplicative-system}
Soit $\mathcal{C}$ une catégorie. Un ensemble de flèches $S$ de $\mathcal{C}$
est appelé un {\it système multiplicatif à gauche} s'il possède les propriétés
suivantes :
\begin{enumerate}
\item[LMS1] L'identité de chaque objet de $\mathcal{C}$ est dans $S$ et la
composition de deux éléments composables de $S$ est dans $S$.
\item[LMS2] Tout diagramme en traits pleins
$$
\xymatrix{
X \ar[d]_t \ar[r]_g & Y \ar@{..>}[d]^s \\
Z \ar@{..>}[r]^f & W
}
$$
avec $t \in S$ peut être complété en un carré commutatif par des flèches pointillées avec $s \in S$.
\item[LMS3] Pour chaque paire de morphismes $f, g : X \to Y$ et $t \in S$ de
but $X$ tel que $f \circ t = g \circ t$, il existe un $s \in S$ de
source $Y$ tel que $s \circ f = s \circ g$.
\end{enumerate}
Un ensemble de flèches $S$ de $\mathcal{C}$ est appelé un {\it système
multiplicatif à droite} s'il possède les propriétés suivantes :
\begin{enumerate}
\item[RMS1] L'identité de chaque objet de $\mathcal{C}$ est dans $S$ et la
composition de deux éléments composables de $S$ est dans $S$.
\item[RMS2] Tout diagramme en traits pleins
$$
\xymatrix{
X \ar@{..>}[d]_t \ar@{..>}[r]_g & Y \ar[d]^s \\
Z \ar[r]^f & W
}
$$
avec $s \in S$ peut être complété en un carré commutatif par des flèches pointillées avec $t \in S$.
\item[RMS3] Pour chaque paire de morphismes $f, g : X \to Y$ et $s \in S$ de
source $Y$ tel que $s \circ f = s \circ g$, il existe un $t \in S$ de
but $X$ tel que $f \circ t = g \circ t$.
\end{enumerate}
Un ensemble de flèches $S$ de $\mathcal{C}$ est appelé un {\it système
multiplicatif} s'il s'agit à la fois d'un système multiplicatif à gauche et d'un
système multiplicatif à droite. En d'autres termes, cela signifie que MS1, MS2,
MS3 sont vérifiées, où MS1 $=$ LMS1 $+$ RMS1, MS2 $=$ LMS2 $+$ RMS2 et MS3 $=$
LMS3 $+$ RMS3. (Cela dit, bien sûr LMS1 $=$ RMS1 $=$ MS1.)
\end{definition}

\noindent
Ces conditions sont utiles pour construire les catégories $S^{-1}\mathcal{C}$
comme suit.

\medskip\noindent
{\bf Calcul des fractions à gauche.} Soit $\mathcal{C}$ une catégorie et $S$
un système multiplicatif à gauche. Nous définissons une nouvelle catégorie
$S^{-1}\mathcal{C}$ comme suit (nous vérifions que cela fonctionne dans la
preuve du lemme \ref{lemma-left-localization}) :
\begin{enumerate}
\item Nous définissons $\Ob(S^{-1}\mathcal{C}) = \Ob(\mathcal{C})$.
\item Les morphismes $X \to Y$ de $S^{-1}\mathcal{C}$ sont représentés par des paires
$(f : X \to Y', s : Y \to Y')$ avec $s \in S$ à équivalence près.
(L'équivalence est définie ci-dessous. Considérons la classe d'équivalence d'une
paire $(f, s)$ comme $s^{-1}f : X \to Y$.)
\item Deux paires $(f_1 : X \to Y_1, s_1 : Y \to Y_1)$ et $(f_2 : X \to Y_2, s_2 : Y \to Y_2)$ sont dites équivalentes s'il existe une troisième paire
$(f_3 : X \to Y_3, s_3 : Y \to Y_3)$ ainsi que des morphismes $u : Y_1 \to Y_3$ et $v : Y_2 \to Y_3$ de $\mathcal{C}$ qui s'insèrent dans le diagramme commutatif
$$
\xymatrix{
 & Y_1 \ar[d]^u & \\
X \ar[ru]^{f_1} \ar[r]^{f_3} \ar[rd]_{f_2} &
Y_3 &
Y \ar[lu]_{s_1} \ar[l]_{s_3} \ar[ld]^{s_2} \\
& Y_2 \ar[u]_v &
}
$$
\item La composition des classes d'équivalence des couples $(f : X \to Y', s : Y \to Y')$ et $(g : Y \to Z', t : Z \to Z')$ est définie comme la classe
d'équivalence d'un couple $(h \circ f : X \to Z'', u \circ t : Z \to Z'')$ où
$h$ et $u \in S$ sont choisis pour s'insérer dans un diagramme commutatif
$$
\xymatrix{
Y \ar[d]_s \ar[r]_g & Z' \ar[d]^u \\
Y' \ar[r]^h & Z''
}
$$
qui existe par hypothèse.
\item Le morphisme identité $X \to X$ dans $S^{-1} \mathcal{C}$ est la
classe d'équivalence de la paire $(\text{id} : X \to X,
\text{id} : X \to X)$.
\end{enumerate}

\begin{lemma}
\label{lemma-left-localization}
Soit $\mathcal{C}$ une catégorie et $S$ un système multiplicatif à gauche.
\begin{enumerate}
\item La relation sur les paires définie ci-dessus est une relation d'équivalence.
\item La règle de composition donnée ci-dessus est bien définie sur les
classes d'équivalence.
\item La composition est associative (et les morphismes identité satisfont
les axiomes d'identité), et donc $S^{-1}\mathcal{C}$ est une catégorie.
\end{enumerate}
\end{lemma}

\begin{proof}
Preuve de (1). Disons que deux paires $p_1 = (f_1 : X \to Y_1, s_1 : Y \to Y_1)$ et $p_2 = (f_2 : X \to Y_2, s_2 : Y \to Y_2)$ sont élémentairement
équivalentes s'il existe un morphisme $a : Y_1 \to Y_2$ de $\mathcal{C}$ tel
que $a \circ f_1 = f_2$ et $a \circ s_1 = s_2$. Le diagramme :
$$
\xymatrix{
X \ar@{=}[d] \ar[r]_{f_1} & Y_1 \ar[d]^a & Y \ar[l]^{s_1} \ar@{=}[d] \\
X \ar[r]^{f_2} & Y_2 & Y \ar[l]_{s_2}
}
$$
Notons cette propriété en disant $p_1Ep_2$. Notez que $pEp$ et $aEb, bEc \Rightarrow aEc$. (Malgré son nom, $E$ n'est pas une relation d'équivalence.)
La partie (1) affirme que la relation $p \sim p' \Leftrightarrow \exists q: pEq \wedge p'Eq$ (où $q$ est une paire satisfaisant les mêmes
conditions que $p$ et $p'$) est une relation d'équivalence. Un simple argument
formel, utilisant les propriétés de $E$ ci-dessus, montre qu'il suffit de
prouver $p_3Ep_1, p_3Ep_2 \Rightarrow p_1 \sim p_2$. Supposons donc qu'on nous
donne un diagramme commutatif
$$
\xymatrix{
 & Y_1 & \\
X \ar[ru]^{f_1} \ar[r]^{f_3} \ar[rd]_{f_2} &
Y_3 \ar[u]_{a_{31}} \ar[d]^{a_{32}} &
Y \ar[lu]_{s_1} \ar[l]_{s_3} \ar[ld]^{s_2} \\
& Y_2 &
}
$$
avec $s_i \in S$. Nous appliquons d’abord LMS2 pour obtenir un diagramme
commutatif
$$
\xymatrix{
Y \ar[d]_{s_1} \ar[r]_{s_2} & Y_2 \ar@{..>}[d]^{a_{24}} \\
Y_1 \ar@{..>}[r]^{a_{14}} & Y_4
}
$$
avec $a_{24} \in S$. Ensuite, nous avons
$$
a_{14} \circ a_{31} \circ s_3 =
a_{14} \circ s_1 =
a_{24} \circ s_2 =
a_{24} \circ a_{32} \circ s_3.
$$
Ainsi, par LMS3, il existe un morphisme $s_{44} : Y_4 \to Y'_4$ tel que
$s_{44} \in S$ et $s_{44} \circ a_{14} \circ a_{31}
= s_{44} \circ a_{24} \circ a_{32}$. Ainsi, après avoir remplacé $Y_4$, $a_{14}$ et $a_{24}$ par
$Y'_4$, $s_{44} \circ a_{14}$ et $s_{44} \circ a_{24}$, nous pouvons supposer
que $a_{14} \circ a_{31} = a_{24} \circ a_{32}$ (et nous avons toujours
$a_{24} \in S$ et $a_{14} \circ s_1 = a_{24} \circ s_2$). Posons
$$
f_4 =
a_{14} \circ f_1 =
a_{14} \circ a_{31} \circ f_3 =
a_{24} \circ a_{32} \circ f_3 =
a_{24} \circ f_2
$$
et $s_4 = a_{14} \circ s_1 = a_{24} \circ s_2$. Ensuite, le diagramme
$$
\xymatrix{
X \ar@{=}[d] \ar[r]_{f_1} & Y_1 \ar[d]^{a_{14}} & Y \ar[l]^{s_1} \ar@{=}[d] \\
X \ar[r]^{f_4} & Y_4 & Y \ar[l]_{s_4}
}
$$
est commutatif, et nous avons $s_4 \in S$ (par LMS1). Ainsi, $p_1 E p_4$, où
$p_4 = (f_4, s_4)$. De même, $p_2 E p_4$. Ces deux faits donnent $p_1 \sim p_2$.

\medskip\noindent
Preuve de (2). Soient $p = (f : X \to Y', s : Y \to Y')$ et $q = (g : Y \to Z', t : Z \to Z')$ des paires comme dans la définition de la composition
ci-dessus. Pour composer on choisit un diagramme
$$
\xymatrix{
Y \ar[d]_s \ar[r]_g & Z' \ar[d]^{u_2} \\
Y' \ar[r]^{h_2} & Z_2
}
$$
avec $u_2 \in S$. Nous montrons d’abord que la classe d’équivalence du couple
$r_2 = (h_2 \circ f : X \to Z_2, u_2 \circ t : Z \to Z_2)$ est indépendante du
choix de $(Z_2, h_2, u_2)$. Supposons en effet que $(Z_3, h_3, u_3)$ soit un
autre choix avec la composition correspondante $r_3 = (h_3 \circ f : X \to Z_3, u_3 \circ t : Z \to Z_3)$. Alors, par LMS2 on peut choisir un diagramme
$$
\xymatrix{
Z' \ar[d]_{u_2} \ar[r]_{u_3} & Z_3 \ar[d]^{u_{34}} \\
Z_2 \ar[r]^{h_{24}} & Z_4
}
$$
avec $u_{34} \in S$. Nous avons $h_2 \circ s = u_2 \circ g$ et de même $h_3 \circ s = u_3 \circ g$. De plus,
$$
u_{34} \circ h_3 \circ s
= u_{34} \circ u_3 \circ g
= h_{24} \circ u_2 \circ g
= h_{24} \circ h_2 \circ s.
$$
Ainsi, LMS3 montre qu'il existe un $Z'_4$ et un $s_{44} : Z_4 \to Z'_4$ tel
que $s_{44} \circ u_{34} \circ h_3 = s_{44} \circ h_{24} \circ h_2$. En
remplaçant $Z_4$, $h_{24}$ et $u_{34}$ par $Z'_4$, $s_{44} \circ h_{24}$ et
$s_{44} \circ u_{34}$, nous pouvons supposer que $u_{34} \circ h_3 = h_{24} \circ h_2$. Pendant ce temps, les relations $u_{34} \circ u_3 = h_{24} \circ u_2$ et $u_{34} \in S$ restent satisfaites. Nous pouvons maintenant définir
$h_4 = u_{34} \circ h_3 = h_{24} \circ h_2$ et $u_4 = u_{34} \circ u_3 = h_{24} \circ u_2$. On a alors un diagramme commutatif
$$
\xymatrix{
X \ar@{=}[d] \ar[r]_{h_2\circ f} &
Z_2 \ar[d]^{h_{24}} &
Z \ar[l]^{u_2 \circ t} \ar@{=}[d] \\
X \ar@{=}[d] \ar[r]^{h_4\circ f} &
Z_4 &
Z \ar@{=}[d] \ar[l]_{u_4 \circ t} \\
X \ar[r]^{h_3 \circ f} &
Z_3 \ar[u]^{u_{34}} &
Z \ar[l]_{u_3 \circ t}
}
$$
Nous obtenons donc une paire $r_4 =
(h_4 \circ f : X \to Z_4, u_4 \circ t : Z \to Z_4)$ et le diagramme ci-dessus montre que nous avons $r_2Er_4$ et
$r_3Er_4$, d'où $r_2 \sim r_3$, comme souhaité. Il est donc désormais légitime
de définir $p \circ q$ comme la classe d'équivalence de toutes les paires
possibles $r$ obtenues comme ci-dessus.

\medskip\noindent
Pour terminer la preuve de (2), nous devons montrer que, pour des paires
$p_1, p_2, q$ telles que $p_1Ep_2$, on a $p_1 \circ q = p_2 \circ q$ et $q \circ p_1 = q \circ p_2$ chaque fois que les compositions ont un sens. Pour ce
faire, écrivons $p_1 = (f_1 : X \to Y_1, s_1 : Y \to Y_1)$ et $p_2 = (f_2 : X \to Y_2, s_2 : Y \to Y_2)$, et soit $a : Y_1 \to Y_2$ un morphisme de
$\mathcal{C}$ tel que $f_2 = a \circ f_1$ et $s_2 = a \circ s_1$. Supposons
d’abord que $q = (g : Y \to Z', t : Z \to Z')$. Dans ce cas, choisissez un
diagramme commutatif comme celui de gauche
$$
\vcenter{
\xymatrix{
Y \ar[d]_{s_2} \ar[r]^g & Z' \ar[d]^u \\
Y_2 \ar[r]^h & Z''
}
}
\quad
\Rightarrow
\quad
\vcenter{
\xymatrix{
Y \ar[d]_{s_1} \ar[r]^g & Z' \ar[d]^u \\
Y_1 \ar[r]^{h \circ a} & Z''
}
}
$$
(avec $u \in S$), ce qui implique que le diagramme de droite est également
commutatif. À l’aide de ces diagrammes, nous voyons que les deux compositions
$q \circ p_1$ et $q \circ p_2$ représentent la classe d’équivalence de $(h \circ a \circ f_1 : X \to Z'', u \circ t : Z \to Z'')$. Ainsi $q \circ p_1 = q \circ p_2$. La preuve de l’autre cas, dans lequel nous devons montrer $p_1 \circ q = p_2 \circ q$, est omise. (C'est similaire au cas que nous avons
traité.)

\medskip\noindent
Preuve de (3). Nous devons prouver l'associativité de la composition.
Considérons un diagramme en traits pleins
$$
\xymatrix{
& & & Z \ar[d] \\
& & Y \ar[d] \ar[r] & Z' \ar@{..>}[d] \\
& X \ar[d] \ar[r] & Y' \ar@{..>}[d] \ar@{..>}[r] & Z'' \ar@{..>}[d] \\
W \ar[r] & X' \ar@{..>}[r] & Y'' \ar@{..>}[r] & Z'''
}
$$
(dont les flèches verticales appartiennent à $S$) qui détermine trois
paires composables. En utilisant LMS2 on peut choisir les flèches pointillées
qui rendent les carrés commutatifs et telles que les flèches verticales soient dans
$S$. Il est alors clair que la composition des trois paires est la classe
d'équivalence de la paire $(W \to Z''', Z \to Z''')$ obtenue en composant les
flèches horizontales sur la ligne du bas et les flèches verticales sur la
colonne de droite.

\medskip\noindent
Nous laissons au lecteur le soin de vérifier les axiomes d'identité.
\end{proof}

\begin{remark}
\label{remark-motivation-localization}
La motivation pour la construction de $S^{-1} \mathcal{C}$ est de « forcer »
les morphismes de $S$ à être inversibles en créant artificiellement des
inverses (au prix de certains morphismes existants qui peuvent s'identifier les
uns aux autres). Ceci est analogue à la localisation d'un anneau commutatif
par rapport à une partie multiplicative, et plus généralement à la
localisation d'un anneau non commutatif par rapport à un ensemble de
dénominateurs à droite (voir \cite[section 10A]{Lam}). C'est plus qu'une simple
similitude : la construction de $S^{-1} \mathcal{C}$ (ou, plus précisément, sa
version pour les catégories additives $\mathcal{C}$) généralise en fait ce
dernier type de localisation. En effet, un anneau non commutatif peut être
considéré comme une catégorie préadditive avec un seul objet (les morphismes
étant les éléments de l'anneau) ; une partie multiplicative de cet anneau
devient alors un ensemble $S$ de morphismes satisfaisant LMS1 (alias RMS1).
Ensuite, les conditions RMS2 et RMS3 pour cette catégorie et cet ensemble
$S$ se traduisent par les deux conditions (« permutable à droite » et
« réversible à droite ») d'un ensemble de dénominateurs à droite (et de même
pour les ensembles de dénominateurs à gauche et LMS), et $S^{-1} \mathcal{C}$
(munie d'une structure additive convenablement définie) est la catégorie à un objet
correspondant à la localisation de l'anneau.
\end{remark}

\begin{definition}
\label{definition-left-localization-as-fraction}
Soit $\mathcal{C}$ une catégorie et soit $S$ un système multiplicatif à gauche
de morphismes de $\mathcal{C}$. Étant donné tout morphisme $f : X \to Y'$ dans
$\mathcal{C}$ et tout morphisme $s : Y \to Y'$ dans $S$, on note {\it $s^{-1} f$} la classe d'équivalence de la paire $(f : X \to Y', s : Y \to Y')$. Il
s'agit d'un morphisme de $X$ à $Y$ dans $S^{-1} \mathcal{C}$.
\end{definition}

\noindent
Cette notation est suggestive, et les choses qu'elle suggère sont vraies :
étant donnés un morphisme $f : X \to Y'$ dans $\mathcal{C}$ et deux
morphismes quelconques $s : Y \to Y'$ et $t : Y' \to Y''$ dans $S$, nous avons
$\left(t \circ s\right)^{-1} \left(t \circ f\right) = s^{-1} f$. De plus, pour
tous les $f : X \to Y'$ et $g : Y' \to Z'$ dans $\mathcal{C}$ et tous les $s : Z \to Z'$ dans $S$, nous avons $s^{-1} \left(g \circ f\right) = \left(s^{-1} g\right) \circ
\left(\text{id}_{Y'}^{-1} f\right)$. Enfin, pour tout $f : X \to Y'$ dans $\mathcal{C}$, tous les $s : Y \to Y'$ dans $S$ et $t : Z \to Y$
dans $S$, nous avons $\left(s \circ t\right)^{-1} f
= \left(t^{-1} \text{id}_Y\right)
\circ \left(s^{-1} f\right)$. Tout cela ressort clairement
de la définition. Nous pouvons « écrire toute famille finie de
morphismes de même but sous forme de fractions à dénominateur commun ».

\begin{lemma}
\label{lemma-morphisms-left-localization}
Soit $\mathcal{C}$ une catégorie et soit $S$ un système multiplicatif à gauche
de morphismes de $\mathcal{C}$. Étant donnée toute famille finie $g_i : X_i \to Y$ de morphismes de $S^{-1}\mathcal{C}$ (indexée par $i$), on peut trouver
un élément $s : Y \to Y'$ de $S$ et une famille de morphismes $f_i : X_i \to Y'$ de $\mathcal{C}$ telle que chaque $g_i$ soit la classe d'équivalence de la
paire $(f_i : X_i \to Y', s : Y \to Y')$.
\end{lemma}

\begin{proof}
Pour chaque $i$, choisissons un représentant $(X_i \to Y_i, s_i : Y \to Y_i)$
de $g_i$. Il suffit de trouver un morphisme $s : Y \to Y'$ dans
$S$ tel que pour chaque $i$ il existe un morphisme $a_i : Y_i \to Y'$ avec
$a_i \circ s_i = s$. Si nous avons deux indices $i = 1, 2$, alors nous pouvons
le faire en complétant le carré
$$
\xymatrix{
Y \ar[d]_{s_1} \ar[r]_{s_2} & Y_2 \ar[d]^{t_2} \\
Y_1 \ar[r]^{a_1} & Y'
}
$$
avec $t_2 \in S$ ce qui est possible par la définition
\ref{definition-multiplicative-system}. Ensuite, $s = t_2 \circ s_2 \in S$
convient. Si nous avons $n > 2$ morphismes, alors nous utilisons
l'argument ci-dessus pour nous ramener au cas de $n - 1$ morphismes, et nous concluons
par induction.
\end{proof}

\noindent
Il existe une caractérisation facile de l'égalité des morphismes s'ils ont le
même dénominateur.

\begin{lemma}
\label{lemma-equality-morphisms-left-localization}
Soit $\mathcal{C}$ une catégorie et soit $S$ un système multiplicatif à gauche
de morphismes de $\mathcal{C}$. Soient $A, B : X \to Y$ des morphismes de
$S^{-1}\mathcal{C}$ qui sont les classes d'équivalence de $(f : X \to Y', s : Y \to Y')$ et $(g : X \to Y', s : Y \to Y')$. Les assertions suivantes sont
équivalentes
\begin{enumerate}
\item $A = B$
\item il existe un morphisme $t : Y' \to Y''$ dans $S$ avec $t \circ f = t \circ g$, et
\item il existe un morphisme $a : Y' \to Y''$ tel que $a \circ f = a \circ g$
et $a \circ s \in S$.
\end{enumerate}
\end{lemma}

\begin{proof}
Nous utiliserons le fait que $S^{-1}\mathcal{C}$ est une catégorie (Lemme
\ref{lemma-left-localization}) et nous utiliserons la notation de la définition
\ref{definition-left-localization-as-fraction} ainsi que la discussion qui suit
cette définition pour identifier certains morphismes dans $S^{-1}\mathcal{C}$.
Ainsi nous écrivons $A = s^{-1}f$ et $B = s^{-1}g$.

\medskip\noindent
Si $A = B$ alors $(\text{id}_{Y'}^{-1}s) \circ A = (\text{id}_{Y'}^{-1}s) \circ B$. Nous avons $(\text{id}_{Y'}^{-1}s) \circ A = \text{id}_{Y'}^{-1}f$
et $(\text{id}_{Y'}^{-1}s) \circ B = \text{id}_{Y'}^{-1}g$. L'égalité de
$\text{id}_{Y'}^{-1}f$ et $\text{id}_{Y'}^{-1}g$ signifie par définition qu'il
existe un diagramme commutatif
$$
\xymatrix{
 & Y' \ar[d]^u & \\
X \ar[ru]^f \ar[r]^h \ar[rd]_g &
Z &
Y' \ar[lu]_{\text{id}_{Y'}} \ar[l]_t \ar[ld]^{\text{id}_{Y'}} \\
& Y' \ar[u]_v &
}
$$
avec $t \in S$. En particulier $u = v = t \in S$ et $t \circ f = t\circ g$.
Ainsi (1) implique (2).

\medskip\noindent
L’implication (2) $\Rightarrow$ (3) est immédiate. Supposons que $a$ soit
comme dans (3). Posons $s' = a \circ s \in S$. Alors $\text{id}_{Y''}^{-1}s'$
est un isomorphisme dans la catégorie $S^{-1}\mathcal{C}$ (avec l'inverse
$(s')^{-1}\text{id}_{Y''}$). Ainsi pour vérifier $A = B$ il suffit de vérifier
que $\text{id}_{Y''}^{-1}s' \circ A = \text{id}_{Y''}^{-1}s' \circ B$. Nous
calculons en utilisant les règles discutées dans le texte suivant la
définition \ref{definition-left-localization-as-fraction} que
$\text{id}_{Y''}^{-1}s' \circ A =
\text{id}_{Y''}^{-1}(a \circ s) \circ s^{-1}f =
\text{id}_{Y''}^{-1}(a \circ f) =
\text{id}_{Y''}^{-1}(a \circ g) =
\text{id}_{Y''}^{-1}(a \circ s) \circ s^{-1}g =
\text{id}_{Y''}^{-1}s' \circ B$ et nous voyons que (1) est vrai.
\end{proof}

\begin{remark}
\label{remark-left-localization-morphisms-colimit}
Soit $\mathcal{C}$ une catégorie. Soit $S$ un système multiplicatif à gauche.
Étant donné un objet $Y$ de $\mathcal{C}$ on note $Y/S$ la catégorie dont les
objets sont $s : Y \to Y'$ avec $s \in S$ et dont les morphismes sont des
diagrammes commutatifs
$$
\xymatrix{
& Y \ar[ld]_s \ar[rd]^t & \\
Y' \ar[rr]^a & & Y''
}
$$
où $a : Y' \to Y''$ est arbitraire. Nous affirmons que la catégorie $Y/S$ est
filtrante (voir définition \ref{definition-directed}). En effet, LMS1 implique
que $\text{id}_Y : Y \to Y$ se trouve dans $Y/S$ ; donc $Y/S$ n'est pas vide.
LMS2 implique qu'étant donnés deux morphismes $s_1 : Y \to Y_1$ et $s_2 : Y \to Y_2$ nous
pouvons trouver un diagramme
$$
\xymatrix{
Y \ar[d]_{s_1} \ar[r]_{s_2} & Y_2 \ar[d]^t \\
Y_1 \ar[r]^a & Y_3
}
$$
avec $t \in S$. Par conséquent, $s_1 : Y \to Y_1$ et $s_2 : Y \to Y_2$ ont
tous deux des morphismes vers $t \circ s_2 : Y \to Y_3$ dans $Y/S$. Enfin, étant
donnés deux morphismes $a, b$ de $s_1 : Y \to Y_1$ à $s_2 : Y \to Y_2$ dans
$Y/S$, nous voyons que $a \circ s_1 = b \circ s_1$ ; donc par LMS3 il existe
un $t : Y_2 \to Y_3$ dans $S$ tel que $t \circ a = t \circ b$. Les
résultats conjoints des lemmes \ref{lemma-morphisms-left-localization} et
\ref{lemma-equality-morphisms-left-localization} nous disent que
\begin{equation}
\label{equation-left-localization-morphisms-colimit}
\Mor_{S^{-1}\mathcal{C}}(X, Y) =
\colim_{(s : Y \to Y') \in Y/S} \Mor_\mathcal{C}(X, Y')
\end{equation}
Cette formule exprimant les ensembles de morphismes dans $S^{-1}\mathcal{C}$
comme limite inductive filtrante d'ensembles de morphismes dans $\mathcal{C}$
est parfois utile.
\end{remark}

\begin{lemma}
\label{lemma-properties-left-localization}
Soit $\mathcal{C}$ une catégorie et soit $S$ un système multiplicatif à gauche
de morphismes de $\mathcal{C}$.
\begin{enumerate}
\item Les règles $X \mapsto X$ et $(f : X \to Y) \mapsto (f : X \to Y, \text{id}_Y : Y \to Y)$ définissent un foncteur $Q : \mathcal{C} \to S^{-1}\mathcal{C}$.
\item Pour tout $s \in S$ le morphisme $Q(s)$ est un isomorphisme dans
$S^{-1}\mathcal{C}$.
\item Si $G : \mathcal{C} \to \mathcal{D}$ est un foncteur tel que $G(s)$ est
inversible pour chaque $s \in S$, alors il existe un foncteur unique $H : S^{-1}\mathcal{C} \to \mathcal{D}$ tel que $H \circ Q = G$.
\end{enumerate}
\end{lemma}

\begin{proof}
Les assertions (1) et (2) sont claires. (Dans (2), l'inverse de $Q(s)$ est la
classe d'équivalence de la paire $(\text{id}_Y, s)$.) Pour voir (3),
définissons simplement $H(X) = G(X)$ et $H((f : X \to Y', s : Y \to Y')) = G(s)^{-1} \circ G(f)$. Détails omis.
\end{proof}

\begin{lemma}
\label{lemma-left-localization-limits}
Soit $\mathcal{C}$ une catégorie et soit $S$ un système multiplicatif à gauche
de morphismes de $\mathcal{C}$. Le foncteur de localisation $Q : \mathcal{C} \to S^{-1}\mathcal{C}$ commute avec les colimites finies.
\end{lemma}

\begin{proof}
Soit $\mathcal{I}$ une catégorie finie et soit $\mathcal{I} \to \mathcal{C}$,
$i \mapsto X_i$ un foncteur dont la colimite existe. Alors, en utilisant
(\ref{equation-left-localization-morphisms-colimit}), le fait que $Y/S$ soit
filtrante et le lemme \ref{lemma-directed-commutes}, nous avons
\begin{align*}
\Mor_{S^{-1}\mathcal{C}}(Q(\colim X_i), Q(Y))
& =
\colim_{(s : Y \to Y') \in Y/S} \Mor_\mathcal{C}(\colim X_i, Y') \\
& =
\colim_{(s : Y \to Y') \in Y/S} \lim_i \Mor_\mathcal{C}(X_i, Y') \\
& =
\lim_i \colim_{(s : Y \to Y') \in Y/S} \Mor_\mathcal{C}(X_i, Y') \\
& =
\lim_i \Mor_{S^{-1}\mathcal{C}}(Q(X_i), Q(Y))
\end{align*}
et cet isomorphisme commute avec les projections des deux côtés vers
l'ensemble $\Mor_{S^{-1}\mathcal{C}}(Q(X_j), Q(Y))$ pour chaque $j \in \Ob(\mathcal{I})$. Ainsi, $Q(\colim X_i)$ satisfait la propriété universelle
de la colimite du foncteur $i \mapsto Q(X_i)$ ; c’est donc cette colimite,
comme souhaité.
\end{proof}

\begin{lemma}
\label{lemma-left-localization-diagram}
Soit $\mathcal{C}$ une catégorie. Soit $S$ un système multiplicatif à gauche.
Si $f : X \to Y$, $f' : X' \to Y'$ sont deux morphismes de $\mathcal{C}$ et si
$$
\xymatrix{
Q(X) \ar[d]_{Q(f)} \ar[r]_a & Q(X') \ar[d]^{Q(f')} \\
Q(Y) \ar[r]^b & Q(Y')
}
$$
est un diagramme commutatif dans $S^{-1}\mathcal{C}$, alors il existe un
morphisme $f'' : X'' \to Y''$ dans $\mathcal{C}$ et un diagramme commutatif
$$
\xymatrix{
X \ar[d]_f \ar[r]_g & X'' \ar[d]^{f''} & X' \ar[d]^{f'} \ar[l]^s \\
Y \ar[r]^h & Y'' & Y' \ar[l]_t
}
$$
dans $\mathcal{C}$ avec $s, t \in S$ et $a = s^{-1}g$, $b = t^{-1}h$.
\end{lemma}

\begin{proof}
Nous choisissons les morphismes et les objets de la manière suivante :
Écrivons d'abord $a = s^{-1}g$ pour un certain $s : X' \to X''$ dans $S$ et un morphisme $g : X \to X''$. Par LMS2 on peut trouver $t : Y' \to Y''$ dans $S$ et $f'' : X'' \to Y''$ de sorte que
$$
\xymatrix{
X \ar[d]_f \ar[r]_g & X'' \ar[d]^{f''} & X' \ar[d]^{f'} \ar[l]^s \\
Y & Y'' & Y' \ar[l]_t
}
$$
est commutatif. À ce stade, dans ce diagramme, nous
allons modifier à plusieurs reprises notre choix de
$$
X'' \xrightarrow{f''} Y'' \xleftarrow{t} Y'
$$
en postcomposant à la fois $t$ et $f''$ par un morphisme $d : Y'' \to Y'''$
avec la propriété $d \circ t \in S$. D'après la remarque
\ref{remark-left-localization-morphisms-colimit}, nous pouvons, après un tel
remplacement, supposer qu'il existe un morphisme $h : Y \to Y''$ tel que $b = t^{-1}h$\footnote{Voici une manière plus terre-à-terre de voir
cela : écrivons $b = q^{-1}i$ pour un certain $q : Y' \to Z$ dans $S$ et un certain
$i : Y \to Z$. Par LMS2, nous pouvons trouver $r : Y'' \to Y'''$ dans $S$ et
$j : Z \to Y'''$ tel que $j \circ q = r \circ t$. Posons alors $d = r$ et $h = j \circ i$.}. À ce stade, nous avons toutes les données du lemme, sauf
que nous ne savons pas que le carré de gauche du diagramme est commutatif.
Mais la définition de composition dans $S^{-1} \mathcal{C}$ montre que $b \circ Q\left(f\right)$ est la classe d'équivalence de la paire $(h \circ f : X \to Y'', t : Y' \to Y'')$ (puisque $b$ est la classe d'équivalence de la paire
$(h : Y \to Y'', t : Y' \to Y'')$, tandis que $Q\left(f\right)$ est la classe
d'équivalence de la paire $(f : X \to Y, \text{id} : Y \to Y)$), tandis que
$Q\left(f'\right) \circ a$ est la classe d'équivalence de la paire $(f'' \circ g : X \to Y'', t : Y' \to Y'')$ (puisque $a$ est la classe d'équivalence de la
paire $(g : X \to X'', s : X' \to X'')$, tandis que $Q\left(f'\right)$ est la
classe d'équivalence de la paire $(f' : X' \to Y', \text{id} : Y' \to Y')$). Puisque l'on sait que $b \circ Q\left(f\right) = Q\left(f'\right) \circ a$, on conclut que les classes d'équivalence des couples $(h \circ f : X \to Y'', t : Y' \to Y'')$ et $(f'' \circ g : X \to Y'', t : Y' \to Y'')$ sont
égales. Ainsi en utilisant le lemme
\ref{lemma-equality-morphisms-left-localization} nous pouvons trouver un
morphisme $d : Y'' \to Y'''$ tel que $d \circ t \in S$ et $d \circ h \circ f = d \circ f'' \circ g$. Nous effectuons donc un remplacement supplémentaire du
type décrit ci-dessus et nous concluons.
\end{proof}

\noindent
{\bf Calcul des fractions à droite.} Soit $\mathcal{C}$ une catégorie et $S$ un
système multiplicatif à droite. Nous définissons une nouvelle catégorie
$S^{-1}\mathcal{C}$ comme suit (nous vérifions que cela fonctionne dans la
preuve du lemme \ref{lemma-right-localization}) :
\begin{enumerate}
\item Nous définissons $\Ob(S^{-1}\mathcal{C}) = \Ob(\mathcal{C})$.
\item Les morphismes $X \to Y$ de $S^{-1}\mathcal{C}$ sont représentés par des couples
$(f : X' \to Y, s : X' \to X)$ avec $s \in S$ à équivalence près.
(L'équivalence est définie ci-dessous. Considérons la classe d'équivalence d'une
paire $(f, s)$ comme $fs^{-1} : X \to Y$.)
\item Deux paires $(f_1 : X_1 \to Y, s_1 : X_1 \to X)$ et $(f_2 : X_2 \to Y, s_2 : X_2 \to X)$ sont dites équivalentes s'il existe une troisième paire
$(f_3 : X_3 \to Y, s_3 : X_3 \to X)$ ainsi que des morphismes $u : X_3 \to X_1$ et $v : X_3 \to X_2$ de $\mathcal{C}$ qui s'insèrent dans le diagramme commutatif
$$
\xymatrix{
 & X_1 \ar[ld]_{s_1} \ar[rd]^{f_1} & \\
X &
X_3 \ar[l]_{s_3} \ar[u]_u \ar[d]^v \ar[r]^{f_3} &
Y \\
& X_2 \ar[lu]^{s_2} \ar[ru]_{f_2} &
}
$$
\item La composition des classes d'équivalence des couples $(f : X' \to Y, s : X' \to X)$ et $(g : Y' \to Z, t : Y' \to Y)$ est définie comme la classe
d'équivalence d'un couple $(g \circ h : X'' \to Z, s \circ u : X'' \to X)$ où
$h$ et $u \in S$ sont choisis pour s'insérer dans un diagramme commutatif
$$
\xymatrix{
X'' \ar[d]_u \ar[r]^h & Y' \ar[d]^t \\
X' \ar[r]^f & Y
}
$$
qui existe par hypothèse.
\item Le morphisme identité $X \to X$ dans $S^{-1} \mathcal{C}$ est la
classe d'équivalence de la paire $(\text{id} : X \to X,
\text{id} : X \to X)$.
\end{enumerate}

\begin{lemma}
\label{lemma-right-localization}
Soit $\mathcal{C}$ une catégorie et $S$ un système multiplicatif à droite.
\begin{enumerate}
\item La relation sur les paires définie ci-dessus est une relation d'équivalence.
\item La règle de composition donnée ci-dessus est bien définie sur les
classes d'équivalence.
\item La composition est associative (et les morphismes identité satisfont
les axiomes d'identité), et donc $S^{-1}\mathcal{C}$ est une catégorie.
\end{enumerate}
\end{lemma}

\begin{proof}
Ce lemme est dual du lemme \ref{lemma-left-localization}. Il découle
formellement de ce lemme en remplaçant $\mathcal{C}$ par sa catégorie opposée
dans laquelle $S$ est un système multiplicatif à gauche.
\end{proof}

\begin{definition}
\label{definition-right-localization-as-fraction}
Soit $\mathcal{C}$ une catégorie et soit $S$ un système multiplicatif à droite
de morphismes de $\mathcal{C}$. Étant donné tout morphisme $f : X' \to Y$ dans
$\mathcal{C}$ et tout morphisme $s : X' \to X$ dans $S$, on note {\it $f s^{-1}$} la classe d'équivalence de la paire $(f : X' \to Y, s : X' \to X)$.
Il s'agit d'un morphisme de $X$ à $Y$ dans $S^{-1} \mathcal{C}$.
\end{definition}

\noindent
On a des identités analogues (en fait, duales) à celles de la définition
\ref{definition-left-localization-as-fraction}. Nous pouvons « écrire
toute famille finie de morphismes de même source sous forme de
fractions à dénominateur commun ».

\begin{lemma}
\label{lemma-morphisms-right-localization}
Soit $\mathcal{C}$ une catégorie et soit $S$ un système multiplicatif à droite
de morphismes de $\mathcal{C}$. Étant donnée toute famille finie $g_i : X \to Y_i$ de morphismes de $S^{-1}\mathcal{C}$ (indexée par $i$), on peut
trouver un élément $s : X' \to X$ de $S$ et une famille de morphismes $f_i : X' \to Y_i$ de $\mathcal{C}$ tels que chaque $g_i$ soit la classe d'équivalence du
couple $(f_i : X' \to Y_i, s : X' \to X)$.
\end{lemma}

\begin{proof}
Ce lemme est le dual du lemme \ref{lemma-morphisms-left-localization} et
découle formellement de ce lemme en remplaçant toutes les catégories considérées
par leurs opposés.
\end{proof}

\noindent
Il existe une caractérisation facile de l'égalité des morphismes s'ils ont le
même dénominateur.

\begin{lemma}
\label{lemma-equality-morphisms-right-localization}
Soit $\mathcal{C}$ une catégorie et soit $S$ un système multiplicatif à droite
de morphismes de $\mathcal{C}$. Soient $A, B : X \to Y$ des morphismes de
$S^{-1}\mathcal{C}$ qui sont les classes d'équivalence de $(f : X' \to Y, s : X' \to X)$ et $(g : X' \to Y, s : X' \to X)$. Les assertions suivantes sont
équivalentes
\begin{enumerate}
\item $A = B$,
\item il existe un morphisme $t : X'' \to X'$ dans $S$ avec $f \circ t = g \circ t$, et
\item il existe un morphisme $a : X'' \to X'$ avec $f \circ a = g \circ a$ et
$s \circ a \in S$.
\end{enumerate}
\end{lemma}

\begin{proof}
Ce lemme est dual du lemme \ref{lemma-equality-morphisms-left-localization}.
\end{proof}

\begin{remark}
\label{remark-right-localization-morphisms-colimit}
Soit $\mathcal{C}$ une catégorie. Soit $S$ un système multiplicatif à droite.
Étant donné un objet $X$ de $\mathcal{C}$ on note $S/X$ la catégorie dont les
objets sont $s : X' \to X$ avec $s \in S$ et dont les morphismes sont des
diagrammes commutatifs
$$
\xymatrix{
X' \ar[rd]_s \ar[rr]_a & & X'' \ar[ld]^t \\
& X
}
$$
où $a : X' \to X''$ est arbitraire. La catégorie $S/X$ est cofiltrante (voir
définition \ref{definition-codirected}). (Ceci est le dual de l'énoncé
correspondant dans la remarque
\ref{remark-left-localization-morphisms-colimit}.) Les résultats
conjoints des lemmes \ref{lemma-morphisms-right-localization} et
\ref{lemma-equality-morphisms-right-localization} nous disent que
\begin{equation}
\label{equation-right-localization-morphisms-colimit}
\Mor_{S^{-1}\mathcal{C}}(X, Y) =
\colim_{(s : X' \to X) \in (S/X)^{opp}} \Mor_\mathcal{C}(X', Y)
\end{equation}
Cette formule exprimant les ensembles de morphismes dans $S^{-1}\mathcal{C}$
comme limite inductive filtrante d'ensembles de morphismes dans $\mathcal{C}$
est parfois utile.
\end{remark}

\begin{lemma}
\label{lemma-properties-right-localization}
Soit $\mathcal{C}$ une catégorie et soit $S$ un système multiplicatif à droite
de morphismes de $\mathcal{C}$.
\begin{enumerate}
\item Les règles $X \mapsto X$ et $(f : X \to Y) \mapsto (f : X \to Y, \text{id}_X : X \to X)$ définissent un foncteur $Q : \mathcal{C} \to S^{-1}\mathcal{C}$.
\item Pour tout $s \in S$ le morphisme $Q(s)$ est un isomorphisme dans
$S^{-1}\mathcal{C}$.
\item Si $G : \mathcal{C} \to \mathcal{D}$ est un foncteur tel que $G(s)$ est
inversible pour chaque $s \in S$, alors il existe un foncteur unique $H : S^{-1}\mathcal{C} \to \mathcal{D}$ tel que $H \circ Q = G$.
\end{enumerate}
\end{lemma}

\begin{proof}
Ce lemme est le dual du lemme \ref{lemma-properties-left-localization} et
découle formellement de ce lemme en remplaçant toutes les catégories considérées
par leurs opposés.
\end{proof}

\begin{lemma}
\label{lemma-right-localization-limits}
Soit $\mathcal{C}$ une catégorie et soit $S$ un système multiplicatif à droite
de morphismes de $\mathcal{C}$. Le foncteur de localisation $Q : \mathcal{C} \to S^{-1}\mathcal{C}$ commute avec les limites finies.
\end{lemma}

\begin{proof}
Ce lemme est dual du lemme \ref{lemma-left-localization-limits}.
\end{proof}

\begin{lemma}
\label{lemma-right-localization-diagram}
Soit $\mathcal{C}$ une catégorie. Soit $S$ un système multiplicatif à droite.
Si $f : X \to Y$, $f' : X' \to Y'$ sont deux morphismes de $\mathcal{C}$ et si
$$
\xymatrix{
Q(X) \ar[d]_{Q(f)} \ar[r]_a & Q(X') \ar[d]^{Q(f')} \\
Q(Y) \ar[r]^b & Q(Y')
}
$$
est un diagramme commutatif dans $S^{-1}\mathcal{C}$, alors il existe un
morphisme $f'' : X'' \to Y''$ dans $\mathcal{C}$ et un diagramme commutatif
$$
\xymatrix{
X \ar[d]_f & X'' \ar[l]^s \ar[d]^{f''} \ar[r]_g & X' \ar[d]^{f'} \\
Y & Y'' \ar[l]_t \ar[r]^h & Y'
}
$$
dans $\mathcal{C}$ avec $s, t \in S$ et $a = gs^{-1}$, $b = ht^{-1}$.
\end{lemma}

\begin{proof}
Ce lemme est dual du lemme \ref{lemma-left-localization-diagram}.
\end{proof}

\noindent
{\bf Systèmes multiplicatifs et calcul bilatéral des fractions.} Si $S$ est un
système multiplicatif, alors les calculs de fractions à gauche et à droite donnent des
catégories canoniquement isomorphes.

\begin{lemma}
\label{lemma-multiplicative-system}
Soit $\mathcal{C}$ une catégorie et $S$ un système multiplicatif. Les catégories
de fractions à gauche et de fractions à droite
$S^{-1}\mathcal{C}$ sont canoniquement isomorphes.
\end{lemma}

\begin{proof}
Notons $\mathcal{C}_{left}$, $\mathcal{C}_{right}$ les deux catégories de
fractions. Par les propriétés universelles des lemmes
\ref{lemma-properties-left-localization} et
\ref{lemma-properties-right-localization} nous obtenons les foncteurs
$\mathcal{C}_{left} \to \mathcal{C}_{right}$ et $\mathcal{C}_{right} \to \mathcal{C}_{left}$. D'après l'énoncé d'unicité dans les propriétés
universelles, ces foncteurs sont inverses les uns des autres.
\end{proof}

\begin{definition}
\label{definition-saturated-multiplicative-system}
Soit $\mathcal{C}$ une catégorie et $S$ un système multiplicatif. On dit que
$S$ est {\it saturé} si, en plus de MS1, MS2, MS3, on a aussi
\begin{enumerate}
\item[MS4] Étant donnés trois morphismes composables $f, g, h$, si $fg, gh \in S$, alors $g \in S$.
\end{enumerate}
\end{definition}

\noindent
Notez qu'un système multiplicatif saturé contient tous les isomorphismes. De
plus, si $f, g, h$ sont des morphismes composables dans une catégorie et que
$fg, gh$ sont des isomorphismes, alors $g$ est un isomorphisme (car alors $g$
a à la fois un inverse gauche et un inverse droit, donc est inversible).

\begin{lemma}
\label{lemma-what-gets-inverted}
Soit $\mathcal{C}$ une catégorie et $S$ un système multiplicatif. Notons $Q : \mathcal{C} \to S^{-1}\mathcal{C}$ le foncteur de localisation. L'ensemble
$$
\hat S = \{f \in \text{Flèches}(\mathcal{C}) \mid
Q(f) \text{ est un isomorphisme}\}
$$
est égal à
$$
S' = \{f \in \text{Flèches}(\mathcal{C}) \mid
\text{il existe }g, h\text{ tels que }gf, fh \in S\}
$$
et est le plus petit système multiplicatif saturé contenant $S$. En
particulier, si $S$ est saturé, alors $\hat S = S$.
\end{lemma}

\begin{proof}
Il est clair que $S \subset S' \subset \hat S$ car les éléments de $S'$
correspondent aux morphismes de $S^{-1}\mathcal{C}$ qui ont à la fois des
inverses à gauche et à droite. Notez que $S'$ satisfait à MS4 et que $\hat S$
satisfait à MS1. Montrons ensuite que $S' = \hat S$.

\medskip\noindent
Soit $f \in \hat S$. Soit $s^{-1}g = ht^{-1}$ le morphisme inverse dans
$S^{-1}\mathcal{C}$. (Nous pouvons utiliser à la fois des fractions gauches et
des fractions droites pour décrire les morphismes dans $S^{-1}\mathcal{C}$,
voir le lemme \ref{lemma-multiplicative-system}.) La relation $\text{id}_X = s^{-1}gf$ dans $S^{-1}\mathcal{C}$ signifie qu'il existe un diagramme
commutatif.
$$
\xymatrix{
 & X' \ar[d]^u & \\
X \ar[ru]^{gf} \ar[r]^{f'} \ar[rd]_{\text{id}_X} &
X'' &
X \ar[lu]_s \ar[l]_{s'} \ar[ld]^{\text{id}_X} \\
& X \ar[u]_v &
}
$$
pour certains morphismes $f', u, v$ et $s' \in S$. D'où $ugf = s' \in S$. De
même, en utilisant l'égalité $\text{id}_Y = fht^{-1}$ on prouve que $fhw \in S$ pour
un certain $w$. Nous concluons que $f \in S'$. Ainsi $S' = \hat S$. Une fois
établi que $S' = \hat S$ est un système multiplicatif, il est
clair que $S' = \hat S$ est le plus petit système
saturé contenant $S$.

\medskip\noindent
Nos remarques ci-dessus établissent MS1 et MS4, donc pour terminer la
preuve du lemme nous devons montrer que LMS2, RMS2, LMS3, RMS3 sont valables
pour $\hat S$. Vérifions que LMS2 est satisfaite pour $\hat S$. Supposons que
nous ayons un diagramme en traits pleins
$$
\xymatrix{
X \ar[d]_t \ar[r]_g & Y \ar@{..>}[d]^s \\
Z \ar@{..>}[r]^f & W
}
$$
avec $t \in \hat S$. Choisissons un morphisme $a : Z \to Z'$ tel que $at \in S$. Nous pouvons alors appliquer LMS2 à $S$ pour trouver un diagramme
commutatif
$$
\xymatrix{
X \ar[d]_t \ar[r]_g & Y \ar[dd]^s \\
Z \ar[d]_a \\
Z' \ar[r]^{f'} & W
}
$$
et, en posant $f = f' \circ a$, nous concluons. La preuve de RMS2 est duale. Enfin,
supposons qu'une paire de morphismes $f, g : X \to Y$ et $t \in \hat S$ soient
donnés, où ce dernier est de but $X$ et vérifie $ft = gt$. Choisissons alors un
morphisme $b$ tel que $tb \in S$. Alors $ftb = gtb$, ce qui implique par LMS3 pour
$S$ qu'il existe un $s \in S$ de source $Y$ tel que $sf = sg$, comme
souhaité. La preuve de RMS3 est duale.
\end{proof}








\section{Propriétés formelles}
\label{section-formal-cat-cat}

\noindent
Dans cette section, nous discutons de certaines propriétés formelles de la
$2$-catégorie des catégories. Cela nous amènera plus tard à la définition d’une
$2$-catégorie (stricte).

\medskip\noindent
Notons $\Ob(\textit{Cat})$ la classe de toutes les catégories. Pour chaque
paire de catégories $\mathcal{A}, \mathcal{B} \in \Ob(\textit{Cat})$, nous
avons la « petite » catégorie des foncteurs $\text{Fun}(\mathcal{A}, \mathcal{B})$. La composition de transformations de foncteurs comme
$$
\xymatrix{
\mathcal{A}
\rruppertwocell^{F''}{t'}
\ar[rr]_(.3){F'}
\rrlowertwocell_F{t}
& &
\mathcal{B}
}
\text{ a pour composé }
\xymatrix{
\mathcal{A}
\rrtwocell^{F''}_F{\ \ t \circ t'}
& &
\mathcal{B}
}
$$
est appelée composition {\it verticale}. Nous utiliserons pour cela le symbole
habituel $\circ$. Nous définirons ensuite la composition {\it horizontale}.
Pour ce faire, nous expliquons un peu plus la structure en question.

\medskip\noindent
En effet, pour chaque triplet de catégories $\mathcal{A}$, $\mathcal{B}$ et
$\mathcal{C}$, il existe une loi de composition
$$
\circ : \Ob(\text{Fun}(\mathcal{B}, \mathcal{C}))
\times
\Ob(\text{Fun}(\mathcal{A}, \mathcal{B}))
\longrightarrow
\Ob(\text{Fun}(\mathcal{A}, \mathcal{C}))
$$
provenant de la composition des foncteurs. Cette loi de composition est
associative et les foncteurs identité agissent comme des unités. Autrement
dit -- en oubliant les transformations des foncteurs -- on voit que
$\textit{Cat}$ forme une catégorie. Comment cette structure interagit-elle
avec les morphismes entre foncteurs ?

\medskip\noindent
Soit $t : F \to F'$ une transformation de foncteurs, où $F, F' : \mathcal{A} \to \mathcal{B}$, et soit $G : \mathcal{B} \to \mathcal{C}$ un foncteur.
Nous pouvons définir une transformation de foncteurs $G\circ F \to G \circ F'$, que nous noterons ${}_Gt$. Elle est donnée par la formule
$({}_Gt)_x = G(t_x) : G(F(x)) \to G(F'(x))$ pour tous les $x \in \mathcal{A}$.
De cette façon, la composition avec $G$ devient un foncteur
$$
\text{Fun}(\mathcal{A}, \mathcal{B})
\longrightarrow
\text{Fun}(\mathcal{A}, \mathcal{C}).
$$
Pour le voir, il suffit de vérifier que ${}_G(\text{id}_F) = \text{id}_{G \circ F}$ et que ${}_G(t_1 \circ t_2) = {}_Gt_1 \circ {}_Gt_2$.
Bien sûr, nous avons aussi ${}_{\text{id}_\mathcal{B}}t = t$.

\medskip\noindent
De même, soit $s : G \to G'$ une transformation de foncteurs, où $G, G' : \mathcal{B} \to \mathcal{C}$, et soit $F : \mathcal{A} \to \mathcal{B}$ un foncteur.
Nous pouvons définir $s_F$ comme étant la transformation de foncteurs $G\circ F \to G' \circ F$ donnée par $(s_F)_x = s_{F(x)} : G(F(x)) \to G'(F(x))$ pour
tous les $x \in \mathcal{A}$. De cette façon, la composition avec $F$ devient
un foncteur
$$
\text{Fun}(\mathcal{B}, \mathcal{C})
\longrightarrow
\text{Fun}(\mathcal{A}, \mathcal{C}).
$$
Pour le voir, il suffit de vérifier que $(\text{id}_G)_F = \text{id}_{G\circ F}$ et que $(s_1 \circ s_2)_F = s_{1, F} \circ s_{2, F}$.
Bien sûr, nous avons aussi $s_{\text{id}_\mathcal{B}} = s$.

\medskip\noindent
Ces constructions satisfont en outre aux propriétés
$$
{}_{G_1}({}_{G_2}t) = {}_{G_1\circ G_2}t,
\ (s_{F_1})_{F_2} = s_{F_1 \circ F_2},
\text{ et }{}_H(s_F) = ({}_Hs)_F
$$
chaque fois que cela a du sens. Enfin, étant donnés les foncteurs $F, F' : \mathcal{A} \to \mathcal{B}$ et $G, G' : \mathcal{B} \to \mathcal{C}$ ainsi que les
transformations $t : F \to F'$ et $s : G \to G'$, le diagramme suivant est
commutatif
$$
\xymatrix{
G \circ F \ar[r]^{{}_Gt} \ar[d]_{s_F}
&
G \circ F' \ar[d]^{s_{F'}} \\
G' \circ F \ar[r]_{{}_{G'}t}
&
G' \circ F'
}
$$
en d’autres termes ${}_{G'}t \circ s_F = s_{F'}\circ {}_Gt$. Pour le prouver,
considérons simplement ce qui se passe sur tout objet $x \in \Ob(\mathcal{A})$ :
$$
\xymatrix{
G(F(x)) \ar[r]^{G(t_x)} \ar[d]_{s_{F(x)}}
&
G(F'(x)) \ar[d]^{s_{F'(x)}} \\
G'(F(x)) \ar[r]_{G'(t_x)}
&
G'(F'(x))
}
$$
qui est commutatif car $s$ est une transformation de foncteurs. Cette relation
de compatibilité permet de définir une composition horizontale.

\begin{definition}
\label{definition-horizontal-composition}
Étant donné un diagramme comme dans le premier diagramme ci-dessous :
$$
\xymatrix{
\mathcal{A}
\rtwocell^F_{F'}{t}
&
\mathcal{B}
\rtwocell^G_{G'}{s}
&
\mathcal{C}
}
\text{ donne }
\xymatrix{
\mathcal{A}
\rrtwocell^{G \circ F} _{G' \circ F'}{\ \ s \star t}
& &
\mathcal{C}
}
$$
nous définissons la composition {\it horizontale} $s \star t$ comme étant la
transformation de foncteurs ${}_{G'}t \circ s_F = s_{F'}\circ {}_Gt$.
\end{definition}

\noindent
Nous voyons maintenant que l'on peut retrouver nos transformations
précédemment construites ${}_Gt$ et $s_F$ sous les noms $ {}_Gt = \text{id}_G \star t $ et $ s_F = s \star \text{id}_F $. De plus, toutes les règles que
nous avons trouvées ci-dessus sont des conséquences des propriétés énoncées
dans le lemme qui suit.

\begin{lemma}
\label{lemma-properties-2-cat-cats}
Les compositions horizontales et verticales ont les propriétés suivantes
\begin{enumerate}
\item $\circ$ et $\star$ sont associatifs,
\item les transformations identité $\text{id}_F$ sont des unités pour
$\circ$,
\item les transformations identité des foncteurs identité
$\text{id}_{\text{id}_\mathcal{A}}$ sont des unités pour $\star$ et $\circ$,
et
\item étant donné un diagramme
$$
\xymatrix{
\mathcal{A}
\rruppertwocell^F{t}
\ar[rr]_(.3){F'}
\rrlowertwocell_{F''}{t'}
& &
\mathcal{B}
\rruppertwocell^G{s}
\ar[rr]_(.3){G'}
\rrlowertwocell_{G''}{s'}
& &
\mathcal{C}
}
$$
nous avons $ (s' \circ s) \star (t' \circ t) = (s' \star t') \circ (s \star t)$.
\end{enumerate}
\end{lemma}

\begin{proof}
La dernière assertion devient, dans la notation précédente, l'équation
suivante
$$
s'_{F''}
\circ
{}_{G'}t'
\circ
s_{F'}
\circ
{}_Gt
=
(s' \circ s)_{F''}
\circ
{}_G(t' \circ t).
$$
D'après notre résultat ci-dessus appliqué à la composition du milieu, nous
pouvons réécrire le côté gauche comme $
s'_{F''}
\circ
s_{F''}
\circ
{}_Gt'
\circ
{}_Gt
$ qui est facilement reconnu égal au côté droit.
\end{proof}

\noindent
Une autre façon de formuler la condition (4) du lemme est que la composition
des foncteurs et la composition horizontale des transformations de foncteurs
définissent un foncteur
$$
(\circ, \star) :
\text{Fun}(\mathcal{B}, \mathcal{C})
\times
\text{Fun}(\mathcal{A}, \mathcal{B})
\longrightarrow
\text{Fun}(\mathcal{A}, \mathcal{C})
$$
dont la source est la catégorie produit, voir Définition
\ref{definition-product-category}.

\section{2-catégories}
\label{section-2-categories}

\noindent
Nous donnerons une définition des $2$-catégories (strictes) telles qu'elles
apparaissent dans le contexte des champs. Avant de poursuivre, consulter
la section \ref{section-formal-cat-cat} et l'exemple
\ref{example-2-1-category-of-categories}. Il s'agit essentiellement de partir de cet
exemple et d'expliciter toutes les règles satisfaites par les objets, les
$1$-morphismes et les $2$-morphismes dans cet exemple.

\begin{definition}
\label{definition-2-category}
Une (stricte) {\it $2$-catégorie} $\mathcal{C}$ se compose des données suivantes
\begin{enumerate}
\item Un ensemble d'objets $\Ob(\mathcal{C})$.
\item Pour chaque paire $x, y \in \Ob(\mathcal{C})$ une catégorie
$\Mor_\mathcal{C}(x, y)$. Les objets de $\Mor_\mathcal{C}(x, y)$ seront
appelés {\it $1$-morphismes} et notés $F : x \to y$. Les morphismes entre ces
$1$-morphismes seront appelés {\it $2$-morphismes} et notés $t : F' \to F$. La
composition des $2$-morphismes dans $\Mor_\mathcal{C}(x, y)$ sera appelée
composition {\it verticale} et sera notée $t \circ t'$ pour $t : F' \to F$ et
$t' : F'' \to F'$.
\item Pour chaque triple $x, y, z\in \Ob(\mathcal{C})$ un foncteur
$$
(\circ, \star) :
\Mor_\mathcal{C}(y, z) \times \Mor_\mathcal{C}(x, y)
\longrightarrow
\Mor_\mathcal{C}(x, z).
$$
L'image du couple de $1$-morphismes $(F, G)$ dans le membre de gauche sera appelée
la {\it composition} de $F$ et $G$, et notée $F\circ G$. L'image de la paire
de $2$-morphismes $(t, s)$ sera appelée composition {\it horizontale} et notée
$t \star s$.
\end{enumerate}
Ces données doivent satisfaire aux règles suivantes :
\begin{enumerate}
\item L'ensemble des objets ainsi que l'ensemble des $1$-morphismes dotés
d'une composition de $1$-morphismes forment une catégorie.
\item La composition horizontale des $2$-morphismes est associative.
\item Le $2$-morphisme identité $\text{id}_{\text{id}_x}$ de l'identité
$1$-morphisme $\text{id}_x$ est une unité pour la composition horizontale.
\end{enumerate}
\end{definition}

\noindent
Ce type de structure n'est évidemment pas très commode à manier. D’un autre
côté, il existe de nombreux exemples où la manière de procéder est
assez claire. Le seul exemple que nous ayons jusqu'à présent est
celui de la $2$-catégorie dont les objets forment une famille donnée de
catégories, les $1$-morphismes sont des foncteurs entre ces catégories, et les
$2$-morphismes sont des transformations naturelles de foncteurs, voir la
section \ref{section-formal-cat-cat}. En ce qui concerne ce texte, toutes les
$2$-catégories seront des sous-$2$-catégories de cet exemple. Voici ce que
signifie être une sous-$2$-catégorie.

\begin{definition}
\label{definition-sub-2-category}
Soit $\mathcal{C}$ une $2$-catégorie. Une {\it sous-$2$-catégorie}
$\mathcal{C}'$ de $\mathcal{C}$ est donnée par un sous-ensemble
$\Ob(\mathcal{C}')$ de $\Ob(\mathcal{C})$ et des sous-catégories
$\Mor_{\mathcal{C}'}(x, y)$ des catégories $\Mor_\mathcal{C}(x, y)$ pour tous
les $x, y \in \Ob(\mathcal{C}')$ de sorte que ces données, avec les opérations
$\circ$ (composition des $1$-morphismes), $\circ$ (composition verticale des
$2$-morphismes) et $\star$ (composition horizontale) forment une
$2$-catégorie.
\end{definition}

\begin{remark}
\label{remark-big-2-categories}
Grandes $2$-catégories. Dans de nombreux textes, une $2$-catégorie est
autorisée à avoir une classe d'objets (mais, espérons-le, une « classe de
classes » n'est pas autorisée). Nous autoriserons également ces « grandes »
$2$-catégories, mais uniquement dans la liste de cas suivante (à mettre à jour
au fur et à mesure) :
\begin{enumerate}
\item La $2$-catégorie des catégories $\textit{Cat}$.
\item La $(2, 1)$-catégorie des catégories $\textit{Cat}$.
\item La $2$-catégorie des groupoïdes $\textit{Groupoïdes}$ ; il s'agit d'une
$(2, 1)$-catégorie.
\item La $2$-catégorie des catégories fibrées au-dessus d'une catégorie fixe.
\item La $(2, 1)$-catégorie des catégories fibrées au-dessus d'une catégorie fixe.
\item La $2$-catégorie des catégories fibrées en groupoïdes au-dessus d'une catégorie
fixe ; il s'agit d'une $(2, 1)$-catégorie.
\item La $2$-catégorie des champs sur un site fixe.
\item La $(2, 1)$-catégorie des champs sur un site fixe.
\item La $2$-catégorie des champs en groupoïdes sur un site fixe ; il s'agit
d'une $(2, 1)$-catégorie.
\item La $2$-catégorie des champs en setoïdes sur un site fixe ; il s'agit d'une
$(2, 1)$-catégorie.
\item La $2$-catégorie des champs algébriques sur un schéma fixe ; il s'agit
d'une $(2, 1)$-catégorie.
\end{enumerate}
Voir la définition \ref{definition-2-1-category}. Notez que dans chaque cas,
la classe d'objets de la $2$-catégorie $\mathcal{C}$ est une classe propre,
mais pour tous les objets $x, y \in \Ob(\mathcal{C})$, la catégorie $\Mor_\mathcal{C}(x, y)$ est « petite » (selon nos conventions).
\end{remark}

\noindent
La notion d'équivalence des catégories que nous avons définie dans la section
\ref{section-definition-categories} s'étend au cadre plus général des
$2$-catégories comme suit.

\begin{definition}
\label{definition-equivalence}
Deux objets $x, y$ d'une $2$-catégorie sont {\it équivalents} s'il existe des
$1$-morphismes $F : x \to y$ et $G : y \to x$ tels que $F \circ G$ est
$2$-isomorphe à $\text{id}_y$ et $G \circ F$ est $2$-isomorphe à
$\text{id}_x$.
\end{definition}

\noindent
Parfois, nous devons dire ce que signifie avoir un foncteur d'une catégorie
dans une $2$-catégorie.

\begin{definition}
\label{definition-functor-into-2-category}
Soit $\mathcal{A}$ une catégorie et $\mathcal{C}$ une $2$-catégorie.
\begin{enumerate}
\item Un {\it foncteur} d'une catégorie ordinaire dans une $2$-catégorie
ignorera les $2$-morphismes, sauf indication contraire. Autrement dit, ce sera
un foncteur « habituel » dans la catégorie obtenue à partir de cette 2-catégorie en oubliant
tous les 2-morphismes.
\item Un {\it foncteur faible}, ou {\it pseudo-foncteur} $\varphi$ de
$\mathcal{A}$ dans la 2-catégorie $\mathcal{C}$ est défini par les données
suivantes :
\begin{enumerate}
\item une application $\varphi : \Ob(\mathcal{A}) \to \Ob(\mathcal{C})$,
\item pour chaque paire $x, y\in \Ob(\mathcal{A})$ et chaque morphisme $f : x \to y$, un $1$-morphisme $\varphi(f) : \varphi(x) \to \varphi(y)$,
\item pour chaque $x\in \Ob(\mathcal{A})$ un $2$-morphisme $\alpha_x : \text{id}_{\varphi(x)} \to \varphi(\text{id}_x)$, et
\item pour chaque paire de morphismes composables $f : x \to y$, $g : y \to z$
de $\mathcal{A}$, un $2$-morphisme $\alpha_{g, f} : \varphi(g \circ f) \to \varphi(g) \circ \varphi(f)$.
\end{enumerate}
Ces données sont soumises aux conditions suivantes :
\begin{enumerate}
\item les $2$-morphismes $\alpha_x$ et $\alpha_{g, f}$ sont tous des
isomorphismes,
\item pour tout morphisme $f : x \to y$ dans $\mathcal{A}$ nous avons
$\alpha_{\text{id}_y, f} = \alpha_y \star \text{id}_{\varphi(f)}$ :
$$
\xymatrix{
\varphi(x)
\rrtwocell^{\varphi(f)}_{\varphi(f)}{\ \ \ \ \text{id}_{\varphi(f)}}
& &
\varphi(y)
\rrtwocell^{\text{id}_{\varphi(y)}}_{\varphi(\text{id}_y)}{\alpha_y}
& &
\varphi(y)
}
=
\xymatrix{
\varphi(x)
\rrtwocell^{\varphi(f)}_{\varphi(\text{id}_y) \circ \varphi(f)}{\ \ \ \ \alpha_{\text{id}_y, f}}
& &
\varphi(y)
}
$$
\item pour tout morphisme $f : x \to y$ dans $\mathcal{A}$ nous avons
$\alpha_{f, \text{id}_x} = \text{id}_{\varphi(f)} \star \alpha_x$,
\item pour tout triple de morphismes composables $f : w \to x$, $g : x \to y$
et $h : y \to z$ de $\mathcal{A}$ nous avons
$$
(\text{id}_{\varphi(h)} \star \alpha_{g, f})
\circ
\alpha_{h, g \circ f}
=
(\alpha_{h, g} \star \text{id}_{\varphi(f)})
\circ
\alpha_{h \circ g, f}
$$
en d'autres termes, le diagramme suivant, dont les objets sont les $1$-morphismes
et les flèches les $2$-morphismes, est commutatif
$$
\xymatrix{
\varphi(h \circ g \circ f)
\ar[d]_{\alpha_{h, g \circ f}}
\ar[rr]_{\alpha_{h \circ g, f}}
& &
\varphi(h \circ g) \circ \varphi(f)
\ar[d]^{\alpha_{h, g} \star \text{id}_{\varphi(f)}} \\
\varphi(h) \circ \varphi(g \circ f)
\ar[rr]^{\text{id}_{\varphi(h)} \star \alpha_{g, f}}
& &
\varphi(h) \circ \varphi(g) \circ \varphi(f)
}
$$
\end{enumerate}
\end{enumerate}
\end{definition}

\noindent
Encore une fois, ce n'est pas une notion très pratique, mais elle revient
parfois. Il existe un théorème qui dit que tout pseudo-foncteur est isomorphe
à un foncteur. Enfin, il y a les notions de {\it foncteur entre
$2$-catégories} et de {\it pseudo-foncteur entre $2$-catégories}. Cette dernière
notion nous amène dans le domaine des $3$-catégories. Nous aimerions
éviter de devoir définir cela à tout prix ou presque !

\section{(2, 1)-catégories}
\label{section-2-1-categories}

\noindent
Certaines $2$-catégories ont la propriété que tous les $2$-morphismes sont des
isomorphismes. Celles-ci joueront un rôle important dans la suite et seront plus
faciles à utiliser.

\begin{definition}
\label{definition-2-1-category}
Une (stricte) {\it $(2, 1)$-catégorie} est une $2$-catégorie dans laquelle
tous les $2$-morphismes sont des isomorphismes.
\end{definition}

\begin{example}
\label{example-2-1-category-of-categories}
La $2$-catégorie $\textit{Cat}$, voir la remarque
\ref{remark-big-2-categories}, peut être transformée en $(2, 1)$-catégorie en
ne conservant comme $2$-morphismes que les isomorphismes de foncteurs.
\end{example}

\noindent
En fait, plus généralement toute $2$-catégorie $\mathcal{C}$ donne une
$(2, 1)$-catégorie en considérant la sous-$2$-catégorie $\mathcal{C}'$ avec
les mêmes objets et $1$-morphismes mais dont les $2$-morphismes sont les
$2$-morphismes inversibles de $\mathcal{C}$. Dans cette situation, nous dirons
« {\it soit $\mathcal{C}'$ la $(2, 1)$-catégorie associée à
$\mathcal{C}$} » ou une formulation analogue. Par exemple, la $(2, 1)$-catégorie de
groupoïdes désigne la $2$-catégorie dont les objets sont des groupoïdes, dont
les $1$-morphismes sont des foncteurs et dont les $2$-morphismes sont des
isomorphismes de foncteurs. Cet exemple est toutefois mal choisi, car une
transformation entre foncteurs entre groupoïdes est automatiquement un
isomorphisme !

\begin{remark}
\label{remark-other-2-categories}
Il existe donc des variantes de la construction de l'exemple
\ref{example-2-1-category-of-categories} ci-dessus dans lesquelles on considère la
$2$-catégorie des groupoïdes, ou des catégories fibrées en groupoïdes au-dessus d'une
catégorie fixe, ou des champs. Et ainsi de suite.
\end{remark}






\section{2-produits fibrés}
\label{section-2-fibre-products}

\noindent
Dans cette section, nous présentons les $2$-produits fibrés. Supposons que
$\mathcal{C}$ soit une 2-catégorie. On dit qu'un diagramme
$$
\xymatrix{
w \ar[r] \ar[d] & y \ar[d] \\
x \ar[r] & z }
$$
est 2-commutatif si les deux 1-morphismes $w \to y \to z$ et $w \to x \to z$ sont
2-isomorphes. Dans une 2-catégorie, il est plus naturel de demander la
2-commutativité des diagrammes que leur commutativité stricte. (En
effet, certains diront peut-être que nous ne devrions pas du tout travailler
avec des 2-catégories strictes, et dans une 2-catégorie « faible » la notion
de diagramme commutatif de 1-morphismes n'a même pas de sens.) En conséquence,
la notion de produit fibré doit être ajustée.

\medskip\noindent
Soit $\mathcal{C}$ une $2$-catégorie. Soient $x, y, z\in \Ob(\mathcal{C})$ et
$f\in \Mor_\mathcal{C}(x, z)$ et $g\in \Mor_{\mathcal C}(y, z)$. Afin de
définir le 2-produit fibré de $f$ et $g$, nous allons examiner les
diagrammes 2-commutatifs
$$
\xymatrix{
& w \ar[r]_a \ar[d]_b & x \ar[d]^{f} \\
& y \ar[r]^{g} & z. }
$$
Dans le cas des catégories ordinaires, le produit fibré est un objet final
dans la catégorie de tels diagrammes. En conséquence, un 2-produit fibré
est un objet final dans une 2-catégorie (voir la définition ci-dessous). La {\it
$2$-catégorie des diagrammes $2$-commutatifs sur $f$ et $g$} est la
$2$-catégorie définie comme suit :
\begin{enumerate}
\item Les objets sont des quadruples $(w, a, b, \phi)$ comme ci-dessus où
$\phi$ est un 2-morphisme inversible $\phi : f \circ a \to g \circ b$,
\item Les 1-morphismes de $(w', a', b', \phi')$ à $(w, a, b, \phi)$ sont
les triplets $(k : w' \to w, \alpha : a' \to a \circ k, \beta : b' \to b \circ k)$ tels que
$$
\xymatrix{
f \circ a'
\ar[rr]_{\text{id}_f \star \alpha}
\ar[d]_{\phi'}
& &
f \circ a \circ k
\ar[d]^{\phi \star \text{id}_k}
\\
g \circ b'
\ar[rr]^{\text{id}_g \star \beta}
& &
g \circ b \circ k
}
$$
est commutatif,
\item étant donné un deuxième $1$-morphisme $(k', \alpha', \beta') : (w'', a'', b'', \phi'') \to
(w', a', b', \phi')$, la composition des
$1$-morphismes est donnée par la règle
$$
(k, \alpha, \beta) \circ (k', \alpha', \beta') =
(k \circ k',
(\alpha \star \text{id}_{k'}) \circ \alpha',
(\beta \star \text{id}_{k'}) \circ \beta'),
$$
\item un 2-morphisme entre les $1$-morphismes $(k_i, \alpha_i, \beta_i)$, $i = 1, 2$ avec la même source et le même but est donné par un 2-morphisme
$\delta : k_1 \to k_2$ tel que
$$
\xymatrix{
a'
\ar[rd]_{\alpha_2}
\ar[r]_{\alpha_1} &
a \circ k_1
\ar[d]^{\text{id}_a \star \delta} &
&
b \circ k_1
\ar[d]_{\text{id}_b \star \delta} &
b'
\ar[l]^{\beta_1}
\ar[ld]^{\beta_2}
\\
&
a \circ k_2
&
&
b \circ k_2
&
}
$$
est commutatif,
\item la composition verticale des $2$-morphismes est donnée par la
composition verticale des morphismes $\delta$ dans $\mathcal{C}$, et
\item la composition horizontale du diagramme
$$
\xymatrix{
(w'', a'', b'', \phi'')
\rrtwocell^{(k'_1, \alpha'_1, \beta'_1)}_{(k'_2, \alpha'_2, \beta'_2)}{\delta'}
& &
(w', a', b', \phi')
\rrtwocell^{(k_1, \alpha_1, \beta_1)}_{(k_2, \alpha_2, \beta_2)}{\delta}
& &
(w, a, b, \phi)
}
$$
est donnée par le diagramme
$$
\xymatrix@C=12pc{
(w'', a'', b'', \phi'')
\rtwocell^{(k_1 \circ k'_1, (\alpha_1 \star \text{id}_{k'_1}) \circ \alpha'_1, (\beta_1 \star \text{id}_{k'_1}) \circ \beta'_1)}_{(k_2 \circ k'_2, (\alpha_2 \star \text{id}_{k'_2}) \circ \alpha'_2, (\beta_2 \star \text{id}_{k'_2}) \circ \beta'_2)}{\ \ \ \delta \star \delta'}
&
(w, a, b, \phi)
}
$$
\end{enumerate}
Notez que si $\mathcal{C}$ est en fait une $(2, 1)$-catégorie, les morphismes
$\alpha$ et $\beta$ dans (2) ci-dessus sont automatiquement également des
isomorphismes\footnote{En fait, il semble que dans le cas d'une $2$-catégorie,
l'on pourrait définir une autre 2-catégorie de diagrammes 2-commutatifs où le
sens des flèches $\alpha$, $\beta$ est inversé, ou encore où le sens d'une
seule d'entre elles est inversé. C'est pourquoi nous nous limiterons plus tard
aux $(2, 1)$-catégories.}. De plus, la $2$-catégorie des diagrammes
$2$-commutatifs est également une $(2, 1)$-catégorie si $\mathcal{C}$ est une
$(2, 1)$-catégorie.

\begin{definition}
\label{definition-final-object-2-category}
Un {\it objet final} d'une $(2, 1)$-catégorie $\mathcal{C}$ est un objet $x$
tel que
\begin{enumerate}
\item pour chaque $y \in \Ob(\mathcal{C})$ il existe un morphisme $y \to x$,
et
\item deux morphismes quelconques $y \to x$ sont isomorphes par un 2-morphisme
unique.
\end{enumerate}
\end{definition}

\noindent
Il est probable que dans le cas plus général des $2$-catégories, il existe
différentes versions d'objets finaux. Nous ne voulons pas entrer dans ce sujet
et c'est pourquoi nous définissons uniquement les $2$-produits fibrés dans
le cas des $(2, 1)$-catégories.

\begin{definition}
\label{definition-2-fibre-products}
Soit $\mathcal{C}$ une $(2, 1)$-catégorie. Soient $x, y, z\in \Ob(\mathcal{C})$ et $f\in \Mor_\mathcal{C}(x, z)$ et $g\in \Mor_{\mathcal C}(y, z)$. Un {\it 2-produit fibré de $f$ et $g$} est un objet final dans
la $2$-catégorie des diagrammes $2$-commutatifs décrite ci-dessus. Si un 2-produit
fibré existe, nous le noterons $x \times_z y\in \Ob(\mathcal{C})$, et
noterons les morphismes correspondants $p\in \Mor_{\mathcal C}(x \times_z y, x)$ et
$q\in \Mor_{\mathcal C}(x \times_z y, y)$ rendant le diagramme
$$
\xymatrix{
& x \times_z y \ar[r]^{p} \ar[d]_q & x \ar[d]^{f} \\
& y \ar[r]^{g} & z }
$$
2-commutatif ; nous noterons le 2-morphisme inversible qui en témoigne
$\psi : f \circ p \to g \circ q$.
\end{definition}

\noindent
On a donc la propriété universelle suivante : pour tout $w\in \Ob(\mathcal{C})$, tous les morphismes $a \in \Mor_{\mathcal C}(w, x)$ et $b \in \Mor_\mathcal{C}(w, y)$ munis d'un 2-isomorphisme $\phi : f \circ a \to g\circ b$
donné, il existe un $\gamma \in \Mor_{\mathcal C}(w, x \times_z y)$ tel que le
diagramme
$$
\xymatrix{
w\ar[rrrd]^a \ar@{-->}[rrd]_\gamma \ar[rrdd]_b & & \\
& & x \times_z y \ar[r]_p \ar[d]_q & x \ar[d]^{f} \\
& & y \ar[r]^{g} & z }
$$
soit 2-commutatif, de telle sorte que, pour des choix appropriés de $a \to p \circ \gamma$ et $b \to q \circ \gamma$, le diagramme
$$
\xymatrix{
f \circ a \ar[r] \ar[d]_\phi &
f \circ p \circ \gamma
\ar[d]^{\psi \star \text{id}_\gamma}
\\
g\circ b
\ar[r]
&
g \circ q \circ \gamma
}
$$
est commutatif. De plus, $\gamma$ est unique à
l'isomorphisme près. Bien entendu, les propriétés précises sont plus fines que
cela. Tous les cas de 2-produits fibrés dont nous aurons besoin plus tard
proviennent de l'exemple suivant de 2-produits fibrés dans la 2-catégorie de
catégories.

\begin{example}
\label{example-2-fibre-product-categories}
Soient $\mathcal{A}$, $\mathcal{B}$ et $\mathcal{C}$ des catégories. Soient $F : \mathcal{A} \to \mathcal{C}$ et $G : \mathcal{B} \to \mathcal{C}$ des
foncteurs. Nous définissons une catégorie $\mathcal{A} \times_\mathcal{C} \mathcal{B}$ comme suit :
\begin{enumerate}
\item un objet de $\mathcal{A} \times_\mathcal{C} \mathcal{B}$ est un triple
$(A, B, f)$, où $A\in \Ob(\mathcal{A})$, $B\in \Ob(\mathcal{B})$ et $f : F(A) \to G(B)$ est un isomorphisme dans $\mathcal{C}$,
\item un morphisme $(A, B, f) \to (A', B', f')$ est donné par une paire $(a, b)$, où $a : A \to A'$ est un morphisme dans $\mathcal{A}$, et $b : B \to B'$
est un morphisme dans $\mathcal{B}$ tel que le diagramme
$$
\xymatrix{
F(A) \ar[r]^f \ar[d]^{F(a)} & G(B) \ar[d]^{G(b)} \\
F(A') \ar[r]^{f'} & G(B')
}
$$
est commutatif.
\end{enumerate}
De plus, nous définissons les foncteurs $p : \mathcal{A} \times_\mathcal{C}\mathcal{B} \to \mathcal{A}$ et $q : \mathcal{A} \times_\mathcal{C}\mathcal{B} \to \mathcal{B}$ en posant
$$
p(A, B, f) = A, \quad q(A, B, f) = B,
$$
en d'autres termes, ce sont les foncteurs d'oubli. Nous définissons une
transformation de foncteurs $\psi : F \circ p \to G \circ q$. Sur l'objet $\xi = (A, B, f)$ elle est donnée par $\psi_\xi = f : F(p(\xi)) = F(A) \to G(B) = G(q(\xi))$.
\end{example}

\begin{lemma}
\label{lemma-2-fibre-product-categories}
Dans la $(2, 1)$-catégorie des catégories, des $2$-produits fibrés existent
et sont donnés par la construction de l'exemple
\ref{example-2-fibre-product-categories}.
\end{lemma}

\begin{proof}
Vérifions la propriété universelle : soit $\mathcal{W}$ une catégorie, soit $a : \mathcal{W} \to \mathcal{A}$ et $b : \mathcal{W} \to \mathcal{B}$ des
foncteurs, et soit $t : F \circ a \to G \circ b$ un isomorphisme de foncteurs.

\medskip\noindent
Considérons le foncteur $\gamma : \mathcal{W} \to \mathcal{A} \times_\mathcal{C}\mathcal{B}$ donné par $W \mapsto (a(W), b(W), t_W)$.
(La vérification qu'il s'agit d'un foncteur est omise.) De plus, considérons $\alpha : a \to p \circ \gamma$ et $\beta : b \to q \circ \gamma$ obtenus à partir des
identités $p \circ \gamma = a$ et $q \circ \gamma = b$. Il est alors clair que
$(\gamma, \alpha, \beta)$ est un morphisme de $(\mathcal{W}, a, b, t)$ vers $(\mathcal{A} \times_\mathcal{C} \mathcal{B}, p, q, \psi)$.

\medskip\noindent
Soit $(k, \alpha', \beta') :
(\mathcal{W}, a, b, t) \to (\mathcal{A} \times_\mathcal{C} \mathcal{B}, p, q, \psi)$ un deuxième morphisme de ce type. Pour un objet $W$
de $\mathcal{W}$ écrivons $k(W) = (a_k(W), b_k(W), t_{k, W})$. D'où $p(k(W)) = a_k(W)$, et ainsi de suite. La transformation $\alpha'$ correspond à une
famille fonctorielle de morphismes $\alpha' : a(W) \to a_k(W)$. Puisque nous
travaillons dans la $(2, 1)$-catégorie des catégories, en fait chacun de ces
morphismes $a(W) \to a_k(W)$ est un isomorphisme. Nous pouvons les utiliser
(ainsi que leurs analogues $b(W) \to b_k(W)$) pour obtenir des isomorphismes
$$
\delta_W :
\gamma(W) = (a(W), b(W), t_W)
\longrightarrow
(a_k(W), b_k(W), t_{k, W}) = k(W).
$$
Il est immédiat de montrer que $\delta$ définit un $2$-isomorphisme entre
$\gamma$ et $k$ dans la $2$-catégorie des diagrammes $2$-commutatifs, comme
souhaité.
\end{proof}

\begin{remark}
\label{remark-other-description-2-fibre-product}
Soient $\mathcal{A}$, $\mathcal{B}$ et $\mathcal{C}$ des catégories. Soient $F : \mathcal{A} \to \mathcal{C}$ et $G : \mathcal{B} \to \mathcal{C}$ des
foncteurs. Une autre construction, légèrement plus symétrique, d'un
$2$-produit fibré $\mathcal{A} \times_\mathcal{C} \mathcal{B}$ est la suivante. Un
objet est un quintuple $(A, B, C, a, b)$ où $A, B, C$ sont des objets de
$\mathcal{A}, \mathcal{B}, \mathcal{C}$ et où $a : F(A) \to C$ et $b : G(B) \to C$ sont des isomorphismes. Un morphisme $(A, B, C, a, b) \to (A', B', C', a', b')$ est donné par un triplet de morphismes $A \to A', B \to B', C \to C'$
compatible avec les morphismes $a, b, a', b'$. Nous pouvons prouver
directement que cette construction donne un $2$-produit fibré. Cependant, il est plus
facile d'observer que le foncteur $(A, B, C, a, b) \mapsto (A, B, b^{-1} \circ a)$ donne une équivalence entre la catégorie des quintuples et la catégorie
construite dans l'exemple \ref{example-2-fibre-product-categories}.
\end{remark}

\begin{lemma}
\label{lemma-functoriality-2-fibre-product}
Soit le diagramme
$$
\xymatrix{
& \mathcal{Y} \ar[d]_I \ar[rd]^K & \\
\mathcal{X} \ar[r]^H \ar[rd]^L &
\mathcal{Z} \ar[rd]^M & \mathcal{B} \ar[d]^G \\
& \mathcal{A} \ar[r]^F & \mathcal{C}
}
$$
$2$-commutatif de catégories. Un choix d'isomorphismes
$\alpha : G \circ K \to M \circ I$ et $\beta : M \circ H \to F \circ L$
détermine un morphisme
$$
\mathcal{X} \times_\mathcal{Z} \mathcal{Y}
\longrightarrow
\mathcal{A} \times_\mathcal{C} \mathcal{B}
$$
entre les $2$-produits fibrés associés à cette situation.
\end{lemma}

\begin{proof}
On utilise simplement le foncteur
$$
(X, Y, \phi) \longmapsto (L(X), K(Y),
\alpha^{-1}_Y \circ M(\phi) \circ \beta^{-1}_X)
$$
sur les objets et
$$
(a, b) \longmapsto (L(a), K(b))
$$
sur les morphismes.
\end{proof}

\begin{lemma}
\label{lemma-equivalence-2-fibre-product}
Plaçons-nous sous les hypothèses du lemme \ref{lemma-functoriality-2-fibre-product}.
\begin{enumerate}
\item Si $K$ et $L$ sont fidèles alors le morphisme $\mathcal{X} \times_\mathcal{Z} \mathcal{Y} \to
\mathcal{A} \times_\mathcal{C} \mathcal{B}$
est fidèle.
\item Si $K$ et $L$ sont pleinement fidèles et que $M$ est fidèle alors le
morphisme $\mathcal{X} \times_\mathcal{Z} \mathcal{Y} \to
\mathcal{A} \times_\mathcal{C} \mathcal{B}$ est pleinement fidèle.
\item Si $K$ et $L$ sont des équivalences et que $M$ est pleinement fidèle
alors le morphisme $\mathcal{X} \times_\mathcal{Z} \mathcal{Y} \to
\mathcal{A} \times_\mathcal{C} \mathcal{B}$ est une équivalence.
\end{enumerate}
\end{lemma}

\begin{proof}
Soient $(X, Y, \phi)$ et $(X', Y', \phi')$ des objets de $\mathcal{X} \times_\mathcal{Z} \mathcal{Y}$. Posons $Z = H(X)$ et identifions-le à
$I(Y)$ au moyen de $\phi$. Identifions également $M(Z)$ à $F(L(X))$ au moyen de $\beta_X$
et $M(Z)$ à $G(K(Y))$ au moyen de $\alpha_Y$. De même pour $Z' = H(X')$
et $M(Z')$. L'application sur les morphismes est l'application
$$
\xymatrix{
\Mor_\mathcal{X}(X, X')
\times_{\Mor_\mathcal{Z}(Z, Z')}
\Mor_\mathcal{Y}(Y, Y')
\ar[d] \\
\Mor_\mathcal{A}(L(X), L(X'))
\times_{\Mor_\mathcal{C}(M(Z), M(Z'))}
\Mor_\mathcal{B}(K(Y), K(Y'))
}
$$
Les parties (1) et (2) suivent donc. De plus, si $K$ et $L$ sont des
équivalences et que $M$ est pleinement fidèle, alors tout objet $(A, B, \phi)$ appartient à l'image essentielle pour les raisons suivantes : choisissons
$X$, $Y$ tels que $L(X) \cong A$ et $K(Y) \cong B$. Alors la pleine fidélité de
$M$ garantit que l’on peut trouver un isomorphisme $H(X) \cong I(Y)$. Certains
détails sont omis.
\end{proof}

\begin{lemma}
\label{lemma-associativity-2-fibre-product}
Soit le diagramme
$$
\xymatrix{
\mathcal{A} \ar[rd] & & \mathcal{C} \ar[ld] \ar[rd] & & \mathcal{E} \ar[ld] \\
& \mathcal{B} & & \mathcal{D}
}
$$
de catégories et de foncteurs. Alors il existe un
isomorphisme canonique
$$
(\mathcal{A} \times_\mathcal{B} \mathcal{C}) \times_\mathcal{D} \mathcal{E}
\cong
\mathcal{A} \times_\mathcal{B} (\mathcal{C} \times_\mathcal{D} \mathcal{E})
$$
de catégories.
\end{lemma}

\begin{proof}
On utilise simplement le foncteur
$$
((A, C, \phi), E, \psi)
\longmapsto
(A, (C, E, \psi), \phi)
$$
ce qui définit le foncteur voulu, si l'on voit ce que cela signifie.
\end{proof}

\noindent
Désormais, nous omettrons les parenthèses lorsqu'il s'agit de produits
fibrés de plus de 2 catégories.

\begin{lemma}
\label{lemma-triple-2-fibre-product-pr02}
Soit le diagramme
$$
\xymatrix{
\mathcal{A} \ar[rd] & & \mathcal{C} \ar[ld] \ar[rd] & & \mathcal{E} \ar[ld] \\
& \mathcal{B} \ar[rd]_F & & \mathcal{D} \ar[ld]^G \\
& & \mathcal{F} &
}
$$
commutatif de catégories et de foncteurs. Alors il existe un
foncteur canonique
$$
\text{pr}_{02} :
\mathcal{A} \times_\mathcal{B} \mathcal{C} \times_\mathcal{D} \mathcal{E}
\longrightarrow
\mathcal{A} \times_\mathcal{F} \mathcal{E}
$$
de catégories.
\end{lemma}

\begin{proof}
Si nous écrivons $\mathcal{A} \times_\mathcal{B} \mathcal{C}
\times_\mathcal{D} \mathcal{E}$ comme $(\mathcal{A} \times_\mathcal{B} \mathcal{C})
\times_\mathcal{D} \mathcal{E}$ alors nous pouvons simplement
utiliser le foncteur
$$
((A, C, \phi), E, \psi)
\longmapsto
(A, E, G(\psi) \circ F(\phi))
$$
ce qui définit le foncteur voulu, si l'on voit ce que cela signifie.
\end{proof}

\begin{lemma}
\label{lemma-2-fibre-product-erase-factor}
Soit le diagramme
$$
\mathcal{A} \to
\mathcal{B} \leftarrow \mathcal{C} \leftarrow \mathcal{D}
$$
de catégories et de foncteurs. Alors il existe un
isomorphisme canonique
$$
\mathcal{A} \times_\mathcal{B} \mathcal{C} \times_\mathcal{C} \mathcal{D}
\cong
\mathcal{A} \times_\mathcal{B} \mathcal{D}
$$
de catégories.
\end{lemma}

\begin{proof}
Omis.
\end{proof}

\noindent
Nous affirmons que cela signifie que l'on peut travailler avec ces
$2$-produits fibrés comme avec des produits fibrés ordinaires. Voici quelques
autres lemmes qui reviendront plus tard.

\begin{lemma}
\label{lemma-diagonal-1}
Soit le diagramme
$$
\xymatrix{
\mathcal{C}_3 \ar[r] \ar[d] & \mathcal{S} \ar[d]^\Delta \\
\mathcal{C}_1 \times \mathcal{C}_2 \ar[r]^{G_1 \times G_2} &
\mathcal{S} \times \mathcal{S}
}
$$
un $2$-produit fibré de catégories. Il existe alors un isomorphisme
canonique $\mathcal{C}_3 \cong
\mathcal{C}_1 \times_{G_1, \mathcal{S}, G_2} \mathcal{C}_2$.
\end{lemma}

\begin{proof}
Nous pouvons supposer que $\mathcal{C}_3$ est la catégorie $(\mathcal{C}_1 \times \mathcal{C}_2)\times_{\mathcal{S} \times \mathcal{S}}
\mathcal{S}$
construite dans l'exemple \ref{example-2-fibre-product-categories}. Un objet
est donc un triple $((X_1, X_2), S, \phi)$ où $\phi = (\phi_1, \phi_2) : (G_1(X_1), G_2(X_2)) \to (S, S)$ est un isomorphisme. On peut ainsi y associer
le triple $(X_1, X_2, \phi_2^{-1} \circ \phi_1)$. À l’inverse, si $(X_1, X_2, \psi)$ est un objet de $\mathcal{C}_1 \times_{G_1, \mathcal{S}, G_2} \mathcal{C}_2$, alors on peut lui associer le triple $((X_1, X_2), G_2(X_2), (\psi, \text{id}_{G_2(X_2)}))$. Nous affirmons que ces constructions
définissent des foncteurs mutuellement inverses. Nous omettons de décrire comment
traiter les morphismes et de montrer qu'ils sont mutuellement inverses.
\end{proof}

\begin{lemma}
\label{lemma-diagonal-2}
Soit le diagramme
$$
\xymatrix{
\mathcal{C}' \ar[r] \ar[d] & \mathcal{S} \ar[d]^\Delta \\
\mathcal{C} \ar[r]^{(G_1, G_2)} &
\mathcal{S} \times \mathcal{S}
}
$$
un $2$-produit fibré de catégories. Alors il existe un isomorphisme
canonique
$$
\mathcal{C}' \cong
(\mathcal{C} \times_{G_1, \mathcal{S}, G_2} \mathcal{C})
\times_{(p, q), \mathcal{C} \times \mathcal{C}, \Delta}
\mathcal{C}.
$$
\end{lemma}

\begin{proof}
Un objet du côté droit est donné par $((C_1, C_2, \phi), C_3, \psi)$ où $\phi : G_1(C_1) \to G_2(C_2)$ est un isomorphisme et $\psi = (\psi_1, \psi_2) : (C_1, C_2) \to (C_3, C_3)$ est un isomorphisme. On peut donc y associer le
triple $(C_3, G_1(C_1), (G_1(\psi_1^{-1}), \phi^{-1} \circ G_2(\psi_2^{-1})))$
qui est un objet de $\mathcal{C}'$. Détails omis.
\end{proof}

\begin{lemma}
\label{lemma-fibre-product-after-map}
Soient $\mathcal{A} \to \mathcal{C}$, $\mathcal{B} \to \mathcal{C}$ et
$\mathcal{C} \to \mathcal{D}$ des foncteurs entre catégories. Alors le
diagramme
$$
\xymatrix{
\mathcal{A} \times_\mathcal{C} \mathcal{B} \ar[d] \ar[r] &
\mathcal{A} \times_\mathcal{D} \mathcal{B} \ar[d] \\
\mathcal{C} \ar[r]^-{\Delta_{\mathcal{C}/\mathcal{D}}} &
\mathcal{C} \times_\mathcal{D} \mathcal{C}
}
$$
est $2$-cartésien.
\end{lemma}

\begin{proof}
Omis.
\end{proof}

\begin{lemma}
\label{lemma-base-change-diagonal}
Soit le diagramme
$$
\xymatrix{
\mathcal{U} \ar[d] \ar[r] & \mathcal{V} \ar[d] \\
\mathcal{X} \ar[r] & \mathcal{Y}
}
$$
un $2$-produit fibré de catégories. Alors le diagramme
$$
\xymatrix{
\mathcal{U} \ar[d] \ar[r] &
\mathcal{U} \times_\mathcal{V} \mathcal{U} \ar[d] \\
\mathcal{X} \ar[r] &
\mathcal{X} \times_\mathcal{Y} \mathcal{X}
}
$$
est $2$-cartésien.
\end{lemma}

\begin{proof}
Il s'agit d'un énoncé purement $2$-catégorique, valable dans
toute $(2, 1)$-catégorie avec des $2$-produits fibrés. Explicitement, cela
découle de la chaîne d’équivalences suivante :
\begin{align*}
\mathcal{X} \times_{(\mathcal{X} \times_\mathcal{Y} \mathcal{X})}
(\mathcal{U} \times_\mathcal{V} \mathcal{U})
& =
\mathcal{X} \times_{(\mathcal{X} \times_\mathcal{Y} \mathcal{X})}
((\mathcal{X} \times_\mathcal{Y} \mathcal{V})
\times_\mathcal{V} (\mathcal{X} \times_\mathcal{Y} \mathcal{V})) \\
& =
\mathcal{X} \times_{(\mathcal{X} \times_\mathcal{Y} \mathcal{X})}
(\mathcal{X} \times_\mathcal{Y} \mathcal{X}
\times_\mathcal{Y} \mathcal{V}) \\
& =
\mathcal{X} \times_\mathcal{Y} \mathcal{V} = \mathcal{U}
\end{align*}
voir les lemmes \ref{lemma-associativity-2-fibre-product} et
\ref{lemma-2-fibre-product-erase-factor}.
\end{proof}








\section{Catégories au-dessus de catégories}
\label{section-categories-over-categories}

\noindent
Dans cette section, nous avons un foncteur $p : \mathcal{S} \to \mathcal{C}$.
Nous considérons $\mathcal{S}$ comme placée au-dessus et $\mathcal{C}$ comme
placée au-dessous. Pour nous
assurer que tout le monde sait de quoi nous parlons, nous définissons la
$2$-catégorie des catégories au-dessus de $\mathcal{C}$.

\begin{definition}
\label{definition-categories-over-C}
Soit $\mathcal{C}$ une catégorie. La {\it $2$-catégorie des catégories au-dessus de
$\mathcal{C}$} est la $2$-catégorie définie comme suit :
\begin{enumerate}
\item Ses objets seront des foncteurs $p : \mathcal{S} \to \mathcal{C}$.
\item Ses $1$-morphismes $(\mathcal{S}, p) \to (\mathcal{S}', p')$ seront des
foncteurs $G : \mathcal{S} \to \mathcal{S}'$ tels que $p' \circ G = p$.
\item Ses $2$-morphismes $t : G \to H$ pour $G, H : (\mathcal{S}, p) \to (\mathcal{S}', p')$ seront des morphismes de foncteurs tels que $p'(t_x) = \text{id}_{p(x)}$ pour tous $x \in \Ob(\mathcal{S})$.
\end{enumerate}
Dans cette situation nous désignerons par
$$
\Mor_{\textit{Cat}/\mathcal{C}}(\mathcal{S}, \mathcal{S}')
$$
la catégorie des $1$-morphismes entre $(\mathcal{S}, p)$ et $(\mathcal{S}', p')$.
\end{definition}

\noindent
Dans cette $2$-catégorie, nous définissons la composition horizontale et
verticale exactement comme cela est fait pour $\textit{Cat}$ dans la section
\ref{section-formal-cat-cat}. Les axiomes d'une $2$-catégorie sont satisfaits
pour la même raison que ceux de $\textit{Cat}$. Pour voir cela, on peut
également utiliser le fait que les axiomes sont valables dans $\textit{Cat}$
et vérifier des assertions telles que « la composition verticale des
$2$-morphismes au-dessus de $\mathcal{C}$ donne un autre $2$-morphisme
au-dessus de $\mathcal{C}$ ». C'est
clair.

\medskip\noindent
Par analogie avec la fibre d'une application d'espaces, nous avons la
notion de catégorie fibre, ainsi que quelques notions de relèvement associées à
cette situation.

\begin{definition}
\label{definition-fibre-category}
Soit $\mathcal{C}$ une catégorie. Soit $p : \mathcal{S} \to \mathcal{C}$ une
catégorie au-dessus de $\mathcal{C}$.
\begin{enumerate}
\item La {\it catégorie fibre} au-dessus d'un objet $U\in \Ob(\mathcal{C})$ est la
catégorie $\mathcal{S}_U$ avec des objets
$$
\Ob(\mathcal{S}_U) = \{x\in \Ob(\mathcal{S}) :
p(x) = U\}
$$
et morphismes
$$
\Mor_{\mathcal{S}_U}(x, y) = \{ \phi \in \Mor_\mathcal{S}(x, y) :
p(\phi) = \text{id}_U\}.
$$
\item Un {\it relèvement} d'un objet $U \in \Ob(\mathcal{C})$ est un objet
$x\in \Ob(\mathcal{S})$ tel que $p(x) = U$, c'est-à-dire que $x\in \Ob(\mathcal{S}_U)$. Nous dirons aussi parfois que {\it $x$ est au-dessus de
$U$}.
\item De même, un {\it relèvement} d'un morphisme $f : V \to U$ dans $\mathcal{C}$
est un morphisme $\phi : y \to x$ dans $\mathcal{S}$ tel que $p(\phi) = f$. On
dit parfois que {\it $\phi$ est au-dessus de $f$}.
\end{enumerate}
\end{definition}

\noindent
Nous pourrions faire ici quelques observations. Par exemple, si $F : (\mathcal{S}, p) \to (\mathcal{S}', p')$ est un $1$-morphisme de catégories
au-dessus de $\mathcal{C}$, alors $F$ induit des foncteurs entre catégories
fibres $F : \mathcal{S}_U \to \mathcal{S}'_U$. Il en va de même pour les
$2$-morphismes.

\medskip\noindent
Voici le lemme incontournable décrivant le $2$-produit fibré dans la
$(2, 1)$-catégorie des catégories au-dessus de $\mathcal{C}$.

\begin{lemma}
\label{lemma-2-product-categories-over-C}
Soit $\mathcal{C}$ une catégorie. La $(2, 1)$-catégorie des catégories
au-dessus de $\mathcal{C}$ admet des 2-produits fibrés. Supposons que $F : \mathcal{X} \to \mathcal{S}$ et $G : \mathcal{Y} \to \mathcal{S}$ soient des
morphismes de catégories au-dessus de $\mathcal{C}$. Un 2-produit fibré explicite
$\mathcal{X} \times_\mathcal{S}\mathcal{Y}$ est donné par la description
suivante :
\begin{enumerate}
\item un objet de $\mathcal{X} \times_\mathcal{S} \mathcal{Y}$ est un
quadruple $(U, x, y, f)$, où $U \in \Ob(\mathcal{C})$, $x\in \Ob(\mathcal{X}_U)$, $y\in \Ob(\mathcal{Y}_U)$ et $f : F(x) \to G(y)$ est un
isomorphisme dans $\mathcal{S}_U$,
\item un morphisme $(U, x, y, f) \to (U', x', y', f')$ est donné par une paire
$(a, b)$, où $a : x \to x'$ est un morphisme dans $\mathcal{X}$, et $b : y \to y'$ est un morphisme dans $\mathcal{Y}$ tel que
\begin{enumerate}
\item $a$ et $b$ induisent le même morphisme $U \to U'$, et
\item le diagramme
$$
\xymatrix{
F(x) \ar[r]^f \ar[d]^{F(a)} & G(y) \ar[d]^{G(b)} \\
F(x') \ar[r]^{f'} & G(y')
}
$$
est commutatif.
\end{enumerate}
\end{enumerate}
Les foncteurs $p : \mathcal{X} \times_\mathcal{S}\mathcal{Y} \to \mathcal{X}$
et $q : \mathcal{X} \times_\mathcal{S}\mathcal{Y} \to \mathcal{Y}$ sont les
foncteurs d'oubli dans ce cas. La transformation $\psi : F \circ p \to
G \circ q$ est donnée sur l'objet $\xi = (U, x, y, f)$ par $\psi_\xi = f : F(p(\xi)) = F(x) \to G(y) = G(q(\xi))$.
\end{lemma}

\begin{proof}
Vérifions la propriété universelle : soit $p_\mathcal{W} : \mathcal{W}\to \mathcal{C}$ une catégorie au-dessus de $\mathcal{C}$, soient $X : \mathcal{W} \to \mathcal{X}$ et $Y : \mathcal{W} \to \mathcal{Y}$ des foncteurs au-dessus de
$\mathcal{C}$, et soit $t : F \circ X \to G \circ Y$ un isomorphisme de
foncteurs au-dessus de $\mathcal{C}$. Le foncteur recherché $\gamma : \mathcal{W} \to \mathcal{X} \times_\mathcal{S} \mathcal{Y}$ est donné par $W \mapsto (p_\mathcal{W}(W), X(W), Y(W), t_W)$. Détails omis ; comparer avec le lemme
\ref{lemma-2-fibre-product-categories}.
\end{proof}

\begin{example}
\label{example-different-2-fibre-products}
Les constructions de $2$-produits fibrés dans la 2-catégorie des catégories
au-dessus d'une catégorie, données dans le lemme \ref{lemma-2-product-categories-over-C},
et dans la 2-catégorie des catégories, données dans le lemme
\ref{lemma-2-fibre-product-categories} (comme dans l'exemple
\ref{example-2-fibre-product-categories}) produisent des résultats non
équivalents en général. Plus précisément, soit $\mathcal{S}$ le groupoïde
avec un objet et deux flèches, et soit $\mathcal{X}$ la catégorie discrète
avec un objet. En prenant le $2$-produit fibré $\mathcal{X} \times_\mathcal{S} \mathcal{X}$ en tant que catégories, on obtient la catégorie
discrète avec deux objets. Cependant, si nous considérons toutes ces catégories
comme des catégories au-dessus de $\mathcal{S}$, le $2$-produit fibré $\mathcal{X} \times_\mathcal{S} \mathcal{X}$ en tant que catégories au-dessus de $\mathcal{S}$ est
la catégorie discrète avec un objet. La différence est que (dans la notation
du lemme \ref{lemma-2-product-categories-over-C}), nous étions autorisés à
choisir n'importe quel isomorphisme de comparaison $f$ dans la première
situation, mais nous ne pouvions choisir la flèche d'identité que dans la
seconde situation.
\end{example}

\begin{lemma}
\label{lemma-fibre-2-fibre-product-categories-over-C}
Soit $\mathcal{C}$ une catégorie. Soient $f : \mathcal{X} \to \mathcal{S}$ et $g : \mathcal{Y} \to \mathcal{S}$ des morphismes de catégories au-dessus de $\mathcal{C}$.
Pour tout objet $U$ de $\mathcal{C}$ nous avons l'identité suivante entre les
catégories fibres
$$
\left(\mathcal{X} \times_\mathcal{S}\mathcal{Y}\right)_U
=
\mathcal{X}_U \times_{\mathcal{S}_U} \mathcal{Y}_U
$$
\end{lemma}

\begin{proof}
Omis.
\end{proof}

\section{Catégories fibrées}
\label{section-fibred-categories}

\noindent
Il convient de donner une très brève présentation des catégories fibrées.

\medskip\noindent
Soit $p : \mathcal{S} \to \mathcal{C}$ une catégorie au-dessus de
$\mathcal{C}$. Étant donné un objet $x \in \mathcal{S}$ tel que $p(x) = U$,
et un morphisme $f : V \to U$, nous pouvons essayer de former une sorte de
« produit fibré $V \times_U x$ » (ou un {\it changement de base} de $x$
par $V \to U$). En effet, un morphisme d'un objet $z \in \mathcal{S}$
vers « $V \times_U x$ » devrait être donné par une paire
$(\varphi, g)$, où
$\varphi : z \to x$, $g : p(z) \to V$ et
$p(\varphi) = f \circ g$. Voici la situation :
$$
\xymatrix{
z \ar@{~>}[d]^p \ar@{-}[r] &
? \ar[r] \ar@{~>}[d]^p &
x \ar@{~>}[d]^p \\
p(z) \ar[r] & V \ar[r]^f & U
}
$$
Si un tel morphisme $V \times_U x \to x$ existe, on l'appelle un morphisme
fortement cartésien.

\begin{definition}
\label{definition-cartesian-over-C}
Soit $\mathcal{C}$ une catégorie.
Soit $p : \mathcal{S} \to \mathcal{C}$ une catégorie au-dessus de
$\mathcal{C}$. Un {\it morphisme fortement cartésien}, ou plus précisément
un {\it morphisme fortement $\mathcal{C}$-cartésien}, est un
morphisme $\varphi : y \to x$ de $\mathcal{S}$ tel que, pour tout
$z \in \Ob(\mathcal{S})$, l'application
$$
\Mor_\mathcal{S}(z, y)
\longrightarrow
\Mor_\mathcal{S}(z, x)
\times_{\Mor_\mathcal{C}(p(z), p(x))}
\Mor_\mathcal{C}(p(z), p(y)),
$$
donnée par $\psi \longmapsto (\varphi \circ \psi, p(\psi))$,
soit bijective.
\end{definition}

\noindent
Notons que, par le lemme de Yoneda \ref{lemma-yoneda}, étant donnés
$x \in \Ob(\mathcal{S})$ au-dessus de $U \in \Ob(\mathcal{C})$
et un morphisme $f : V \to U$ de $\mathcal{C}$, s'il existe un
morphisme fortement cartésien $\varphi : y \to x$ tel que
$p(\varphi) = f$, alors $(y, \varphi)$ est unique à isomorphisme unique près.
Cela résulte clairement de la définition ci-dessus, puisque le foncteur
$$
z
\longmapsto
\Mor_\mathcal{S}(z, x)
\times_{\Mor_\mathcal{C}(p(z), U)}
\Mor_\mathcal{C}(p(z), V)
$$
ne dépend que des données $(x, U, f : V \to U)$. Ainsi, nous emploierons
parfois $V \times_U x \to x$ ou $f^*x \to x$ pour désigner un morphisme
fortement cartésien qui relève $f$.

\begin{lemma}
\label{lemma-composition-cartesian}
Soit $\mathcal{C}$ une catégorie.
Soit $p : \mathcal{S} \to \mathcal{C}$ une catégorie au-dessus de
$\mathcal{C}$.
\begin{enumerate}
\item Le composé de deux morphismes fortement cartésiens est fortement
cartésien.
\item Tout isomorphisme de $\mathcal{S}$ est fortement cartésien.
\item Tout morphisme fortement cartésien $\varphi$ tel que $p(\varphi)$
soit un isomorphisme est un isomorphisme.
\end{enumerate}
\end{lemma}

\begin{proof}
Démonstration de (1). Soient $\varphi : y \to x$ et
$\psi : z \to y$ des morphismes fortement cartésiens. Soit $t$ un objet
quelconque de $\mathcal{S}$. Alors
\begin{align*}
& \Mor_\mathcal{S}(t, z) \\
& =
\Mor_\mathcal{S}(t, y)
\times_{\Mor_\mathcal{C}(p(t), p(y))}
\Mor_\mathcal{C}(p(t), p(z)) \\
& =
\Mor_\mathcal{S}(t, x)
\times_{\Mor_\mathcal{C}(p(t), p(x))}
\Mor_\mathcal{C}(p(t), p(y))
\times_{\Mor_\mathcal{C}(p(t), p(y))}
\Mor_\mathcal{C}(p(t), p(z)) \\
& =
\Mor_\mathcal{S}(t, x)
\times_{\Mor_\mathcal{C}(p(t), p(x))}
\Mor_\mathcal{C}(p(t), p(z))
\end{align*}
et donc $z \to x$ est fortement cartésien.

\medskip\noindent
Démonstration de (2). Soit $y \to x$ un isomorphisme. Alors
$p(y) \to p(x)$ est également un isomorphisme. Par conséquent,
$\Mor_\mathcal{C}(p(z), p(y)) \to
\Mor_\mathcal{C}(p(z), p(x))$
est une bijection. Ainsi,
$\Mor_\mathcal{S}(z, x)
\times_{\Mor_\mathcal{C}(p(z), p(x))}
\Mor_\mathcal{C}(p(z), p(y))$ est en bijection avec
$\Mor_\mathcal{S}(z, x)$.
L'application affichée dans la
définition \ref{definition-cartesian-over-C}
est donc bijective puisque $y \to x$ est un isomorphisme, et nous en
concluons que $y \to x$ est fortement cartésien.

\medskip\noindent
Démonstration de (3). Supposons que $\varphi : y \to x$ soit fortement
cartésien et que $p(\varphi) : p(y) \to p(x)$ soit un isomorphisme.
En appliquant la définition avec $z = x$, on voit que
$(\text{id}_x, p(\varphi)^{-1})$ provient d'un unique morphisme
$\chi : x \to y$. Nous omettons la vérification que $\chi$ est l'inverse
de $\varphi$.
\end{proof}

\begin{lemma}
\label{lemma-cartesian-over-cartesian}
Soient $F : \mathcal{A} \to \mathcal{B}$ et
$G : \mathcal{B} \to \mathcal{C}$ des foncteurs composables entre
catégories. Soit $x \to y$ un morphisme de $\mathcal{A}$. Si
$x \to y$ est fortement $\mathcal{B}$-cartésien et si
$F(x) \to F(y)$ est fortement $\mathcal{C}$-cartésien, alors
$x \to y$ est fortement $\mathcal{C}$-cartésien.
\end{lemma}

\begin{proof}
Cela résulte directement de la définition.
\end{proof}

\begin{lemma}
\label{lemma-strongly-cartesian-fibre-product}
Soit $\mathcal{C}$ une catégorie.
Soit $p : \mathcal{S} \to \mathcal{C}$ une catégorie au-dessus de
$\mathcal{C}$. Soient $x \to y$ et $z \to y$ des morphismes de
$\mathcal{S}$. Supposons que
\begin{enumerate}
\item $x \to y$ soit fortement cartésien,
\item $p(x) \times_{p(y)} p(z)$ existe, et
\item il existe un morphisme fortement cartésien $a : w \to z$ dans
$\mathcal{S}$ tel que $p(w) = p(x) \times_{p(y)} p(z)$ et
$p(a) = \text{pr}_2 : p(x) \times_{p(y)} p(z) \to p(z)$.
\end{enumerate}
Alors le produit fibré $x \times_y z$ existe et est isomorphe à $w$.
\end{lemma}

\begin{proof}
Comme $x \to y$ est fortement cartésien, il existe un unique morphisme
$b : w \to x$ tel que $p(b) = \text{pr}_1$. Pour voir que $w$ est le
produit fibré, calculons
\begin{align*}
& \Mor_\mathcal{S}(t, w) \\
& = \Mor_\mathcal{S}(t, z)
\times_{\Mor_\mathcal{C}(p(t), p(z))}
\Mor_\mathcal{C}(p(t), p(w)) \\
& = \Mor_\mathcal{S}(t, z)
\times_{\Mor_\mathcal{C}(p(t), p(z))}
(\Mor_\mathcal{C}(p(t), p(x))
\times_{\Mor_\mathcal{C}(p(t), p(y))}
\Mor_\mathcal{C}(p(t), p(z))) \\
& = \Mor_\mathcal{S}(t, z)
\times_{\Mor_\mathcal{C}(p(t), p(y))}
\Mor_\mathcal{C}(p(t), p(x)) \\
& = \Mor_\mathcal{S}(t, z)
\times_{\Mor_\mathcal{S}(t, y)}
\Mor_\mathcal{S}(t, y)
\times_{\Mor_\mathcal{C}(p(t), p(y))}
\Mor_\mathcal{C}(p(t), p(x)) \\
& = \Mor_\mathcal{S}(t, z)
\times_{\Mor_\mathcal{S}(t, y)}
\Mor_\mathcal{S}(t, x)
\end{align*}
comme voulu. La première égalité vaut parce que $a : w \to z$ est
fortement cartésien, et la dernière parce que $x \to y$ est fortement
cartésien.
\end{proof}

\begin{definition}
\label{definition-fibred-category}
Soit $\mathcal{C}$ une catégorie.
Soit $p : \mathcal{S} \to \mathcal{C}$ une catégorie au-dessus de
$\mathcal{C}$. On dit que $\mathcal{S}$ est une
{\it catégorie fibrée au-dessus de $\mathcal{C}$} si, pour tout
$x \in \Ob(\mathcal{S})$ au-dessus de
$U \in \Ob(\mathcal{C})$ et tout morphisme $f : V \to U$ de
$\mathcal{C}$, il existe un morphisme fortement cartésien
$f^*x \to x$ au-dessus de $f$.
\end{definition}

\noindent
Supposons que $p : \mathcal{S} \to \mathcal{C}$ soit une catégorie fibrée.
Pour tous $f : V \to U$ et $x\in \Ob(\mathcal{S}_U)$ comme dans la
définition, nous pouvons choisir un morphisme fortement cartésien
$f^\ast x \to x$ au-dessus de $f$. Par l'axiome du choix, nous pouvons
choisir simultanément $f^*x \to x$ pour tous les
$f: V \to U = p(x)$. Nous affirmons que, pour tout morphisme
$\phi : x \to x'$ dans $\mathcal{S}_U$ et tout $f : V \to U$, il existe
un unique morphisme $f^\ast \phi : f^\ast x \to f^\ast x'$ dans
$\mathcal{S}_V$ tel que
$$
\xymatrix{
f^\ast x \ar[r]_{f^\ast \phi} \ar[d] & f^\ast x' \ar[d] \\
x \ar[r]^{\phi} & x' }
$$
soit commutatif. En effet, la flèche existe et est unique parce que
$f^*x' \to x'$ est fortement cartésien. L'unicité de cette flèche garantit
que $f^\ast$ (désormais également défini sur les morphismes) est un
foncteur $ f^\ast : \mathcal{S}_U \to \mathcal{S}_V$.

\begin{definition}
\label{definition-pullback-functor-fibred-category}
Supposons que $p : \mathcal{S} \to \mathcal{C}$ soit une catégorie fibrée.
\begin{enumerate}
\item Un {\it choix de changements de base}\footnote{Cette terminologie
n'est probablement pas standard. Dans certains textes, on parle de
« clivage », mais ce terme évoque une image inadéquate. Peut-être
« opération de clivage » conviendrait-il mieux. Une notion connexe est
celle de « scindage », mais, dans beaucoup de textes, un « scindage »
désigne un choix de changements de base tel que
$g^*f^* = (f \circ g)^*$ pour toute paire de morphismes composables.
Comparer aussi à la définition \ref{definition-split-fibred-category}.}
pour $p : \mathcal{S} \to \mathcal{C}$
consiste à choisir un morphisme fortement cartésien
$f^\ast x \to x$ au-dessus de $f$ pour tout morphisme
$f: V \to U$ de $\mathcal{C}$ et tout $x \in \Ob(\mathcal{S}_U)$.
\item Étant donné un choix de changements de base, pour tout morphisme
$f : V \to U$ de $\mathcal{C}$, le foncteur
$f^* : \mathcal{S}_U \to \mathcal{S}_V$ décrit ci-dessus est appelé
{\it foncteur de changement de base} (associé aux choix
$f^*x \to x$ faits ci-dessus).
\end{enumerate}
\end{definition}

\noindent
Nous pouvons évidemment toujours supposer que notre choix de changements
de base vérifie $\text{id}_U^*x = x$, bien qu'en pratique cette propriété
soit inutile si l'on n'impose pas d'autres hypothèses aux changements de base.

\begin{lemma}
\label{lemma-fibred}
Supposons que $p : \mathcal{S} \to \mathcal{C}$ soit une catégorie fibrée.
Supposons que l'on ait choisi des changements de base pour
$p : \mathcal{S} \to \mathcal{C}$.
\begin{enumerate}
\item Pour toute paire de morphismes composables $f : V \to U$,
$g : W \to V$, il existe un unique isomorphisme
$$
\alpha_{g, f} :
(f \circ g)^\ast
\longrightarrow
g^\ast \circ f^\ast
$$
de foncteurs $\mathcal{S}_U \to \mathcal{S}_W$ tel que, pour tout
$y\in \Ob(\mathcal{S}_U)$, le diagramme suivant soit commutatif :
$$
\xymatrix{
g^\ast f^\ast y \ar[r]
&
f^\ast y \ar[d] \\
(f \circ g)^\ast y \ar[r]
\ar[u]^{(\alpha_{g, f})_y}
&
y
}
$$
\item Si $f = \text{id}_U$, il existe un isomorphisme canonique
$\alpha_U : \text{id} \to (\text{id}_U)^*$ de foncteurs
$\mathcal{S}_U \to \mathcal{S}_U$.
\item Le quadruplet
$(U \mapsto \mathcal{S}_U, f \mapsto f^*, \alpha_{g, f}, \alpha_U)$
définit un pseudo-foncteur de $\mathcal{C}^{opp}$ vers la
$(2, 1)$-catégorie des catégories, voir la
définition \ref{definition-functor-into-2-category}.
\end{enumerate}
\end{lemma}

\begin{proof}
En fait, il est clair que le diagramme commutatif de la partie (1)
détermine de manière unique le morphisme
$(\alpha_{g, f})_y$ dans la catégorie fibre
$\mathcal{S}_W$. C'est un isomorphisme, puisque tant le morphisme
$(f \circ g)^*y \to y$ que le composé
$g^*f^*y \to f^*y \to y$ sont des morphismes fortement cartésiens
relevant $f \circ g$ (voir la discussion qui suit la
définition \ref{definition-cartesian-over-C} et le
lemme \ref{lemma-composition-cartesian}). De même, comme
$\text{id}_x : x \to x$ est manifestement fortement cartésien
au-dessus de $\text{id}_U$ (avec $U = p(x)$), nous voyons qu'il existe
un isomorphisme $(\alpha_U)_x : x \to (\text{id}_U)^*x$.
(Nous aurions évidemment pu supposer au préalable que $f^*x = x$
chaque fois que $f$ est un morphisme identité, mais il vaut mieux,
par souci de généralité, ne pas faire cette hypothèse.)
Nous omettons la vérification que $\alpha_{g, f}$ et
$\alpha_U$ ainsi obtenus sont des transformations de foncteurs.
Nous omettons aussi la vérification de (3).
\end{proof}

\begin{lemma}
\label{lemma-fibred-equivalent}
Soit $\mathcal{C}$ une catégorie.
Soient $\mathcal{S}_1$, $\mathcal{S}_2$ des catégories au-dessus de
$\mathcal{C}$. Supposons que $\mathcal{S}_1$ et $\mathcal{S}_2$ soient
équivalentes comme catégories au-dessus de $\mathcal{C}$.
Alors $\mathcal{S}_1$ est fibrée au-dessus de $\mathcal{C}$ si et seulement
si $\mathcal{S}_2$ est fibrée au-dessus de $\mathcal{C}$.
\end{lemma}

\begin{proof}
Notons $p_i : \mathcal{S}_i \to \mathcal{C}$ les foncteurs donnés.
Soient $F : \mathcal{S}_1 \to \mathcal{S}_2$,
$G : \mathcal{S}_2 \to \mathcal{S}_1$ des foncteurs au-dessus de
$\mathcal{C}$, et soient
$i : F \circ G \to \text{id}_{\mathcal{S}_2}$,
$j : G \circ F \to \text{id}_{\mathcal{S}_1}$ des isomorphismes de
foncteurs au-dessus de $\mathcal{C}$.
Nous affirmons que, dans ce cas, $F$ envoie les morphismes fortement
cartésiens sur des morphismes fortement cartésiens. En effet, supposons que
$\varphi : y \to x$ soit fortement cartésien dans $\mathcal{S}_1$.
Posons $f : V \to U$ ; il s'agit de $p_1(\varphi)$. Supposons que
$z' \in \Ob(\mathcal{S}_2)$, avec $W = p_2(z')$, et que l'on se soit donné
$g : W \to V$ et $\psi' : z' \to F(x)$ tels que
$p_2(\psi') = f \circ g$. Alors
$$
\psi = j \circ G(\psi') : G(z') \to G(F(x)) \to x
$$
est un morphisme dans $\mathcal{S}_1$ tel que
$p_1(\psi) = f \circ g$.
Par hypothèse, il existe donc un unique morphisme $\xi : G(z') \to y$
au-dessus de $g$ tel que $\psi = \varphi \circ \xi$. Celui-ci fournit à son
tour un morphisme
$$
\xi' = F(\xi) \circ i^{-1} : z' \to F(G(z')) \to F(y)
$$
au-dessus de $g$ tel que $\psi' = F(\varphi) \circ \xi'$. Nous omettons
la vérification que $\xi'$ est unique.
\end{proof}

\noindent
Le lemme \ref{lemma-fibred-equivalent} permet de conclure que les équivalences
envoient les morphismes fortement cartésiens sur des morphismes fortement
cartésiens. Cela peut toutefois être faux pour un foncteur quelconque entre
catégories fibrées au-dessus de $\mathcal{C}$. Nous définissons donc la
$2$-catégorie des catégories fibrées comme suit.

\begin{definition}
\label{definition-fibred-categories-over-C}
Soit $\mathcal{C}$ une catégorie.
La {\it $2$-catégorie des catégories fibrées au-dessus de $\mathcal{C}$}
est la sous-$2$-catégorie de la $2$-catégorie des catégories
au-dessus de $\mathcal{C}$ (voir la
définition \ref{definition-categories-over-C}) définie comme suit :
\begin{enumerate}
\item Ses objets sont les catégories fibrées
$p : \mathcal{S} \to \mathcal{C}$.
\item Ses $1$-morphismes $(\mathcal{S}, p) \to (\mathcal{S}', p')$
sont les foncteurs $G : \mathcal{S} \to \mathcal{S}'$ tels que
$p' \circ G = p$ et tels que $G$ envoie les morphismes fortement cartésiens
sur des morphismes fortement cartésiens.
\item Ses $2$-morphismes $t : G \to H$, pour
$G, H : (\mathcal{S}, p) \to (\mathcal{S}', p')$, sont les morphismes de
foncteurs tels que $p'(t_x) = \text{id}_{p(x)}$
pour tout $x \in \Ob(\mathcal{S})$.
\end{enumerate}
Dans cette situation, nous noterons
$$
\Mor_{\textit{Fib}/\mathcal{C}}(\mathcal{S}, \mathcal{S}')
$$
la catégorie des $1$-morphismes entre
$(\mathcal{S}, p)$ et $(\mathcal{S}', p')$.
\end{definition}

\noindent
Notons la condition imposée aux $1$-morphismes.
Notons aussi qu'il s'agit d'une véritable $2$-catégorie et non d'une
$(2, 1)$-catégorie. Ainsi, lorsque nous prenons des $2$-produits fibrés,
nous passons d'abord à la $(2, 1)$-catégorie associée.

\begin{lemma}
\label{lemma-2-product-fibred-categories-over-C}
Soit $\mathcal{C}$ une catégorie.
La $(2, 1)$-catégorie des catégories fibrées au-dessus de $\mathcal{C}$
possède des 2-produits fibrés, qui sont décrits comme dans le
lemme \ref{lemma-2-product-categories-over-C}.
\end{lemma}

\begin{proof}
Il s'agit essentiellement de montrer que, si
$F : \mathcal{X} \to \mathcal{S}$ et
$G : \mathcal{Y} \to \mathcal{S}$ sont des morphismes de catégories
fibrées au-dessus de $\mathcal{C}$, alors la catégorie
$\mathcal{X} \times_\mathcal{S} \mathcal{Y}$
décrite dans le lemme \ref{lemma-2-product-categories-over-C} est fibrée.
Montrons que $\mathcal{X} \times_\mathcal{S} \mathcal{Y}$
possède suffisamment de morphismes fortement cartésiens.
Soit donc $(U, x, y, \phi)$ un objet de
$\mathcal{X} \times_\mathcal{S} \mathcal{Y}$.
Soit $f : V \to U$ un morphisme de $\mathcal{C}$.
Choisissons des morphismes fortement cartésiens $a : f^*x \to x$ dans
$\mathcal{X}$ au-dessus de $f$ et $b : f^*y \to y$ dans $\mathcal{Y}$
au-dessus de $f$.
Par hypothèse, $F(a)$ et $G(b)$ sont fortement cartésiens.
Comme $\phi : F(x) \to G(y)$ est un isomorphisme, l'unicité des morphismes
fortement cartésiens fournit un unique isomorphisme
$f^*\phi : F(f^*x) \to G(f^*y)$ tel que
$G(b) \circ f^*\phi = \phi \circ F(a)$. Autrement dit,
$(a, b) : (V, f^*x, f^*y, f^*\phi) \to (U, x, y, \phi)$
est un morphisme dans $\mathcal{X} \times_\mathcal{S} \mathcal{Y}$.
Nous omettons la vérification que c'est un morphisme fortement cartésien
(et que ce sont en fait les seuls morphismes fortement cartésiens).
\end{proof}

\begin{lemma}
\label{lemma-cute}
Soit $\mathcal{C}$ une catégorie. Soit $U \in \Ob(\mathcal{C})$.
Si $p : \mathcal{S} \to \mathcal{C}$ est une catégorie fibrée
et si $p$ se factorise par $p' : \mathcal{S} \to \mathcal{C}/U$,
alors $p' : \mathcal{S} \to \mathcal{C}/U$ est une catégorie fibrée.
\end{lemma}

\begin{proof}
Supposons que $\varphi : x' \to x$ soit fortement cartésien relativement
à $p$. Nous affirmons que $\varphi$ est également fortement cartésien
relativement à $p'$. Posons $g = p'(\varphi)$, de sorte que
$g : V'/U \to V/U$ pour certains morphismes $f : V \to U$ et
$f' : V' \to U$.
Soit $z \in \Ob(\mathcal{S})$. Posons $p'(z) = (W \to U)$.
Pour montrer que $\varphi$ est fortement cartésien pour $p'$, nous devons
montrer que l'application
$$
\Mor_\mathcal{S}(z, x')
\longrightarrow
\Mor_\mathcal{S}(z, x)
\times_{\Mor_{\mathcal{C}/U}(W/U, V/U)}
\Mor_{\mathcal{C}/U}(W/U, V'/U),
$$
donnée par $\psi' \longmapsto (\varphi \circ \psi', p'(\psi'))$,
est bijective. Soit $(\psi, h)$ un élément du membre de droite. Alors, en
particulier, $g \circ h = p'(\psi)$ et, puisque $\varphi$ est fortement
cartésien, nous obtenons un unique morphisme $\psi' : z \to x'$ tel que
$\psi = \varphi \circ \psi'$ et tel que $p(\psi')$ égale l'application
sous-jacente à $h$. Or $p'(\psi') : W/U \to V'/U$ est un morphisme
dont l'application sous-jacente $W \to V'$ est l'application sous-jacente
de $h$ ; il est donc égal à $h$ comme morphisme dans $\mathcal{C}/U$.
Ainsi, $\psi'$ est l'unique morphisme $z \to x'$ qui s'envoie sur la paire
donnée $(\psi, h)$. Cela démontre l'affirmation.

\medskip\noindent
Enfin, soient $g : V'/U \to V/U$ et $x$ tel que $p'(x) = V/U$.
Puisque $p : \mathcal{S} \to \mathcal{C}$ est une catégorie fibrée,
il existe un morphisme fortement cartésien $\varphi : x' \to x$
tel que $p(\varphi)$ soit l'application sous-jacente de $g$.
Le même argument que ci-dessus montre que
$p'(\varphi) = g : V'/U \to V/U$. Comme nous l'avons vu plus haut,
le morphisme $\varphi$ est fortement cartésien. Les conditions de la
définition \ref{definition-fibred-category} sont donc satisfaites, ce qui
achève la preuve.
\end{proof}

\begin{lemma}
\label{lemma-fibred-over-fibred}
Soient $\mathcal{A} \to \mathcal{B} \to \mathcal{C}$ des foncteurs entre
catégories. Si $\mathcal{A}$ est fibrée au-dessus de $\mathcal{B}$ et si
$\mathcal{B}$ est fibrée au-dessus de $\mathcal{C}$, alors $\mathcal{A}$
est fibrée au-dessus de $\mathcal{C}$.
\end{lemma}

\begin{proof}
Cela résulte des définitions et du
lemme \ref{lemma-cartesian-over-cartesian}.
\end{proof}

\begin{lemma}
\label{lemma-fibred-category-representable-goes-up}
Soit $p : \mathcal{S} \to \mathcal{C}$ une catégorie fibrée.
Soient $x \to y$ et $z \to y$ des morphismes de $\mathcal{S}$,
où $x \to y$ est fortement cartésien. Si
$p(x) \times_{p(y)} p(z)$ existe, alors $x \times_y z$ existe,
$p(x \times_y z) = p(x) \times_{p(y)} p(z)$, et
$x \times_y z \to z$ est fortement cartésien.
\end{lemma}

\begin{proof}
Choisissons un morphisme fortement cartésien
$\text{pr}_2^*z \to z$ au-dessus de
$\text{pr}_2 : p(x) \times_{p(y)} p(z) \to p(z)$. Alors
$\text{pr}_2^*z = x \times_y z$ par le
lemme \ref{lemma-strongly-cartesian-fibre-product}.
\end{proof}

\begin{lemma}
\label{lemma-ameliorate-morphism-fibred-categories}
Soit $\mathcal{C}$ une catégorie. Soit
$F : \mathcal{X} \to \mathcal{Y}$ un $1$-morphisme de catégories fibrées
au-dessus de $\mathcal{C}$.
Il existe des $1$-morphismes de catégories fibrées au-dessus de
$\mathcal{C}$
$$
\xymatrix{
\mathcal{X} \ar@<1ex>[r]^u &
\mathcal{X}' \ar[r]^v \ar@<1ex>[l]^w & \mathcal{Y}
}
$$
tels que $F = v \circ u$ et tels que
\begin{enumerate}
\item $u : \mathcal{X} \to \mathcal{X}'$ soit pleinement fidèle,
\item $w$ soit adjoint à gauche de $u$, et
\item le foncteur $v : \mathcal{X}' \to \mathcal{Y}$ définisse une
catégorie fibrée.
\end{enumerate}
\end{lemma}

\begin{proof}
Notons $p : \mathcal{X} \to \mathcal{C}$ et
$q : \mathcal{Y} \to \mathcal{C}$ les foncteurs structuraux.
Nous construisons explicitement $\mathcal{X}'$ comme suit.
Un objet de $\mathcal{X}'$ est un quadruplet $(U, x, y, f)$ où
$x \in \Ob(\mathcal{X}_U)$, $y \in \Ob(\mathcal{Y}_U)$
et $f : y \to F(x)$ est un morphisme dans $\mathcal{Y}_U$.
Un morphisme $(a, b) : (U, x, y, f) \to (U', x', y', f')$ est donné
par $a : x \to x'$ et $b : y \to y'$ tels que
$p(a) = q(b) : U \to U'$ et tels que
$f' \circ b = F(a) \circ f$.

\medskip\noindent
Faisons un choix de changements de base pour $p$ et pour $q$, et employons
la même notation pour les désigner.
Soit $(U, x, y, f)$ un objet et soit $h : V \to U$ un morphisme.
Considérons le morphisme
$c : (V, h^*x, h^*y, h^*f) \to (U, x, y, f)$
provenant des morphismes fortement cartésiens donnés
$h^*x \to x$ et $h^*y \to y$.
Nous affirmons que $c$ est fortement cartésien dans $\mathcal{X}'$
au-dessus de $\mathcal{C}$.
Supposons en effet que l'on se soit donné un objet
$(W, x', y', f')$ de $\mathcal{X}'$, un morphisme
$(a, b) : (W, x', y', f') \to (U, x, y, f)$ au-dessus de
$W \to U$, et une factorisation $W \to V \to U$ de $W \to U$ par $h$.
Comme $h^*x \to x$ et $h^*y \to y$ sont fortement cartésiens, nous obtenons
des morphismes $a' : x' \to h^*x$ et $b' : y' \to h^*y$ au-dessus du
morphisme donné $W \to V$. Considérons le diagramme
$$
\xymatrix{
y' \ar[d]_{f'} \ar[r] & h^*y \ar[r] \ar[d]_{h^*f} & y \ar[d]_f \\
F(x') \ar[r] & F(h^*x) \ar[r] & F(x)
}
$$
Le rectangle extérieur et le carré de droite sont commutatifs.
Puisque $F$ est un $1$-morphisme de catégories fibrées, le morphisme
$F(h^*x) \to F(x)$ est fortement cartésien.
Le carré de gauche est donc commutatif par la propriété universelle des
morphismes fortement cartésiens. Cela montre que $\mathcal{X}'$
est fibrée au-dessus de $\mathcal{C}$.

\medskip\noindent
Le foncteur $u : \mathcal{X} \to \mathcal{X}'$ est donné par
$x \mapsto (p(x), x, F(x), \text{id})$. Il est pleinement fidèle.
Le foncteur $\mathcal{X}' \to \mathcal{Y}$ est donné par
$(U, x, y, f) \mapsto y$. Le foncteur
$w : \mathcal{X}' \to \mathcal{X}$ est donné par
$(U, x, y, f) \mapsto x$. Chacun de ces foncteurs est un
$1$-morphisme de catégories fibrées au-dessus de $\mathcal{C}$ d'après
notre description des morphismes fortement cartésiens de $\mathcal{X}'$
au-dessus de $\mathcal{C}$. Dire que $w$ est adjoint à gauche de $u$ signifie que
$$
\Mor_\mathcal{X}(x, x') =
\Mor_{\mathcal{X}'}((U, x, y, f), (p(x'), x', F(x'), \text{id})),
$$
ce qui résulte immédiatement des définitions.

\medskip\noindent
Enfin, nous devons montrer que $\mathcal{X}' \to \mathcal{Y}$ est une
catégorie fibrée. Soit $c : y' \to y$ un morphisme dans $\mathcal{Y}$
et soit $(U, x, y, f)$ un objet de $\mathcal{X}'$ au-dessus de $y$.
Posons $V = q(y')$ et $h = q(c) : V \to U$. Soient
$a : h^*x \to x$ et $b : h^*y \to y$ les morphismes fortement cartésiens
au-dessus de $h$.
Comme $F$ est un $1$-morphisme de catégories fibrées, nous pouvons identifier
$h^*F(x) = F(h^*x)$, le morphisme fortement cartésien correspondant étant
$F(a) : F(h^*x) \to F(x)$. Par la propriété universelle de
$b : h^*y \to y$, il existe un morphisme $c' : y' \to h^*y$ dans
$\mathcal{Y}_V$ tel que $c = b \circ c'$. Nous affirmons que
$$
(a, c) : (V, h^*x, y', h^*f \circ c') \longrightarrow (U, x, y, f)
$$
est fortement cartésien dans $\mathcal{X}'$ au-dessus de $\mathcal{Y}$.
Pour le voir, soit $(W, x_1, y_1, f_1)$ un objet de $\mathcal{X}'$, soit
$(a_1, b_1) : (W, x_1, y_1, f_1) \to (U, x, y, f)$ un morphisme,
et supposons que $b_1 = c \circ b_1'$ pour un morphisme
$b_1' : y_1 \to y'$. Alors
$$
(a_1', b_1') : (W, x_1, y_1, f_1) \longrightarrow (V, h^*x, y', h^*f \circ c')
$$
(où $a_1' : x_1 \to h^*x$ est l'unique morphisme au-dessus du morphisme donné
$q(b_1') : W \to V$ tel que $a_1 = a \circ a_1'$)
est le morphisme recherché.
\end{proof}
\section{Inertie}
\label{section-inertia}

\noindent
Étant données des catégories fibrées $p : \mathcal{S} \to \mathcal{C}$ et
$p' : \mathcal{S}' \to \mathcal{C}$ au-dessus d'une catégorie $\mathcal{C}$
ainsi qu'un $1$-morphisme $F : \mathcal{S} \to \mathcal{S}'$,
nous disposons du morphisme diagonal
$$
\Delta = \Delta_{\mathcal{S}/\mathcal{S}'} :
\mathcal{S} \longrightarrow \mathcal{S} \times_{\mathcal{S}'} \mathcal{S}
$$
dans la $(2, 1)$-catégorie des catégories fibrées au-dessus de $\mathcal{C}$.

\begin{lemma}
\label{lemma-inertia-fibred-category}
Soit $\mathcal{C}$ une catégorie. Soient
$p : \mathcal{S} \to \mathcal{C}$ et
$p' : \mathcal{S}' \to \mathcal{C}$ des catégories fibrées.
Soit $F : \mathcal{S} \to \mathcal{S}'$ un $1$-morphisme de
catégories fibrées au-dessus de $\mathcal{C}$. Considérons la catégorie
$\mathcal{I}_{\mathcal{S}/\mathcal{S}'}$ au-dessus de $\mathcal{C}$ dont
\begin{enumerate}
\item les objets sont les paires $(x, \alpha)$ où $x \in \Ob(\mathcal{S})$
et $\alpha : x \to x$ est un automorphisme tel que $F(\alpha) = \text{id}$,
\item les morphismes $(x, \alpha) \to (y, \beta)$ sont donnés par les morphismes
$\phi : x \to y$ tels que le diagramme
$$
\xymatrix{
x\ar[r]_\phi\ar[d]_\alpha &
y\ar[d]^{\beta} \\
x\ar[r]^\phi &
y \\
}
$$
soit commutatif, et
\item le foncteur $\mathcal{I}_{\mathcal{S}/\mathcal{S}'} \to \mathcal{C}$
est donné par $(x, \alpha) \mapsto p(x)$.
\end{enumerate}
Alors
\begin{enumerate}
\item il existe une équivalence
$$
\mathcal{I}_{\mathcal{S}/\mathcal{S}'} \longrightarrow
\mathcal{S}
\times_{\Delta, (\mathcal{S} \times_{\mathcal{S}'} \mathcal{S}), \Delta}
\mathcal{S}
$$
dans la $(2, 1)$-catégorie des catégories au-dessus de $\mathcal{C}$, et
\item $\mathcal{I}_{\mathcal{S}/\mathcal{S}'}$ est une catégorie fibrée au-dessus de
$\mathcal{C}$.
\end{enumerate}
\end{lemma}
\begin{proof}
Remarquons que (2) découle de (1) par les
lemmes \ref{lemma-2-product-fibred-categories-over-C} et
\ref{lemma-fibred-equivalent}. Il suffit donc de prouver (1).
Nous utiliserons sans autre mention la construction du $2$-produit fibré
donnée dans le
lemme \ref{lemma-2-product-fibred-categories-over-C}.
En particulier, un objet de
$\mathcal{S}
\times_{\Delta, (\mathcal{S} \times_{\mathcal{S}'} \mathcal{S}), \Delta}
\mathcal{S}$
est un triplet $(x, y, (\iota, \kappa))$ où $x$ et $y$ sont des objets de
$\mathcal{S}$, et
$(\iota, \kappa) : (x, x, \text{id}_{F(x)}) \to (y, y, \text{id}_{F(y)})$
est un isomorphisme dans $\mathcal{S} \times_{\mathcal{S}'} \mathcal{S}$.
Cela signifie simplement que $\iota, \kappa : x \to y$ sont des isomorphismes et que
$F(\iota) = F(\kappa)$. Considérons le foncteur
$$
\mathcal{I}_{\mathcal{S}/\mathcal{S}'}
\longrightarrow
\mathcal{S}
\times_{\Delta, (\mathcal{S} \times_{\mathcal{S}'} \mathcal{S}), \Delta}
\mathcal{S}
$$
qui, à un objet $(x, \alpha)$ du membre de gauche, associe l'objet
$(x, x, (\alpha, \text{id}_x))$ du membre de droite
et, à un morphisme $\phi$ du membre de gauche,
associe le morphisme $(\phi, \phi)$ du membre de droite.
Nous affirmons qu'un quasi-inverse de ce foncteur est donné par le
foncteur
$$
\mathcal{S}
\times_{\Delta, (\mathcal{S} \times_{\mathcal{S}'} \mathcal{S}), \Delta}
\mathcal{S}
\longrightarrow
\mathcal{I}_{\mathcal{S}/\mathcal{S}'}
$$
qui, à un objet $(x, y, (\iota, \kappa))$ du membre de gauche,
associe l'objet $(x, \kappa^{-1} \circ \iota)$ du membre de droite
et, à un morphisme
$(\phi, \phi') : (x, y, (\iota, \kappa)) \to (z, w, (\lambda, \mu))$
du membre de gauche, associe le morphisme $\phi$.
En effet, l'endofoncteur de $\mathcal{I}_{\mathcal{S}/\mathcal{S}'}$ obtenu
en composant les deux foncteurs ci-dessus est strictement l'identité, et
l'endofoncteur induit sur
$\mathcal{S}
\times_{\Delta, (\mathcal{S} \times_{\mathcal{S}'} \mathcal{S}), \Delta}
\mathcal{S}$
est isomorphe à
l'identité par l'isomorphisme naturel
$$
(\text{id}_x, \kappa) :
(x, x, (\kappa^{-1} \circ \iota, \text{id}_x))
\longrightarrow
(x, y, (\iota, \kappa)).
$$
Certains détails sont omis.
\end{proof}

\begin{definition}
\label{definition-inertia-fibred-category}
Soit $\mathcal{C}$ une catégorie.
\begin{enumerate}
\item Soit $F : \mathcal{S} \to \mathcal{S}'$ un $1$-morphisme de
catégories fibrées au-dessus de $\mathcal{C}$. L'{\it inertie relative
de $\mathcal{S}$ au-dessus de $\mathcal{S}'$} est la catégorie fibrée
$\mathcal{I}_{\mathcal{S}/\mathcal{S}'} \to \mathcal{C}$ du
lemme \ref{lemma-inertia-fibred-category}.
\item Par {\it catégorie fibrée d'inertie $\mathcal{I}_\mathcal{S}$
de $\mathcal{S}$}, nous entendons
$\mathcal{I}_\mathcal{S} = \mathcal{I}_{\mathcal{S}/\mathcal{C}}$.
\end{enumerate}
\end{definition}

\noindent
Remarquons qu'il existe des $1$-morphismes canoniques
\begin{equation}
\label{equation-inertia-structure-map}
\mathcal{I}_{\mathcal{S}/\mathcal{S}'} \longrightarrow \mathcal{S}
\quad\text{et}\quad
\mathcal{I}_\mathcal{S} \longrightarrow \mathcal{S}
\end{equation}
de catégories fibrées au-dessus de $\mathcal{C}$. Dans la description du
lemme \ref{lemma-inertia-fibred-category},
ils envoient simplement l'objet $(x, \alpha)$ sur l'objet $x$ et le morphisme
$\phi : (x, \alpha) \to (y, \beta)$ sur le morphisme $\phi : x \to y$.
Il existe aussi une {\it section neutre}
\begin{equation}
\label{equation-neutral-section}
e : \mathcal{S} \to \mathcal{I}_{\mathcal{S}/\mathcal{S}'}
\quad\text{et}\quad
e : \mathcal{S} \to \mathcal{I}_\mathcal{S}
\end{equation}
définie par les règles $x \mapsto (x, \text{id}_x)$ et
$(\phi : x \to y) \mapsto \phi$. C'est un inverse à droite de
(\ref{equation-inertia-structure-map}). Étant donné un carré $2$-commutatif
$$
\xymatrix{
\mathcal{S}_1 \ar[d]_{F_1} \ar[r]_G & \mathcal{S}_2 \ar[d]^{F_2} \\
\mathcal{S}'_1 \ar[r]^{G'} & \mathcal{S}'_2
}
$$
il existe des {\it morphismes de fonctorialité}
\begin{equation}
\label{equation-functorial}
\mathcal{I}_{\mathcal{S}_1/\mathcal{S}'_1}
\longrightarrow
\mathcal{I}_{\mathcal{S}_2/\mathcal{S}'_2}
\quad\text{et}\quad
\mathcal{I}_{\mathcal{S}_1}
\longrightarrow
\mathcal{I}_{\mathcal{S}_2}
\end{equation}
définis par les règles $(x, \alpha) \mapsto (G(x), G(\alpha))$
et $\phi \mapsto G(\phi)$. En particulier, il existe toujours un
morphisme de comparaison
\begin{equation}
\label{equation-comparison}
\mathcal{I}_{\mathcal{S}/\mathcal{S}'}
\longrightarrow
\mathcal{I}_\mathcal{S}
\end{equation}
et tous les morphismes ci-dessus sont compatibles avec celui-ci.

\begin{lemma}
\label{lemma-relative-inertia-as-fibre-product}
Soit $F : \mathcal{S} \to \mathcal{S}'$ un $1$-morphisme de catégories
fibrées au-dessus d'une catégorie $\mathcal{C}$. Alors le diagramme
$$
\xymatrix{
\mathcal{I}_{\mathcal{S}/\mathcal{S}'}
\ar[d]_{F \circ (\ref{equation-inertia-structure-map})}
\ar[rr]_{(\ref{equation-comparison})} & &
\mathcal{I}_\mathcal{S} \ar[d]^{(\ref{equation-functorial})} \\
\mathcal{S}' \ar[rr]^e & &
\mathcal{I}_{\mathcal{S}'}
}
$$
est un $2$-produit fibré.
\end{lemma}

\begin{proof}
Omis.
\end{proof}





\section{Catégories fibrées en groupoïdes}
\label{section-fibred-groupoids}

\noindent
Dans cette section, nous expliquons comment concevoir les catégories fibrées en groupoïdes
et nous voyons qu'elles reviennent essentiellement aux foncteurs à valeurs dans
la $(2, 1)$-catégorie des groupoïdes.

\begin{definition}
\label{definition-fibred-groupoids}
Soit $p : \mathcal{S} \to \mathcal{C}$ un foncteur.
On dit que $\mathcal{S}$ est une {\it catégorie fibrée en groupoïdes} au-dessus de $\mathcal{C}$ si
les deux conditions suivantes sont satisfaites :
\begin{enumerate}
\item Pour tout morphisme $f : V \to U$ dans $\mathcal{C}$ et tout
relèvement $x$ de $U$, il existe un relèvement $\phi : y \to x$ de $f$
de but $x$.
\item Pour toute paire de morphismes $\phi : y \to x$ et $ \psi : z \to x$
et tout morphisme $f : p(z) \to p(y)$ tel que $p(\phi) \circ f = p(\psi)$,
il existe un unique relèvement $\chi : z \to y$ de $f$ tel que
$\phi \circ \chi = \psi$.
\end{enumerate}
\end{definition}

\noindent
Formulée autrement, la condition (2) dit que
l'application induite par le foncteur $p$ est une bijection entre les ensembles
de flèches pointillées dans le diagramme commutatif suivant :
$$
\xymatrix{
y \ar[r] & x & p(y) \ar[r] & p(x) \\
z \ar@{-->}[u] \ar[ru] & & p(z) \ar@{-->}[u]\ar[ru] & \\
}
$$
Voici une autre manière d'interpréter la deuxième condition.
Supposons que $g : W \to V$ et $f : V \to U$ soient des morphismes de $\mathcal{C}$.
Soit $x \in \Ob(\mathcal{S}_U)$. Par la première condition, nous pouvons relever
$f$ en $ \phi : y \to x$, puis relever $g$ en $\psi : z \to y$.
Au lieu de procéder en deux étapes, nous pouvons relever directement $f \circ g$ en
$\gamma : z' \to x$. On obtient ainsi les flèches pleines du diagramme
\begin{equation}
\label{equation-fibred-groupoids}
\vcenter{
\xymatrix{
z' \ar@{-->}[d]\ar[rrd]^\gamma & & \\
z \ar@{-->}[u] \ar[r]^\psi \ar@{~>}[d]^p &
y \ar[r]^\phi \ar@{~>}[d]^p &
x \ar@{~>}[d]^p
\\
W \ar[r]^g & V \ar[r]^f & U \\
}
}
\end{equation}
où les flèches ondulées représentent non pas des morphismes, mais le foncteur $p$.
En appliquant la deuxième condition aux flèches $\phi \circ \psi$, $\gamma$
et $\text{id}_W$, nous concluons qu'il existe un unique morphisme
$\chi : z \to z'$ dans $\mathcal{S}_W$ tel que
$\gamma \circ \chi = \phi \circ \psi$. De même, il existe un unique morphisme
$z' \to z$. L'unicité implique que les morphismes $z' \to z$ et
$z\to z'$ sont inverses l'un de l'autre, autrement dit des isomorphismes.

\medskip\noindent
Cette discussion montre clairement qu'une
catégorie fibrée en groupoïdes est très étroitement liée
à une catégorie fibrée. Voici le résultat.

\begin{lemma}
\label{lemma-fibred-groupoids}
Soit $p : \mathcal{S} \to \mathcal{C}$ un foncteur.
Les assertions suivantes sont équivalentes :
\begin{enumerate}
\item $p : \mathcal{S} \to \mathcal{C}$ est une catégorie
fibrée en groupoïdes, et
\item toutes les catégories fibres sont des groupoïdes et
$\mathcal{S}$ est une catégorie fibrée au-dessus de $\mathcal{C}$.

\end{enumerate}
De plus, dans ce cas, tout morphisme de $\mathcal{S}$ est
fortement cartésien. En outre, étant donnés des $f^\ast x \to x$
au-dessus de $f$ pour tous les $f: V \to U = p(x)$, les données
$(U \mapsto \mathcal{S}_U, f \mapsto f^*, \alpha_{g, f}, \alpha_U)$
construites dans le lemme \ref{lemma-fibred}
définissent un pseudo-foncteur de $\mathcal{C}^{opp}$ vers
la $(2, 1)$-catégorie des groupoïdes.
\end{lemma}

\begin{proof}
Supposons que le foncteur $p : \mathcal{S} \to \mathcal{C}$ définisse une catégorie fibrée en groupoïdes.
Pour montrer que toutes les catégories fibres $\mathcal{S}_U$, pour
$U \in \Ob(\mathcal{C})$,
sont des groupoïdes, nous devons exhiber, pour tout $f : y \to x$ dans $\mathcal{S}_U$, un
morphisme inverse. Le diagramme de gauche (dans $\mathcal{S}_U$) est envoyé par
$p$ sur le diagramme de droite :
$$
\xymatrix{
y \ar[r]^f & x & U \ar[r]^{\text{id}_U} & U \\
x \ar@{-->}[u] \ar[ru]_{\text{id}_x} & &
U \ar@{-->}[u]\ar[ru]_{\text{id}_U} & \\
}
$$
Puisque seul $\text{id}_U$ rend commutatif le diagramme de droite, il existe un
unique $g : x \to y$ rendant commutatif le diagramme de gauche, donc
$fg = \text{id}_x$. Par un argument analogue, il existe un unique $h : y \to x$ tel
que $gh = \text{id}_y$. Alors $fgh = f : y \to x$. Comme $fg = \text{id}_x$,
on a $h = f$. La condition (2) de la définition \ref{definition-fibred-groupoids} dit
exactement que tout morphisme de $\mathcal{S}$ est fortement cartésien. Ainsi,
la condition (1) de la définition \ref{definition-fibred-groupoids} implique que
$\mathcal{S}$ est une catégorie fibrée au-dessus de $\mathcal{C}$.
\medskip\noindent
Réciproquement, supposons que toutes les catégories fibres soient des groupoïdes et que
$\mathcal{S}$ soit une catégorie fibrée au-dessus de $\mathcal{C}$.
Il faut vérifier les conditions (1) et (2) de la
définition \ref{definition-fibred-groupoids}.
La première condition est immédiate. Soient $\phi : y \to x$,
$\psi : z \to x$ et $f : p(z) \to p(y)$ tels que
$p(\phi) \circ f = p(\psi)$, comme dans la condition (2) de la
définition \ref{definition-fibred-groupoids}.
Écrivons $U = p(x)$, $V = p(y)$, $W = p(z)$, $p(\phi) = g : V \to U$,
$p(\psi) = h : W \to U$. Choisissons un morphisme fortement cartésien $g^*x \to x$
au-dessus de $g$. Nous obtenons alors un morphisme $i : y \to g^*x$ dans
$\mathcal{S}_V$, qui est donc un isomorphisme. Nous
obtenons aussi un morphisme $j : z \to g^*x$ correspondant à
la paire $(\psi, f)$ puisque $g^*x \to x$ est fortement cartésien.
On vérifie alors que $\chi = i^{-1} \circ j$ est une solution.

\medskip\noindent
Nous avons vu dans la preuve de (1) $\Rightarrow$ (2) que
tout morphisme de $\mathcal{S}$ est fortement cartésien.
La dernière assertion découle directement du lemme \ref{lemma-fibred}.
\end{proof}

\begin{lemma}
\label{lemma-fibred-gives-fibred-groupoids}
Soit $\mathcal{C}$ une catégorie.
Soit $p : \mathcal{S} \to \mathcal{C}$ une catégorie fibrée.
Soit $\mathcal{S}'$ la sous-catégorie de $\mathcal{S}$ définie
comme suit :
\begin{enumerate}
\item $\Ob(\mathcal{S}') = \Ob(\mathcal{S})$, et
\item pour $x, y \in \Ob(\mathcal{S}')$, l'ensemble des morphismes entre $x$
et $y$ dans $\mathcal{S}'$ est l'ensemble des morphismes fortement cartésiens entre
$x$ et $y$ dans $\mathcal{S}$.
\end{enumerate}
Soit $p' : \mathcal{S}' \to \mathcal{C}$ la restriction de $p$
à $\mathcal{S}'$. Alors $p' : \mathcal{S}' \to \mathcal{C}$ est une catégorie
fibrée en groupoïdes.
\end{lemma}

\begin{proof}
Remarquons que la construction a un sens puisque, d'après le
lemme \ref{lemma-composition-cartesian},
le morphisme identité de tout objet de $\mathcal{S}$ est fortement cartésien,
et la composée de morphismes fortement cartésiens est fortement cartésienne.
La première propriété de relèvement de la
définition \ref{definition-fibred-groupoids}
découle de la condition selon laquelle, dans une catégorie fibrée,
pour tout morphisme $f : V \to U$ et tout $x$ au-dessus de $U$, il existe
un morphisme fortement cartésien $\varphi : y \to x$ au-dessus de $f$.
Vérifions la deuxième propriété de relèvement de la
définition \ref{definition-fibred-groupoids}
pour la catégorie $p' : \mathcal{S}' \to \mathcal{C}$ au-dessus de $\mathcal{C}$.
Pour ce faire, raisonnons comme dans la discussion qui suit la
définition \ref{definition-fibred-groupoids}.
Ainsi, dans le diagramme \ref{equation-fibred-groupoids}, les
morphismes $\phi$, $\psi$ et $\gamma$ sont des morphismes fortement cartésiens
de $\mathcal{S}$.
Par conséquent, $\gamma$ et $\phi \circ \psi$ sont des morphismes fortement cartésiens
de $\mathcal{S}$ au-dessus de la même flèche de $\mathcal{C}$ et
ayant le même but dans $\mathcal{S}$. D'après la discussion qui suit la
définition \ref{definition-cartesian-over-C},
ces deux flèches sont donc isomorphes, comme voulu (nous utilisons aussi ici
le fait que tout isomorphisme de $\mathcal{S}$ est fortement cartésien, d'après le
lemme \ref{lemma-composition-cartesian}).
\end{proof}

\begin{example}
\label{example-group-homomorphism-fibreedingroupoids}
Un homomorphisme de groupes $p : G \to H$ donne naissance à un foncteur
$p : \mathcal{S}\to\mathcal{C}$ comme dans l'exemple
\ref{example-group-homomorphism-functor}. Ce foncteur
$p : \mathcal{S}\to\mathcal{C}$ définit une catégorie fibrée en groupoïdes si et seulement si
$p$ est surjectif. La catégorie fibre $\mathcal{S}_U$ au-dessus de l'unique
objet $U\in \Ob(\mathcal{C})$ est la catégorie associée au
noyau de $p$, comme dans l'exemple \ref{example-group-groupoid}.
\end{example}

\noindent
Étant donné $p : \mathcal{S} \to \mathcal{C}$, on peut se demander : si la catégorie
fibre $\mathcal{S}_U$ est un groupoïde pour tout $U \in \Ob(\mathcal{C})$,
$\mathcal{S}$ est-elle nécessairement fibrée en groupoïdes au-dessus de $\mathcal{C}$ ?
La réponse est négative, comme on peut le voir ainsi. Partons d'une catégorie fibrée en
groupoïdes $p : \mathcal{S} \to \mathcal{C}$. Modifier les morphismes de $\mathcal{S}$
qui ne s'envoient pas sur le morphisme identité d'un objet ne modifie pas les
catégories $\mathcal{S}_U$. On peut donc faire échouer les conditions d'existence et
d'unicité des relèvements. Un exemple est fourni par le foncteur de l'exemple
\ref{example-group-homomorphism-fibreedingroupoids} lorsque $G \to H$ n'est pas
surjectif. Voici un autre exemple.

\begin{example}
\label{example-not-fibred-in-groupoids-but-fibre-cats-are}
Posons $\Ob(\mathcal{C}) = \{A, B, T\}$ et
$\Mor_\mathcal{C}(A, B) = \{f\}$, $\Mor_\mathcal{C}(B, T) = \{g\}$,
$\Mor_\mathcal{C}(A, T) = \{h\} = \{gf\}, $ auxquels s'ajoute le morphisme identité
de chaque objet. Le diagramme ci-dessous représente cette catégorie. Posons maintenant
$\Ob(\mathcal{S}) = \{A', B', T'\}$ et
$\Mor_\mathcal{S}(A', B') = \emptyset$,
$\Mor_\mathcal{S}(B', T') = \{g'\}$,
$\Mor_\mathcal{S}(A', T') = \{h'\}, $ auxquels s'ajoutent les morphismes identité. Le
foncteur $p : \mathcal{S} \to \mathcal{C}$ est évident. Alors, pour tout
$U \in \Ob(\mathcal{C})$, $\mathcal{S}_U$ est la catégorie ayant un seul
objet et le morphisme identité de cet objet, donc un groupoïde, mais le
morphisme $f: A \to B$ ne peut être relevé. De même, si l'on pose
$\Mor_\mathcal{S}(A', B') = \{f'_1, f'_2\}$ et
$ \Mor_\mathcal{S}(A', T') = \{h'\} = \{g'f'_1 \} = \{g'f'_2\}$, alors
les catégories fibres sont les mêmes et $f: A \to B$ dans le diagramme ci-dessous
possède deux relèvements.
$$
\xymatrix{
B' \ar[r]^{g'} & T' & & B \ar[r]^g & T & \\
A' \ar@{-->}[u]^{??} \ar[ru]_{h'} & & \ar@{}[u]^{\text{ci-dessus}} &
A \ar[u]^f \ar[ru]_{gf = h} & \\
}
$$
\end{example}

\noindent
Plus loin, nous voudrons formuler des assertions telles que « toute catégorie fibrée en
groupoïdes au-dessus de $\mathcal{C}$ est équivalente à une catégorie scindée », ou
« toute catégorie fibrée en groupoïdes dont les catégories fibres sont assimilables à des ensembles
est équivalente à une catégorie fibrée en ensembles ». La notion d'équivalence
dépend de la $2$-catégorie dans laquelle nous travaillons.

\begin{definition}
\label{definition-categories-fibred-in-groupoids-over-C}
Soit $\mathcal{C}$ une catégorie.
La {\it $2$-catégorie des catégories fibrées en groupoïdes au-dessus de $\mathcal{C}$}
est la sous-$2$-catégorie de la $2$-catégorie des catégories fibrées
au-dessus de $\mathcal{C}$ (voir la définition \ref{definition-fibred-categories-over-C})
définie comme suit :
\begin{enumerate}
\item Ses objets sont les catégories
$p : \mathcal{S} \to \mathcal{C}$ fibrées en groupoïdes.
\item Ses $1$-morphismes $(\mathcal{S}, p) \to (\mathcal{S}', p')$
sont les foncteurs $G : \mathcal{S} \to \mathcal{S}'$ tels que
$p' \circ G = p$ (puisque tout morphisme est fortement cartésien,
$G$ les préserve automatiquement).
\item Ses $2$-morphismes $t : G \to H$ pour
$G, H : (\mathcal{S}, p) \to (\mathcal{S}', p')$
sont les morphismes de foncteurs
tels que $p'(t_x) = \text{id}_{p(x)}$
pour tout $x \in \Ob(\mathcal{S})$.
\end{enumerate}
\end{definition}

\noindent
Remarquons que tout $2$-morphisme est automatiquement un isomorphisme !
Il s'agit donc en fait d'une $(2, 1)$-catégorie, et non simplement d'une
$2$-catégorie. Voici l'inévitable lemme sur les $2$-produits fibrés.

\begin{lemma}
\label{lemma-2-product-fibred-categories}
Soit $\mathcal{C}$ une catégorie.
La $2$-catégorie des catégories fibrées en groupoïdes
au-dessus de $\mathcal{C}$ admet des $2$-produits fibrés, et ceux-ci sont décrits comme dans le
lemme \ref{lemma-2-product-categories-over-C}.
\end{lemma}

\begin{proof}
D'après le lemme \ref{lemma-2-product-fibred-categories-over-C},
le produit fibré décrit dans le
lemme \ref{lemma-2-product-categories-over-C} est une catégorie fibrée.
Il suffit donc de prouver que les catégories fibres sont des
groupoïdes ; voir le lemme \ref{lemma-fibred-groupoids}.
D'après le lemme \ref{lemma-fibre-2-fibre-product-categories-over-C},
il suffit de montrer que le 2-produit fibré de groupoïdes
est un groupoïde, ce qui est clair (par exemple d'après la construction du
lemme \ref{lemma-2-fibre-product-categories}).
\end{proof}

\begin{remark}
\label{remark-2-product-fibred-categories-same}
Soit $\mathcal{C}$ une catégorie. Soient $f : \mathcal{X} \to \mathcal{S}$
et $g : \mathcal{Y} \to \mathcal{S}$ des $1$-morphismes de
catégories fibrées en groupoïdes au-dessus de $\mathcal{C}$.
Soit $p : \mathcal{S} \to \mathcal{C}$ le foncteur donné.
Nous affirmons que le $2$-produit fibré du lemme \ref{lemma-2-product-fibred-categories}
est canoniquement équivalent (comme catégorie) à celui de
l'exemple \ref{example-2-fibre-product-categories}.
Les objets du premier sont les quadruplets $(U, x, y, \alpha)$
où $p(\alpha) = \text{id}_U$
(voir le lemme \ref{lemma-2-product-categories-over-C}),
et les objets du second sont les triplets $(x, y, \alpha)$ (voir
l'exemple \ref{example-2-fibre-product-categories}).
L'équivalence entre les deux catégories est donnée par
les règles $(U, x, y, \alpha) \mapsto (x, y, \alpha)$
et $(x, y, \alpha) \mapsto (p(f(x)), x, y', \alpha')$,
où $\alpha' = g(\gamma)^{-1} \circ \alpha$ et $\gamma : y' \to y$
est un relèvement de la flèche $p(\alpha) : p(f(x)) \to p(g(y))$.
Les détails sont omis.
\end{remark}

\begin{lemma}
\label{lemma-equivalence-fibred-categories}
Soient $p : \mathcal{S}\to \mathcal{C}$ et
$p' : \mathcal{S'}\to \mathcal{C}$ des catégories fibrées en groupoïdes, et
supposons que $G : \mathcal{S}\to \mathcal {S}'$ soit un foncteur au-dessus de
$\mathcal{C}$.
\begin{enumerate}
\item Alors $G$ est fidèle (resp.\ pleinement fidèle, resp.\ une équivalence)
si et seulement si, pour chaque $U\in\Ob(\mathcal{C})$, le foncteur induit
$G_U : \mathcal{S}_U\to \mathcal{S}'_U$ est fidèle
(resp.\ pleinement fidèle, resp.\ une équivalence).
\item Si $G$ est une équivalence, alors $G$ est une équivalence dans la
$2$-catégorie des catégories fibrées en groupoïdes au-dessus de $\mathcal{C}$.
\end{enumerate}
\end{lemma}

\begin{proof}
Soient $x, y$ des objets de $\mathcal{S}$ au-dessus du même objet $U$.
Considérons le diagramme commutatif
$$
\xymatrix{
\Mor_\mathcal{S}(x, y) \ar[rd]_p \ar[rr]_G & &
\Mor_{\mathcal{S}'}(G(x), G(y)) \ar[ld]^{p'} \\
& \Mor_\mathcal{C}(U, U) &
}
$$
Ce diagramme montre clairement que, si $G$ est fidèle (resp.\ pleinement fidèle),
alors chaque $G_U$ l'est aussi.

\medskip\noindent
Supposons que $G$ soit une équivalence. Pour tout objet
$x'$ de $\mathcal{S}'$, il existe un objet $x$ de $\mathcal{S}$
tel que $G(x)$ soit isomorphe à $x'$. Supposons que $x'$ soit
au-dessus de $U'$ et que $x$ soit au-dessus de $U$. Il existe alors un isomorphisme
$f : U' \to U$ dans $\mathcal{C}$, à savoir l'image par $p'$ de
l'isomorphisme $x' \to G(x)$. Par les axiomes d'une catégorie fibrée
en groupoïdes, il existe une flèche $f^*x \to x$ de $\mathcal{S}$
au-dessus de $f$. Il existe donc un isomorphisme
$\alpha : x' \to G(f^*x)$ tel que $p'(\alpha) = \text{id}_{U'}$
(cette fois par les axiomes de $\mathcal{S}'$). En définitive, nous concluons
que, pour tout objet $x'$ de $\mathcal{S}'$, nous pouvons choisir
une paire $(o_{x'}, \alpha_{x'})$ formée d'un objet
$o_{x'}$ de $\mathcal{S}$ et d'un isomorphisme $\alpha_{x'} : x' \to G(o_{x'})$
tel que $p'(\alpha_{x'}) = \text{id}_{p'(x')}$. À partir de là,
nous procédons comme d'habitude (voir la preuve du
lemme \ref{lemma-equivalence-categories}) pour construire un foncteur
inverse $F : \mathcal{S}' \to \mathcal{S}$, en prenant
$x' \mapsto o_{x'}$ et $\varphi' : x' \to y'$ pour l'unique
flèche $\varphi_{\varphi'} : o_{x'} \to o_{y'}$ telle que
$\alpha_{y'}^{-1} \circ G(\varphi_{\varphi'}) \circ \alpha_{x'} = \varphi'$.
Avec ces choix, $F$ est un foncteur au-dessus de $\mathcal{C}$.
Nous omettons de vérifier que $G \circ F$ et $F \circ G$ sont
$2$-isomorphes aux foncteurs identité respectifs
(dans la $2$-catégorie des catégories fibrées en groupoïdes au-dessus de $\mathcal{C}$).

\medskip\noindent
Supposons que $G_U$ soit fidèle (resp.\ pleinement fidèle)
pour tout $U\in\Ob(\mathcal C)$. Pour
montrer que $G$ est fidèle (resp.\ pleinement fidèle),
il faut montrer, pour tous objets
$x, y\in\Ob(\mathcal{S})$, que $G$ induit une
injection (resp.\ bijection) entre
$\Mor_\mathcal{S}(x, y)$ et
$\Mor_{\mathcal{S}'}(G(x), G(y))$.
Posons $U = p(x)$ et $V = p(y)$.
Il suffit de prouver que $G$
induit une injection (resp.\ bijection) entre les morphismes
$x \to y$ au-dessus de $f$ et les morphismes $G(x) \to G(y)$ au-dessus de $f$,
pour tout morphisme $f : U \to V$.
Fixons maintenant $f : U \to V$. Choisissons un morphisme fortement cartésien
$f^*y \to y$ au-dessus de ce morphisme. Alors $G(f^*y) \to G(y)$ est lui aussi
fortement cartésien.
L'ensemble des morphismes de $x$ vers $y$ au-dessus de $f$
est en bijection avec l'ensemble des morphismes entre
$x$ et $f^*y$ au-dessus de $\text{id}_U$. (Par le deuxième axiome
d'une catégorie fibrée en groupoïdes.) De même,
l'ensemble des morphismes de $G(x)$ vers $G(y)$ au-dessus de $f$
est en bijection avec l'ensemble des morphismes entre
$G(x)$ et $G(f^*y)$ au-dessus de $\text{id}_U$.
Le fait que $G_U$ soit fidèle (resp.\ pleinement fidèle)
donne donc le résultat voulu.

\medskip\noindent
Enfin, supposons que tous les $G_U$ soient des équivalences pour tout $U$ ; ils sont donc
pleinement fidèles et essentiellement surjectifs. Nous avons vu que cela implique que $G$ est
pleinement fidèle ; pour montrer que c'est une équivalence, il reste donc à montrer qu'il
est essentiellement surjectif. C'est clair : si $z'\in
\Ob(\mathcal{S}')$, alors $z'\in \Ob(\mathcal{S}'_U)$, où
$U = p'(z')$. Puisque $G_U$ est essentiellement surjectif, nous savons que
$z'$ est isomorphe, dans $\mathcal{S}'_U$, à un objet de la forme
$G_U(z)$ pour un certain $z\in \Ob(\mathcal{S}_U)$. Or les morphismes
de $\mathcal{S}'_U$ sont des morphismes de $\mathcal{S}'$ ; ainsi $z'$ est
isomorphe à $G(z)$ dans $\mathcal{S}'$.
\end{proof}

\begin{lemma}
\label{lemma-fully-faithful-diagonal-equivalence}
Soit $\mathcal{C}$ une catégorie. Soient $p : \mathcal{S}\to \mathcal{C}$ et
$p' : \mathcal{S'}\to \mathcal{C}$ des catégories fibrées en groupoïdes.
Soit $G : \mathcal{S}\to \mathcal {S}'$ un foncteur au-dessus de $\mathcal{C}$.
Alors $G$ est pleinement fidèle si et seulement si la diagonale
$$
\Delta_G :
\mathcal{S}
\longrightarrow
\mathcal{S} \times_{G, \mathcal{S}', G} \mathcal{S}
$$
est une équivalence.
\end{lemma}

\begin{proof}
D'après le
lemme \ref{lemma-equivalence-fibred-categories},
il suffit de considérer les catégories fibres au-dessus d'un objet $U$ de $\mathcal{C}$.
Un objet du membre de droite est un triplet $(x, x', \alpha)$ où
$\alpha : G(x) \to G(x')$ est un morphisme de $\mathcal{S}'_U$.
Le foncteur $\Delta_G$ envoie l'objet $x$ de $\mathcal{S}_U$
sur le triplet $(x, x, \text{id}_{G(x)})$. Remarquons que $(x, x', \alpha)$
appartient à l'image essentielle de $\Delta_G$ si et seulement si $\alpha = G(\beta)$
pour un certain morphisme $\beta : x \to x'$ dans $\mathcal{S}_U$ (détails omis).
Par conséquent, pour que $\Delta_G$ soit une équivalence, tout $\alpha$ doit
être l'image d'un morphisme $\beta : x \to x'$, et deux morphismes
distincts $\beta, \beta' : x \to x'$ doivent aussi donner des morphismes distincts
$G(\beta), G(\beta')$. Cela prouve le lemme.
\end{proof}

\begin{lemma}
\label{lemma-morphisms-equivalent-fibred-groupoids}
Soit $\mathcal{C}$ une catégorie.
Soient $\mathcal{S}_i$, $i = 1, 2, 3, 4$, des catégories fibrées en
groupoïdes au-dessus de $\mathcal{C}$.
Supposons que $\varphi : \mathcal{S}_1 \to \mathcal{S}_2$ et
$\psi : \mathcal{S}_3 \to \mathcal{S}_4$ soient des équivalences
au-dessus de $\mathcal{C}$. Alors
$$
\Mor_{\textit{Cat}/\mathcal{C}}(\mathcal{S}_2, \mathcal{S}_3)
\longrightarrow
\Mor_{\textit{Cat}/\mathcal{C}}(\mathcal{S}_1, \mathcal{S}_4),
\quad \alpha \longmapsto \psi \circ \alpha \circ \varphi
$$
est une équivalence de catégories.
\end{lemma}

\begin{proof}
C'est un fait général, valable dans toute $2$-catégorie.
\end{proof}

\begin{lemma}
\label{lemma-inertia-fibred-groupoids}
Soit $\mathcal{C}$ une catégorie.
Si $p : \mathcal{S} \to \mathcal{C}$ est une catégorie fibrée en groupoïdes, alors
la catégorie fibrée d'inertie $\mathcal{I}_\mathcal{S} \to \mathcal{C}$ l'est aussi.
\end{lemma}

\begin{proof}
Cela ressort de la construction du
lemme \ref{lemma-inertia-fibred-category},
ou bien du fait (établi dans le même lemme) que
$\mathcal{I}_\mathcal{S} \to \mathcal{S}
\times_{\Delta, \mathcal{S} \times_\mathcal{C} \mathcal{S}, \Delta}\mathcal{S}$
est une équivalence, combiné au
lemme \ref{lemma-2-product-fibred-categories}.
\end{proof}

\begin{lemma}
\label{lemma-cute-groupoids}
Soit $\mathcal{C}$ une catégorie. Soit $U \in \Ob(\mathcal{C})$.
Si $p : \mathcal{S} \to \mathcal{C}$ est une catégorie fibrée en groupoïdes
et si $p$ se factorise par $p' : \mathcal{S} \to \mathcal{C}/U$,
alors $p' : \mathcal{S} \to \mathcal{C}/U$ est une catégorie fibrée en groupoïdes.
\end{lemma}

\begin{proof}
Nous avons déjà vu dans le lemme \ref{lemma-cute} que $p'$ est une catégorie
fibrée. Il suffit donc de prouver que les catégories fibres sont des groupoïdes ;
voir le lemme \ref{lemma-fibred-groupoids}.
Pour $V \in \Ob(\mathcal{C})$, on a
$$
\mathcal{S}_V = \coprod\nolimits_{f : V \to U} \mathcal{S}_{(f : V \to U)}
$$
où le membre de gauche est la catégorie fibre de $p$ et le membre de droite
est la réunion disjointe des catégories fibres de $p'$.
D'où le résultat.
\end{proof}

\begin{lemma}
\label{lemma-fibred-in-groupoids-over-fibred-in-groupoids}
Soient $\mathcal{A} \to \mathcal{B} \to \mathcal{C}$ des foncteurs entre
catégories. Si $\mathcal{A}$ est fibrée en groupoïdes au-dessus de $\mathcal{B}$
et $\mathcal{B}$ est fibrée en groupoïdes au-dessus de $\mathcal{C}$, alors
$\mathcal{A}$ est fibrée en groupoïdes au-dessus de $\mathcal{C}$.
\end{lemma}

\begin{proof}
On peut le prouver directement à partir de la définition. Toutefois, nous raisonnerons
à l'aide du critère du lemme \ref{lemma-fibred-groupoids}.
D'après le lemme \ref{lemma-fibred-over-fibred}, $\mathcal{A}$
est fibrée au-dessus de $\mathcal{C}$. Pour achever la preuve, montrons que la catégorie
fibre $\mathcal{A}_U$ est un groupoïde pour tout objet $U$ de $\mathcal{C}$.
En effet, si $x \to y$ est un morphisme de $\mathcal{A}_U$, alors son
image dans $\mathcal{B}$ est un isomorphisme puisque $\mathcal{B}_U$ est
un groupoïde. Mais alors $x \to y$ est un isomorphisme, par exemple d'après le
lemme \ref{lemma-composition-cartesian} et le fait que tout
morphisme de $\mathcal{A}$ est fortement $\mathcal{B}$-cartésien
(voir le lemme \ref{lemma-fibred-groupoids}).
\end{proof}

\begin{lemma}
\label{lemma-fibred-groupoids-fibre-product-goes-up}
Soit $p : \mathcal{S} \to \mathcal{C}$ une catégorie fibrée en groupoïdes.
Soient $x \to y$ et $z \to y$ des morphismes de $\mathcal{S}$.
Si $p(x) \times_{p(y)} p(z)$ existe, alors
$x \times_y z$ existe et $p(x \times_y z) = p(x) \times_{p(y)} p(z)$.
\end{lemma}

\begin{proof}
Cela découle du
lemme \ref{lemma-fibred-category-representable-goes-up}.
\end{proof}

\begin{lemma}
\label{lemma-ameliorate-morphism-categories-fibred-groupoids}
Soit $\mathcal{C}$ une catégorie. Soit $F : \mathcal{X} \to \mathcal{Y}$
un $1$-morphisme de catégories fibrées en groupoïdes au-dessus de $\mathcal{C}$.
Il existe une factorisation $\mathcal{X} \to \mathcal{X}' \to \mathcal{Y}$
par des $1$-morphismes de catégories fibrées en groupoïdes au-dessus de $\mathcal{C}$,
telle que $\mathcal{X} \to \mathcal{X}'$ soit une équivalence au-dessus de $\mathcal{C}$
et que $\mathcal{X}'$ soit une catégorie fibrée en groupoïdes au-dessus de
$\mathcal{Y}$.
\end{lemma}

\begin{proof}
Notons $p : \mathcal{X} \to \mathcal{C}$ et $q : \mathcal{Y} \to \mathcal{C}$
les foncteurs structuraux. Construisons explicitement $\mathcal{X}'$ comme suit.
Un objet de $\mathcal{X}'$ est un quadruplet $(U, x, y, f)$ où
$x \in \Ob(\mathcal{X}_U)$, $y \in \Ob(\mathcal{Y}_U)$
et $f : F(x) \to y$ est un isomorphisme dans $\mathcal{Y}_U$.
Un morphisme $(a, b) : (U, x, y, f) \to (U', x', y', f')$ est donné
par $a : x \to x'$ et $b : y \to y'$, avec $p(a) = q(b)$ et
tels que $f' \circ F(a) = b \circ f$. Autrement dit,
$\mathcal{X}' = \mathcal{X} \times_{F, \mathcal{Y}, \text{id}} \mathcal{Y}$
avec la construction du $2$-produit fibré du
lemme \ref{lemma-2-product-categories-over-C}.
D'après le
lemme \ref{lemma-2-product-fibred-categories},
$\mathcal{X}'$ est une catégorie fibrée en groupoïdes au-dessus de
$\mathcal{C}$ et $\mathcal{X}' \to \mathcal{Y}$ est un morphisme de
catégories au-dessus de $\mathcal{C}$. Comme foncteur $\mathcal{X} \to \mathcal{X}'$, prenons
$x \mapsto (p(x), x, F(x), \text{id}_{F(x)})$ sur les objets et
$(a : x \to x') \mapsto (a, F(a))$ sur les morphismes. Il est clair que
la composée $\mathcal{X} \to \mathcal{X}' \to \mathcal{Y}$
est égale à $F$. Nous omettons de vérifier que
$\mathcal{X} \to \mathcal{X}'$ est une équivalence de catégories fibrées au-dessus de
$\mathcal{C}$.

\medskip\noindent
Enfin, il reste à montrer que $\mathcal{X}' \to \mathcal{Y}$ est une catégorie
fibrée en groupoïdes. Soit $b : y' \to y$ un morphisme de $\mathcal{Y}$
et soit $(U, x, y, f)$ un objet de $\mathcal{X}'$ au-dessus de $y$.
Puisque $\mathcal{X}$ est fibrée en groupoïdes au-dessus de $\mathcal{C}$, nous
pouvons trouver un morphisme $a : x' \to x$ au-dessus de $U' = q(y') \to q(y) = U$.
Puisque $\mathcal{Y}$ est fibrée en groupoïdes au-dessus de $\mathcal{C}$ et que
$F(x') \to F(x)$ et $y' \to y$ sont tous deux au-dessus du même morphisme $U' \to U$,
nous pouvons trouver $f' : F(x') \to y'$ au-dessus de $\text{id}_{U'}$ tel que
$f \circ F(a) = b \circ f'$. Nous obtenons ainsi
$(a, b) : (U', x', y', f') \to (U, x, y, f)$.
Ceci vérifie la première condition (1) de la
définition \ref{definition-fibred-groupoids}.
Pour vérifier (2), soient
$(a, b) : (U', x', y', f') \to (U, x, y, f)$ et
$(a', b') : (U'', x'', y'', f'') \to (U, x, y, f)$ des morphismes de
$\mathcal{X}'$, et soit $b'' : y' \to y''$ un morphisme de $\mathcal{Y}$
tel que $b' \circ b'' = b$. Il faut montrer qu'il existe
un unique morphisme $a'' : x' \to x''$ tel que
$f'' \circ F(a'') = b'' \circ f'$ et tel que
$(a', b') \circ (a'', b'') = (a, b)$. Puisque $\mathcal{X}$ est fibrée
en groupoïdes, nous savons qu'il existe un unique morphisme
$a'' : x' \to x''$ tel que $a' \circ a'' = a$ et $p(a'') = q(b'')$.
Puisque $\mathcal{Y}$ est fibrée en groupoïdes, nous voyons que
$F(a'')$ est l'unique morphisme $F(x') \to F(x'')$ tel que
$F(a') \circ F(a'') = F(a)$ et $q(F(a'')) = q(b'')$. La relation
$f'' \circ F(a'') = b'' \circ f'$ découle alors de cette unicité et des relations données
$f \circ F(a) = b \circ f'$ et $f \circ F(a') = b' \circ f''$.
\end{proof}

\begin{lemma}
\label{lemma-amelioration-unique}
Soit $\mathcal{C}$ une catégorie. Soit $F : \mathcal{X} \to \mathcal{Y}$
un $1$-morphisme de catégories fibrées en groupoïdes au-dessus de $\mathcal{C}$.
Supposons donné un diagramme $2$-commutatif
$$
\xymatrix{
\mathcal{X}' \ar[rd]_f &
\mathcal{X} \ar[l]^a \ar[d]^F \ar[r]_b &
\mathcal{X}'' \ar[ld]^g \\
& \mathcal{Y}
}
$$
où $a$ et $b$ sont des équivalences de catégories au-dessus de $\mathcal{C}$
et où $f$ et $g$ définissent des catégories fibrées en groupoïdes. Il existe alors
une équivalence $h : \mathcal{X}'' \to \mathcal{X}'$ de catégories au-dessus de
$\mathcal{Y}$ telle que $h \circ b$ soit $2$-isomorphe à $a$ comme $1$-morphismes
de catégories au-dessus de $\mathcal{C}$. Si le diagramme ci-dessus est effectivement
commutatif, on peut faire en sorte que $h \circ b$ soit $2$-isomorphe à $a$ comme
$1$-morphismes de catégories au-dessus de $\mathcal{Y}$.
\end{lemma}

\begin{proof}
Nous montrerons que $\mathcal{X}'$ et $\mathcal{X}''$, au-dessus de $\mathcal{Y}$,
sont toutes deux équivalentes à la catégorie fibrée en groupoïdes
$\mathcal{X} \times_{F, \mathcal{Y}, \text{id}} \mathcal{Y}$
au-dessus de $\mathcal{Y}$ ; voir la preuve du
lemme \ref{lemma-ameliorate-morphism-categories-fibred-groupoids}.
Choisissons un quasi-inverse $b^{-1} : \mathcal{X}'' \to \mathcal{X}$ dans la
$2$-catégorie des catégories au-dessus de $\mathcal{C}$.
Puisque le triangle de droite du diagramme est $2$-commutatif, le diagramme
$$
\xymatrix{
\mathcal{X} \ar[d]_F & \mathcal{X}'' \ar[l]^{b^{-1}} \ar[d]^g \\
\mathcal{Y} & \mathcal{Y} \ar[l]
}
$$
est $2$-commutatif. Nous obtenons donc un $1$-morphisme
$c : \mathcal{X}'' \to
\mathcal{X} \times_{F, \mathcal{Y}, \text{id}} \mathcal{Y}$
par la propriété universelle du $2$-produit fibré. De plus, $c$
est un morphisme de catégories au-dessus de $\mathcal{Y}$ (!) et une équivalence
(puisque $b$ est supposé être une équivalence ; voir le
lemme \ref{lemma-equivalence-2-fibre-product}).
Ainsi, $c$ est une équivalence dans la $2$-catégorie des catégories fibrées
en groupoïdes au-dessus de $\mathcal{Y}$ d'après le
lemme \ref{lemma-equivalence-fibred-categories}.

\medskip\noindent
Il reste à construire un $2$-isomorphisme entre $c \circ b$ et
le foncteur $d : \mathcal{X} \to
\mathcal{X} \times_{F, \mathcal{Y}, \text{id}} \mathcal{Y}$,
$x \mapsto (p(x), x, F(x), \text{id}_{F(x)})$,
construit dans la preuve du
lemme \ref{lemma-ameliorate-morphism-categories-fibred-groupoids}.
Soient $\alpha : F \to g \circ b$ et $\beta : b^{-1} \circ b \to \text{id}$
des $2$-isomorphismes entre $1$-morphismes de catégories au-dessus de $\mathcal{C}$.
Remarquons que $c \circ b$ est donné sur les objets par la règle
$$
x \mapsto (p(x), b^{-1}(b(x)), g(b(x)), \alpha_x \circ F(\beta_x))
$$
On voit alors que
$$
(\beta_x, \alpha_x) :
(p(x), x, F(x), \text{id}_{F(x)})
\longrightarrow
(p(x), b^{-1}(b(x)), g(b(x)), \alpha_x \circ F(\beta_x))
$$
est un isomorphisme fonctoriel qui donne notre $2$-morphisme
$d \to c \circ b$. Enfin, si le diagramme est commutatif, alors
$\alpha_x$ est l'identité pour tout $x$, et ce
$2$-morphisme est donc un $2$-morphisme dans la $2$-catégorie des catégories
au-dessus de $\mathcal{Y}$.
\end{proof}


































\section{Préfaisceaux de catégories}
\label{section-presheaves-categories}

\noindent
Dans cette section, nous comparons la notion de catégorie fibrée
à la notion étroitement apparentée de « préfaisceau de catégories ».
La construction de base est expliquée dans l'exemple suivant.

\begin{example}
\label{example-functor-categories}
Soit $\mathcal{C}$ une catégorie.
Supposons que $F : \mathcal{C}^{opp} \to \textit{Cat}$ soit un foncteur
vers la $2$-catégorie des catégories ; voir la
définition \ref{definition-functor-into-2-category}.
Pour $f : V \to U$ dans $\mathcal{C}$, nous écrirons de manière suggestive
$F(f) = f^\ast$ pour le foncteur de $F(U)$ vers $F(V)$.
À partir de là, nous pouvons construire une catégorie fibrée $\mathcal{S}_F$ au-dessus de
$\mathcal{C}$ comme suit. Posons
$$
\Ob(\mathcal{S}_F) =
\{(U, x) \mid U\in \Ob(\mathcal{C}), x\in \Ob(F(U))\}.
$$
Pour $(U, x), (V, y) \in \Ob(\mathcal{S}_F)$, posons
\begin{align*}
\Mor_{\mathcal{S}_F}((V, y), (U, x)) & =
\{ (f, \phi) \mid f \in \Mor_\mathcal{C}(V, U),
\phi \in \Mor_{F(V)}(y, f^\ast x)\} \\
& =
\coprod\nolimits_{f \in \Mor_\mathcal{C}(V, U)}
\Mor_{F(V)}(y, f^\ast x)
\end{align*}
Pour définir la composition, nous utilisons l'égalité $g^\ast \circ f^\ast =
(f \circ g)^\ast$ pour une paire de morphismes composables de $\mathcal{C}$
(par définition d'un foncteur vers une $2$-catégorie).
Plus précisément, nous définissons la composée de $\psi : z \to g^\ast y$ et
$ \phi : y \to f^\ast x$ comme $ g^\ast(\phi) \circ \psi$. Le foncteur
$p_F : \mathcal{S}_F \to \mathcal{C}$ est donné par la règle
$(U, x) \mapsto U$.
Vérifions qu'il s'agit bien d'une catégorie fibrée.
Étant donnés $f: V \to U$ dans $\mathcal{C}$ et $(U, x)$, objet au-dessus de $U$, nous
affirmons que $(f, \text{id}_{f^\ast x}): (V, {f^\ast x}) \to (U, x)$ est un
relèvement fortement cartésien de $f$.
Il faut montrer qu'un $h$ dans le diagramme de gauche
détermine $(h, \nu)$ dans celui de droite :
$$
\xymatrix{
V \ar[r]^f &
U &
(V, f^*x) \ar[r]^{(f, \text{id}_{f^*x})} &
(U, x) \\
W \ar@{-->}[u]^h \ar[ru]_g & &
(W, z) \ar@{-->}[u]^{(h, \nu)} \ar[ru]_{(g, \psi)} &
}
$$
Il suffit de prendre $\nu = \psi$, ce qui convient puisque $f \circ h = g$
et donc $g^*x = h^*f^*x$. De plus, c'est l'unique relèvement
qui rende commutatif le diagramme (de droite).
\end{example}

\begin{definition}
\label{definition-split-fibred-category}
Soit $\mathcal{C}$ une catégorie.
Supposons que $F : \mathcal{C}^{opp} \to \textit{Cat}$ soit un foncteur
vers la $2$-catégorie des catégories.
Nous noterons $p_F : \mathcal{S}_F \to \mathcal{C}$ la
catégorie fibrée construite dans
l'exemple \ref{example-functor-categories}.
Une {\it catégorie fibrée scindée} est une catégorie fibrée isomorphe (!),
au-dessus de $\mathcal{C}$, à l'une de ces catégories {\it $\mathcal{S}_F$}.
\end{definition}

\begin{lemma}
\label{lemma-when-split}
Soit $\mathcal{C}$ une catégorie.
Soit $\mathcal{S}$ une catégorie fibrée au-dessus de $\mathcal{C}$.
Alors $\mathcal{S}$ est scindée si et seulement si, pour un certain clivage
(voir la définition \ref{definition-pullback-functor-fibred-category}),
les foncteurs de changement de base
$(f \circ g)^*$ et $g^* \circ f^*$ sont égaux.
\end{lemma}

\begin{proof}
Cela découle immédiatement des définitions.
\end{proof}

\begin{lemma}
\label{lemma-fibred-strict}
Soit $ p : \mathcal{S} \to \mathcal{C}$ une catégorie fibrée.
Il existe un foncteur contravariant $F : \mathcal{C} \to \textit{Cat}$
tel que $\mathcal{S}$ soit équivalente à $\mathcal{S}_F$
dans la $2$-catégorie des catégories fibrées au-dessus de $\mathcal{C}$. Autrement
dit, toute catégorie fibrée est équivalente à une catégorie scindée.
\end{lemma}

\begin{proof}
Choisissons un clivage (voir la
définition \ref{definition-pullback-functor-fibred-category}).
Le lemme \ref{lemma-fibred} fournit des foncteurs de changement de base $f^*$ pour
tout morphisme $f$ de $\mathcal{C}$.

\medskip\noindent
Construisons une nouvelle catégorie $\mathcal{S}'$ comme suit.
Les objets de $\mathcal{S}'$ sont les paires $(x, f)$
formées d'un morphisme $f : V \to U$ de $\mathcal{C}$
et d'un objet $x$ de $\mathcal{S}$ au-dessus de $U$, c'est-à-dire
$x\in \Ob(\mathcal{S}_U)$. Le foncteur
$p' : \mathcal{S}' \to \mathcal{C}$ envoie la paire $(x, f)$ sur la source
du morphisme $f$ ; autrement dit, $p'(x, f : V\to U) = V$. Un morphisme
$\varphi : (x_1, f_1: V_1 \to U_1) \to (x_2, f_2 : V_2 \to U_2)$ est donné par une
paire $(\varphi, g)$ formée d'un morphisme $g : V_1 \to V_2$ et d'un morphisme
$\varphi : f_1^\ast x_1 \to f_2^\ast x_2$ tel que $p(\varphi) = g$. Il est
immédiat de définir la loi de composition : $(\varphi, g) \circ (\psi, h) =
(\varphi \circ \psi, g\circ h)$ pour toute paire de morphismes composables.
Il existe un foncteur naturel $\mathcal{S} \to \mathcal{S}'$ qui envoie simplement
$x$ au-dessus de $U$ sur la paire $(x, \text{id}_U)$.

\medskip\noindent
Il nous faut à présent vérifier que $p'$ fait de $\mathcal{S}'$ une
catégorie fibrée au-dessus de $\mathcal{C}$, et que
$\mathcal{S} \to \mathcal{S}'$ est une équivalence de catégories au-dessus de
$\mathcal{C}$ qui envoie les morphismes fortement cartésiens sur des morphismes fortement
cartésiens. Nous omettons ces vérifications.

\medskip\noindent
Enfin, nous pouvons définir des foncteurs de changement de base sur $\mathcal{S}'$
en posant $g^\ast(x, f) = (x, f \circ g)$ sur les objets si
$g : V' \to V$ et $f : V \to U$. Sur les morphismes
$(\varphi, \text{id}_V) : (x_1, f_1) \to (x_2, f_2)$
de $\mathcal{S}'_V$, nous posons $g^\ast(\varphi, \text{id}_V) =
(g^\ast\varphi, \text{id}_{V'})$, où nous utilisons les identifications uniques
$g^\ast f_i^\ast x_i = (f_i \circ g)^\ast x_i$ du lemme
\ref{lemma-fibred} pour considérer $g^\ast\varphi$ comme un morphisme de
$(f_1 \circ g)^\ast x_1$ vers $(f_2 \circ g)^\ast x_2$. Manifestement, ces foncteurs
de changement de base $g^\ast$ vérifient
$g_1^\ast \circ g_2^\ast = (g_2\circ g_1)^\ast$ ; autrement dit, $\mathcal{S}'$
est scindée, comme voulu.
\end{proof}



















\section{Préfaisceaux de groupoïdes}
\label{section-presheaves-groupoids}

\noindent
Dans cette section, nous comparons la notion de catégorie fibrée en groupoïdes
à la notion étroitement apparentée de « préfaisceau de groupoïdes ». La construction
de base est expliquée dans l'exemple suivant.

\begin{example}
\label{example-functor-groupoids}
Cet exemple est l'analogue de
l'exemple \ref{example-functor-categories},
pour les « préfaisceaux de groupoïdes » plutôt que pour les « préfaisceaux de catégories ».
Le résultat sera une catégorie fibrée en groupoïdes plutôt qu'une catégorie fibrée.
Supposons que $F : \mathcal{C}^{opp} \to \textit{Groupoïdes}$ soit un foncteur
vers la $2$-catégorie de groupoïdes ; voir la
définition \ref{definition-functor-into-2-category}.
Pour $f : V \to U$ dans $\mathcal{C}$, nous écrirons de manière suggestive
$F(f) = f^\ast$ pour le foncteur de $F(U)$ vers $F(V)$.
Construisons une catégorie $\mathcal{S}_F$ fibrée en groupoïdes au-dessus de $\mathcal{C}$
comme suit. Posons
$$
\Ob(\mathcal{S}_F) =
\{(U, x) \mid U\in \Ob(\mathcal{C}), x\in \Ob(F(U))\}.
$$
Pour $(U, x), (V, y) \in \Ob(\mathcal{S}_F)$, posons
\begin{align*}
\Mor_{\mathcal{S}_F}((V, y), (U, x))
& =
\{ (f, \phi) \mid f \in \Mor_\mathcal{C}(V, U),
\phi \in \Mor_{F(V)}(y, f^\ast x)\} \\
& =
\coprod\nolimits_{f \in \Mor_\mathcal{C}(V, U)}
\Mor_{F(V)}(y, f^\ast x)
\end{align*}
Pour définir la composition, nous utilisons l'égalité $g^\ast \circ f^\ast =
(f \circ g)^\ast$ pour une paire de morphismes composables de $\mathcal{C}$
(par définition d'un foncteur vers une $2$-catégorie).
Plus précisément, nous définissons la composée de $\psi : z \to g^\ast y$ et
$ \phi : y \to f^\ast x$ comme $ g^\ast(\phi) \circ \psi$. Le foncteur
$p_F : \mathcal{S}_F \to \mathcal{C}$ est donné par la règle $(U, x) \mapsto U$.
La condition selon laquelle $F(U)$ est un groupoïde pour tout $U$ garantit que
$\mathcal{S}_F$ est fibrée en groupoïdes au-dessus de $\mathcal{C}$, puisque nous avons
déjà vu dans
l'exemple \ref{example-functor-categories}
que $\mathcal{S}_F$ est une catégorie fibrée ; voir le
lemme \ref{lemma-fibred-groupoids}.
Mais nous pouvons aussi démontrer directement les conditions (1), (2) de la
définition \ref{definition-fibred-groupoids}
comme suit : (1) Les relèvements de
morphismes existent : étant donnés $f: V \to U$ dans $\mathcal{C}$ et $(U, x)$,
objet de $\mathcal{S}_F$ au-dessus de $U$,
$(f, \text{id}_{f^\ast x}): (V, {f^\ast x}) \to (U, x)$ est un relèvement de $f$.
(2) Supposons données les flèches pleines des diagrammes suivants :
$$
\xymatrix{
V \ar[r]^f & U & (V, y) \ar[r]^{(f, \phi)} & (U, x) \\
W \ar@{-->}[u]^h \ar[ru]_g & &
(W, z) \ar@{-->}[u]^{(h, \nu)} \ar[ru]_{(g, \psi)} & \\
}
$$
Alors, pour les flèches pointillées, on a $\nu = (h^\ast \phi)^{-1} \circ \psi$ ;
ainsi, $h$ étant donné, il existe un $\nu$, unique par unicité des inverses.
\end{example}

\begin{definition}
\label{definition-split-category-fibred-in-groupoids}
Soit $\mathcal{C}$ une catégorie.
Supposons que $F : \mathcal{C}^{opp} \to \textit{Groupoïdes}$ soit un foncteur
vers la $2$-catégorie des groupoïdes.
Nous noterons $p_F : \mathcal{S}_F \to \mathcal{C}$ la
catégorie fibrée en groupoïdes construite dans
l'exemple \ref{example-functor-groupoids}.
Une {\it catégorie fibrée en groupoïdes scindée} est une
catégorie fibrée en groupoïdes isomorphe (!),
au-dessus de $\mathcal{C}$, à l'une de ces catégories {\it $\mathcal{S}_F$}.
\end{definition}

\begin{lemma}
\label{lemma-fibred-groupoids-strict}
Soit $ p : \mathcal{S} \to \mathcal{C}$ une catégorie fibrée en groupoïdes.
Il existe un foncteur contravariant $F : \mathcal{C} \to \textit{Groupoïdes}$
tel que $\mathcal{S}$ soit équivalente à $\mathcal{S}_F$ au-dessus de $\mathcal{C}$.
Autrement dit, toute catégorie fibrée en groupoïdes est équivalente à une catégorie scindée.
\end{lemma}

\begin{proof}
Choisissons un clivage (voir la
définition \ref{definition-pullback-functor-fibred-category}).
Les lemmes \ref{lemma-fibred} et \ref{lemma-fibred-groupoids}
fournissent des foncteurs de changement de base $f^*$ pour
tout morphisme $f$ de $\mathcal{C}$.

\medskip\noindent
Construisons une nouvelle catégorie $\mathcal{S}'$ comme suit.
Les objets de $\mathcal{S}'$ sont les paires $(x, f)$
formées d'un morphisme $f : V \to U$ de $\mathcal{C}$
et d'un objet $x$ de $\mathcal{S}$ au-dessus de $U$, c'est-à-dire
$x\in \Ob(\mathcal{S}_U)$. Le foncteur
$p' : \mathcal{S}' \to \mathcal{C}$ envoie la paire $(x, f)$ sur la source
du morphisme $f$ ; autrement dit, $p'(x, f : V\to U) = V$. Un morphisme
$\varphi : (x_1, f_1: V_1 \to U_1) \to (x_2, f_2 : V_2 \to U_2)$ est donné par une
paire $(\varphi, g)$ formée d'un morphisme $g : V_1 \to V_2$ et d'un morphisme
$\varphi : f_1^\ast x_1 \to f_2^\ast x_2$ tel que $p(\varphi) = g$. Il est
immédiat de définir la loi de composition : $(\varphi, g) \circ (\psi, h) =
(\varphi \circ \psi, g\circ h)$ pour toute paire de morphismes composables.
Il existe un foncteur naturel $\mathcal{S} \to \mathcal{S}'$ qui envoie simplement
$x$ au-dessus de $U$ sur la paire $(x, \text{id}_U)$.

\medskip\noindent
Il nous faut à présent vérifier que $p'$ fait de $\mathcal{S}'$ une catégorie
fibrée en groupoïdes au-dessus de $\mathcal{C}$, et que
$\mathcal{S} \to \mathcal{S}'$ est une équivalence de catégories au-dessus de
$\mathcal{C}$. Nous omettons ces vérifications.

\medskip\noindent
Enfin, nous pouvons définir des foncteurs de changement de base sur $\mathcal{S}'$
en posant $g^\ast(x, f) = (x, f \circ g)$ sur les objets si
$g : V' \to V$ et $f : V \to U$. Sur les morphismes
$(\varphi, \text{id}_V) : (x_1, f_1) \to (x_2, f_2)$
de $\mathcal{S}'_V$, nous posons $g^\ast(\varphi, \text{id}_V) =
(g^\ast\varphi, \text{id}_{V'})$, où nous utilisons les identifications uniques
$g^\ast f_i^\ast x_i = (f_i \circ g)^\ast x_i$ du lemme
\ref{lemma-fibred-groupoids} pour considérer $g^\ast\varphi$ comme un morphisme de
$(f_1 \circ g)^\ast x_1$ vers $(f_2 \circ g)^\ast x_2$. Manifestement, ces foncteurs
de changement de base $g^\ast$ vérifient
$g_1^\ast \circ g_2^\ast = (g_2\circ g_1)^\ast$ ; autrement dit, $\mathcal{S}'$
est scindée, comme voulu.
\end{proof}

\noindent
Nous verrons une autre preuve de ce lemme dans la
section \ref{section-representable-1-morphisms}.





\section{Catégories fibrées en ensembles}
\label{section-fibred-in-sets}

\begin{definition}
\label{definition-discrete}
Une catégorie est dite {\it discrète} si ses seuls morphismes sont les
identités.
\end{definition}

\noindent
Une catégorie discrète ne possède qu'une seule donnée intéressante :
son ensemble d'objets. Ainsi, nous confondons parfois les catégories discrètes
avec les ensembles.

\begin{definition}
\label{definition-category-fibred-sets}
Soit $\mathcal{C}$ une catégorie.
Une {\it catégorie fibrée en ensembles}, ou une {\it catégorie fibrée
en catégories discrètes}, est une catégorie fibrée en groupoïdes dont
toutes les catégories fibres sont discrètes.
\end{definition}

\noindent
Nous voulons préciser la relation entre les catégories fibrées en ensembles
et les préfaisceaux (voir la définition \ref{definition-presheaf}).
Pour ce faire, il est naturel de commencer par la définition suivante.

\begin{definition}
\label{definition-categories-fibred-in-sets-over-C}
Soit $\mathcal{C}$ une catégorie.
La {\it $2$-catégorie des catégories fibrées en ensembles au-dessus de
$\mathcal{C}$} est la sous-$2$-catégorie de la $2$-catégorie des catégories
fibrées en groupoïdes au-dessus de $\mathcal{C}$ (voir la
définition \ref{definition-categories-fibred-in-groupoids-over-C})
définie comme suit :
\begin{enumerate}
\item Ses objets sont les catégories
$p : \mathcal{S} \to \mathcal{C}$ fibrées en ensembles.
\item Ses $1$-morphismes $(\mathcal{S}, p) \to (\mathcal{S}', p')$
sont les foncteurs $G : \mathcal{S} \to \mathcal{S}'$ tels que
$p' \circ G = p$ (puisque tout morphisme est fortement cartésien,
$G$ les préserve automatiquement).
\item Ses $2$-morphismes $t : G \to H$ pour
$G, H : (\mathcal{S}, p) \to (\mathcal{S}', p')$
sont les morphismes de foncteurs
tels que $p'(t_x) = \text{id}_{p(x)}$
pour tout $x \in \Ob(\mathcal{S})$.
\end{enumerate}
\end{definition}

\noindent
Notons que tout $2$-morphisme est automatiquement un isomorphisme.
Cette $2$-catégorie est donc en fait une $(2, 1)$-catégorie.
Voici le lemme obligé sur l'existence des $2$-produits fibrés.

\begin{lemma}
\label{lemma-2-product-categories-fibred-sets}
Soit $\mathcal{C}$ une catégorie.
La 2-catégorie des catégories fibrées en ensembles au-dessus de $\mathcal{C}$
admet des 2-produits fibrés. Plus précisément, le 2-produit fibré décrit dans le
lemme \ref{lemma-2-product-categories-over-C}
est une catégorie fibrée en ensembles lorsque les données de départ le sont.
\end{lemma}

\begin{proof}
Omis.
\end{proof}

\begin{example}
\label{example-presheaf}
Cet exemple est l'analogue des
exemples \ref{example-functor-categories} et
\ref{example-functor-groupoids}
pour les préfaisceaux au lieu des « préfaisceaux de catégories ».
Le résultat sera une catégorie fibrée en ensembles plutôt qu'une catégorie fibrée.
Supposons que $F : \mathcal{C}^{opp} \to \textit{Ensembles}$ soit un préfaisceau.
Pour $f : V \to U$ dans $\mathcal{C}$, nous écrirons de manière suggestive
$F(f) = f^\ast : F(U) \to F(V)$.
Nous construisons une catégorie $\mathcal{S}_F$ fibrée en ensembles au-dessus de
$\mathcal{C}$ comme suit. Posons
$$
\Ob(\mathcal{S}_F) =
\{(U, x) \mid U \in \Ob(\mathcal{C}), x \in F(U)\}.
$$
Pour $(U, x), (V, y) \in \Ob(\mathcal{S}_F)$, nous posons
\begin{align*}
\Mor_{\mathcal{S}_F}((V, y), (U, x))
& =
\{f \in \Mor_\mathcal{C}(V, U) \mid f^*x = y\}
\end{align*}
La composition est héritée de celle de $\mathcal{C}$ ; cela fonctionne car
$g^\ast \circ f^\ast = (f \circ g)^\ast$
pour une paire de morphismes composables de $\mathcal{C}$.
Le foncteur $p_F : \mathcal{S}_F \to \mathcal{C}$
est donné par la règle $(U, x) \mapsto U$.
Puisque toute catégorie fibre $\mathcal{S}_{F, U}$ est discrète, d'ensemble
sous-jacent $F(U)$, et que nous avons déjà vu dans
l'exemple \ref{example-functor-groupoids}
que $\mathcal{S}_F$ est une catégorie fibrée en groupoïdes,
nous concluons que $\mathcal{S}_F$ est fibrée en ensembles.
\end{example}

\begin{lemma}
\label{lemma-2-category-fibred-sets}
\begin{slogan}
Les catégories fibrées en ensembles sont exactement les préfaisceaux.
\end{slogan}
Soit $\mathcal{C}$ une catégorie.
Les seuls $2$-morphismes entre catégories fibrées en ensembles sont les identités.
Autrement dit, la $2$-catégorie des catégories fibrées en ensembles est une catégorie.
De plus, il existe une équivalence de catégories
$$
\left\{
\begin{matrix}
\text{la catégorie des préfaisceaux}\\
\text{d'ensembles sur }\mathcal{C}
\end{matrix}
\right\}
\leftrightarrow
\left\{
\begin{matrix}
\text{la catégorie des catégories}\\
\text{fibrées en ensembles au-dessus de }\mathcal{C}
\end{matrix}
\right\}
$$
Le foncteur de gauche à droite est la construction
$F \to \mathcal{S}_F$ étudiée dans
l'exemple \ref{example-presheaf}.
Le foncteur de droite à gauche associe à $p : \mathcal{S} \to \mathcal{C}$
le préfaisceau d'objets $U \mapsto \Ob(\mathcal{S}_U)$.
\end{lemma}

\begin{proof}
La première assertion est claire, puisque les seuls morphismes des catégories
fibres sont les identités.

\medskip\noindent
Supposons que $p :
\mathcal{S} \to \mathcal{C}$ soit fibrée en ensembles. Soit $f : V \to U$
un morphisme de $\mathcal{C}$ et soit $x \in \Ob(\mathcal{S}_U)$.
Il existe alors exactement un choix pour l'objet $f^\ast x$. Nous voyons donc que
$(f \circ g)^\ast x = g^\ast(f^\ast x)$ pour $f, g$ comme dans le lemme
\ref{lemma-fibred-groupoids}. Il s'ensuit que l'on peut considérer les
applications $U \mapsto \Ob(\mathcal{S}_U)$ et $f \mapsto f^\ast$
comme un préfaisceau sur $\mathcal{C}$.
\end{proof}

\noindent
Voici un exemple important de catégorie fibrée en ensembles.

\begin{example}
\label{example-fibred-category-from-functor-of-points}
Soit $\mathcal{C}$ une catégorie. Soit $X \in \Ob(\mathcal{C})$.
Considérons le préfaisceau représentable $h_X = \Mor_\mathcal{C}(-, X)$
(voir l'exemple \ref{example-hom-functor}).
D'autre part, considérons la catégorie $p : \mathcal{C}/X \to \mathcal{C}$
de l'exemple \ref{example-category-over-X}.
La catégorie fibre $(\mathcal{C}/X)_U$ a pour objets les morphismes
$h : U \to X$, et pour seuls morphismes les identités. Nous voyons donc que,
sous la correspondance du
lemme \ref{lemma-2-category-fibred-sets},
nous avons
$$
h_X \longleftrightarrow \mathcal{C}/X.
$$
Autrement dit, la catégorie $\mathcal{C}/X$ est canoniquement équivalente
à la catégorie $\mathcal{S}_{h_X}$ associée
à $h_X$ dans
l'exemple \ref{example-presheaf}.
\end{example}

\noindent
Pour cette raison, il est tentant de définir un objet « représentable » de la
2-catégorie des catégories fibrées en groupoïdes comme une catégorie fibrée en
ensembles dont le préfaisceau associé est représentable. Ce ne serait toutefois
pas une bonne définition, car nous préférons disposer d'une notion qui soit
invariante par équivalences. Pour préciser ce point, nous étudions exactement
quelles catégories fibrées en groupoïdes sont équivalentes à des catégories
fibrées en ensembles.








\section{Catégories fibrées en setoïdes}
\label{section-fibred-in-setoids}

\begin{definition}
\label{definition-setoid}
Appelons {\it setoïde}\footnote{Un ensemble sous stéroïdes !?} toute catégorie
qui est un groupoïde dans lequel tout objet
possède exactement un automorphisme : l'identité.
\end{definition}

\noindent
Si $C$ est un ensemble muni d'une relation d'équivalence $\sim$, on peut en faire
un setoïde $\mathcal{C}$ comme suit : $\Ob(\mathcal{C}) = C$ et
$\Mor_\mathcal{C}(x, y) = \emptyset$ sauf si $x \sim y$, auquel
cas on pose $\Mor_\mathcal{C}(x, y) = \{1\}$. La transitivité de
$\sim$ signifie que l'on peut composer les morphismes. Réciproquement, tout
setoïde définit sur ses objets la relation d'équivalence « être isomorphes »,
à partir de laquelle on retrouve la catégorie (à isomorphisme
unique près, et non simplement à équivalence près) par le procédé qui vient d'être décrit.

\medskip\noindent
Les catégories discrètes sont des setoïdes. Pour tout setoïde $\mathcal{C}$,
il existe un procédé canonique qui produit une catégorie
discrète équivalente : on remplace $\Ob(\mathcal{C})$ par l'ensemble des classes
d'isomorphisme (et l'on ajoute les identités). En termes d'ensembles
munis d'une relation d'équivalence, cela correspond au passage au quotient
par la relation d'équivalence.

\begin{definition}
\label{definition-category-fibred-setoids}
Soit $\mathcal{C}$ une catégorie. Une {\it catégorie fibrée en setoïdes} est
une catégorie fibrée en groupoïdes dont toutes les catégories fibres sont des
setoïdes.
\end{definition}

\noindent
Nous préciserons ci-dessous la relation entre les catégories fibrées en
setoïdes et les catégories fibrées en ensembles.

\begin{definition}
\label{definition-categories-fibred-in-setoids-over-C}
Soit $\mathcal{C}$ une catégorie.
La {\it $2$-catégorie des catégories fibrées en setoïdes
au-dessus de $\mathcal{C}$} est la sous-$2$-catégorie de la $2$-catégorie des
catégories fibrées en groupoïdes au-dessus de $\mathcal{C}$ (voir la
définition \ref{definition-categories-fibred-in-groupoids-over-C})
définie comme suit :
\begin{enumerate}
\item Ses objets sont les catégories
$p : \mathcal{S} \to \mathcal{C}$ fibrées en setoïdes.
\item Ses $1$-morphismes $(\mathcal{S}, p) \to (\mathcal{S}', p')$
sont les foncteurs $G : \mathcal{S} \to \mathcal{S}'$ tels que
$p' \circ G = p$ (puisque tout morphisme est fortement cartésien,
$G$ les préserve automatiquement).
\item Ses $2$-morphismes $t : G \to H$ pour
$G, H : (\mathcal{S}, p) \to (\mathcal{S}', p')$
sont les morphismes de foncteurs
tels que $p'(t_x) = \text{id}_{p(x)}$
pour tout $x \in \Ob(\mathcal{S})$.
\end{enumerate}
\end{definition}

\noindent
Notons que tout $2$-morphisme est automatiquement un isomorphisme.
Cette $2$-catégorie est donc en fait une $(2, 1)$-catégorie.

\noindent
Voici le lemme obligé sur l'existence des $2$-produits fibrés.

\begin{lemma}
\label{lemma-2-product-categories-fibred-setoids}
Soit $\mathcal{C}$ une catégorie.
La 2-catégorie des catégories fibrées en setoïdes au-dessus de
$\mathcal{C}$ admet des 2-produits fibrés. Plus précisément, le 2-produit fibré
décrit dans le lemme \ref{lemma-2-product-categories-over-C} est une catégorie
fibrée en setoïdes lorsque les données de départ le sont.
\end{lemma}

\begin{proof}
Omis.
\end{proof}

\begin{lemma}
\label{lemma-setoid-fibres}
Soit $\mathcal{C}$ une catégorie. Soit $\mathcal{S}$ une catégorie
au-dessus de $\mathcal{C}$.
\begin{enumerate}
\item Si $\mathcal{S} \to \mathcal{S}'$ est une équivalence
au-dessus de $\mathcal{C}$, avec $\mathcal{S}'$ fibrée en ensembles au-dessus
de $\mathcal{C}$, alors
\begin{enumerate}
\item $\mathcal{S}$ est fibrée en setoïdes au-dessus de
$\mathcal{C}$, et
\item pour tout $U \in \Ob(\mathcal{C})$, l'application
$\Ob(\mathcal{S}_U) \to \Ob(\mathcal{S}'_U)$
identifie le but à l'ensemble des classes d'isomorphisme de la source.
\end{enumerate}
\item Si $p : \mathcal{S} \to \mathcal{C}$ est une catégorie fibrée en
setoïdes, alors il existe une catégorie fibrée en ensembles
$p' : \mathcal{S}' \to \mathcal{C}$ et une équivalence
$\text{can} : \mathcal{S} \to \mathcal{S}'$ au-dessus de $\mathcal{C}$.
\end{enumerate}
\end{lemma}

\begin{proof}
Démontrons (2).
Un objet de la catégorie $\mathcal{S}'$ est une paire $(U, \xi)$, où
$U \in \Ob(\mathcal{C})$ et $\xi$ est une classe d'isomorphisme d'objets
de $\mathcal{S}_U$. Un morphisme $(U, \xi) \to (V , \psi)$ est donné par un
morphisme $x \to y$, où $x \in \xi$ et $y \in \psi$. Ici, nous identifions
deux morphismes $x \to y$ et $x' \to y'$ s'ils induisent le même morphisme
$U \to V$, et si, pour certains choix d'isomorphismes $x \to x'$ dans
$\mathcal{S}_U$ et $y \to y'$ dans $\mathcal{S}_V$, les composées
$x \to x' \to y'$ et $x \to y \to y'$ coïncident. Par construction, les
applications sur les objets et les morphismes induites par $\mathcal{S} \to
\mathcal{S}'$ sont surjectives. Nous définissons la composition des morphismes dans
$\mathcal{S}'$ comme l'unique loi qui fasse de $\mathcal{S} \to \mathcal{S}'$
un foncteur. Certains détails sont omis.
\end{proof}

\noindent
Ainsi, les catégories fibrées en setoïdes sont exactement les
catégories fibrées en groupoïdes qui sont équivalentes à des catégories fibrées
en ensembles. De plus, une équivalence de catégories fibrées en ensembles est
un isomorphisme d'après le lemme \ref{lemma-2-category-fibred-sets}.

\begin{lemma}
\label{lemma-2-category-fibred-setoids}
Soit $\mathcal{C}$ une catégorie. La construction du
lemme \ref{lemma-setoid-fibres},
partie (2), donne un foncteur
$$
F :
\left\{
\begin{matrix}
\text{la 2-catégorie des catégories}\\
\text{fibrées en setoïdes au-dessus de }\mathcal{C}
\end{matrix}
\right\}
\longrightarrow
\left\{
\begin{matrix}
\text{la catégorie des catégories}\\
\text{fibrées en ensembles au-dessus de }\mathcal{C}
\end{matrix}
\right\}
$$
(voir la
définition \ref{definition-functor-into-2-category}).
Ce foncteur est une équivalence au sens suivant :
\begin{enumerate}
\item pour deux 1-morphismes quelconques
$f, g : \mathcal{S}_1 \to \mathcal{S}_2$
tels que $F(f) = F(g)$, il existe un unique 2-isomorphisme $f \to g$,
\item pour tout morphisme $h : F(\mathcal{S}_1) \to F(\mathcal{S}_2)$,
il existe un 1-morphisme $f : \mathcal{S}_1 \to \mathcal{S}_2$
tel que $F(f) = h$, et
\item toute catégorie fibrée en ensembles $\mathcal{S}$ est égale à
$F(\mathcal{S})$.
\end{enumerate}
En particulier, si l'on définit $F_i \in \textit{PSh}(\mathcal{C})$ par la
règle $F_i(U) = \Ob(\mathcal{S}_{i, U})/\cong$, on a
$$
\Mor_{\textit{Cat}/\mathcal{C}}(\mathcal{S}_1, \mathcal{S}_2)
\Big/
2\text{-isomorphisme}
=
\Mor_{\textit{PSh}(\mathcal{C})}(F_1, F_2)
$$
Plus précisément, pour tout morphisme $\phi : F_1 \to F_2$, il existe un
$1$-morphisme $f : \mathcal{S}_1 \to \mathcal{S}_2$ qui induit
$\phi$ sur les classes d'isomorphisme d'objets et
qui est unique à un unique $2$-isomorphisme près.
\end{lemma}

\begin{proof}
D'après le
lemme \ref{lemma-2-category-fibred-sets},
le but de $F$ est une catégorie ; l'assertion a donc un sens.
La construction du
lemme \ref{lemma-setoid-fibres}, partie (2),
associe à $\mathcal{S}$ la catégorie fibrée en ensembles dont la valeur au-dessus
de $U$ est l'ensemble des classes d'isomorphisme dans $\mathcal{S}_U$. Il est
donc clair qu'elle définit un foncteur comme indiqué.
Soient $f, g : \mathcal{S}_1 \to \mathcal{S}_2$
tels que $F(f) = F(g)$ comme en (1). Pour tout objet $U$ de $\mathcal{C}$
et tout objet $x$ de $\mathcal{S}_{1, U}$, l'hypothèse donne
$f(x) \cong g(x)$. Comme $\mathcal{S}_2$ est fibrée en setoïdes,
il existe un unique isomorphisme $t_x : f(x) \to g(x)$ dans
$\mathcal{S}_{2, U}$.
Il est clair que la règle $x \mapsto t_x$ donne le $2$-isomorphisme souhaité
$f \to g$. Nous omettons les démonstrations de (2) et (3).
Pour voir la dernière assertion, appliquons le
lemme \ref{lemma-2-category-fibred-sets}
pour identifier le membre de droite à
$\Mor_{\textit{Cat}/\mathcal{C}}(F(\mathcal{S}_1), F(\mathcal{S}_2))$
et appliquons (1) et (2) ci-dessus.
\end{proof}

\noindent
Voici une autre caractérisation des catégories fibrées en setoïdes parmi
toutes les catégories fibrées en groupoïdes.

\begin{lemma}
\label{lemma-characterize-fibred-setoids-inertia}
Soit $\mathcal{C}$ une catégorie.
Soit $p : \mathcal{S} \to \mathcal{C}$ une catégorie fibrée en groupoïdes.
Les conditions suivantes sont équivalentes :
\begin{enumerate}
\item $p : \mathcal{S} \to \mathcal{C}$ est une catégorie fibrée en setoïdes,
et
\item le $1$-morphisme canonique $\mathcal{I}_\mathcal{S} \to \mathcal{S}$,
voir (\ref{equation-inertia-structure-map}), est une équivalence (de catégories
au-dessus de $\mathcal{C}$).
\end{enumerate}
\end{lemma}

\begin{proof}
Supposons (2). La catégorie $\mathcal{I}_\mathcal{S}$ a pour objets
les $(x, \alpha)$ où $x \in \mathcal{S}$, avec, disons, $p(x) = U$, et où
$\alpha : x \to x$ est un morphisme de $\mathcal{S}_U$. Ainsi, si
$\mathcal{I}_\mathcal{S} \to \mathcal{S}$ est une équivalence au-dessus de
$\mathcal{C}$, alors les deux objets de toute paire
$(x, \alpha)$, $(x, \alpha')$ sont isomorphes dans la catégorie fibre de
$\mathcal{I}_\mathcal{S}$ au-dessus de $U$.
En examinant la définition des morphismes de $\mathcal{I}_\mathcal{S}$,
nous concluons que $\alpha$, $\alpha'$ sont conjugués dans le groupe
des automorphismes de $x$. En prenant $\alpha' = \text{id}_x$, nous concluons
donc que tout automorphisme de $x$ est égal à l'identité.
Puisque $\mathcal{S} \to \mathcal{C}$ est fibrée en groupoïdes, cela implique
que $\mathcal{S} \to \mathcal{C}$ est fibrée en setoïdes.
Nous omettons la démonstration de (1) $\Rightarrow$ (2).
\end{proof}

\begin{lemma}
\label{lemma-category-fibred-setoids-presheaves-products}
Soit $\mathcal{C}$ une catégorie.
La construction du
lemme \ref{lemma-2-category-fibred-setoids},
qui associe un préfaisceau à une catégorie fibrée en setoïdes,
est compatible avec les produits, au sens où le préfaisceau associé à un
$2$-produit fibré $\mathcal{X} \times_\mathcal{Y} \mathcal{Z}$
est le produit fibré des préfaisceaux associés à
$\mathcal{X}, \mathcal{Y}, \mathcal{Z}$.
\end{lemma}

\begin{proof}
Soit $U \in \Ob(\mathcal{C})$. Le lemme affirme simplement que
$$
\Ob((\mathcal{X} \times_\mathcal{Y} \mathcal{Z})_U)/\!\cong
\quad \text{est égal à} \quad
\Ob(\mathcal{X}_U)/\!\cong
\ \times_{\Ob(\mathcal{Y}_U)/\!\cong}
\ \Ob(\mathcal{Z}_U)/\!\cong
$$
et nous en omettons la démonstration. (Notons cependant que ce ne serait pas
vrai en général si la catégorie $\mathcal{Y}_U$ n'était pas un setoïde.)
\end{proof}










\section{Catégories fibrées en groupoïdes représentables}
\label{section-representable-fibred-groupoids}

\noindent
Voici notre définition d'une catégorie fibrée en groupoïdes représentable.
Comme promis, elle est invariante par équivalences.

\begin{definition}
\label{definition-representable-fibred-category}
Soit $\mathcal{C}$ une catégorie.
Une catégorie fibrée en groupoïdes $p : \mathcal{S} \to \mathcal{C}$ est dite
{\it représentable} s'il existe un objet
$X$ de $\mathcal{C}$ et une équivalence $j : \mathcal{S} \to \mathcal{C}/X$
(dans la $2$-catégorie des catégories fibrées en groupoïdes au-dessus de
$\mathcal{C}$).
\end{definition}

\noindent
L'abus de notation habituel consiste à dire que
{\it $X$ représente $\mathcal{S}$} sans mentionner l'équivalence $j$.
Explicitons ce que cela implique.

\begin{lemma}
\label{lemma-characterize-representable-fibred-category}
Soit $\mathcal{C}$ une catégorie.
Soit $p : \mathcal{S} \to \mathcal{C}$ une catégorie fibrée en groupoïdes.
\begin{enumerate}
\item $\mathcal{S}$ est représentable si et seulement si
les conditions suivantes sont satisfaites :
\begin{enumerate}
\item $\mathcal{S}$ est fibrée en setoïdes, et
\item le préfaisceau $U \mapsto \Ob(\mathcal{S}_U)/\cong$ est
représentable.
\end{enumerate}
\item Si $\mathcal{S}$ est représentable, la paire $(X, j)$, où $j$ est
l'équivalence $j : \mathcal{S} \to \mathcal{C}/X$, est déterminée de manière
unique à isomorphisme près.
\end{enumerate}
\end{lemma}

\begin{proof}
La première assertion découle immédiatement du
lemme \ref{lemma-setoid-fibres}.
Pour la seconde, supposons que $(X', j')$ soit
une seconde paire de ce type. Choisissons un $1$-morphisme
$t' : \mathcal{C}/X' \to \mathcal{S}$ tel que
$j' \circ t' \cong \text{id}_{\mathcal{C}/X'}$ et
$t' \circ j' \cong \text{id}_\mathcal{S}$. Alors
$j \circ t' : \mathcal{C}/X' \to \mathcal{C}/X$ est une équivalence.
C'est donc un isomorphisme, voir le lemme \ref{lemma-2-category-fibred-sets}.
Par le lemme de Yoneda \ref{lemma-yoneda} (via, par exemple,
l'exemple \ref{example-fibred-category-from-functor-of-points}),
il est donc donné par un isomorphisme $X' \to X$.
\end{proof}

\begin{lemma}
\label{lemma-morphisms-representable-fibred-categories}
Soit $\mathcal{C}$ une catégorie.
Soient $\mathcal{X}$, $\mathcal{Y}$ des catégories fibrées en groupoïdes
au-dessus de $\mathcal{C}$. Supposons que $\mathcal{X}$, $\mathcal{Y}$
soient représentées par des objets $X$, $Y$ de $\mathcal{C}$.
Alors
$$
\Mor_{\textit{Cat}/\mathcal{C}}(\mathcal{X}, \mathcal{Y})
\Big/
2\text{-isomorphisme}
=
\Mor_\mathcal{C}(X, Y)
$$
Plus précisément, pour tout $\phi : X \to Y$, il existe un
$1$-morphisme $f : \mathcal{X} \to \mathcal{Y}$ qui induit
$\phi$ sur les classes d'isomorphisme d'objets et
qui est unique à un unique $2$-isomorphisme près.
\end{lemma}

\begin{proof}
D'après
l'exemple \ref{example-fibred-category-from-functor-of-points},
nous avons $\mathcal{C}/X = \mathcal{S}_{h_X}$ et
$\mathcal{C}/Y = \mathcal{S}_{h_Y}$. D'après le
lemme \ref{lemma-2-category-fibred-setoids},
nous avons
$$
\Mor_{\textit{Cat}/\mathcal{C}}(\mathcal{X}, \mathcal{Y})
\Big/
2\text{-isomorphisme}
=
\Mor_{\textit{PSh}(\mathcal{C})}(h_X, h_Y)
$$
Par le lemme de Yoneda
\ref{lemma-yoneda},
nous avons $\Mor_{\textit{PSh}(\mathcal{C})}(h_X, h_Y)
= \Mor_\mathcal{C}(X, Y)$.
\end{proof}








\section{Le lemme de 2-Yoneda}
\label{section-2-yoneda}

\noindent
Soit $\mathcal{C}$ une catégorie. La $2$-catégorie des catégories fibrées
au-dessus de $\mathcal{C}$ a été définie dans la
définition \ref{definition-fibred-categories-over-C}.
Si $\mathcal{S}$, $\mathcal{S}'$ sont des catégories fibrées
au-dessus de $\mathcal{C}$, alors
$$
\Mor_{\textit{Fib}/\mathcal{C}}(\mathcal{S}, \mathcal{S}')
$$
désigne la catégorie des $1$-morphismes dans cette $2$-catégorie.
Voici l'analogue $2$-catégorique du lemme de Yoneda
dans le cadre des catégories fibrées.

\begin{lemma}[Lemme de 2-Yoneda pour les catégories fibrées]
\label{lemma-yoneda-2category-fibred}
Soit $\mathcal{C}$ une catégorie.
Soit $\mathcal{S} \to \mathcal{C}$ une catégorie fibrée au-dessus de
$\mathcal{C}$.
Soit $U \in \Ob(\mathcal{C})$.
Le foncteur
$$
\Mor_{\textit{Fib}/\mathcal{C}}(\mathcal{C}/U, \mathcal{S})
\longrightarrow
\mathcal{S}_U
$$
donné par $G \mapsto G(\text{id}_U)$ est une équivalence.
\end{lemma}

\begin{proof}
Choisissons un clivage de $\mathcal{S}$
(voir la définition \ref{definition-pullback-functor-fibred-category}).
Nous définissons un foncteur
$$
\mathcal{S}_U
\longrightarrow
\Mor_{\textit{Fib}/\mathcal{C}}(\mathcal{C}/U, \mathcal{S})
$$
comme suit. Étant donné
$x \in \Ob(\mathcal{S}_U)$,
le foncteur associé est défini
\begin{enumerate}
\item sur les objets par $(f : V \to U) \mapsto f^*x$, et
\item sur les morphismes en envoyant la flèche
$(g : V'/U \to V/U)$ sur la composée
$$
(f \circ g)^*x \xrightarrow{(\alpha_{g, f})_x} g^*f^*x \rightarrow f^*x
$$
où $\alpha_{g, f}$ est comme dans le lemme \ref{lemma-fibred}.
\end{enumerate}
Nous omettons la vérification que ce foncteur est un inverse de celui du lemme.
\end{proof}

\noindent
Soit $\mathcal{C}$ une catégorie. La $2$-catégorie des catégories
fibrées en groupoïdes au-dessus de $\mathcal{C}$ est une sous-$2$-catégorie
« pleine » de la $2$-catégorie des catégories au-dessus de
$\mathcal{C}$ (voir la
définition \ref{definition-categories-fibred-in-groupoids-over-C}).
Ainsi, si $\mathcal{S}$, $\mathcal{S}'$ sont fibrées en groupoïdes
au-dessus de $\mathcal{C}$, alors
$$
\Mor_{\textit{Cat}/\mathcal{C}}(\mathcal{S}, \mathcal{S}')
$$
désigne la catégorie des $1$-morphismes dans cette $2$-catégorie
(voir la définition \ref{definition-categories-over-C}).
Ce sont tous des groupoïdes, voir les remarques qui suivent la
définition \ref{definition-categories-fibred-in-groupoids-over-C}.
Voici l'analogue $2$-catégorique du lemme de Yoneda.

\begin{lemma}[Lemme de 2-Yoneda]
\label{lemma-yoneda-2category}
Soit $\mathcal{S}\to \mathcal{C}$ fibrée en groupoïdes.
Soit $U \in \Ob(\mathcal{C})$.
Le foncteur
$$
\Mor_{\textit{Cat}/\mathcal{C}}(\mathcal{C}/U, \mathcal{S})
\longrightarrow
\mathcal{S}_U
$$
donné par $G \mapsto G(\text{id}_U)$ est une équivalence.
\end{lemma}

\begin{proof}
Choisissons un clivage de $\mathcal{S}$
(voir la définition \ref{definition-pullback-functor-fibred-category}).
Nous définissons un foncteur
$$
\mathcal{S}_U
\longrightarrow
\Mor_{\textit{Cat}/\mathcal{C}}(\mathcal{C}/U, \mathcal{S})
$$
comme suit. Étant donné
$x \in \Ob(\mathcal{S}_U)$,
le foncteur associé est défini
\begin{enumerate}
\item sur les objets par $(f : V \to U) \mapsto f^*x$, et
\item sur les morphismes en envoyant la flèche
$(g : V'/U \to V/U)$ sur la composée
$$
(f \circ g)^*x \xrightarrow{(\alpha_{g, f})_x} g^*f^*x \rightarrow f^*x
$$
où $\alpha_{g, f}$ est comme dans le lemme \ref{lemma-fibred-groupoids}.
\end{enumerate}
Nous omettons la vérification que ce foncteur est un inverse de celui du lemme.
\end{proof}

\begin{remark}
\label{remark-alternative-fibred-groupoids-strict}
On peut utiliser le lemme de $2$-Yoneda pour donner une autre démonstration du
lemme \ref{lemma-fibred-groupoids-strict}.
Soit $p : \mathcal{S} \to \mathcal{C}$ une catégorie fibrée en groupoïdes.
Nous définissons un foncteur contravariant $F$ de $\mathcal{C}$ vers la
catégorie des groupoïdes comme suit : pour $U\in \Ob(\mathcal{C})$,
posons
$$
F(U) = \Mor_{\textit{Cat}/\mathcal{C}}(\mathcal{C}/U, \mathcal{S}).
$$
Si $f : U \to V$, le foncteur induit $\mathcal{C}/U \to \mathcal{C}/V$
induit le morphisme $F(f) : F(V) \to F(U)$. Il est clair que $F$ est un foncteur.
Soit $\mathcal{S}'$ la catégorie fibrée en groupoïdes associée par
l'exemple \ref{example-functor-groupoids}.
Il existe un foncteur évident $G : \mathcal{S}' \to \mathcal{S}$
au-dessus de $\mathcal{C}$, qui envoie la paire $(U, x)$, où
$U \in \Ob(\mathcal{C})$ et $x \in F(U)$, sur
$x(\text{id}_U) \in \mathcal{S}$. Le
lemme \ref{lemma-yoneda-2category}
implique maintenant que, pour tout $U$,
$$
G_U : \mathcal{S}'_U = F(U)=
\Mor_{\textit{Cat}/\mathcal{C}}(\mathcal{C}/U, \mathcal{S})
\to
\mathcal{S}_U
$$
est une équivalence ; ainsi, $G$ est une équivalence entre $\mathcal{S}$ et
$\mathcal{S}'$ d'après le lemme \ref{lemma-equivalence-fibred-categories}.
\end{remark}





\section{1-flèches représentables}
\label{section-representable-1-morphisms}

\noindent
Soit $\mathcal{C}$ une catégorie.
Dans cette section, nous expliquons ce que signifie, pour une $1$-flèche
entre catégories fibrées en groupoïdes au-dessus de $\mathcal{C}$,
d'être représentable.

\medskip\noindent
Soit $\mathcal{C}$ une catégorie.
Soient $\mathcal{X}$, $\mathcal{Y}$ des catégories fibrées en groupoïdes
au-dessus de $\mathcal{C}$.
Soit $U \in \Ob(\mathcal{C})$.
Soient $F : \mathcal{X} \to \mathcal{Y}$ et
$G : \mathcal{C}/U \to \mathcal{Y}$ des $1$-morphismes de catégories
fibrées en groupoïdes au-dessus de $\mathcal{C}$.
Nous voulons décrire
le $2$-produit fibré
$$
\xymatrix{
(\mathcal{C}/U) \times_\mathcal{Y} \mathcal{X} \ar[r] \ar[d] &
\mathcal{X} \ar[d]^F \\
\mathcal{C}/U \ar[r]^G &
\mathcal{Y}
}
$$
Soit $y = G(\text{id}_U) \in \mathcal{Y}_U$.
Choisissons un clivage de $\mathcal{Y}$
(voir la définition \ref{definition-pullback-functor-fibred-category}).
Alors $G$ est isomorphe au foncteur $(f : V \to U) \mapsto f^*y$,
voir le lemme \ref{lemma-yoneda-2category} et sa démonstration.
On peut considérer un objet de
$(\mathcal{C}/U) \times_\mathcal{Y} \mathcal{X}$
comme un quadruplet $(V, f : V \to U, x, \phi)$, voir le
lemme \ref{lemma-2-product-categories-over-C}.
Avec la description de $G$ ci-dessus, on peut considérer $\phi$ comme
un isomorphisme $\phi : f^*y \to F(x)$ dans $\mathcal{Y}_V$.

\begin{lemma}
\label{lemma-identify-fibre-product}
Dans la situation ci-dessus, la catégorie fibre de
$(\mathcal{C}/U) \times_\mathcal{Y} \mathcal{X}$ au-dessus
d'un objet $f : V \to U$ de $\mathcal{C}/U$
est la catégorie décrite comme suit :
\begin{enumerate}
\item ses objets sont les paires $(x, \phi)$,
où $x \in \Ob(\mathcal{X}_V)$ et
$\phi : f^*y \to F(x)$ est un morphisme de $\mathcal{Y}_V$ ;
\item l'ensemble des morphismes entre $(x, \phi)$ et $(x', \phi')$
est l'ensemble des morphismes $\psi : x \to x'$ dans $\mathcal{X}_V$
tels que $F(\psi) = \phi' \circ \phi^{-1}$.
\end{enumerate}
\end{lemma}

\begin{proof}
Voir la discussion ci-dessus.
\end{proof}

\begin{lemma}
\label{lemma-prepare-representable-map-stack-in-groupoids}
Soit $\mathcal{C}$ une catégorie.
Soient $\mathcal{X}$, $\mathcal{Y}$ des catégories fibrées en groupoïdes
au-dessus de $\mathcal{C}$.
Soit $F : \mathcal{X} \to \mathcal{Y}$ un $1$-morphisme.
Soit $G : \mathcal{C}/U \to \mathcal{Y}$ un $1$-morphisme.
Alors
$$
(\mathcal{C}/U) \times_\mathcal{Y} \mathcal{X}
\longrightarrow
\mathcal{C}/U
$$
est une catégorie fibrée en groupoïdes.
\end{lemma}

\begin{proof}
Nous avons déjà vu dans le lemme \ref{lemma-2-product-fibred-categories}
que la composée
$$
(\mathcal{C}/U) \times_\mathcal{Y} \mathcal{X}
\longrightarrow
\mathcal{C}/U
\longrightarrow
\mathcal{C}
$$
est une catégorie fibrée en groupoïdes. Le lemme découle alors du
lemme \ref{lemma-cute-groupoids}.
\end{proof}

\begin{definition}
\label{definition-representable-map-categories-fibred-in-groupoids}
Soit $\mathcal{C}$ une catégorie.
Soient $\mathcal{X}$, $\mathcal{Y}$ des catégories fibrées en groupoïdes
au-dessus de $\mathcal{C}$.
Soit $F : \mathcal{X} \to \mathcal{Y}$ un $1$-morphisme.
Nous disons que $F$ est {\it représentable}, ou que
{\it $\mathcal{X}$ est relativement représentable au-dessus de $\mathcal{Y}$},
si, pour tout $U \in \Ob(\mathcal{C})$
et tout $G : \mathcal{C}/U \to \mathcal{Y}$,
la catégorie fibrée en groupoïdes
$$
(\mathcal{C}/U) \times_\mathcal{Y} \mathcal{X}
\longrightarrow
\mathcal{C}/U
$$
est représentable.
\end{definition}

\begin{lemma}
\label{lemma-spell-out-representable-map-stack-in-groupoids}
Soit $\mathcal{C}$ une catégorie.
Soient $\mathcal{X}$, $\mathcal{Y}$ des catégories fibrées en groupoïdes
au-dessus de $\mathcal{C}$.
Soit $F : \mathcal{X} \to \mathcal{Y}$ un $1$-morphisme.
Si $F$ est représentable, alors chacun des foncteurs
$$
F_U : \mathcal{X}_U \longrightarrow \mathcal{Y}_U
$$
entre catégories fibres est fidèle.
\end{lemma}

\begin{proof}
Cela résulte de la description des catégories fibres du
lemme \ref{lemma-identify-fibre-product} et de la caractérisation
des catégories fibrées représentables du
lemme \ref{lemma-characterize-representable-fibred-category}.
\end{proof}

\begin{lemma}
\label{lemma-criterion-representable-map-stack-in-groupoids}
Soit $\mathcal{C}$ une catégorie.
Soient $\mathcal{X}$, $\mathcal{Y}$ des catégories fibrées en groupoïdes
au-dessus de $\mathcal{C}$.
Soit $F : \mathcal{X} \to \mathcal{Y}$ un $1$-morphisme.
Choisissons un clivage de $\mathcal{Y}$.
Supposons que
\begin{enumerate}
\item tout foncteur $F_U : \mathcal{X}_U \longrightarrow \mathcal{Y}_U$
entre catégories fibres soit fidèle, et
\item pour tout $U$ et tout $y \in \mathcal{Y}_U$, le préfaisceau
$$
(f : V \to U)
\longmapsto
\{(x, \phi) \mid x \in \mathcal{X}_V, \phi : f^*y \to F(x)\}/\cong
$$
soit un préfaisceau représentable sur $\mathcal{C}/U$.
\end{enumerate}
Alors $F$ est représentable.
\end{lemma}

\begin{proof}
Cela résulte de la description des catégories fibres du
lemme \ref{lemma-identify-fibre-product} et de la caractérisation
des catégories fibrées représentables du
lemme \ref{lemma-characterize-representable-fibred-category}.
\end{proof}

\noindent
Avant d'énoncer le lemme suivant, signalons que la $2$-catégorie
des catégories fibrées en groupoïdes est une $(2, 1)$-catégorie, et que nous
savons donc ce que signifie pour elle d'avoir un objet final (voir la
définition \ref{definition-final-object-2-category}). Elle possède effectivement
un objet final, à savoir $\text{id} : \mathcal{C} \to \mathcal{C}$.
Nous définissons donc les {\it $2$-produits} de catégories fibrées en groupoïdes
au-dessus de $\mathcal{C}$ comme les $2$-produits fibrés
$$
\mathcal{X} \times \mathcal{Y} :=
\mathcal{X} \times_\mathcal{C} \mathcal{Y}.
$$
Avec cette définition, le lemme suivant a un sens.

\begin{lemma}
\label{lemma-representable-diagonal-groupoids}
Soit $\mathcal{C}$ une catégorie.
Soit $\mathcal{S} \to \mathcal{C}$ une catégorie fibrée en groupoïdes.
Supposons que $\mathcal{C}$ admette les produits de paires d'objets et les
produits fibrés.
Les conditions suivantes sont équivalentes :
\begin{enumerate}
\item La diagonale $\mathcal{S} \to \mathcal{S} \times \mathcal{S}$
est représentable.
\item Pour tout $U$ dans $\mathcal{C}$, tout
$G : \mathcal{C}/U \to \mathcal{S}$ est représentable.
\end{enumerate}
\end{lemma}

\begin{proof}
Supposons la diagonale représentable, et fixons $U, G$.
Considérons un objet quelconque $V \in \Ob(\mathcal{C})$ et un
$G' : \mathcal{C}/V \to \mathcal{S}$ quelconque.
Notons que $\mathcal{C}/U \times \mathcal{C}/V = \mathcal{C}/(U \times V)$
est représentable. Par hypothèse, le produit fibré
$$
\xymatrix{
(\mathcal{C}/(U \times V))
\times_{(\mathcal{S} \times \mathcal{S})}
\mathcal{S}
\ar[r] \ar[d] &
\mathcal{S} \ar[d] \\
\mathcal{C}/(U \times V) \ar[r]^{(G, G')} &
\mathcal{S} \times \mathcal{S}
}
$$
est représentable.
Cela signifie qu'il existe $W \to U \times V$ dans $\mathcal{C}$,
tel que
$$
\xymatrix{
\mathcal{C}/W \ar[d] \ar[r] & \mathcal{S} \ar[d] \\
\mathcal{C}/U \times \mathcal{C}/V \ar[r] & \mathcal{S} \times \mathcal{S}
}
$$
soit cartésien. Cela implique que
$\mathcal{C}/W \cong \mathcal{C}/U \times_\mathcal{S} \mathcal{C}/V$
(voir le lemme \ref{lemma-diagonal-1} et la
remarque \ref{remark-2-product-fibred-categories-same}),
comme voulu.

\medskip\noindent
Supposons (2). Considérons un objet quelconque $V \in \Ob(\mathcal{C})$
et un $(G, G') : \mathcal{C}/V \to \mathcal{S} \times \mathcal{S}$
quelconque. Nous devons montrer que
$\mathcal{C}/V \times_{\mathcal{S} \times \mathcal{S}} \mathcal{S}$
est représentable. Nous savons que
$\mathcal{C}/V \times_{G, \mathcal{S}, G'} \mathcal{C}/V$
est représentable, disons par $a : W \to V$ dans $\mathcal{C}/V$.
L'équivalence
$$
\mathcal{C}/W \to \mathcal{C}/V \times_{G, \mathcal{S}, G'} \mathcal{C}/V
$$
suivie de la seconde projection vers $\mathcal{C}/V$ donne un
second morphisme $a' : W \to V$. Considérons
$W' = W \times_{(a, a'), V \times V} V$.
Il existe une équivalence
$$
\mathcal{C}/W' \cong
\mathcal{C}/V \times_{\mathcal{S} \times \mathcal{S}} \mathcal{S}
$$
à savoir
\begin{eqnarray*}
\mathcal{C}/W' & \cong &
\mathcal{C}/W \times_{(\mathcal{C}/V \times \mathcal{C}/V)} \mathcal{C}/V \\
& \cong &
\left(\mathcal{C}/V \times_{(G, \mathcal{S}, G')} \mathcal{C}/V\right)
\times_{(\mathcal{C}/V \times \mathcal{C}/V)} \mathcal{C}/V \\
& \cong &
\mathcal{C}/V \times_{(\mathcal{S} \times \mathcal{S})} \mathcal{S}
\end{eqnarray*}
(pour le dernier isomorphisme, voir le lemme \ref{lemma-diagonal-2} et la
remarque \ref{remark-2-product-fibred-categories-same}),
ce qui démontre le lemme.
\end{proof}

\noindent
{\bf Notes bibliographiques :}
Certaines parties de ce qui précède sont tirées des notes de Vistoli \cite{Vis2}.





\section{Catégories monoïdales}
\label{section-monoidal}

\noindent
Soit $\mathcal{C}$ une catégorie. Supposons donné un foncteur
$$
\otimes : \mathcal{C} \times \mathcal{C} \longrightarrow \mathcal{C}
$$
On cherche souvent à savoir si $\otimes$ satisfait une loi d'associativité
et s'il existe une unité pour $\otimes$.

\medskip\noindent
Une {\it contrainte d'associativité} pour $(\mathcal{C}, \otimes)$ est
un isomorphisme fonctoriel
$$
\phi_{X, Y, Z} : X \otimes (Y \otimes Z) \to (X \otimes Y) \otimes Z
$$
tel que, pour tous objets $X, Y, Z, W$, le diagramme
\begin{equation}
\label{equation-pentagon}
\vcenter{
\xymatrix{
X \otimes (Y \otimes ( Z \otimes W)) \ar[r] \ar[d] &
X \otimes ((Y \otimes Z) \otimes W) \ar[d] \\
(X \otimes Y) \otimes (Z \otimes W) \ar[d] &
(X \otimes (Y \otimes Z)) \otimes W \ar[ld] \\
((X \otimes Y) \otimes Z) \otimes W
}
}
\end{equation}
soit commutatif, toute flèche étant obtenue en appliquant convenablement
$\phi$ et en utilisant la fonctorialité de $\otimes$.

\medskip\noindent
Dans \cite[Théorème 3.1]{associativity}, il est expliqué en quel sens un triplet
$(\mathcal{C}, \otimes, \phi)$ comme ci-dessus est cohérent. Une reformulation
consiste à dire qu'un tel triplet donne lieu à un système de foncteurs bien définis
$$
\mathcal{C} \times \ldots \times \mathcal{C} \longrightarrow \mathcal{C},
\quad
(X_1, \ldots, X_n) \longmapsto X_1 \otimes \ldots \otimes X_n
$$
pour tout $n \geq 1$, ainsi qu'à des isomorphismes fonctoriels, pour
$n, m \geq 1$,
$$
(X_1 \otimes \ldots \otimes X_n) \otimes (Y_1 \otimes \ldots \otimes Y_m)
\to
X_1 \otimes \ldots \otimes X_n \otimes Y_1 \otimes \ldots \otimes Y_m
$$
tels que tous les diagrammes que l'on peut former à partir de ces isomorphismes
commutent (le diagramme du pentagone ci-dessus en est un exemple). Pour $n = 5$,
cela signifie que les expressions
$$
\begin{matrix}
(((X_1 \otimes X_2) \otimes X_3) \otimes X_4) \otimes X_5 &
((X_1 \otimes (X_2 \otimes X_3)) \otimes X_4) \otimes X_5 \\
((X_1 \otimes X_2) \otimes (X_3 \otimes X_4)) \otimes X_5 &
(X_1 \otimes ((X_2 \otimes X_3) \otimes X_4)) \otimes X_5 \\
(X_1 \otimes (X_2 \otimes (X_3 \otimes X_4))) \otimes X_5 &
((X_1 \otimes X_2) \otimes X_3) \otimes (X_4 \otimes X_5) \\
(X_1 \otimes (X_2 \otimes X_3)) \otimes (X_4 \otimes X_5) &
(X_1 \otimes X_2) \otimes ((X_3 \otimes X_4) \otimes X_5) \\
(X_1 \otimes X_2) \otimes (X_3 \otimes (X_4 \otimes X_5)) &
X_1 \otimes (((X_2 \otimes X_3) \otimes X_4) \otimes X_5) \\
X_1 \otimes ((X_2 \otimes (X_3 \otimes X_4)) \otimes X_5) &
X_1 \otimes ((X_2 \otimes X_3) \otimes (X_4 \otimes X_5)) \\
X_1 \otimes (X_2 \otimes ((X_3 \otimes X_4) \otimes X_5)) &
X_1 \otimes (X_2 \otimes (X_3 \otimes (X_4 \otimes X_5)))
\end{matrix}
$$
sont toutes fonctoriellement isomorphes à
$X_1 \otimes \ldots \otimes X_5$ et que ces isomorphismes sont tous
compatibles avec tous les isomorphismes possibles que l'on obtient entre ces
foncteurs en appliquant $\phi$ à trois positions consécutives (lorsque cela est
possible). Nous utiliserons ce fait sans autre mention dans la suite et
n'écrirons plus les parenthèses dans les itérations de $\otimes$.

\medskip\noindent
Une {\it unité} pour un triplet $(\mathcal{C}, \otimes, \phi)$ comme ci-dessus
est un objet $\mathbf{1}$ de $\mathcal{C}$ muni d'isomorphismes fonctoriels
$$
l : \mathbf{1} \otimes X \to X
\quad\text{et}\quad
r : X \otimes \mathbf{1} \to X
$$
tels que, pour tous objets $X, Y$, le diagramme
\begin{equation}
\label{equation-triangle}
\vcenter{
\xymatrix{
X \otimes (\mathbf{1} \otimes Y) \ar[rr]_\phi \ar[rd]_{\text{id} \otimes l} & &
(X \otimes \mathbf{1}) \otimes Y \ar[ld]^{r \otimes \text{id}} \\
& X \otimes Y
}
}
\end{equation}
soit commutatif. Nous considérerons souvent les unités comme des paires
$(\mathbf{1}, 1)$, comme dans le lemme suivant.

\begin{lemma}
\label{lemma-monoidal-unit}
Soit $(\mathcal{C}, \otimes, \phi)$ comme ci-dessus. Il existe une
correspondance bijective entre les unités $(\mathbf{1}, l, r)$ de $\mathcal{C}$
et les paires $(\mathbf{1}, 1)$ où $\mathbf{1}$ est un objet de $\mathcal{C}$ et
$1 : \mathbf{1} \otimes \mathbf{1} \to \mathbf{1}$
est un isomorphisme tel que les foncteurs
$L : X \mapsto \mathbf{1} \otimes X$
et $R : X \mapsto X \otimes \mathbf{1}$ soient des équivalences.
\end{lemma}

\begin{proof}
Étant donnée une unité $(\mathbf{1}, l, r)$, on obtient un isomorphisme
$r : \mathbf{1} \otimes \mathbf{1} \to \mathbf{1}$, et $L$ et $R$ sont
des équivalences puisqu'ils sont isomorphes au foncteur identité.
Réciproquement, supposons donnée $(\mathbf{1}, 1)$ telle que
$L$ et $R$ soient des équivalences. On obtient des isomorphismes fonctoriels
$l_X : \mathbf{1} \otimes X \to X$ et $r_X : X \otimes \mathbf{1} \to X$
caractérisés par $L(l_X) = 1 \otimes \text{id}_X$ et
$R(r_X) = \text{id}_X \otimes 1$. Il nous faut alors montrer que les deux
flèches $\text{id}_X \otimes l_Y$ et $r_X \otimes \text{id}_Y$ de
$X \otimes \mathbf{1} \otimes Y$ vers $X \otimes Y$ sont égales
pour tous $X$ et $Y$. Cette propriété ne dépend que des classes d'isomorphisme
de $X$ et $Y$.
Puisque $R$ et $L$ sont des équivalences, il suffit de
le vérifier pour $X = Z \otimes \mathbf{1}$ et $Y = \mathbf{1} \otimes W$
pour certains objets $Z$ et $W$.
Autrement dit, nous devons montrer que
$$
\text{id}_Z \otimes \text{id}_\mathbf{1} \otimes l_{\mathbf{1} \otimes W}
=
r_{Z \otimes \mathbf{1}} \otimes \text{id}_\mathbf{1} \otimes \text{id}_W
$$
Par construction, ces morphismes sont respectivement égaux à
$\text{id}_Z \otimes 1 \otimes \text{id}_\mathbf{1} \otimes \text{id}_W$ et
$\text{id}_Z \otimes \text{id}_\mathbf{1} \otimes 1 \otimes \text{id}_W$.
Il suffit donc de montrer que
$1 \otimes \text{id}_\mathbf{1} = \text{id}_\mathbf{1} \otimes 1$.

\medskip\noindent
Nous avons
$\text{id}_\mathbf{1} \otimes 1 = r_\mathbf{1} \otimes \text{id}_\mathbf{1}$
et
$l_\mathbf{1} \otimes \text{id}_\mathbf{1} = \text{id}_\mathbf{1} \otimes 1$.
On peut écrire $r_\mathbf{1} = a \circ 1$ et $l_\mathbf{1} = b \circ 1$
pour certains automorphismes $a, b$ de $\mathbf{1}$. Ainsi, on a
$\text{id}_\mathbf{1} \otimes 1 =
(a \otimes \text{id}_\mathbf{1}) \circ (1 \otimes \text{id}_\mathbf{1})$
et
$1 \otimes \text{id}_\mathbf{1} =
(\text{id}_\mathbf{1} \otimes b) \circ (\text{id}_\mathbf{1} \otimes 1)$.
On peut alors écrire
\begin{align*}
(1 \otimes \text{id}_\mathbf{1}) \circ
(\text{id}_\mathbf{1} \otimes 1 \otimes \text{id}_\mathbf{1})
& =
(1 \otimes \text{id}_\mathbf{1}) \circ
(\text{id}_\mathbf{1} \otimes \text{id}_\mathbf{1} \otimes b)
\circ
(\text{id}_\mathbf{1} \otimes \text{id}_\mathbf{1} \otimes 1) \\
& =
(\text{id}_\mathbf{1} \otimes b) \circ
(1 \otimes \text{id}_\mathbf{1}) \circ
(\text{id}_\mathbf{1} \otimes \text{id}_\mathbf{1} \otimes 1) \\
& =
(\text{id}_\mathbf{1} \otimes b) \circ
(1 \otimes 1)
\end{align*}
et l'on a aussi
\begin{align*}
(1 \otimes \text{id}_\mathbf{1}) \circ
(\text{id}_\mathbf{1} \otimes 1 \otimes \text{id}_\mathbf{1})
& =
(\text{id}_\mathbf{1} \otimes b) \circ
(\text{id}_\mathbf{1} \otimes 1) \circ
(a \otimes \text{id}_\mathbf{1} \otimes \text{id}_\mathbf{1}) \circ
(1 \otimes \text{id}_\mathbf{1} \otimes \text{id}_\mathbf{1}) \\
& =
(\text{id}_\mathbf{1} \otimes b) \circ
(a \otimes \text{id}_\mathbf{1}) \circ
(1 \otimes 1)
\end{align*}
Cela montre que $a \otimes \text{id}_\mathbf{1}$ est l'identité, et donc
que $a$ est l'identité, comme voulu.

\medskip\noindent
Pour achever la démonstration, notons que les règles ci-dessus déterminent
des équivalences inverses de catégories entre la catégorie des unités
(convenablement définie) et la catégorie des paires $(\mathbf{1}, 1)$.
\end{proof}

\begin{lemma}
\label{lemma-monoidal-unit-unique}
Soit $(\mathcal{C}, \otimes, \phi)$ comme ci-dessus.
Soit $(\mathbf{1}, 1)$ une unité (voir le lemme \ref{lemma-monoidal-unit}).
Alors
\begin{enumerate}
\item $1 \otimes \text{id}_\mathbf{1} = \text{id}_\mathbf{1} \otimes 1$
\item $\Gamma = \text{Mor}(\mathbf{1}, \mathbf{1})$ est un monoïde commutatif,
\item  $a = 1 \circ (a \otimes \text{id}_\mathbf{1}) \circ 1^{-1} =
1 \circ (\text{id}_\mathbf{1} \otimes a) \circ 1^{-1}$ pour tout
$a \in \Gamma$,
\item toute autre unité est isomorphe à $(\mathbf{1}, 1)$
par un unique isomorphisme.
\end{enumerate}
\end{lemma}

\begin{proof}
La partie (1) a été démontrée dans la preuve du
lemme \ref{lemma-monoidal-unit}.
Pour $a \in \Gamma$, on a
\begin{align*}
(1 \circ (a \otimes \text{id}_\mathbf{1})) \otimes \text{id}_\mathbf{1}
& =
(1 \otimes \text{id}_\mathbf{1}) \circ
(a \otimes \text{id}_\mathbf{1} \otimes \text{id}_\mathbf{1}) \\
& =
(\text{id}_\mathbf{1} \otimes 1) \circ
(a \otimes \text{id}_\mathbf{1} \otimes \text{id}_\mathbf{1}) \\
& =
(a \otimes \text{id}_\mathbf{1}) \circ
(\text{id}_\mathbf{1} \otimes 1) \\
& =
(a \otimes \text{id}_\mathbf{1}) \circ
(1 \otimes \text{id}_\mathbf{1}) \\
& =
(a \circ 1) \otimes \text{id}_\mathbf{1}
\end{align*}
Ainsi, $1 \circ (a \otimes \text{id}_\mathbf{1}) = a \circ 1$,
ce qui donne une moitié de (3). L'autre moitié s'obtient
de la même manière. Pour voir (2), observons que, pour $a, b \in \Gamma$,
on a
$$
a \circ b =
1 \circ (a \otimes \text{id}_\mathbf{1}) \circ 1^{-1}
\circ
1 \circ (\text{id}_\mathbf{1} \otimes b) \circ 1^{-1}
=
1 \circ (a \otimes b) \circ 1^{-1}
$$
et la même expression donne $b \circ a$. Ainsi, (2) est satisfaite.
Pour voir (4), supposons que $(\mathbf{1}', 1')$ soit une seconde unité.
En utilisant $r$ et $l'$, on obtient des isomorphismes
$\mathbf{1}' \otimes \mathbf{1} \to \mathbf{1}'$ et
$\mathbf{1}' \otimes \mathbf{1} \to \mathbf{1}$. Il existe donc
un isomorphisme $t : \mathbf{1}' \to \mathbf{1}$.
Alors le diagramme
$$
\xymatrix{
\mathbf{1}' \otimes \mathbf{1}' \ar[r]_{t \otimes t} \ar[d]_{1'} &
\mathbf{1} \otimes \mathbf{1} \ar[d]^1 \\
\mathbf{1}' \ar[r]^t & \mathbf{1}
}
$$
commute à un automorphisme $a$ de $\mathbf{1}$ près. Après avoir remplacé
$t$ par $a \circ t$, le diagramme est commutatif (indication : utiliser (3) pour voir
que $1 \circ (a \otimes a) = a^2 \circ 1$).
\end{proof}

\begin{definition}
\label{definition-monoidal-category}
Un triplet $(\mathcal{C}, \otimes, \phi)$ où $\mathcal{C}$ est une catégorie,
$\otimes : \mathcal{C} \times \mathcal{C} \to \mathcal{C}$ est un foncteur,
et $\phi$ est une contrainte d'associativité, est appelé
{\it catégorie monoïdale} s'il existe une unité $\mathbf{1}$.
\end{definition}

\noindent
Nous écrivons toujours $\mathbf{1}$ pour désigner une unité d'une catégorie
monoïdale et nous notons $1 : \mathbf{1} \otimes \mathbf{1} \to \mathbf{1}$
un isomorphisme choisi ; puisque la paire $(\mathbf{1}, 1)$ est déterminée
à un unique isomorphisme près (lemme \ref{lemma-monoidal-unit-unique}),
on peut en choisir une sans inconvénient.

\begin{lemma}
\label{lemma-monoidal-coherence}
Dans une catégorie monoïdale $\mathcal{C}, \otimes, \phi, \mathbf{1}, 1$,
et avec les notations de la démonstration du lemme \ref{lemma-monoidal-unit},
on a
\begin{enumerate}
\item les flèches $1, r_\mathbf{1}, l_\mathbf{1} :
\mathbf{1} \otimes \mathbf{1} \to \mathbf{1}$ coïncident ;
\item les flèches
$l_X \otimes \text{id}_Y, l_{X \otimes Y} :
\mathbf{1} \otimes X \otimes Y \to X \otimes Y$ coïncident ; et
\item les flèches
$\text{id}_X \otimes r_Y , r_{X \otimes Y} :
X \otimes Y \otimes \mathbf{1} \to X \otimes Y$ coïncident.
\end{enumerate}
Une catégorie monoïdale satisfait les hypothèses de
\cite[Théorème 5.2]{associativity}.
\end{lemma}

\begin{proof}
Nous avons établi (1) dans la démonstration du
lemme \ref{lemma-monoidal-unit}.
Nous avons vu dans la démonstration du lemme \ref{lemma-monoidal-unit}
que $l_X$ et $l_{X \otimes Y}$ sont
les uniques morphismes tels que
$\text{id}_\mathbf{1} \otimes l_{X \otimes Y} =
1 \otimes \text{id}_{X \otimes Y}$ et
$\text{id}_\mathbf{1} \otimes l_X = 1 \otimes \text{id}_X$.
La partie (2) en découle immédiatement. La partie (3) se démontre de manière
analogue. Avec la commutativité de (\ref{equation-pentagon})
et de (\ref{equation-triangle}), cela donne la dernière assertion du lemme.
\end{proof}

\noindent
Dans \cite[Théorème 5.2]{associativity}, il est expliqué en quel sens un
quintuplet $(\mathcal{C}, \otimes, \phi, \mathbf{1}, 1)$ comme dans le lemme
ci-dessus est cohérent. Une reformulation est que, dans une catégorie monoïdale
(au sens où nous l'entendons), pour tous $n \geq i \geq 0$, on dispose
d'isomorphismes fonctoriels
$$
X_1 \otimes \ldots \otimes X_i \otimes \mathbf{1}
\otimes X_{i + 1} \otimes \ldots \otimes X_n
\to
X_1 \otimes \ldots \otimes X_n
$$
tels que tous les diagrammes possibles formés à partir de ces isomorphismes
commutent. Ainsi, par exemple, si l'on part de deux insertions de $\mathbf{1}$
et que l'on utilise les isomorphismes dans des ordres différents, on obtient le
même morphisme. Outre la convention consistant à supprimer les parenthèses dans
les itérations de $\otimes$, nous identifions désormais $X \otimes \mathbf{1}$
et $\mathbf{1} \otimes X$ à $X$ sans autre mention.
En outre, nous dirons « soit $\mathcal{C}$ une catégorie monoïdale »,
les données $\otimes, \phi, \mathbf{1}$ étant sous-entendues.

\begin{definition}
\label{definition-functor-monoidal-categories}
Soient $\mathcal{C}$ et $\mathcal{C}'$ des catégories monoïdales.
Un {\it foncteur monoïdal} $F : \mathcal{C} \to \mathcal{C}'$
est la donnée du foncteur $F$ indiqué et d'un isomorphisme
$$
F(X) \otimes F(Y) \to F(X \otimes Y)
$$
fonctoriel en $X$ et $Y$, tel que, pour tous objets $X$, $Y$ et $Z$,
le diagramme
$$
\xymatrix{
F(X) \otimes (F(Y) \otimes F(Z)) \ar[r] \ar[d] &
F(X) \otimes F(Y \otimes Z) \ar[r] &
F(X \otimes (Y \otimes Z)) \ar[d] \\
(F(X) \otimes F(Y)) \otimes F(Z) \ar[r] &
F(X \otimes Y) \otimes F(Z) \ar[r] &
F((X \otimes Y) \otimes Z)
}
$$
est commutatif et tel que $F(\mathbf{1})$ soit une unité de $\mathcal{C}'$.
\end{definition}

\noindent
Selon nos conventions sur les unités, on peut toujours supposer que
$F(\mathbf{1}) = \mathbf{1}$ si $F$ est un foncteur monoïdal.
À titre d'exemple, si $A \to B$ est un homomorphisme d'anneaux, alors
le foncteur $M \mapsto M \otimes_A B$ est un foncteur monoïdal de
$\text{Mod}_A$ vers $\text{Mod}_B$.

\begin{lemma}
\label{lemma-invertible}
Soit $\mathcal{C}$ une catégorie monoïdale. Soit $X$ un objet de
$\mathcal{C}$. Les conditions suivantes sont équivalentes :
\begin{enumerate}
\item le foncteur $L : Y \mapsto X \otimes Y$ est une équivalence ;
\item le foncteur $R : Y \mapsto Y \otimes X$ est une équivalence ;
\item il existe un objet $X'$ tel que
$X \otimes X' \cong X' \otimes X \cong \mathbf{1}$.
\end{enumerate}
\end{lemma}

\begin{proof}
Supposons (1). Choisissons $X'$ tel que $L(X') = \mathbf{1}$, c'est-à-dire
$X \otimes X' \cong \mathbf{1}$. Notons $L'$ et $R'$ les foncteurs
correspondant à $X'$. L'égalité $X \otimes X' \cong \mathbf{1}$
implique $L \circ L' \cong \text{id}$. Ainsi, $L'$ est nécessairement le
quasi-inverse de $L$ (qui existe par hypothèse). Par conséquent,
$L' \circ L \cong \text{id}$.
Il s'ensuit que $X' \otimes X \cong \mathbf{1}$. Ainsi, (3) est satisfaite.

\medskip\noindent
La démonstration de (2) $\Rightarrow$ (3) est duale de celle qui précède.

\medskip\noindent
Supposons (3). Il est alors clair que $L'$ et $L$ sont quasi-inverses
l'un de l'autre, et que $R'$ et $R$ le sont également. Ainsi, (1) et (2)
sont satisfaites.
\end{proof}

\begin{definition}
\label{definition-invertible}
Soit $\mathcal{C}$ une catégorie monoïdale. Un objet $X$ de $\mathcal{C}$
est dit {\it inversible} si l'une quelconque (donc toutes) des conditions
équivalentes du lemme \ref{lemma-invertible} est satisfaite.
\end{definition}

\noindent
Observons que, si $F : \mathcal{C} \to \mathcal{C}'$ est un foncteur
monoïdal, alors $F$ envoie les objets inversibles sur des objets
inversibles.

\begin{definition}
\label{definition-dual}
Étant donnés une catégorie monoïdale $(\mathcal{C}, \otimes, \phi)$
et un objet $X$, un {\it dual à gauche} de celui-ci est un objet $Y$ muni de
morphismes $\eta : \mathbf{1} \to X \otimes Y$ et
$\epsilon : Y \otimes X \to \mathbf{1}$
tels que les diagrammes
$$
\vcenter{
\xymatrix{
X \ar[rd]_1 \ar[r]_-{\eta \otimes 1} &
X \otimes Y \otimes X \ar[d]^{1 \otimes \epsilon} \\
& X
}
}
\quad\text{et}\quad
\vcenter{
\xymatrix{
Y \ar[rd]_1 \ar[r]_-{1 \otimes \eta} &
Y \otimes X \otimes Y \ar[d]^{\epsilon \otimes 1} \\
& Y
}
}
$$
commutent. Dans cette situation, on dit que $X$ est un
{\it dual à droite} de $Y$.
\end{definition}

\noindent
Observons que, si $F : \mathcal{C} \to \mathcal{C}'$ est un foncteur
monoïdal, alors $F(Y)$ est un dual à gauche de $F(X)$ si
$Y$ est un dual à gauche de $X$.

\begin{lemma}
\label{lemma-left-dual}
Soit $\mathcal{C}$ une catégorie monoïdale. Si $Y$ est un dual à gauche de $X$,
alors
$$
\Mor(Z' \otimes X, Z) = \Mor(Z', Z \otimes Y)
\quad\text{et}\quad
\Mor(Y \otimes Z', Z) = \Mor(Z', X \otimes Z)
$$
fonctoriellement en $Z$ et $Z'$.
\end{lemma}

\begin{proof}
Considérons les applications
$$
\Mor(Z' \otimes X, Z) \to
\Mor(Z' \otimes X \otimes Y, Z \otimes Y) \to
\Mor(Z', Z \otimes Y)
$$
où l'on utilise $\eta$ dans la seconde flèche,
ainsi que la suite d'applications
$$
\Mor(Z', Z \otimes Y) \to
\Mor(Z' \otimes X, Z \otimes Y \otimes X) \to
\Mor(Z' \otimes X, Z)
$$
où l'on utilise $\epsilon$ dans la seconde flèche.
Pour montrer que ces flèches sont inverses l'une de l'autre, considérons une
morphisme $a : Z' \to Z \otimes Y$. Nous devons montrer que
$$
Z' \xrightarrow{\text{id}_{Z'} \otimes \eta}
Z' \otimes X \otimes Y \xrightarrow{a \otimes \text{id}_{X \otimes Y}}
Z \otimes Y \otimes X \otimes Y
\xrightarrow{\text{id}_Z \otimes \epsilon \otimes \text{id}_Y}
Z \otimes Y
$$
est égale à $a$. La composée des deux premières flèches est égale à
$(\text{id}_{Z \otimes Y} \otimes \eta) \circ a :
Z' \to Z \otimes Y \to Z \otimes Y \otimes X \otimes Y$.
Puis la composée de $\text{id}_{Z \otimes Y} \otimes \eta$
et de $\text{id}_Z \otimes \epsilon \otimes \text{id}_Y$ est
l'identité par définition du dual. Il en va de même pour l'autre composée.
Nous omettons la démonstration de la seconde égalité.
\end{proof}

\begin{remark}
\label{remark-left-dual-adjoint}
Le lemme \ref{lemma-left-dual} affirme en particulier que
$Z \mapsto Z \otimes Y$ est l'adjoint à droite de
$Z' \mapsto Z' \otimes X$. En particulier, l'unicité des foncteurs adjoints
garantit qu'un dual à gauche de $X$, s'il existe, est unique à un unique
isomorphisme près.
Réciproquement, supposons que le foncteur $Z \mapsto Z \otimes Y$ soit un
adjoint à droite du foncteur $Z' \mapsto Z' \otimes X$, c'est-à-dire que l'on
se donne une bijection
$$
\Mor(Z' \otimes X, Z) \longrightarrow \Mor(Z', Z \otimes Y)
$$
fonctorielle à la fois en $Z$ et en $Z'$. L'unité de l'adjonction fournit
des morphismes
$$
\eta_Z : Z \to Z \otimes X \otimes Y
$$
fonctoriels en $Z$, et la counité de l'adjonction fournit des morphismes
$$
\epsilon_{Z'} : Z' \otimes Y \otimes X \to Z'
$$
fonctoriels en $Z'$. En particulier, on obtient
$\eta = \eta_\mathbf{1} : \mathbf{1} \to X \otimes Y$ et
$\epsilon = \epsilon_\mathbf{1} : Y \otimes X \to \mathbf{1}$.
En guise d'exercice sur la relation entre les unités, les counités et l'isomorphisme
d'adjonction, le lecteur peut montrer que l'on a
$$
(\epsilon \otimes \text{id}_Y) \circ \eta_Y = \text{id}_Y
\quad\text{et}\quad
\epsilon_X \circ (\eta \otimes \text{id}_X) = \text{id}_X
$$
Cela ne suffit toutefois pas à montrer que
$(\epsilon \otimes \text{id}_Y) \circ (\text{id}_Y \otimes \eta) =
\text{id}_Y$ et
$(\text{id}_X \otimes \epsilon) \circ (\eta \otimes \text{id}_X) =
\text{id}_X$, car on ne sait pas en général que
$\eta_Y = \text{id}_Y \otimes \eta$, ni que
$\epsilon_X = \epsilon \otimes \text{id}_X$. Pour cela, il suffirait
de savoir que notre isomorphisme d'adjonction possède la propriété suivante :
pour tous $W, Z, Z'$, le diagramme
$$
\xymatrix{
\Mor(Z' \otimes X, Z) \ar[r] \ar[d]_{\text{id}_W \otimes -} &
\Mor(Z', Z \otimes Y) \ar[d]^{\text{id}_W \otimes -} \\
\Mor(W \otimes Z' \otimes X, W \otimes Z) \ar[r] &
\Mor(W \otimes Z', W \otimes Z \otimes Y)
}
$$
est commutatif.
Si tel est le cas, nous dirons que {\it l'adjonction est compatible avec
la structure tensorielle donnée}. Ainsi, exiger que
$Z \mapsto Z \otimes Y$ soit l'adjoint à droite de $Z' \mapsto Z' \otimes X$
compatible avec la structure tensorielle donnée constitue une formulation
équivalente de la propriété d'être un dual à gauche.
\end{remark}

\begin{lemma}
\label{lemma-tensor-dual}
Soit $\mathcal{C}$ une catégorie monoïdale. Si $Y_i$, $i = 1, 2$,
sont des duaux à gauche de $X_i$, $i = 1, 2$, alors $Y_2 \otimes Y_1$ est
un dual à gauche de $X_1 \otimes X_2$.
\end{lemma}

\begin{proof}
Cela découle de l'unicité des adjoints et de la
remarque \ref{remark-left-dual-adjoint}.
\end{proof}

\noindent
Une {\it contrainte de commutativité} pour $(\mathcal{C}, \otimes)$ est un
isomorphisme fonctoriel
$$
\psi : X \otimes Y \longrightarrow Y \otimes X
$$
tel que la composée
$$
X \otimes Y \xrightarrow{\psi} Y \otimes X \xrightarrow{\psi} X \otimes Y
$$
soit l'identité. On dit que $\psi$ est {\it compatible} avec une contrainte
d'associativité donnée $\phi$ si, pour tous objets $X, Y, Z$, le diagramme
\begin{equation}
\label{equation-hexagon}
\vcenter{
\xymatrix{
X \otimes (Y \otimes Z) \ar[r]_\phi \ar[d]^\psi &
(X \otimes Y) \otimes Z \ar[r]_\psi &
Z \otimes (X \otimes Y) \ar[d]^\phi \\
X \otimes (Z \otimes Y) \ar[r]^\phi &
(X \otimes Z) \otimes Y \ar[r]^\psi &
(Z \otimes X) \otimes Y
}
}
\end{equation}
est commutatif.

\begin{definition}
\label{definition-symmetric-monoidal-category}
Un quadruplet $(\mathcal{C}, \otimes, \phi, \psi)$ où
$\mathcal{C}$ est une catégorie,
$\otimes : \mathcal{C} \times \mathcal{C} \to \mathcal{C}$ est un foncteur,
$\phi$ est une contrainte d'associativité et
$\psi$ est une contrainte de commutativité compatible avec $\phi$,
est appelé {\it catégorie monoïdale symétrique} s'il existe
une unité.
\end{definition}

\noindent
Bien entendu, si $(\mathcal{C}, \otimes, \phi, \psi)$ est une
catégorie monoïdale symétrique, alors $(\mathcal{C}, \otimes, \phi)$
est une catégorie monoïdale et nous pouvons employer le langage et les
notations introduits ci-dessus.

\begin{lemma}
\label{lemma-symmetric-monoidal-coherence}
Dans une catégorie monoïdale symétrique
$\mathcal{C}, \otimes, \phi, \psi, \mathbf{1}, 1$,
on a
\begin{enumerate}
\item les flèches $1 \circ \psi, 1 :
\mathbf{1} \otimes \mathbf{1} \to \mathbf{1}$ coïncident ;
\item les flèches $\text{id}_X \otimes l_Y,
(l_X \otimes \text{id}) \circ (\psi \otimes \text{id}_Y):
X \otimes \mathbf{1} \otimes Y \to X \otimes Y$ coïncident,
\end{enumerate}
Une catégorie monoïdale symétrique satisfait les hypothèses de
\cite[Théorème 5.1]{associativity}.
\end{lemma}

\begin{proof}
On peut écrire $\psi = a \otimes \text{id}_\mathbf{1}$ pour un unique
isomorphisme $a : \mathbf{1} \to \mathbf{1}$.
Le lemme \ref{lemma-monoidal-unit-unique} implique que
$a \otimes \text{id}_\mathbf{1} = \text{id}_\mathbf{1} \otimes a$.
La fonctorialité de $\psi$ affirme que le diagramme
$$
\xymatrix{
(\mathbf{1} \otimes \mathbf{1}) \otimes \mathbf{1}
\ar[r]_\psi
\ar[d]_{1 \otimes \text{id}_\mathbf{1}} &
\mathbf{1} \otimes (\mathbf{1} \otimes \mathbf{1})
\ar[d]^{\text{id}_\mathbf{1} \otimes 1} \\
\mathbf{1} \otimes \mathbf{1}  \ar[r]^\psi &
\mathbf{1} \otimes \mathbf{1}
}
$$
est commutatif. Ainsi, la flèche du haut est égale à
$a \otimes \text{id}_\mathbf{1} \otimes \text{id}_\mathbf{1}$.
Par conséquent, (\ref{equation-hexagon}) pour
$X = Y = Z = \mathbf{1}$ affirme que
$a \otimes \text{id}_\mathbf{1} \otimes \text{id}_\mathbf{1}$
est égal à son carré. Ainsi, $a = \text{id}_\mathbf{1}$.
Cela démontre (1).

\medskip\noindent
La partie (2) affirme que $\psi : X \otimes \mathbf{1} \to \mathbf{1} \otimes X$
est l'identité lorsque l'on identifie la source et le but à $X$
dans notre catégorie monoïdale. Cela résulte du fait que, pour
(\ref{equation-hexagon}) et $\mathbf{1}, X, \mathbf{1}$, le diagramme
$$
\xymatrix{
\mathbf{1} \otimes (X \otimes \mathbf{1}) \ar[r]_\phi \ar[d]^\psi &
(\mathbf{1} \otimes X) \otimes \mathbf{1} \ar[r]_\psi &
\mathbf{1} \otimes (\mathbf{1} \otimes X) \ar[d]^\phi \\
\mathbf{1} \otimes (\mathbf{1} \otimes X) \ar[r]^\phi &
(\mathbf{1} \otimes \mathbf{1}) \otimes X \ar[r]^\psi &
(\mathbf{1} \otimes \mathbf{1}) \otimes X
}
$$
est commutatif et que tous les morphismes $\psi$ sauf un
sont égaux à l'identité.

\medskip\noindent
Outre (1) et (2), la commutativité des diagrammes
(\ref{equation-pentagon}), (\ref{equation-triangle}), (\ref{equation-hexagon})
et les résultats du lemme \ref{lemma-monoidal-coherence}
impliquent la dernière assertion du lemme.
\end{proof}

\noindent
Dans \cite[Théorème 5.1]{associativity}, il est expliqué en quel sens un
sextuplet $(\mathcal{C}, \otimes, \phi, \psi, \mathbf{1}, 1)$ comme dans le
lemme ci-dessus est cohérent. Une reformulation est que, dans une catégorie
monoïdale symétrique (au sens où nous l'entendons), pour tout $n \geq 1$ et
toute permutation $\sigma$ de $\{1, \ldots, n\}$, on dispose d'isomorphismes
fonctoriels
$$
X_1 \otimes \ldots \otimes X_n
\to
X_{\sigma(1)} \otimes \ldots \otimes X_{\sigma(n)}
$$
tels que tous les diagrammes possibles formés à partir de ces isomorphismes
commutent et que ces isomorphismes soient compatibles avec la structure de
catégorie monoïdale (notamment avec les isomorphismes obtenus lorsque l'on insère
un $\mathbf{1}$ dans une position).

\begin{lemma}
\label{lemma-dual-symmetric}
Soit $(\mathcal{C}, \otimes, \phi, \psi)$ une catégorie monoïdale symétrique.
Soit $X$ un objet de $\mathcal{C}$, et supposons que $Y$,
$\eta : \mathbf{1} \to X \otimes Y$, et
$\epsilon : Y \otimes X \to \mathbf{1}$
définissent un dual à gauche de $X$ comme dans la définition \ref{definition-dual}.
Alors $\eta' = \psi \circ \eta : \mathbf{1} \to Y \otimes X$
et $\epsilon' = \epsilon \circ \psi : X \otimes Y \to \mathbf{1}$
font de $X$ un dual à gauche de $Y$.
\end{lemma}

\begin{proof}
Omis. Indication : un exercice agréable sur les définitions.
\end{proof}

\begin{definition}
\label{definition-functor-symmetric-monoidal-categories}
Soient $\mathcal{C}$ et $\mathcal{C}'$ des catégories monoïdales symétriques.
Un {\it foncteur monoïdal symétrique}
$F : \mathcal{C} \to \mathcal{C}'$
est la donnée du foncteur $F$ indiqué et d'un isomorphisme
$$
F(X) \otimes F(Y) \to F(X \otimes Y)
$$
fonctoriel en $X$ et $Y$, tel que $F$ soit un foncteur monoïdal et tel que,
pour tous objets $X$ et $Y$, le diagramme
$$
\xymatrix{
F(X) \otimes F(Y) \ar[r] \ar[d] &
F(X \otimes Y) \ar[d] \\
F(Y) \otimes F(X) \ar[r] &
F(Y \otimes X)
}
$$
est commutatif.
\end{definition}

\begin{remark}
\label{remark-internal-hom-monoidal}
Soit $\mathcal{C}$ une catégorie monoïdale. Nous disons que $\mathcal{C}$
possède un {\it Hom interne} si, pour toute paire d'objets $X, Y$ de
$\mathcal{C}$, il existe un objet $hom(X, Y)$ de $\mathcal{C}$ tel que l'on ait
$$
\Mor(X, hom(Y, Z)) = \Mor(X \otimes Y, Z)
$$
fonctoriellement en $X, Y, Z$. Par le lemme de Yoneda, le bifoncteur
$(X, Y) \mapsto hom(X, Y)$ est déterminé à un unique isomorphisme près,
s'il existe. Étant donné un Hom interne, on obtient des morphismes canoniques
\begin{enumerate}
\item $hom(X, Y) \otimes X \to Y$,
\item $hom(Y, Z) \otimes hom(X, Y) \to hom(X, Z)$,
\item $Z \otimes hom(X, Y) \to hom(X, Z \otimes Y)$,
\item $Y \to hom(X, Y \otimes X)$, et
\item $hom(Y, Z) \otimes X \to hom(hom(X, Y), Z)$ lorsque
$\mathcal{C}$ est monoïdale symétrique.
\end{enumerate}
En effet, le morphisme de (1) est l'image de $\text{id}_{hom(X, Y)}$
par $\Mor(hom(X, Y), hom(X, Y)) \to \Mor(hom(X, Y) \otimes X, Y)$.
Pour construire le morphisme de (2), la propriété caractéristique de
$hom(X, Z)$ ramène à construire un morphisme
$$
hom(Y, Z) \otimes hom(X, Y) \otimes X \longrightarrow Z
$$
et un tel morphisme existe puisque, d'après (1), on dispose des
morphismes $hom(X, Y) \otimes X \to Y$ et $hom(Y, Z) \otimes Y \to Z$.
Pour construire le morphisme de (3), la propriété caractéristique de
$hom(X, Z \otimes Y)$ ramène à construire un morphisme
$$
Z \otimes hom(X, Y) \otimes X \to Z \otimes Y
$$
pour laquelle on utilise $\text{id}_Z \otimes a$, où
$a$ est le morphisme de (1). Pour construire le morphisme de (4),
notons que nous disposons déjà du morphisme
$Y \otimes hom(X, X) \to hom(X, Y \otimes X)$ donnée par (3).
Il suffit donc de construire un morphisme $\mathbf{1} \to hom(X, X)$ ;
pour cela, on prend l'élément de $\Mor(\mathbf{1}, hom(X, X))$
qui correspond à l'isomorphisme canonique $\mathbf{1} \otimes X \to X$
dans $\Mor(\mathbf{1} \otimes X, X)$.
Enfin, considérons (5). Par la propriété universelle de
$hom(hom(X, Y), Z)$, il suffit de construire un morphisme
$$
hom(Y, Z) \otimes X \otimes hom(X, Y) \longrightarrow Z
$$
Pour ce faire, on permute les deux derniers facteurs tensoriels au moyen de la
contrainte de commutativité, puis on utilise les morphismes
$hom(X, Y) \otimes X \to Y$ et $hom(Y, Z) \otimes Y \to Z$.
\end{remark}





\section{Catégories de flèches pointillées}
\label{section-dotted-arrows}

\noindent
Nous étudions certaines « catégories de flèches pointillées » dans les
$(2,1)$-catégories. Elles apparaîtront dans la formulation de divers critères de
relèvement pour les champs algébriques ; voir, par exemple, Morphismes de champs,
section \ref{stacks-morphisms-section-valuative} et
Compléments sur les morphismes de champs, section
\ref{stacks-more-morphisms-section-formally-smooth}.

\begin{definition}
\label{definition-dotted-arrows}
Soit $\mathcal{C}$ une $(2,1)$-catégorie. Considérons un diagramme
$2$-commutatif formé par les flèches pleines
\begin{equation}
\label{equation-dotted-arrows}
\vcenter{
\xymatrix{
S \ar[r]_-x \ar[d]_j & X \ar[d]^f \\
T \ar[r]^-y \ar@{..>}[ru] & Y
}
}
\end{equation}
dans $\mathcal{C}$. Fixons un $2$-isomorphisme
$$
\gamma : y \circ j \rightarrow f \circ x
$$
attestant la $2$-commutativité du diagramme.
Étant donnés (\ref{equation-dotted-arrows}) et $\gamma$, une
\emph{flèche pointillée} est un triplet $(a, \alpha, \beta)$ constitué d'un
morphisme $a \colon T \to X$ et de $2$-isomorphismes
$\alpha : a \circ j \to x$, $\beta : y \to f \circ a$
tels que
$\gamma = (\text{id}_f \star \alpha) \circ (\beta \star \text{id}_j)$,
autrement dit tels que
$$
\xymatrix{
& f \circ a \circ j \ar[rd]^{\text{id}_f \star \alpha} \\
y \circ j \ar[ru]^{\beta \star \text{id}_j} \ar[rr]^\gamma & &
f \circ x
}
$$
soit commutatif. Un {\it morphisme de flèches pointillées}
$(a, \alpha, \beta) \to (a', \alpha', \beta')$ est une
$2$-flèche $\theta : a \to a'$ telle que
$\alpha = \alpha' \circ (\theta \star \text{id}_j)$ et
$\beta' = (\text{id}_f \star \theta) \circ \beta$.
\end{definition}

\noindent
Dans la situation de la définition \ref{definition-dotted-arrows}, il existe
une \emph{catégorie de flèches pointillées} associée.
Cette catégorie est un groupoïde. Elle peut en général dépendre de $\gamma$.
Les deux lemmes suivants affirment que les catégories de flèches pointillées
sont compatibles avec le changement de base et avec la composition du
morphisme $f$.

\begin{lemma}
\label{lemma-cat-dotted-arrows-base-change}
Soit $\mathcal{C}$ une $(2,1)$-catégorie. Supposons donné un diagramme
$2$-commutatif
$$
\xymatrix{
S \ar[r]_-{x'} \ar[d]_j &
X' \ar[d]^p \ar[r]_q &
X \ar[d]^f \\
T \ar[r]^-{y'} &
Y' \ar[r]^g &
Y
}
$$
dans $\mathcal{C}$, où le carré de droite est $2$-cartésien
relativement à un $2$-isomorphisme $\phi \colon g \circ p \to f \circ q$.
Choisissons une $2$-flèche
$\gamma' : y' \circ j \to p \circ x'$. Posons
$x = q \circ x'$, $y = g \circ y'$ et soit
$\gamma : y \circ j \to f \circ x$ le $2$-isomorphisme
$\gamma = (\phi \star \text{id}_{x'}) \circ (\text{id}_g \star \gamma')$.
Alors la catégorie $\mathcal{D}'$ des flèches pointillées
pour le carré de gauche et $\gamma'$ est équivalente à la catégorie
$\mathcal{D}$ des flèches pointillées
pour le rectangle extérieur et $\gamma$.
\end{lemma}

\begin{proof}
Il existe un foncteur $\mathcal{D}' \to \mathcal{D}$ qui est donné sur les
objets par $(a, \alpha, \beta) \mapsto (q \circ a, \text{id}_q \star \alpha,
(\phi \star \text{id}_a) \circ (\text{id}_g \star \beta))$
et sur les flèches par $\theta \mapsto \text{id}_q \star \theta$.
Le fait que ce foncteur $\mathcal{D}' \to \mathcal{D}$ soit une équivalence
découle formellement de la propriété universelle des $2$-produits fibrés,
comme dans la section \ref{section-2-fibre-products}. Les détails sont omis.
\end{proof}

\begin{lemma}
\label{lemma-cat-dotted-arrows-composition}
Soit $\mathcal{C}$ une $(2,1)$-catégorie. Supposons donné un diagramme
$2$-commutatif aux flèches pleines
$$
\xymatrix{
S \ar[r]_-x \ar[dd]_j & X \ar[d]^f \\
& Y \ar[d]^g \\
T \ar[r]^-z \ar@{..>}[ruu] & Z
}
$$
dans $\mathcal{C}$.
Choisissons un $2$-isomorphisme
$\gamma \colon z \circ j \to g \circ f \circ x$.
Soit $\mathcal{D}$ la catégorie des flèches pointillées pour
le rectangle extérieur et $\gamma$. Soit $\mathcal{D}'$
la catégorie des flèches pointillées pour le carré aux flèches pleines
$$
\xymatrix{
S \ar[r]_-{f \circ x} \ar[d]_j & Y \ar[d]^g \\
T \ar[r]^-z \ar@{..>}[ru] & Z
}
$$
et $\gamma$. Alors $\mathcal{D}$ est équivalente à
une catégorie $\mathcal{D}''$ qui possède la propriété suivante :
il existe un foncteur $\mathcal{D}'' \to \mathcal{D}'$ qui fait de
$\mathcal{D}''$ une catégorie fibrée en groupoïdes au-dessus de
$\mathcal{D}'$ et dont les catégories fibres sont isomorphes à des catégories
de flèches pointillées pour certains carrés aux flèches pleines de la forme
$$
\xymatrix{
S \ar[r]_-x \ar[d]_j & X \ar[d]^f \\
T \ar[r]^-y \ar@{..>}[ru] & Y
}
$$
et certains choix de $2$-isomorphisme $y \circ j \to f \circ x$.
\end{lemma}

\begin{proof}
Construisons la catégorie $\mathcal{D}''$ dont les objets sont les quintuplets
$(a,\alpha,\beta,b,\eta)$,
où $(a,\alpha,\beta)$ est un objet de $\mathcal{D}$,
$b \colon T \rightarrow Y$ est un $1$-morphisme
et $\eta \colon b \rightarrow f \circ a$ est un $2$-isomorphisme.
Les morphismes
$(a,\alpha,\beta,b,\eta) \rightarrow (a',\alpha',\beta',b',\eta')$
dans $\mathcal{D}''$
sont les paires $(\theta_1,\theta_2)$, où
$\theta_1 \colon a \rightarrow a'$ définit une flèche
$(a, \alpha, \beta) \rightarrow (a', \alpha', \beta')$
dans $\mathcal{D}$ et où $\theta_2 \colon b \rightarrow b'$ est un
$2$-isomorphisme satisfaisant à la condition de compatibilité
$\eta' \circ \theta_2 = (\text{id}_f \star \theta_1) \circ \eta$.

\medskip\noindent
Il existe un foncteur $\mathcal{D}'' \rightarrow \mathcal{D}'$ qui est donné
sur les objets par
$(a, \alpha, \beta, b, \eta) \mapsto
(b, (\text{id}_f \star \alpha) \circ (\eta \star \text{id}_j),
(\text{id}_g \star \eta^{-1}) \circ \beta)$
et sur les flèches par $(\theta_1,\theta_2) \mapsto \theta_2$.
Alors $\mathcal{D}'' \rightarrow \mathcal{D}'$ est fibrée en groupoïdes.

\medskip\noindent
Si $(y, \delta, \epsilon)$ est un objet de $\mathcal{D}'$, notons
$\mathcal{D}_{y,\delta}$
la catégorie des flèches pointillées pour le dernier diagramme affiché, avec
$y \circ j \rightarrow f \circ x$ donné par $\delta$.
Il existe un foncteur $\mathcal{D}_{y,\delta} \rightarrow \mathcal{D}''$ donné
sur les objets par
$(a, \alpha, \eta) \mapsto
(a, \alpha, (\text{id}_g \star \eta) \circ \epsilon, y, \eta)$
et sur les flèches par $\theta \mapsto (\theta, \text{id}_y)$.
Ce foncteur donne un isomorphisme de $\mathcal{D}_{y,\delta}$ vers
la catégorie fibre de $\mathcal{D}'' \rightarrow \mathcal{D}'$ au-dessus de
$(y,\delta,\epsilon)$.

\medskip\noindent
Il existe aussi un foncteur $\mathcal{D} \rightarrow \mathcal{D}''$ qui est
donné sur les objets par
$(a,\alpha,\beta) \mapsto (a,\alpha,\beta,f \circ a, \text{id}_{f \circ a})$
et sur les flèches par
$\theta \mapsto (\theta, \text{id}_f \star \theta)$.
Ce foncteur est pleinement fidèle et essentiellement surjectif ; c'est donc une
équivalence. Les détails sont omis.
\end{proof}








\begin{multicols}{2}[\section{Autres chapitres}]
\noindent
Préliminaires
\begin{enumerate}
\item \hyperref[introduction-section-phantom]{Introduction}
\item \hyperref[conventions-section-phantom]{Conventions}
\item \hyperref[sets-section-phantom]{Théorie des ensembles}
\item \hyperref[categories-section-phantom]{Catégories}
\item \hyperref[topology-section-phantom]{Topologie}
\item \hyperref[sheaves-section-phantom]{Faisceaux sur les espaces}
\item \hyperref[sites-section-phantom]{Sites et faisceaux}
\item \hyperref[stacks-section-phantom]{Champs}
\item \hyperref[fields-section-phantom]{Corps}
\item \hyperref[algebra-section-phantom]{Algèbre commutative}
\item \hyperref[brauer-section-phantom]{Groupes de Brauer}
\item \hyperref[homology-section-phantom]{Algèbre homologique}
\item \hyperref[derived-section-phantom]{Catégories dérivées}
\item \hyperref[simplicial-section-phantom]{Méthodes simpliciales}
\item \hyperref[more-algebra-section-phantom]{Compléments d'algèbre}
\item \hyperref[smoothing-section-phantom]{Lissage des morphismes d'anneaux}
\item \hyperref[modules-section-phantom]{Faisceaux de modules}
\item \hyperref[sites-modules-section-phantom]{Modules sur les sites}
\item \hyperref[injectives-section-phantom]{Objets injectifs}
\item \hyperref[cohomology-section-phantom]{Cohomologie des faisceaux}
\item \hyperref[sites-cohomology-section-phantom]{Cohomologie sur les sites}
\item \hyperref[dga-section-phantom]{Algèbre différentielle graduée}
\item \hyperref[dpa-section-phantom]{Algèbre à puissances divisées}
\item \hyperref[sdga-section-phantom]{Faisceaux différentiels gradués}
\item \hyperref[hypercovering-section-phantom]{Hyperrecouvrements}
\end{enumerate}
Schémas
\begin{enumerate}
\setcounter{enumi}{25}
\item \hyperref[schemes-section-phantom]{Schémas}
\item \hyperref[constructions-section-phantom]{Constructions de schémas}
\item \hyperref[properties-section-phantom]{Propriétés des schémas}
\item \hyperref[morphisms-section-phantom]{Morphismes de schémas}
\item \hyperref[coherent-section-phantom]{Cohomologie des schémas}
\item \hyperref[divisors-section-phantom]{Diviseurs}
\item \hyperref[limits-section-phantom]{Limites de schémas}
\item \hyperref[varieties-section-phantom]{Variétés}
\item \hyperref[topologies-section-phantom]{Topologies sur les schémas}
\item \hyperref[descent-section-phantom]{Descente}
\item \hyperref[perfect-section-phantom]{Catégories dérivées de schémas}
\item \hyperref[more-morphisms-section-phantom]{Compléments sur les morphismes}
\item \hyperref[flat-section-phantom]{Compléments sur la platitude}
\item \hyperref[groupoids-section-phantom]{Schémas en groupoïdes}
\item \hyperref[more-groupoids-section-phantom]{Compléments sur les schémas en groupoïdes}
\item \hyperref[etale-section-phantom]{Morphismes étales de schémas}
\end{enumerate}
Thèmes de théorie des schémas
\begin{enumerate}
\setcounter{enumi}{41}
\item \hyperref[chow-section-phantom]{Homologie de Chow}
\item \hyperref[intersection-section-phantom]{Théorie de l'intersection}
\item \hyperref[pic-section-phantom]{Schémas de Picard des courbes}
\item \hyperref[weil-section-phantom]{Théories cohomologiques de Weil}
\item \hyperref[adequate-section-phantom]{Modules adéquats}
\item \hyperref[dualizing-section-phantom]{Complexes dualisants}
\item \hyperref[duality-section-phantom]{Dualité pour les schémas}
\item \hyperref[discriminant-section-phantom]{Discriminants et différentes}
\item \hyperref[derham-section-phantom]{Cohomologie de de Rham}
\item \hyperref[local-cohomology-section-phantom]{Cohomologie locale}
\item \hyperref[algebraization-section-phantom]{Géométrie algébrique et formelle}
\item \hyperref[curves-section-phantom]{Courbes algébriques}
\item \hyperref[resolve-section-phantom]{Résolution des surfaces}
\item \hyperref[models-section-phantom]{Réduction semi-stable}
\item \hyperref[functors-section-phantom]{Foncteurs et morphismes}
\item \hyperref[equiv-section-phantom]{Catégories dérivées de variétés}
\item \hyperref[pione-section-phantom]{Groupes fondamentaux de schémas}
\item \hyperref[etale-cohomology-section-phantom]{Cohomologie étale}
\item \hyperref[crystalline-section-phantom]{Cohomologie cristalline}
\item \hyperref[proetale-section-phantom]{Cohomologie pro-étale}
\item \hyperref[relative-cycles-section-phantom]{Cycles relatifs}
\item \hyperref[more-etale-section-phantom]{Compléments de cohomologie étale}
\item \hyperref[trace-section-phantom]{La formule des traces}
\end{enumerate}
Espaces algébriques
\begin{enumerate}
\setcounter{enumi}{64}
\item \hyperref[spaces-section-phantom]{Espaces algébriques}
\item \hyperref[spaces-properties-section-phantom]{Propriétés des espaces algébriques}
\item \hyperref[spaces-morphisms-section-phantom]{Morphismes d'espaces algébriques}
\item \hyperref[decent-spaces-section-phantom]{Espaces algébriques décents}
\item \hyperref[spaces-cohomology-section-phantom]{Cohomologie des espaces algébriques}
\item \hyperref[spaces-limits-section-phantom]{Limites d'espaces algébriques}
\item \hyperref[spaces-divisors-section-phantom]{Diviseurs sur les espaces algébriques}
\item \hyperref[spaces-over-fields-section-phantom]{Espaces algébriques sur des corps}
\item \hyperref[spaces-topologies-section-phantom]{Topologies sur les espaces algébriques}
\item \hyperref[spaces-descent-section-phantom]{Descente et espaces algébriques}
\item \hyperref[spaces-perfect-section-phantom]{Catégories dérivées d'espaces}
\item \hyperref[spaces-more-morphisms-section-phantom]{Compléments sur les morphismes d'espaces}
\item \hyperref[spaces-flat-section-phantom]{Platitude sur les espaces algébriques}
\item \hyperref[spaces-groupoids-section-phantom]{Groupoïdes dans les espaces algébriques}
\item \hyperref[spaces-more-groupoids-section-phantom]{Compléments sur les groupoïdes dans les espaces}
\item \hyperref[bootstrap-section-phantom]{Amorçage}
\item \hyperref[spaces-pushouts-section-phantom]{Sommes amalgamées d'espaces algébriques}
\end{enumerate}
Thèmes de géométrie
\begin{enumerate}
\setcounter{enumi}{81}
\item \hyperref[spaces-chow-section-phantom]{Groupes de Chow des espaces}
\item \hyperref[groupoids-quotients-section-phantom]{Quotients de groupoïdes}
\item \hyperref[spaces-more-cohomology-section-phantom]{Compléments sur la cohomologie des espaces}
\item \hyperref[spaces-simplicial-section-phantom]{Espaces simpliciaux}
\item \hyperref[spaces-duality-section-phantom]{Dualité pour les espaces}
\item \hyperref[formal-spaces-section-phantom]{Espaces algébriques formels}
\item \hyperref[restricted-section-phantom]{Algébrisation des espaces formels}
\item \hyperref[spaces-resolve-section-phantom]{Résolution des surfaces revisitée}
\end{enumerate}
Théorie des déformations
\begin{enumerate}
\setcounter{enumi}{89}
\item \hyperref[formal-defos-section-phantom]{Théorie formelle des déformations}
\item \hyperref[defos-section-phantom]{Théorie des déformations}
\item \hyperref[cotangent-section-phantom]{Le complexe cotangent}
\item \hyperref[examples-defos-section-phantom]{Problèmes de déformation}
\end{enumerate}
Champs algébriques
\begin{enumerate}
\setcounter{enumi}{93}
\item \hyperref[algebraic-section-phantom]{Champs algébriques}
\item \hyperref[examples-stacks-section-phantom]{Exemples de champs}
\item \hyperref[stacks-sheaves-section-phantom]{Faisceaux sur les champs algébriques}
\item \hyperref[criteria-section-phantom]{Critères de représentabilité}
\item \hyperref[artin-section-phantom]{Axiomes d'Artin}
\item \hyperref[quot-section-phantom]{Espaces de Quot et de Hilbert}
\item \hyperref[stacks-properties-section-phantom]{Propriétés des champs algébriques}
\item \hyperref[stacks-morphisms-section-phantom]{Morphismes de champs algébriques}
\item \hyperref[stacks-limits-section-phantom]{Limites de champs algébriques}
\item \hyperref[stacks-cohomology-section-phantom]{Cohomologie des champs algébriques}
\item \hyperref[stacks-perfect-section-phantom]{Catégories dérivées de champs}
\item \hyperref[stacks-introduction-section-phantom]{Introduction aux champs algébriques}
\item \hyperref[stacks-more-morphisms-section-phantom]{Compléments sur les morphismes de champs}
\item \hyperref[stacks-geometry-section-phantom]{La géométrie des champs}
\end{enumerate}
Thèmes de théorie des espaces de modules
\begin{enumerate}
\setcounter{enumi}{107}
\item \hyperref[moduli-section-phantom]{Champs de modules}
\item \hyperref[moduli-curves-section-phantom]{Espaces de modules de courbes}
\end{enumerate}
Divers
\begin{enumerate}
\setcounter{enumi}{109}
\item \hyperref[examples-section-phantom]{Exemples}
\item \hyperref[exercises-section-phantom]{Exercices}
\item \hyperref[guide-section-phantom]{Guide bibliographique}
\item \hyperref[desirables-section-phantom]{Desiderata}
\item \hyperref[coding-section-phantom]{Style de programmation}
\item \hyperref[obsolete-section-phantom]{Obsolète}
\item \hyperref[fdl-section-phantom]{Licence de documentation libre GNU}
\item \hyperref[index-section-phantom]{Index généré automatiquement}
\end{enumerate}
\end{multicols}


\bibliography{my}
\bibliographystyle{amsalpha}

\end{document}
